\phantomsection\label{def:shadow-class}%
\phantomsection\label{def:shadow-invariants}%
\providecommand{\Br}{\mathrm{Br}}%
\providecommand{\Drinfdouble}[1]{D(#1)}%
\providecommand{\PhiFA}{\Phi^{\mathrm{FA}}}%
\providecommand{\hCS}{\mathrm{hCS}}%
\providecommand{\Wonepinf}{\mathcal{W}_{1+\infty}}%
\providecommand{\protectedPfaff}[1]{\mathrm{Pf}^{\mathrm{prot}}(#1)}%

\chapter{Quantum Chiral Algebras}
\label{ch:quantum-chiral-algebras}

\begin{remark}[Two-stage locus of the holomorphic Chern--Simons observables]
\label{rem:qca-two-stage-locus}
The six-dimensional holomorphic Chern--Simons observables constructed in
this chapter are the local BV/factorisation-algebra model for the
Stage-$1$ output $\Phi^{\mathrm{FA}}_3(\cA)$. They are not, by
themselves, the Hall algebra or the curve-specialised chiral algebra.
The chiral-algebra-on-a-curve output is the Stage-$2$ specialisation
along chosen data $(\Sigma_2, C)$.  The finite comparison presently
available is the DWR/Ran Rees comparison at fixed truncation
$(N,r,L,m)$: it compares the renormalised hCS factorisation algebra with
the oriented Rees critical Hall layer.  Passing from that finite Rees
layer to the ordinary compact critical CoHA additionally requires the
vanishing-cycle realisation map and pro-compatible inverse-limit descent.
\end{remark}

\begin{theorem}[Holomorphic Chern--Simons observables as \texorpdfstring{$\Ethree$}{Ethree}-Dolbeault--chiral Chevalley--Eilenberg algebra]
\label{thm:hcs-obs-e3-ce-dolbeault}
\[
 \mathrm{Obs}^{\mathrm{cl}}_{\mathrm{hCS}}(X) \;\simeq\; \mathrm{CE}^\bullet_{\bar\partial, \mathrm{chir}}(\cE_{\mathrm{hCS},c}, \cO_X)
\]
as completed classical functions on the compactly supported local dg Lie
algebra. The quantum observables are the Costello--Gwilliam
perturbative factorisation algebra carrying a scale-dependent
effective action, BV operator, renormalisation counterterms, and
anomaly class; the classical CE differential supplies the
$\hbar = 0$ associated graded of the full quantum differential. The
ambient $E_3$ structure is higher factorisation coherence
(commutation up to specified homotopy at each configuration arity);
the connected components of $S^5$ are trivial, so the $\pi_1(S^5) = 0$
contribution is the trivial braid monodromy, and the $E_3$ data lives
in the higher homotopy of $\Conf_n(\C^3)$.
\end{theorem}

The central objects of this chapter are $\Eone$-chiral algebras whose
centres or doubles carry the braided representation theory of quantum
groups.  At $d=2$ the CY-to-chiral output is natively $\Etwo$.  At
$d\geq 3$ the curve-specialised output is natively $\Eone$; the first
$\Etwo$ object is the Drinfeld centre of its representation category, or
equivalently the completed Drinfeld double when a Hall pairing and
completion have been fixed.  The finite DWR/Ran Rees hCS--Hall
comparison lives before this centre step.

\section{Definition and structure}
\label{sec:qca-definition}

\begin{definition}[Quantum chiral algebra in dimension stratified form]
\label{def:qca-dimension-stratified}
A \emph{quantum chiral algebra} in this chapter is a chiral algebra $A$
with finite-dimensional weight spaces together with the following
dimension-dependent structure.
\begin{enumerate}[label=\textup{(\roman*)}]
\item For CY$_2$ input, $A=\Phi_2(\cC)$ is an $\Etwo$-chiral algebra.
Its collision braiding gives an $R$-matrix acting on tensor products of
representations.
\item For CY$_d$ input with $d\geq3$ and admissible specialisation data,
\[
 A=A_{\cC}^{(\Sigma_{d-1},C)}
   =\SpCh_{\Sigma_{d-1},C}\bigl(\PhiFA_d(\cC)\bigr)
\]
is an $\Eone$-chiral algebra on $C$.  The braided $\Etwo$ object
attached to $A$ is the Drinfeld centre
$\cZ(\Rep^{\Eone}(A))$.
\item An expression \(R(z)\in A\otimes A\otimes k((z))\) is used only
after a Hopf, double, or matrix-coefficient realisation of the centre has
been chosen.  Intrinsically, the centre carries a half-braiding natural
transformation (equivalently a universal intertwiner); \(R(z)\) is the
meromorphic matrix-coefficient or evaluation family obtained from that
datum after the chosen realisation and spectral parameters.
\end{enumerate}
\end{definition}

\begin{definition}[Five objects attached to an \texorpdfstring{$\Eone$}{E1}-chiral algebra]
\label{def:qca-five-objects}
Let \(A\) be an \(\Eone\)-chiral algebra on a curve.
\begin{enumerate}[label=\textup{(\roman*)}]
\item \(A\) is the associative chiral algebra, for CY$_3$ obtained after
Stage-$2$ specialisation.
\item \(B(A)\) is its conilpotent chiral bar coalgebra.  On a Koszul
locus the counit \(\Omega B(A)\to A\) is a bar--cobar inversion: $A$
is recovered as the cobar $\Omega$ of $B(A)$, with $B(A)$ a distinct
coalgebra object carrying the twisting/coupling structure on the
universal arrow.
\item \(A^i:=H^\bullet(B(A))\) is the bar cohomology coalgebra, or a
chosen bar-coalgebra model before passing to cohomology.
\item \(A^!:=\mathbb D_{\Ran}(A^i)\), or the corresponding continuous
linear dual in a finite model, is the Koszul dual algebra.
\item \(Z^{\mathrm{der}}_{\mathrm{ch}}(A)\) is the derived chiral centre,
the Hochschild cochain object \(\mathrm{RHom}^{\mathrm{ch}}_{A-A}(A,A)\).
\end{enumerate}
The Drinfeld centre \(\cZ(\Rep^{\Eone}(A))\) is a braided monoidal
category.  It is related to, but not equal to, the derived chiral centre
as a chiral algebra.
\end{definition}

\begin{proposition}[Finite Rees hCS--Hall layer separation]
\label{prop:qca-finite-rees-layer-separation}
Let \(X\) be a CY$_3$ equipped with a finite DWR-good atlas and finite
truncation data \((N,r,L,m)\).  Assume the cyclic contractions,
orientation, grading, trace/symmetry, and face-compatibility hypotheses
of Theorem~\ref{thm:finite-dwr-ran-oriented-hcs-hall-comparison} and
Theorem~\ref{thm:finite-dwr-ran-rees-hcs-hall-compatibility}.  Then:
\begin{enumerate}[label=\textup{(\roman*)}]
\item \(\PhiFA_3(\Perf(X))\) is represented locally by the holomorphic
\(E_3\)-factorisation algebra of hCS observables, while every
curve-specialised output
\[
 A_X^{(\Sigma_2,C)}
   =\SpCh_{\Sigma_2,C}\bigl(\PhiFA_3(\Perf(X))\bigr)
\]
is only \(\Eone\)-chiral.
\item At the fixed bounds \((N,r,L,m)\) there is an oriented Rees
comparison
\[
 \Theta_{\hCS\to\Hall}^{\mathrm{Rees};N,r,L,m}\colon
 \Obs_{\hCS}^{q,\leq(N,r,L,m)}
 \longrightarrow
 \Hall_{\mathrm{crit}}^{\mathrm{Rees,or};N,r,L,m}.
\]
It is a morphism of finite ordered Hall-valued factorisation cosheaves.
It is not yet a morphism to the ordinary compact critical CoHA.
\item The ordinary compact critical CoHA requires a monoidal
vanishing-cycle realisation \(R_{\mathrm{vc}}\).  The completed compact
comparison further requires pro-compatible transition data in
\((N,r,L,m)\); the remaining obstruction is the corresponding
\(\varprojlim^1\)-class.
\item On a single affine chart \(X=\C^3\),
\[
 \CoHA(\C^3)=Y^+(\widehat{\fgl}_1),
\]
the positive half.  The full affine Yangian, its universal \(R\)-matrix,
and the \(W_{1+\infty}\) Fock action arise only after passing to the
Drinfeld double or to \(\cZ(\Rep^{\Eone}(Y^+))\).
\item For a multi-chart CY$_3$, DWR/Ran descent glues the finite Rees
Hall layer chart by chart; the further constructions — the global
Drinfeld double, the global \(W_{1+\infty}\) vertex algebra, and the
\(\Etwo\)-chiral structure on the Drinfeld centre
\(\cZ(\Rep^{\Eone}(A_X^{(\Sigma_2,C)}))\) — are obtained by composing
DWR/Ran descent with the centre/double step.
\end{enumerate}
\end{proposition}

\begin{proof}
The two-stage assertion is Theorem~\ref{thm:phi-two-stage-factorisation}
and the native-level statement is the \(d=3\) case of
Theorem~\ref{thm:cy-to-chiral-d3}.  The finite comparison map is the
DWR/Ran Rees construction of
Theorem~\ref{thm:finite-dwr-ran-oriented-hcs-hall-comparison}, in the
ordered \(E_1\) form recorded by
Theorem~\ref{thm:finite-dwr-ran-rees-hcs-hall-compatibility}.  Those
theorems produce a finite Rees critical Hall target; their hypotheses do
not include the monoidal vanishing-cycle realisation to ordinary
critical CoHA or the inverse-limit compatibility required to leave fixed
\((N,r,L,m)\).  The local \(\C^3\) identification is the
Schiffmann--Vasserot/RSYZ positive-half computation; the full Yangian is
obtained by the Drinfeld double, equivalently by the Drinfeld centre of
the representation category.  Drinfeld centre and completion are
functorial operations after the ordered Hall layer has been supplied, so
they cannot be inserted into the finite Rees comparison map.
\end{proof}

\begin{problem}[Quantum vertex chiral group]
\label{prob:qvcg-construction}
For a smooth proper CY$_3$ category $\cC$ with $X$ the underlying
manifold, construct a \emph{quantum vertex chiral group} \(G(X)\) with:
\begin{enumerate}[label=(\alph*)]
 \item A generalized BKM superalgebra $\mathfrak{g}_X$ with root datum $\mathcal{R}(X)$ extracted from the CY$_3$ geometry (Chapter~\ref{ch:k3-times-e});
 \item A vertex algebra structure $V(\mathfrak{g}_X)$: the vacuum module of $\mathfrak{g}_X$ equipped with the state-field correspondence;
 \item An $E_2$-braided monoidal category of representations $\Rep^{E_2}(V(\mathfrak{g}_X))$;
 \item A denominator identity: the Weyl--Kac--Borcherds character formula for $\mathfrak{g}_X$, which equals an automorphic form $\Phi_X$ (e.g., $\Delta_5$ for $K3 \times E$).
\end{enumerate}
\end{problem}

\begin{remark}[Construction locus of \texorpdfstring{$G(X)$}{G(X)}]
\label{rem:qvcg-construction-locus}
Problem~\ref{prob:qvcg-construction} is a construction problem: a
definition of $G(X)$ awaits an explicit construction, or an axiomatic
characterisation implying existence. The currently constructed pieces
are these:
\begin{itemize}
 \item For toric CY$_3$ without compact $4$-cycles, the CoHA
 \(\mathcal H(Q_X,W_X)\) gives the positive half; its Drinfeld double
 supplies the braided representation theory.
 \item For \(d=2\), \(\Phi_2\) constructs the native \(\Etwo\)-chiral
 algebra.
 \item For \(K3\times E\), the BKM superalgebra
 \(\mathfrak g_{\Delta_5}\) and root datum are supplied by
 Gritsenko--Nikulin.  The curve-specialised chiral algebra is the
 framed object-level \(A_{K3\times E}^{(K3,E)}\) of
 Theorem~\ref{thm:cy-to-chiral-d3}; its braided representation theory
 is obtained only after the centre/double operation.
\end{itemize}
Thus a sentence about \(G(X)\) is a theorem only after it factors
through one of the constructed pieces above, or after
Problem~\ref{prob:qvcg-construction} is solved.
\end{remark}

An $\Etwo$-structure on \(A\), or on \(\cZ(\Rep^{\Eone}(A))\) when
\(A\) is only \(\Eone\), determines a universal $R$-matrix
\[
 R(z) \in A \otimes A \otimes k((z))
\]
satisfying the quantum Yang--Baxter equation after a universal
matrix-coefficient realisation has been chosen.
This $R$-matrix is the $E_2$ analogue of the collision $r$-matrix $r(z) = \mathrm{Res}^{\mathrm{coll}}_{0,2}(\Theta_A)$ from Volume~I (whose pole orders are one less than the OPE poles, due to the $d\log$ extraction; see Vol~I).

\section{The CY \texorpdfstring{$R$}{R}-matrix}
\label{sec:cy-r-matrix}

\begin{construction}[CY \texorpdfstring{$R$}{R}-matrix]
\label{constr:cy-r-matrix}
Given a CY category $\cC$ of dimension $2$ and its quantum chiral algebra $A_\cC = \Phi_2(\cC)$, the CY $R$-matrix is
\[
 R_\cC(z) = \mathrm{Res}^{\mathrm{coll}}_{0,2}(\Theta_{A_\cC})
\]
where $\Theta_{A_\cC}$ is the universal MC element from Volume~I. The pole structure of $R_\cC(z)$ encodes the CY shadow depth classification.
\end{construction}

\section{Drinfeld center and the \texorpdfstring{$E_1 \to E_2$}{E-1 -> E-2} passage}
\label{sec:drinfeld-center-yangian}

The passage from \(E_1\) to \(E_2\) is realised by the Drinfeld centre construction. For \(X=\C^3\), the positive CoHA is \(Y^+(\widehat{\mathfrak{gl}}_1)\); after the Hall pairing and Drinfeld double are imposed, the centre of its \(E_1\)-representation category carries the braided category of the full affine Yangian.

The full quantum group attached to a Calabi--Yau target $X$ is not the
positive half $\Yplus{X}$ alone: the positive half is constructed
intrinsically from the equivariant geometry of moduli of effective
objects, while the full quantum group $G(X)$ is its Drinfeld double, a
strictly higher datum requiring a Hall pairing, a completion, an
integral form, stable-envelope transport, and Drinfeld-double descent.
Conflation of $\Yplus{X}$ with $G(X)$ (e.g.\ ``$\CoHA(\C^3) =
\Wonepinf$'') collapses the two-stage assembly. The next definition and
theorem record the assembly explicitly; their toric specialisation is
\cref{thm:drinfeld-center-coha} below.

\begin{definition}[\label{def:Y-plus-equiv-cohomology}Positive half via
equivariant cohomology of moduli of effective objects]
Let $\cA$ be a Calabi--Yau-$d$ category with target $X$, equipped with
the geometric data \textup{(}equivariant torus action $T$ on
$\cM^+_{\mathrm{eff}}(X)$ with isolated $T$-fixed loci; potential $W
\colon \cM^+_{\mathrm{eff}}(X) \to \C$ from the Calabi--Yau
superpotential; orientation data and Serre-compatible NCCR chart on the
critical locus\textup{)} required for the construction below to be
well posed. The \emph{positive half} of $\cA$ is the equivariant
critical cohomology
\[
 \Yplus{X} \;:=\; H^\bullet_{T}\!\bigl(\cM^+_{\mathrm{eff}}(X),
 \phi_W\bigr)
\]
of the moduli of effective objects with respect to the vanishing-cycle
sheaf $\phi_W$, equipped with the Schiffmann--Vasserot Hall product
arising from extension correspondences
\[
 \cM^+_{\mathrm{eff}}(X) \times \cM^+_{\mathrm{eff}}(X)
 \;\xleftarrow{\;p\;}\;
 \cE\!xt^+(X)
 \;\xrightarrow{\;q\;}\;
 \cM^+_{\mathrm{eff}}(X),
 \qquad
 m \;=\; q_\star \circ p^\star.
\]
At $d = 3$, $\Yplus{X}$ is the cohomological Hall algebra
$\CoHA(X)$ of Schiffmann--Vasserot \cite{SchiffmannVasserot2013}; on
$\C^3$ it gives $\CoHA(\C^3) = \Yplus{\widehat{\fgl}_1}$
\textup{(}\cite{SchiffmannVasserot2013}, identified with the positive
half of the affine Yangian, \emph{not} with $\Wonepinf$\textup{)}.

The comparison data are $(p^\star,q_\star,\phi_W)$. The ambient category is
$T$-equivariant cohomology with critical coefficients, in
$\mathrm{Ch}(\mathrm{Vect}_\Q)$ at the chain level and in the perverse
$\infty$-category of constructible sheaves on
$\cM^+_{\mathrm{eff}}(X)$ at the $(\infty,1)$-categorical level. The construction is
unconditional on $\C^3$ and standard toric CY$_3$ charts; on compact
non-toric CY$_3$ targets it is conditional on a compact critical atlas
and orientation data.
\end{definition}

\begin{theorem}[\label{thm:G-eq-D-Yplus}Drinfeld-double identification of
the full quantum group]
\emph{Hypotheses.} Let $\cA$ be a Calabi--Yau-$d$ category with target
$X$ for which $\Yplus{X}$ is constructed in the sense of
\cref{def:Y-plus-equiv-cohomology}. Assume the following five
data:
\begin{enumerate}[label=\textup{(LD\arabic*)}]
\item \emph{Hall pairing.}
A non-degenerate bilinear pairing $\langle -, - \rangle_{\mathrm{Hall}}
\colon \Yplus{X} \otimes \Yplus{X} \to \kk$ obtained from
Serre/Calabi--Yau duality on $\cM^+_{\mathrm{eff}}(X)$
\textup{(}Schiffmann--Vasserot pairing on critical
cohomology\textup{)}.
\item \emph{Completion.}
A topology / completion of $\Yplus{X}$ \textup{(}grading-completion,
$J$-adic completion, or weight-completion\textup{)} for which
\textup{(LD1)} converges and the antipode is well-defined.
\item \emph{Integral form.}
A subring $\Yplus{X}_{\Z} \subset \Yplus{X}$ \textup{(}or
$\Yplus{X}_{\Z[\hbar]}$\textup{)} compatible with the Hall product,
coproduct, and pairing.
\item \emph{Stable-envelope transport.}
A Maulik--Okounkov stable-envelope realisation
$\mathrm{Stab} \colon \Yplus{X} \to \End(\cF_X)$ on the appropriate
Fock-type representation $\cF_X$, satisfying the chamber-by-chamber
factorisation that intertwines the Hall product with the stable-basis
matrix coefficient \cite{MaulikOkounkov2012}.
\item \emph{Drinfeld-double descent.}
A descent datum identifying the algebra
$\Yplus{X} \bowtie (\Yplus{X})^{\ast,\mathrm{cop}}$,
constructed from \textup{(LD1)--(LD4)}, with the chamber-independent
Drinfeld double, well-defined on the underlying critical cohomology
independent of the chamber and orientation choices used in
\textup{(LD4)}.
\end{enumerate}

\emph{Statement.} Under \textup{(LD1)--(LD5)}, the full quantum group
$G(X)$ associated to $\cA$ is the Drinfeld double of the positive
half $\Yplus{X}$:
\begin{equation}
\label{eq:G-eq-D-Yplus}
 G(X) \;=\; \Drinfdouble{\Yplus{X}}
 \;=\; \Yplus{X} \;\bowtie\; (\Yplus{X})^{\ast,\mathrm{cop}},
\end{equation}
where $\bowtie$ is the Hall-pairing biproduct of \textup{(LD1)}
completed in the topology of \textup{(LD2)} on the integral form
\textup{(LD3)}, and the chamber-independence is supplied by
\textup{(LD4)--(LD5)}. The braided category
$\Rep^{E_1}(\Drinfdouble{\Yplus{X}}) = \Rep(G(X))$ realises the
Drinfeld centre of $\Rep^{E_1}(\Yplus{X})$
\textup{(}Majid 1991; Etingof--Gelaki--Nikshych--Ostrik,
Proposition~7.13.8\textup{)}, and the universal $R$-matrix of
$G(X)$ is the half-braiding inherited from this centre.

\emph{Licensing.} Tag $\beta$ \textup{(}comparison--functor:
$\Drinfdouble{\,\cdot\,}$ and the Hall--pairing biproduct\textup{)}.
The full hypothesis set \textup{(LD1)--(LD5)} is the comparison datum;
the theorem is the equality of two algebras over the same critical
cohomology, not a definition of $G(X)$.

\emph{Constructed loci.}
\begin{enumerate}[label=\textup{(\roman*)}]
\item $X = \C^3$: $G(\C^3) = Y(\widehat{\fgl}_1)$ is the affine Yangian
of $\widehat{\fgl}_1$, with $\Yplus{\C^3} = \CoHA(\C^3) =
\Yplus{\widehat{\fgl}_1}$
\textup{(}\cite{SchiffmannVasserot2013}; cf.\
\cref{thm:drinfeld-center-coha}\textup{(i)}\textup{)}.
\item Toric CY$_3$ without compact $4$-cycles: $G(X)$ is the
quantum-toric Yangian $Y_{\mathrm{tor}}(X)$, constructed from the toric
quiver with potential $(Q_X, W_X)$
\textup{(}cf.\ Chapter~\ref{ch:toric-coha}\textup{)}.
\item $X = $ flat affine $\C^2 \times \C$: \textup{(LD1)--(LD5)} are
all chain-level constructible from
$\cM^+_{\mathrm{eff}}(\C^2 \times \C) \cong \mathrm{Hilb}(\C^2)
\times \mathrm{pt}$ and the Maulik--Okounkov stable-envelope
construction on Hilbert schemes \cite{MaulikOkounkov2012}.
\end{enumerate}

\emph{Compact non-toric loci.}
On compact non-toric CY$_3$ \textup{(}quintic, generic Hartshorne
threefold\textup{)} and on $K3 \times E$, the construction
\textup{(LD1)--(LD5)} of \cref{eq:G-eq-D-Yplus} requires the named
data of Kontsevich--Soibelman \cite{KontsevichSoibelman2008}: critical
CoHA on a compact critical atlas, monoidal vanishing-cycle realisation
$R_{\mathrm{vc}}$, pro-compatible inverse-limit transition data, and
PBW/Hall integral form. On these loci the equality \cref{eq:G-eq-D-Yplus}
is the explicit comparison map, not a black-box definition.
\end{theorem}

\begin{proof}[Sketch on the constructed loci]
On $X = \C^3$: by Schiffmann--Vasserot
\cite{SchiffmannVasserot2013}, $\Yplus{\C^3}$ is identified with the
positive half $\Yplus{\widehat{\fgl}_1}$ of the affine Yangian, with
the Hall product matching the Yangian coproduct on shuffle generators.
The Hall pairing \textup{(LD1)} is the Schiffmann--Vasserot pairing on
shuffle algebras, non-degenerate on each weight space; the completion
\textup{(LD2)} is the standard $\hbar$-adic / weight-completion of the
Yangian; the integral form \textup{(LD3)} is the
$\Z[h_1, h_2, h_3]$-form of the shuffle algebra; the stable-envelope
transport \textup{(LD4)} is supplied by Maulik--Okounkov on
$\mathrm{Hilb}^n(\C^2)$ \cite{MaulikOkounkov2012}; chamber-independence
\textup{(LD5)} is the wall-crossing / $R$-matrix factorisation of
\cite[\S 3]{MaulikOkounkov2012}. The Drinfeld-double construction
$\Yplus{X} \bowtie (\Yplus{X})^{\ast,\mathrm{cop}}$ then identifies
with $Y(\widehat{\fgl}_1)$: the negative generators
$Y^-(\widehat{\fgl}_1)$ are the graded dual of $\Yplus{\widehat{\fgl}_1}$
with reversed coproduct, and the cross-relations are read off from the
Hall pairing \textup{(LD1)} \textup{(}Schiffmann--Vasserot
\cite{SchiffmannVasserot2013}, Theorem~9.1\textup{)}. The toric and
flat cases reduce to $\C^3$ on each toric chart by the same argument
applied chartwise.

The non-toric and $K3 \times E$ cases require
\cite{KontsevichSoibelman2008}'s motivic-/critical-CoHA framework and
are recorded as conditional on the compact-critical atlas, which is the
content of Construction Problem~1 in the Vol~II frontier
\textup{(}\texttt{chiral-bar-cobar-vol2:CLAUDE.md}~\S 9, item 1: the
$K3 \times E$ Drinfeld-double construction with $\protectedPfaff{
\fD_X} = \Delta_5$\textup{)}.
\end{proof}

\begin{remark}[\label{rem:G-eq-D-Yplus-not-Wonepinf}$\Wonepinf$ is a
Fock representation of $G(\C^3)$, not equal to $\CoHA(\C^3)$]
The equality $\CoHA(\C^3) = \Wonepinf$ is forbidden
\textup{(}Vol~II voice table, item 8\textup{)}. The correct relation
is the chain
\[
 \CoHA(\C^3) \;=\; \Yplus{\widehat{\fgl}_1}
 \;\xrightarrow[\textup{(LD1)--(LD5)}]{\Drinfdouble{\,\cdot\,}}\;
 G(\C^3) \;=\; Y(\widehat{\fgl}_1)
 \;\xrightarrow[\text{Gaiotto--Rap\v{c}\'ak}]{\text{vacuum action}}\;
 \Wonepinf,
\]
where the last arrow is the Fock-module action of the affine Yangian on
$\Wonepinf$ \textup{(}Proch\'azka--Rap\v{c}\'ak; Gaiotto--Rap\v{c}\'ak;
cf.\ \cref{thm:drinfeld-center-coha}\textup{(ii)}\textup{)}.
$\Wonepinf$ is a representation of $G(\C^3)$, not $G(\C^3)$ itself.
\end{remark}

\begin{theorem}[Drinfeld center of the CoHA]
\label{thm:drinfeld-center-coha}
Let $Y^+(\widehat{\fgl}_1)$ denote the positive half of the affine Yangian (Schiffmann--Vasserot). Then:
\begin{enumerate}[label=(\roman*)]
 \item The Drinfeld center of its representation category is the braided monoidal category of restricted modules for the Drinfeld double, hence for the full affine Yangian:
 \[
 \cZ(\Rep^{E_1}(Y^+(\widehat{\fgl}_1))) \;\simeq\; \Rep^{E_2}_{\mathrm{res}}(Y(\widehat{\fgl}_1)).
 \]
 \item The completed affine Yangian acts on the $W_{1+\infty}$ vacuum/Fock modules by the Gaiotto--Rap\v{c}\'ak/Proch\'azka--Rap\v{c}\'ak evaluation. This is an action/evaluation statement, not an equivalence of the whole center with an unspecified $\Rep(W_{1+\infty})$ category.
 \item The completed Drinfeld double has graded dimension
 \[
 \mathrm{ch}\,D(Y^+) = M(q)^2,
 \]
where $M(q) = \prod_{n \geq 1} 1/(1-q^n)^n$ is the MacMahon
function.  A chosen $W_{1+\infty}$ vacuum/Fock evaluation has
separate trace factor
\[
 \mathrm{ch}\,\cF_{W_{1+\infty}}=P(q)=\prod_{n \geq 1}(1-q^n)^{-1}.
\]
The product \(M(q)^2P(q)\) is therefore an evaluated trace, not an
intrinsic character of the Drinfeld centre.
\end{enumerate}
\end{theorem}

\begin{proof}
Part~(i) proceeds in two steps. First, the general categorical fact (Majid, 1991; see also Etingof--Gelaki--Nikshych--Ostrik, \emph{Tensor Categories}, Proposition~7.13.8): for a bialgebra $H$ over a field $k$, the Drinfeld center $\cZ(\Rep(H))$ is equivalent as a braided monoidal category to $\Rep(D(H))$, where $D(H) = H \bowtie H^{*\mathrm{cop}}$ is the Drinfeld double. The equivalence sends a half-braiding $(V, \sigma_{V,-})$ to the $D(H)$-module structure determined by $\sigma$. Second, for $H = Y^+(\widehat{\fgl}_1)$ the positive half of the affine Yangian, the Drinfeld double $D(Y^+) = Y^+ \bowtie (Y^+)^{*\mathrm{cop}}$ is identified with the full affine Yangian $Y(\widehat{\fgl}_1)$ (Schiffmann--Vasserot, 2013: the negative generators $Y^-$ are the graded dual of $Y^+$ with reversed coproduct). The $E_2$ braiding on $\Rep(Y)$ is induced by the universal $R$-matrix of the Yangian, which satisfies the YBE (Theorem~\ref{thm:yang-r-matrix-ybe} below). The passage from $\Eone$ to $\Etwo$ is the Drinfeld center functor $\cZ \colon \Eone\text{-}\mathrm{Cat} \to \Etwo\text{-}\mathrm{Cat}$ (this is the categorified analogue of the derived center; Drinfeld center $\cZ(\Rep^{\Eone}(A))$ is a braided monoidal \emph{category}, not a chiral algebra).

Part~(ii) is the standard affine-Yangian/$W_{1+\infty}$ evaluation
dictionary: the completed mode algebra acts on the
$W_{1+\infty}[\lambda]$ vacuum/Fock module. The theorem uses only this
module-level action; it does not assert a categorical equivalence with
all $W$-modules.

Part~(iii): the Drinfeld double character is
\(\mathrm{ch}(Y^+) \cdot \mathrm{ch}(Y^-) = M(q)^2\) (the MacMahon
function squared, since both \(Y^+\) and \(Y^-\) have character \(M(q)\)
by the Schiffmann--Vasserot computation of the graded dimension of the
shuffle algebra).  The vacuum module of \(W_{1+\infty}\) contributes
\(P(q)\) only after the evaluation representation has been chosen.
\end{proof}

\begin{theorem}[Yang \texorpdfstring{$R$}{R}-matrix: YBE and unitarity]
\label{thm:yang-r-matrix-ybe}
Let $h_1, h_2, h_3$ satisfy $h_1 + h_2 + h_3 = 0$ and let
\[
 g(u) = \frac{(u - h_1)(u - h_2)(u - h_3)}{(u + h_1)(u + h_2)(u + h_3)}
\]
be the structure function of the affine Yangian $Y(\widehat{\fgl}_1)$.
The \(R\)-matrix on
\(\mathrm{Fock}_2 \otimes \mathrm{Fock}_2\), where
\(\mathrm{Fock}_2=\langle |\mathrm{row}\rangle,|\mathrm{col}\rangle\rangle\)
is the charge-$2$ space indexed by the partitions \((2)\) and
\((1,1)\), is a \(4\times4\) matrix
\(R(u) \in \End(\C^2 \otimes \C^2)((u))\) satisfying:
\begin{enumerate}[label=(\roman*)]
 \item \emph{Quantum Yang--Baxter equation}:
 \[
 R_{12}(u-v)\, R_{13}(u)\, R_{23}(v) = R_{23}(v)\, R_{13}(u)\, R_{12}(u-v)
 \]
 in $\End(\C^2 \otimes \C^2 \otimes \C^2)((u,v))$. This is verified
 entry-by-entry on the \(8\)-dimensional triple tensor product.
 \item \emph{Unitarity}: $R_{12}(u)\, R_{21}(-u) = \mathrm{id}$.
 \item \emph{$E_2$ nontriviality}: $R(u) \neq \mathrm{id}$ whenever $h_1 h_2 h_3 \neq 0$, i.e., whenever the Calabi--Yau deformation parameters are generic. At the self-dual point $(h_1, h_2, h_3) = (1, 0, -1)$, the structure function $g(u) = 1$ identically ($\sigma_3 = h_1 h_2 h_3 = 0$), the $R$-matrix degenerates to the identity, and the $E_2$ braiding becomes symmetric (the braiding data is trivial, though the operadic level remains $E_2$).
 \item \emph{Classical limit}: $R(u) = 1 - \sigma_2/u + O(1/u^2)$ as $u \to \infty$, where $\sigma_2 = h_1 h_2 + h_1 h_3 + h_2 h_3$. The classical $r$-matrix is $r(u) = -\sigma_2/u$.
\end{enumerate}
\end{theorem}

\begin{proof}
The \(4\times4\) matrix \(R(u)\) gives, after placement in the three
tensor factors, an equality of the \(2^6 = 64\) rational matrix entries of
$R_{12}R_{13}R_{23}$ and $R_{23}R_{13}R_{12}$ in
$u,v,h_1,h_2,h_3$.  Unitarity is ordinary matrix multiplication.  For
the classical limit the structure function satisfies
$g(u) = 1 + O(u^{-2})$, because the $u^{-1}$ coefficient is
$-2(h_1+h_2+h_3)=0$; the $R$-matrix on
$\mathrm{Fock}_2\otimes \mathrm{Fock}_2$ nevertheless has a $u^{-1}$
term from the off-diagonal exchange of boxes.  Hence
$R(u)=1-\sigma_2/u+O(u^{-2})$ and $r(u)=-\sigma_2/u$.
\end{proof}


\begin{remark}[Dimofte's \texorpdfstring{$K$}{K}-matrix and CY shadow depth]
\label{rem:dimofte-k-matrix-cy}
In the PIRSA lecture series (Dimofte, PIRSA 25110067), Dimofte isolates a single chiral operator $K_\cA(z)$ which controls the deviation of the boundary coproduct of an HT boundary algebra from the primitive (Hopf) coproduct. Its explicit form is
\begin{equation}
\label{eq:dimofte-k-matrix-cy}
K_\cA(z) \;=\; z^{\alpha_0}\,\exp\!\Biggl(\sum_{n>0} \frac{\alpha_{-n}}{-n}\, z^{n}\Biggr)\,\exp\!\Biggl(\sum_{n>0} \frac{\alpha_{n}}{-n}\, z^{-n}\Biggr),
\end{equation}
where $\{\alpha_n\}_{n \in \Z}$ are the modes of a free boson identified with the $\mathfrak{u}(1)^r$ Cartan factor of the boundary algebra and $\alpha_0$ is the zero-mode charge. The operator $K_\cA(z)$ lives inside $\cA$, not over it. The canonical Vol~II discussion of equation~\eqref{eq:dimofte-k-matrix-cy}, its diagnostic interpretation, and its interaction with the shadow-depth classification is in the ordered bar chapter of Vol~II (\emph{Vol~II}, remark \ref{rem:dimofte-k-matrix-cy} in \cite{VolII}), together with the BPS-shadow-depth parallel in the HT bulk-boundary chapter. This remark records the CY specialization of that Vol~II content.

Let $\cC$ be a smooth proper Calabi--Yau category of dimension $d \in \{2, 3\}$. For $d = 2$ the CY-to-chiral correspondence $\Phi_2$ of Theorem~CY-A$_2$ produces an $E_2$-chiral algebra $A_\cC = \Phi_2(\cC)$ and the $K$-matrix $K_{A_\cC}(z)$ is well defined by equation~\eqref{eq:dimofte-k-matrix-cy} applied to the boundary algebra of the associated 3d HT theory. For $d = 3$ the same prescription is available only when Theorem~CY-A$_3$ supplies the framed object-level specialisation $A_\cC^{(\Sigma_2,C)}=\Phi_3^{(\Sigma_2,C)}(\cC)$ under H1--H4 and fixed admissible data; the chain-level $\mathbb{S}^3$-framing is an assumption or a verified toric/formal witness, not automatic.

The diagnostic interpretation of $K_{A_\cC}(z)$ refines on the CY side into a statement about the BPS hull of $\cC$. Write $\cC^{\mathrm{BPS}} \subset \cC$ for the subcategory of BPS objects (semistable objects at a chosen stability condition), and let $\Phi_d(\cC^{\mathrm{BPS}})$ denote the image of this subcategory under the CY-to-chiral correspondence $\Phi_d$ (at the CY dimension $d$ of $\cC$). Then $K_{A_\cC}(z)$ is determined by $\Phi_d$ applied to the BPS hull: the Cartan zero-mode $\alpha_0$ records the rank of the $\mathfrak{u}(1)^r$ factor, and the higher modes $\alpha_{\pm n}$ ($n \geq 1$) encode the BPS generating-function corrections. Here ``$\Phi$ applied to the BPS hull'' refers to the CY-to-chiral correspondence $\{\Phi_d\}$ at the appropriate $d$.
\begin{itemize}
\item For toric CY$_3$ without compact $4$-cycles (Chapter~\ref{ch:toric-coha}), $K_{A_\cC}(z)$ is the generating function of Donaldson--Thomas invariants on the fiber of the toric quiver: the mode $\alpha_{-n}$ scales as the DT partition function restricted to classes with $n$ BPS quanta, and the Schiffmann--Vasserot/RSYZ identification of Theorem~\ref{thm:drinfeld-center-coha} realizes the identification $K_{A_\cC}(z) = Z^{\mathrm{DT}}_{\mathrm{fiber}}(z)$.
\item For $K3 \times E$ (Chapter~\ref{ch:k3-times-e}), the $K$-matrix specializes to the Fourier coefficients of the Igusa cusp form $\Phi_{10}^{\mathrm{un}}$: the Fourier--Jacobi expansion $\Phi_{10}^{\mathrm{un}} = \sum_{m \geq 1} \phi_m(\tau,z)\, p^m$ records the BKM root multiplicities on successive modes, and these coefficients are precisely the data encoded by equation~\eqref{eq:dimofte-k-matrix-cy}. The weight of $\Phi_{10}^{\mathrm{un}}$ is $10 = 2 \cdot \kappa_{\mathrm{BKM}}(\Delta_5)$, where $\kappa_{\mathrm{BKM}}(\Delta_5) = 5$ is the Borcherds--Kac--Moody modular characteristic (see Table of Chapter~\ref{ch:k3-times-e}).
\end{itemize}

The Vol~I shadow-depth classification $\{\mathbf{G}, \mathbf{L}, \mathbf{C}, \mathbf{M}\}$ coincides, under $\Phi$, with the $K$-matrix complexity of the CY boundary algebra:
\begin{center}
\renewcommand{\arraystretch}{1.2}
\begin{tabular}{lll}
\toprule
Class & $K_{A_\cC}(z)$ & $d_{\mathrm{alg}}(A_\cC)$ \\
\midrule
$\mathbf{G}$ (trivial CY fiber, $\mathfrak{u}(1)^r$) & $1$ & $0$ \\
$\mathbf{L}$ (rank-one CY Heisenberg deformation) & polynomial, simple pole & $1$ \\
$\mathbf{C}$ (toric CY$_3$, $\beta\gamma$-type) & polynomial, double pole & $2$ \\
$\mathbf{M}$ ($K3$, $W_N$-type) & infinite expansion & $\infty$ \\
\bottomrule
\end{tabular}
\end{center}
Equivalently: $K_{A_\cC}(z) = 1$ iff the boundary coproduct is primitive iff $d_{\mathrm{alg}} = 0$ iff $\cC$ is class $\mathbf{G}$. The correspondence $K_{A_\cC}(z) = 1 \Longleftrightarrow \mathbf{G} \Longleftrightarrow d_{\mathrm{alg}} = 0$ is the CY analogue of the Vol~II biconditional.

The $K$-matrix modifies the coproduct, not the product, so the $r$-matrix vanishing check does not apply directly to $K_{A_\cC}(z)$. The corresponding $r$-matrix check is the affine one: the classical $r$-matrix $k_{\mathrm{ch}}\,\Omega/z$ vanishes at $k_{\mathrm{ch}} = 0$, in which case $A_\cC$ collapses to the trivial Heisenberg and $K_{A_\cC}(z) = 1$, consistent with class $\mathbf{G}$.
\end{remark}


\section{Comparison with the Maulik--Okounkov \texorpdfstring{$R$}{R}-matrix}
\label{sec:mo-rmatrix-comparison}

The Maulik--Okounkov stable-envelope construction provides an
independent geometric source of \(R\)-matrices on symplectic
resolutions, in particular on Nakajima quiver varieties such as
\(\mathrm{Hilb}^n(\C^2)\).  A Calabi--Yau threefold CoHA uses
critical Hall and shuffle data instead.  Thus the MO comparison is a
surface-chart comparison; it is not a stable-envelope construction on
\(\mathrm{Hilb}^n(\C^3)\).

\begin{theorem}[MO/chiral \texorpdfstring{$R$}{R}-matrix comparison on symplectic surface charts]
\label{thm:mo-chiral-rmatrix}
Let \(S\) be a toric symplectic surface chart, an ADE surface chart, or
a Nakajima quiver variety chart for which the Maulik--Okounkov
stable-envelope \(R\)-matrix \(R^{\mathrm{MO}}(u)\) is defined.  Let
\(R^{\mathrm{ch}}(u)\) be the chiral \(R\)-matrix from the \(\Etwo\)-bar
complex of the Drinfeld double \(D(Y^+)\) of the corresponding Hall
positive half (Construction~\ref{constr:cy-r-matrix}).  The \(\Etwo\)-bar
complex here is constructed on the Drinfeld double, not on the CoHA
\(Y^+\) itself, which is natively \(\Eone\); the \(\Etwo\)-braiding
arises from the Drinfeld centre, Theorem~\ref{thm:drinfeld-center-coha}.
Then:
\begin{enumerate}[label=(\roman*)]
 \item Both \(R\)-matrices act on the same fixed-point Fock space
 \[
   \cF_S=\bigoplus_n K_T(\mathrm{Hilb}^n(S)^T)
 \]
 in the Hilbert-scheme surface case, and on the analogous Nakajima
 fixed-point representation in the quiver-variety case.
 \item On each partition $\lambda$, both produce the same diagonal eigenvalue:
 \[
 R^{\mathrm{MO}}_{\lambda\lambda}(u) = R^{\mathrm{ch}}_{\lambda\lambda}(u) = \prod_{s \in \lambda} g(u + c(s)),
 \]
 where $c(s) = h_1 i + h_2 j$ is the content of box $s = (i,j) \in \lambda$.
 \item Both satisfy the same Yang--Baxter equation, unitarity, and crossing symmetry.
\end{enumerate}
\end{theorem}

\begin{proof}
Seven independent verification paths confirm the identification:
\begin{enumerate}
 \item \emph{Direct diagonal computation}: the eigenvalues $\prod_{s \in \lambda} g(u + c(s))$ are computed from the structure function in both constructions.
 \item \emph{$\mathrm{Hilb}^2(\C^2)$ check}: at $n = 2$, the MO construction uses stable envelopes on $\mathrm{Hilb}^2(\C^2)$ (two chambers in the equivariant torus); the resulting $2 \times 2$ $R$-matrix matches the Yangian $R$-matrix.
 \item \emph{Unitarity}: $R_{12}(u) R_{21}(-u) = \mathrm{id}$ holds for both.
 \item \emph{YBE}: both satisfy the same quantum Yang--Baxter equation (Theorem~\ref{thm:yang-r-matrix-ybe}).
 \item \emph{Crossing symmetry}: the crossing relation $R_{12}(u)^{t_1} R_{12}(-u - \kappa_{\mathrm{cr}})^{t_1} = f(u) \cdot \mathrm{id}$ holds for both, with the same scalar function $f(u)$; here $\kappa_{\mathrm{cr}}$ is the Yangian crossing parameter (distinct from the modular characteristic $\kappa_{\mathrm{ch}}$).
 \item \emph{Classical $r$-matrix}: both yield $r(u) = -\sigma_2 / u$ in the $u \to \infty$ limit.
 \item \emph{Schiffmann--Vasserot specialization}: at the SV values $(h_1, h_2, h_3) = (1, -1-\epsilon, \epsilon)$, both $R$-matrices specialize to the CoHA shuffle algebra braiding.
\end{enumerate}
\end{proof}

\begin{remark}
\label{rem:mo-comparison-scope}
For \(\C^3\), the framed Donaldson--Thomas space
\(\mathrm{Hilb}^n(\C^3)\) supplies a Fock module for the critical
CoHA/shuffle algebra; it is not a Nakajima quiver variety carrying the
Maulik--Okounkov stable envelope.  In the \(K3\times E\) lane the MO
computation is local: it is made on ADE/Kummer surface charts and then
enters the compact theory only through the Hall/DT comparison datum.
For a general CY threefold without torus action, the MO construction is
not directly available, and the chiral \(R\)-matrix from
Construction~\ref{constr:cy-r-matrix} is the primary object.
\end{remark}

\section{Quantum chiral algebras from holomorphic Chern--Simons theory}
\label{sec:qca-from-qg}

Holomorphic Chern--Simons theory on a CY$_d$ manifold $Y$ is the physical realisation of stage~1 of the two-stage factorisation
\[
 \mathrm{CY\text{-}cat}_d
 \;\xrightarrow{\;\PhiFA_d\;}
 \EdHolFA(Y)
 \;\xrightarrow{\;\SpCh_{\Sigma_{d-1},C}\;}
 \EnHolFA(C)
\]
of Definition~\ref{def:phi-fa-and-sp} and Theorem~\ref{thm:phi-two-stage-factorisation}. Stage~1 ($\PhiFA_d$) outputs the canonical $E_d$-holomorphic factorization algebra on $Y$; stage~2 ($\SpCh_{\Sigma_{d-1},C}$) specialises it to an $E_n$-chiral algebra on a curve $C \subset Y$ by pushing forward along the normal cycle $\Sigma_{d-1} = N_{C/Y}$. The Omega-background parameters $(h_1,\ldots,h_d)$ satisfying $\sum h_i = 0$ are equivariant weights of the chosen torus chart and of the normal directions; they deform the factorisation model and survive its Stage-$2$ specialisation. Three regimes of $d$ realise this two-stage picture concretely.

\begin{construction}[Holomorphic Chern--Simons hierarchy]
\label{constr:hol-cs-hierarchy}
Let $\frakg$ be a finite-dimensional Lie algebra equipped with an invariant bilinear form $\langle -, - \rangle_{\frakg} \colon \frakg \otimes \frakg \to \C$. Holomorphic Chern--Simons theory with gauge algebra $\frakg$ on a complex manifold $M$ has action
\[
 S_{\mathrm{hCS}}(A) \;=\; \frac{1}{2} \int_M \Omega \wedge \langle A \wedge \dbar A + \tfrac{2}{3} A \wedge A \wedge A \rangle_{\frakg},
\]
where $A \in \Omega^{0,1}(M; \frakg)$ is a $(0,1)$-connection and $\Omega$ is a holomorphic volume form on $M$ (when $M$ is Calabi--Yau) or a partial volume form (when $M$ is a product). The three regimes:
\begin{enumerate}[label=\textup{(\roman*)}]
 \item \emph{3d holomorphic CS on $\Sigma \times \R$} (Costello~2007, Costello--Gwilliam~2021):
 the boundary on $\Sigma$ produces the affine Kac--Moody vertex algebra $V_k(\frakg)$ at level $k$ determined by $\langle -, - \rangle_{\frakg}$. Its Beilinson--Drinfeld factorization is symmetric only on disjoint configurations; the collision OPE carries the non-commutative current algebra. The ordered $\Eone$ refinement used by Vol.~I and the Vol.~II Swiss-cheese structure records that collision sector, while PVA descent applies only after the Li-graded classical limit and the finite-type hypotheses named in Vol.~II.

 \item \emph{5d holomorphic CS on $\C^2 \times \R$} (Costello~2013, ``Supersymmetric gauge theory and the Yangian''):
 the Omega-background parameters on $\C^2$ produce the Yangian
 vertex-algebra boundary object \(Y^{\mathrm{VOA}}_\varepsilon(\frakg)\)
 in the Costello--Gaiotto--Yagi simply-laced setting.  The affine
 Yangian \(Y(\widehat{\fgl}_1)\) and its higher-rank affine or toroidal
 analogues require the additional Hall, stable-envelope, or
 double/centre input appropriate to the chosen chart.  The positive
 half $Y^+(\widehat{\fgl}_1)$ is the CoHA; it is associative, not a
 chiral algebra.  The full affine Yangian is recovered by the Drinfeld
 double where the Hall pairing is available.

 \item \emph{6d holomorphic theory on $\C^3$} (Costello--Li holomorphic BV locality and Costello's Omega-background construction):
 the observables on $\C^3$, when the one-loop anomaly is cancelled, carry the holomorphic $E_3$ factorisation structure of Definition~\ref{def:cy3-hcs-bv-complex}. Projection to $\C^2 \subset \C^3$ gives an $\Etwo$-chiral shadow; projection to $C \subset \C^3$ gives an $\Eone$-chiral shadow. For $\frakg = \fgl_1$, the finite Rees hCS--Hall comparison identifies the ordered Hall layer at fixed DWR/Ran bounds.  A quantum toroidal algebra \(U_{q,t}(\widehat{\widehat{\fgl}}_1)\) is an expected curve-level projection after completion and centre/double operations, not a direct consequence of Costello--Francis--Gwilliam~\cite{CFG2026}.
\end{enumerate}
\end{construction}

\begin{remark}[6d theory is not standard holomorphic CS]
\label{rem:6d-not-lagrangian}
The naive field $A\in\Omega^{0,1}(\C^3,\mathfrak g)$ alone does not contain all BV components: $\Omega\wedge A\wedge\dbar A$ has type $(3,3)$ only after the superfield multiplication is read in the full BV complex $\Omega^{0,\bullet}(\C^3,\mathfrak g)[1]$ and the $\Omega^{0,3}$ component is taken. Thus the correct local hCS presentation is the BV superfield action of Definition~\ref{def:cy3-hcs-bv-complex}, not a bare $(0,1)$ connection calculation. Costello--Francis--Gwilliam~\cite{CFG2026} proves the ordinary $3$d Chern--Simons $E_3$ factorisation-homology theorem. The $6$d hCS--Hall input used here is instead the finite DWR/Ran Rees comparison of Proposition~\ref{prop:qca-finite-rees-layer-separation}; ordinary compact critical CoHA still requires the vanishing-cycle realisation and pro-descent beyond fixed bounds. The mixed-holomorphic-topological local model on $\R^2_{\mathrm{top}} \times \C^2_{\mathrm{hol}}$ with Hamiltonian BF sector~\cite{LorgatMixedHTStrings} is the level-$1$ physical realisation lane of Stage-$1$ at toric loci; the global de Rham obstruction $\mathrm{Ob}^{\mathrm{BMK\text{-}curr}}_N$ named in~\cite{LorgatMixedHTStrings} Chapter~\ref{ch:cy-to-chiral} lines~$3207$--$3266$ blocks the formal-Darboux $\Rightarrow$ global compact theory promotion outside the verified locus.
\end{remark}

\begin{theorem}[Three Atiyah-sourced cocycles on compact CY\texorpdfstring{$_3$}{-3}]
\label{thm:three-atiyah-cocycles-qca}
For $X$ a compact Calabi--Yau threefold with $\mathrm{At}(T_X) \in H^1(X, \Omega^1_X \otimes \End T_X)$ the Atiyah class of the holomorphic tangent bundle, three cocycles on $X$ arise and enter the hCS presentation of $\PhiFA_3(\Perf(X))$ at distinct orders:
\begin{enumerate}
\item \emph{Generator.} $\At(T_X)$ is the binary bracket $\ell_2$ of the Kapranov $L_\infty$-algebra on $T_X[-1]$ (Kapranov 1999 Proposition~4.4). Under the C\u ald\u araru--C\u ald\u araru--Tu twisted HKR~\cite{CaldararuCaldararuTu2014}, $\At(T_X)$ is cohomologically represented by the defect from HKR being a Gerstenhaber morphism.
\item \emph{Tree-level Yukawa.} $Y_3 \in H^{0,3}(X) \simeq \mathrm{Sym}^3 H^1(X, T_X)^\vee$ is the Kapranov $\ell_3^{\min}$-bracket,
\[
Y_3(v_1, v_2, v_3) \;=\; \int_X \Omega_X \wedge \mathrm{At}(v_1) \cup \mathrm{At}(v_2) \cup \mathrm{At}(v_3),
\]
and enters the classical hCS action as the cubic correction $\int_X \Omega_X \wedge Y_3 \wedge \langle \cA, \cA, \cA \rangle$ beyond the $\bar\partial$-kinetic plus Chern--Simons cubic. Physically: the Bershadsky--Cecotti--Ooguri--Vafa tree-level three-point Yukawa coupling on the moduli of complex structures.
\item \emph{One-loop determinant-line/gravitational datum.} The BCOV
one-loop contribution is a determinant-line or gravitational anomaly
datum, with coefficient controlled by characteristic-class integrals
such as the Euler characteristic. It is not a class
$\alpha_{\mathrm{BCOV}}\in H^{0,1}(X)$ and it is not
$\frac{\chi(X)}{24}[\Omega_X]^{0,1}$; a holomorphic volume form on a
CY$_3$ has no $(0,1)$ component. This datum is also separate from the
local hCS gauge anomaly, whose CY$_3$ slot is quartic as in
Proposition~\ref{prop:cy3-hcs-quartic-anomaly-slot}.
\end{enumerate}
The tree-level Atiyah/Yukawa data, the local quartic hCS gauge anomaly,
and the BCOV determinant-line/gravitational datum coexist on compact
non-flat $X$, but they live in different complexes; the three are
distinct cohomology classes.
\end{theorem}

\begin{remark}[Geometry at \texorpdfstring{$X = K3 \times E$}{X = K3 x E}]
\label{rem:three-cocycles-at-K3xE-qca}
At $X = K3 \times E$ the K\"unneth decomposition gives
$\dim H^{0,3}(K3 \times E) = \dim H^{0,2}(K3) \cdot \dim H^{0,1}(E) = 1$
and $\dim H^{0,1}(K3 \times E) = 1$. These are ambient Hodge slots,
not automatic obstruction classes. On the formal $K3\times E$ locus the
elliptic tangent Atiyah class is zero, the K3 factor has no
$\Omega^3_{K3}$ receptacle for the Kuranishi cubic, and the coefficient
$\chi(K3 \times E)=\chi(K3)\chi(E)=24\cdot0=0$ sets the BCOV one-loop
coefficient to zero; this matches the K\"unneth-multiplicative
total-space reading $\kcat(K3 \times E)=\chi(\cO_{K3\times E})=0$.
\end{remark}

\begin{theorem}[Boundary chiral algebra at \texorpdfstring{$d = 5$}{d = 5}: the affine Yangian]
\label{thm:5d-boundary-yangian}
Let $h_1, h_2, h_3 \in \C$ satisfy $h_1 + h_2 + h_3 = 0$. The observables of 5d holomorphic Chern--Simons theory on $\C^2_{h_1, h_2} \times \R$ with gauge algebra $\fgl_1$, projected to a holomorphic curve $C \subset \C^2$, form an $\Eone$-chiral algebra on $C$ isomorphic to the affine Yangian $Y(\widehat{\fgl}_1)$ with structure function
\[
 g(u) = \frac{(u - h_1)(u - h_2)(u - h_3)}{(u + h_1)(u + h_2)(u + h_3)}.
\]
The completed affine Yangian acts on the $W_{1+\infty}$ vacuum/Fock
modules by the Proch\'azka--Rap\v{c}\'ak/Gaiotto--Rap\v{c}\'ak
evaluation dictionary; the statement used here is this completed
mode-action, not an identification of $Y^+$ or its full representation
category with the vertex algebra. References: Costello~2013;
Proch\'azka--Rap\v{c}\'ak~2019.
\end{theorem}

\begin{conjecture}[Boundary chiral algebra at \texorpdfstring{$d = 6$}{d = 6}: quantum toroidal from factorization]
\label{conj:6d-boundary-toroidal}
\ClaimStatusConjectured
The $E_3$-holomorphic factorization algebra of 6d hCS observables on
$\C^3_{h_1,h_2,h_3}$, with $h_1+h_2+h_3=0$, projected to an
$\Eone$-chiral algebra on a curve $C\subset\C^3$, is expected to be the
quantum toroidal algebra $U_{q,t}(\widehat{\widehat{\fgl}}_1)$ with
parameters $q=e^{2\pi i h_1/h_3}$ and $t=e^{-2\pi i h_2/h_3}$. The
construction is conditional on the 6d hCS BV--BRST complex and the
Costello--Gwilliam factorisation-algebra comparison. The intermediate
$\Etwo$-projection to $\C^2 \subset \C^3$
should recover the affine Yangian $Y(\widehat{\fgl}_1)$ of
Theorem~\ref{thm:5d-boundary-yangian}; this compatibility between the 6d
projection and the independent 5d construction is itself part of the
conjecture. The $E_3$ structure on $\C^3$ is the \emph{source} of the
second affinization (the second hat in $\widehat{\widehat{\fgl}}_1$); the
first affinization comes from the $\Etwo$ factorization on $\C^2$.
\end{conjecture}

\begin{proposition}[Non-abelian hCS pushforward on \texorpdfstring{$K3\times E$}{K3 x E} and the Gritsenko \texorpdfstring{$\fg_{\Delta_5}$}{g-Delta-5}]
\label{prop:hcs-pushforward-k3e-delta5}
Let $X = K3 \times E$ with holomorphic volume form $\Omega_X = \omega_{K3} \wedge d\tau_E$ and projection $\pi \colon X \to E$. Let $\cE_{\hCS}^{\fg_{\Delta_5},\mathrm{re}}$ denote the BV field complex of $6$D holomorphic Chern--Simons on $X$ with gauge datum restricted to the real-root subalgebra $\fg_{\Delta_5}^{\mathrm{re}} = \bigoplus_{\alpha \in \Delta^{\mathrm{re}}} \fsl_2^{(\alpha)}$ (Gritsenko--Nikulin 1998 Theorem~2.1, rank-$3$ hyperbolic lattice $\Lambda^{2,1}_{II}\subset\Lambda^{3,2}$). Then:
\begin{enumerate}[label=\textup{(\roman*)}]
 \item \emph{Anomaly split.} The $1$-loop anomaly factorises as
 \[
 A_{\hCS}^{(1)}(\fg_{\Delta_5}^{\mathrm{re}}) \;=\; A_{\mathrm{w.f.}} + A_{\mathrm{anom}}, \qquad A_{\mathrm{w.f.}} = -\frac{C_2|_{\mathrm{re}}}{(2\pi)^3},
 \]
 with $A_{\mathrm{anom}}$ the four-vertex box wheel weighted by the
 symmetric quartic adjoint trace
 $q^{abcd} = \operatorname{tr}_{\mathrm{ad}}(T^{(a}T^bT^cT^{d)})$
 (Definition~\ref{def:6d-admissible-gauge-algebra}). The
 kinetic-renormalisation piece $A_{\mathrm{w.f.}}$ is scheme-dependent
 and absorbed by the Costello--Li propagator counter-term. For each
 real-root $\fsl_2^{(\alpha)}$ taken as gauge datum, the primitive
 quartic part of $q^{abcd}$ vanishes --- $A_1$ admits no primitive
 quartic Casimir
 (Lemma~\ref{lem:quartic-casimir-classification}) --- leaving a
 multiple of $(C_2)^2$, which on $K3 \times E$ is cancelled by the
 integrated vanishing of item~(ii); the cross-block couplings between
 distinct real-root copies belong to the compact extension of
 Problem~\ref{op:hcs-pushforward-k3e-delta5-compact-bv}.
 \item \emph{Quartic obstruction coefficient.} The compact-CY$_3$ quartic anomaly polynomial $I_4 \in \Sym^4(\fg_{\Delta_5}^\vee)^{\fg_{\Delta_5}}$ of Conjecture~\ref{conj:quintic-hcs-kanom} integrates to
 \[
 \int_X \mathrm{td}(T_X) \wedge I_4(F_{\nabla}) \;=\; \frac{\chi_{\mathrm{top}}(X)}{24} \cdot \mathrm{tr}_{\mathrm{ad}}(F^4) \;=\; 0,
 \]
 since $\chi_{\mathrm{top}}(K3 \times E) = \chi_{\mathrm{top}}(K3)\,\chi_{\mathrm{top}}(E) = 24 \cdot 0 = 0$. The obstruction coefficient vanishes on $K3 \times E$ without invoking a Green--Schwarz mechanism.
 \item \emph{Pushforward identification.} Under the fibre integration $\pi_* \colon \PhiFA_3(\Perf(X)) \to \PhiFA_1(\Perf(E))$ specialised along $\Sigma_2 = K3 \subset X$, the renormalised hCS factorisation algebra descends to the $\Eone$-chiral algebra on $E$
 \[
 \pi_* \Obs_{\hCS}^{\fg_{\Delta_5}^{\mathrm{re}}}(X) \;\simeq\; V^{\mathrm{ch}}(\fg_{\Delta_5}^{\mathrm{re}}),
 \]
 where the right-hand side is the real-root subenvelope of the
Gritsenko--Nikulin BKM superalgebra. Its Weyl factor is
 \[
 e^{-\rho_W}\prod_{\alpha\in\Delta^{\mathrm{re}}_+}
 \bigl(1-e^{-\alpha}\bigr)^{-1},
 \]
and its Borcherds-weight normalisation is fixed by the $N=1$ identity
$\kBKM(\Delta_5)=c_1(0)/2=5$. The full imaginary-root multiplicity
extension is Theorem~\ref{thm:hcs-pushforward-imaginary-roots-singular-theta}.
 \item \emph{Composite resolvent.} The Costello--Li propagator on $X$ factorises as $P_X = P_{K3}^{(0,1)} \boxtimes P_E^{(0,1)}$; fibre-integrating $P_{K3}^{(0,1)}$ against $\omega_{K3}$ extracts the K3-twined elliptic genus, and the composite $\pi_* P_X = \eta(\tau_E)^{24}$ reproduces the weight-$12$ factor of the unnormalised square $\Phi_{10}^{\mathrm{un}} = \Delta_5^2$ via the Kawai--Nam/DMVV product.
\end{enumerate}
\end{proposition}

\begin{proof}[Sketch]
(i): The splitting $A^{(1)} = A_{\mathrm{w.f.}} + A_{\mathrm{anom}}$ is Costello--Li 2016 Proposition~3.7; the kinetic term is scheme-dependent rather than cohomological. On a $3$-fold the anomalous wheel is the four-vertex box with the quartic adjoint trace, not the three-vertex triangle (the triangle is the anomalous wheel in holomorphic dimension~$2$). Restriction to $\fg_{\Delta_5}^{\mathrm{re}}$: real-root subalgebras are simply-laced rank-$1$ ($\fsl_2$) copies because $\fg_{\Delta_5}$ has real roots of squared length~$2$ in the hyperbolic core $\Lambda^{2,1}_{II}$ (Gritsenko--Nikulin 1998 §2). For $\fsl_2 = A_1$, $(\Sym^4 \fsl_2^\vee)^{\fsl_2} = \C\,(C_2)^2$ with no primitive quartic generator (exponents $\{1\}$; Lemma~\ref{lem:quartic-casimir-classification}), so the per-block box tensor is a pure $(C_2)^2$-multiple; its integrated coefficient on $K3 \times E$ vanishes by item~(ii).

(ii): Grothendieck--Riemann--Roch for the anomaly polynomial (Costello 2013 §9) yields the total coefficient $\chi_{\mathrm{top}}(X)/24$; Kunneth gives $\chi_{\mathrm{top}}(K3 \times E) = 24 \cdot 0 = 0$.

(iii)--(iv): The pushforward $\pi_*$ is the Costello--Gwilliam fibre-integration functor on factorisation algebras (Costello--Gwilliam 2021 Vol II §3.5). Fibre-integrating the BV action $\int_X \Omega_X \wedge \mathrm{tr}(A \bar\partial A + \tfrac23 A^3)$ over $K3$ produces the elliptic-genus-weighted 2d chiral theory on $E$ at the real-root subsystem; the denominator identity is Gritsenko--Nikulin 1998 Theorem~2.1 with Borcherds-lift coefficient extraction. The Borcherds-weight component is the $N=1$ case of Theorem~\ref{thm:bl-master-kappa-bkm-chain-level}; it gives $\kBKM(\Delta_5)=c_1(0)/2=5$. The composite-resolvent identity (iv) is the unnormalised Borcherds square $\Phi_{10}^{\mathrm{un}}(Z)=\Delta_5(Z)^2$ with product expansion $\prod_{(n,l,m) > 0}(1 - p^n r^l q^m)^{c(4nm - l^2)}$; the $\eta^{24}$ factor is the weight-$12$ prefactor from the elliptic direction, matching Dijkgraaf--Moore--Verlinde--Verlinde 1997.
\end{proof}

\begin{problem}[Compact BV anomaly cancellation beyond the real-root subsystem]
\label{op:hcs-pushforward-k3e-delta5-compact-bv}
Extend Proposition~\ref{prop:hcs-pushforward-k3e-delta5} from
\(\fg_{\Delta_5}^{\mathrm{re}}\) to the full compact hCS field theory
with imaginary-root sectors and prove the corresponding anomaly
cancellation as a BV chain-level statement. The Borcherds imaginary-root
multiplicities are supplied by
Theorem~\ref{thm:hcs-pushforward-imaginary-roots-singular-theta}; the
missing datum here is the compact BV cancellation and OPE-compatible
transport, not the automorphic denominator formula.
\end{problem}

\begin{conjecture}[Three-path chain-level triangulation at \texorpdfstring{$N = 1$}{N=1}]
\label{cor:k3e-delta5-three-path-triangulation}
\ClaimStatusConjectured
Assume the BV pushforward, the Maulik--Okounkov stable-envelope residue, and the modular-operadic trace have been constructed on the same K3-fibre branch and compared through the projected period correspondence \(\pi_{\mathrm{Sp}_4}\circ\mathscr P_{K3}\). Then the three Borcherds-scalar readings are expected to satisfy
\[
 \kBKM^{\mathrm{BV}}(\fg_{\Delta_5})
 \;=\;
 \kBKM^{\mathrm{MO}}(\fg_{\Delta_5})
 \;=\;
 \mathrm{STr}^{\mathrm{Borch}}_{\mathrm{mod.op}}\!\bigl[B^{E_3}\bigl(\PhiFA_3(\Perf(K3 \times E))\bigr)\bigr]
 \;=\; \frac{c_1(0)}{2} \;=\; 5.
\]
The compact Hodge supertrace of \(K3\times E\) is a different scalar:
\[
 \kappa_{\mathrm{ch}}^{\mathrm{Hodge}}(K3\times E)
 =
 \sum_q(-1)^q h^{0,q}(K3\times E)
 =
 0.
\]
It is not one of the equal quantities above.
\end{conjecture}

\begin{proof}[Comparison calculation]
Hodge path. The chain-level Hodge supertrace on \(\mathrm{HH}^\bullet(D^b\Coh(K3\times E))=\mathrm{PV}^\bullet(K3\times E)\) gives the total-space value \(0\) by Künneth and Serre duality. It supplies a consistency check on compactness; it does not produce the weight-\(5\) Siegel form.
BV path. Proposition~\ref{prop:hcs-pushforward-k3e-delta5}(ii)--(iv) integrates the Costello--Li--Yagi propagator $P_X = P_{K3}^{(0,1)} \boxtimes P_E^{(0,1)}$ at one loop; the wheel-graph sum localises at the Humbert locus $H_1$. In the divisor normalisation $\mathrm{div}(\Delta_5)=H_1+2H_4$, one has $\operatorname{ord}_{H_1}(\Delta_5)=1$ and $\operatorname{ord}_{H_4}(\Delta_5)=2$; the number $5$ entering this path is the Borcherds weight $\kBKM(\Delta_5)=c_1(0)/2$, not an order of vanishing (Borcherds 1998 Theorem 13.3, Gritsenko 1999 Theorem 1.2).
MO path. The Schiffmann--Vasserot K3 cohomological Hall algebra carries a Maulik--Okounkov stable-envelope $R$-matrix; its residue at the Mukai-lattice wall is given by the Bruinier theta-integrand restriction to $H_1$ (Bruinier 2002 Theorem 4.23) with Weyl-denominator normalisation factor $1/2$.
Pairwise compatibility. The required chain-homotopies must be built in the same completed ambient category. Until this is done, the equality is a conjectural triangulation of the Borcherds scalar, not an unconditional Hodge computation.
Master theorem. The universal Borcherds identity $\kBKM(\Phi_N)=c_N(0)/2$ fixes the terminal scalar. It does not by itself construct the modular-operadic trace on \(B^{E_3}(\PhiFA_3(\cC_N))\).
Primary-literature anchors: Gritsenko--Nikulin 1995 Part~II Theorem~2.1; Borcherds 1998 Theorem 13.3; Gritsenko 1999 Theorem 6.1; Bruinier 2002 Proposition 5.1; Kapranov 1999 Theorem 2.6; Costello--Li 2016 Proposition 5.2; Lorgat 2020.
\end{proof}

\begin{remark}[Real-root and compact-CY\texorpdfstring{$_3$}{3} loci]
\label{rem:hcs-pushforward-scope}
Proposition~\ref{prop:hcs-pushforward-k3e-delta5} is rigorous at the real-root subalgebra level and on the $K3 \times E$ formal locus. The extension to all (real + imaginary) roots of $\fg_{\Delta_5}$ is supplied by the Borcherds imaginary-root multiplicity theorem, whose chiral-avatar transport is the content of Theorem~\ref{thm:hcs-pushforward-imaginary-roots-singular-theta} below. For compact CY$_3$ other than $K3 \times E$ (quintic, generic), the quartic obstruction coefficient $\chi_{\mathrm{top}}(X)/24 \neq 0$ and the analogue of (ii) fails; this is Conjecture~\ref{conj:quintic-hcs-kanom}. The anomaly-split discipline of (i) is universal. The $\fsl_{N \geq 3}$ extension is blocked by the non-zero cubic tensor; the non-simply-laced folding problem is open for the different reason that the twisted comparison is not constructed by the rank-one real-root argument.
\end{remark}

\begin{theorem}[Imaginary-root extension of the \texorpdfstring{$K3 \times E$}{K3 x E} hCS pushforward via the chiral singular theta correspondence]
\label{thm:hcs-pushforward-imaginary-roots-singular-theta}
Let $X = K3 \times E$, $\pi \colon X \to E$, \(L=\Lambda^{3,2}\) the
paramodular lattice of signature \((3,2)\) with Grassmannian
\(\mathcal G(L)\), and
\[
F=\phi_{0,1}^{\mathrm{EZ}}\in
J_{0,1}^{\mathrm{wk}}(\mathrm{SL}_2(\mathbb Z))
\]
the Eichler--Zagier generator of weight \(0\) and index \(1\), with
\(c(-1)=1\), \(c(0)=10\), \(c(3)=-64\), \(c(4)=108\). The physical K3
elliptic genus is \(Z_{K3}=2\phi_{0,1}^{\mathrm{EZ}}\); the
\(\Delta_5\) Borcherds denominator uses the Eichler--Zagier seed \(F\).
Extend
Proposition~\ref{prop:hcs-pushforward-k3e-delta5} past the real-root
subsystem as follows.
\begin{enumerate}[label=\textup{(\roman*)}]
 \item \emph{Chiral avatar of the singular theta integral.} The $K3$-fibre integration of the hCS factorisation algebra is representable as the regularised Siegel theta integral of $F$ against the lattice $L$:
 \[
 \pi_*^{K3} \Obs_{\hCS}^{\fg_{\Delta_5}}(X)(Z) \;\simeq\; \int_{\mathcal{F}}^{\mathrm{reg}} F(\tau) \, \Theta_L(\tau, \bar\tau, Z) \, \frac{d\tau \wedge d\bar\tau}{y^2}, \qquad Z \in \Pi^+_L + i\,\mathcal{G}(L),
 \]
 where $\Theta_L$ is the Siegel theta kernel, $\mathcal{F} = \mathbb{H} / \mathrm{SL}_2(\mathbb{Z})$ the modular fundamental domain, and $\mathrm{reg}$ is the Harvey--Moore regularisation (Borcherds~$1998$ \S6). The identification is the Costello--Gwilliam holomorphic-topological tensor theorem (Costello--Gwilliam~$2017$ Vol~II Theorem~$4.10.1$) applied to the holomorphic $K3$-fibre coupled to the topological $\R^2$-direction of the $E_3$-factorisation datum on $X$.
 \item \emph{Denominator-formula consistency.} Expanding \textup{(i)} near the level-$0$ cusp $F_\rho$ with Weyl vector $\rho_W \in L$ yields the Borcherds product
 \[
 \pi_* \Obs_{\hCS}^{\fg_{\Delta_5}}(X)(Z) \;=\; e^{2\pi i (\rho_W, Z)} \prod_{\substack{\lambda \in L^+ \\ (\lambda, \rho_W) > 0}} \bigl(1 - e^{2\pi i (\lambda, Z)}\bigr)^{c\bigl(-\lambda^2/2\bigr)} \;=\; \Delta_5(Z),
 \]
 with $c\bigl(-\lambda^2/2\bigr) = c_{\phi_{0,1}}(n, \ell)$ for $\lambda = (n, \ell, m)$ and $\lambda^2 = 2(nm - \ell^2/4)$. The root multiplicity of $\alpha \in \Delta_+(\fg_{\Delta_5})$ equals $\mathrm{mult}(\alpha) = c\bigl(-\alpha^2/2\bigr)$ for all $\alpha$, real and imaginary, matching Gritsenko--Nikulin~$1998$ Theorem~$2.1$.
 \item \emph{Imaginary-root multiplicities at low norm.} At the null roots $\alpha \in \{(1,0,0), (0,0,1)\}$ with $\alpha^2 = 0$: $\mathrm{mult}(\alpha) = c(0) = 10$. At the first strictly imaginary root stratum $\alpha^2 = -3$ (so $-\alpha^2/2 = 3/2$, realised at $(n, \ell, m) = (1, 1, 1)$ with $D = 4nm - \ell^2 = 3$): $\mathrm{mult}(\alpha) = c(3) = -64$, interpreted as the signed difference $\dim \Delta^{\mathrm{im}}_0(\alpha) - \dim \Delta^{\mathrm{im}}_1(\alpha)$ of even- and odd-imaginary-root multiplicities per Construction~\ref{constr:k3e-roots}. At $\alpha^2 = -4$ (so $D = 4$): $\mathrm{mult}(\alpha) = c(4) = 108$.
 \item \emph{Weight formula.} The weight of the chiral-avatar automorphic form equals $\mathrm{wt}(\Delta_5) = c(0)/2 = 5 = \kBKM(\Delta_5)$, specialising the Borcherds weight theorem \textup{(}Theorem~\ref{thm:borcherds-lift-universal}\textup{)} and the chain-level master identity \textup{(}Theorem~\ref{thm:bl-master-kappa-bkm-chain-level}\textup{)} to $(V, L) = (V^{K3\mathrm{-VOA}}, \Lambda^{3,2})$. The composite identification $\pi_*^{K3} \Obs_{\hCS}^{\fg_{\Delta_5}}(X) = \Delta_5^{-1} \cdot V^{\mathrm{ch}}(\fg_{\Delta_5})$ extends the real-root identification of Proposition~\ref{prop:hcs-pushforward-k3e-delta5}(iii) to the full root system.
\end{enumerate}
\end{theorem}

\begin{proof}[Sketch]
\emph{(i) Chiral avatar.} The $K3 \times E$ hCS factorisation algebra is the tensor product of a holomorphic factorisation algebra on $K3$ and a holomorphic factorisation algebra on $E$, twisted by the $\fg_{\Delta_5}$ gauge datum. Factorisation homology over the compact complex surface $K3$ is the Getzler--Kapranov modular-operadic supertrace on the $E_3^{\mathrm{hol}}$-observable algebra (Costello--Gwilliam~$2021$ Vol~II Theorem~$4.10.1$; Costello--Gaiotto~$2018$ \S4). For $K3$ rational elliptic, the supertrace is expressible as the regularised integral over the $\mathrm{SL}_2(\mathbb{Z})$-modular base of the Eichler--Zagier seed \(\phi_{0,1}^{\mathrm{EZ}}\) paired with the Siegel--Narain kernel determined by a positive four-plane in the Mukai lattice; the Borcherds product appears after projection to the \(\Lambda^{3,2}\) theta-lift datum. The remaining $\R^2$-factor of the $6\mathrm{d} \to 4\mathrm{d}$ holomorphic-topological reduction supplies the weak modular covariance required for the Harvey--Moore regularisation to converge.
\emph{(ii) Borcherds product.} Borcherds~$1998$ Theorem~$13.3$ expands the regularised theta integral at a $0$-dimensional cusp as the product displayed; the exponents are the Fourier coefficients of $F$ at negative argument $-\lambda^2/2$. For \(L=\Lambda^{3,2}\) and \(F=\phi_{0,1}^{\mathrm{EZ}}\) the product is \(\Delta_5\) up to the \(pqt\) Weyl-vector prefactor (Gritsenko--Nikulin~$1998$ \S$2$, Theorem~$2.1$).
\emph{(iii) Low-discriminant coefficients.} Fourier coefficients of \(\phi_{0,1}^{\mathrm{EZ}}\) depend only on the discriminant \(D=4n-\ell^2\): \(c(-1)=1\) (the polar term driving the \(y^{\pm1}\) head), \(c(0)=10\) (lightlike roots), \(c(3)=-64\) (first strict imaginary), \(c(4)=108\) (second strict imaginary). The generating series opens as \(y^{-1} + 10 + y + q(10 y^{-2} - 64 y^{-1} + 108 - 64 y + 10 y^2) + O(q^2)\) (Eichler--Zagier~$1985$ Table~$1$), reproducing \(c(3)=-64\), \(c(4)=108\). The signed reading \(-64=0-64\) at \(\alpha^2=-3\) is Gritsenko--Nikulin~$1998$ Example~$2$: the root space is a purely odd \(64\)-dimensional superspace.
\emph{(iv) Weight.} $\mathrm{wt}(\Phi_F) = c(0)/2 = 5$ is Borcherds~$1998$ Theorem~$13.3$(iv); the specialisation $V = V^{K3\mathrm{-VOA}}$, $L = \Lambda^{3,2}$ matches Theorem~\ref{thm:borcherds-lift-universal}(iv) and Theorem~\ref{thm:bl-master-kappa-bkm-chain-level} at $N=1$. The composite factorisation of the pushforward is Costello--Gwilliam holomorphic-topological tensor applied after the Weyl-group averaging of (ii).
\end{proof}

\begin{remark}[What the chiral avatar adds to Borcherds $1998$]
\label{rem:chiral-avatar-adds-to-borcherds}
Borcherds~$1998$ Theorem~$13.3$ and Gritsenko--Nikulin~$1998$ establish the automorphic-form identity at the level of formal power series on the Siegel upper half-space. Theorem~\ref{thm:hcs-pushforward-imaginary-roots-singular-theta} states that this identity is the Fourier expansion, at a $0$-dimensional cusp, of the K3-fibre factorisation-homology pushforward of the chiral $\fg_{\Delta_5}$-gauge observables on $K3 \times E$. The addition is not a new proof of the Borcherds product: it is the identification of the Borcherds singular theta integrand with the Costello--Gaiotto $5\mathrm{d}$-hCS factorisation-algebra pairing (a working note \S $E_3$ holomorphic algebra of observables), normalised by the modular-operadic supertrace of Theorem~\ref{thm:bl-master-kappa-bkm-chain-level}. The imaginary-root multiplicities $\mathrm{mult}(\alpha) = c(-\alpha^2/2)$ are the Fourier coefficients of that pushforward automorphic form; the compact BV anomaly problem outside the $K3\times E$ singular-theta locus is not used here.
\end{remark}

\begin{remark}[Why the $b^+ \geq 2$ hypothesis is satisfied for $K3 \times E$]
\label{rem:bplus-geq-two-k3e}
The Borcherds singular theta lift \textup{(}Theorem~\ref{thm:borcherds-lift-universal}\textup{)} requires $b^+ \geq 2$. For $L = \Lambda^{3,2}$ the signature is $(3, 2)$, so $b^+ = 3 \geq 2$. The $N = 1$ Monster case $L = \mathrm{II}_{1,1}$ sits outside this hypothesis; the $K3 \times E$ case sits inside it. The chiral avatar of Theorem~\ref{thm:hcs-pushforward-imaginary-roots-singular-theta} therefore transports the full Borcherds theta correspondence, including the imaginary-root multiplicity identity, without the metaplectic-cover complications of the Monster case \textup{(}Remark~\ref{rem:borcherds-three-cases} item~\textup{(1)}\textup{)}.
\end{remark}

\begin{conjecture}[Green--Schwarz closure for 6D \texorpdfstring{$\hCS$}{hCS} on compact CY\texorpdfstring{$_3$}{-3}]
\label{conj:quintic-hcs-kanom}
\ClaimStatusConjectured
Let $X$ be a compact Calabi--Yau threefold and $\frakg$ simple. The
one-loop gauge anomaly of 6D holomorphic Chern--Simons is controlled
locally by the degree-$4$ invariant polynomial
$I_4\in\mathrm{Sym}^4(\frakg^\vee)^{\frakg}$, not by the cubic Casimir.
A Green--Schwarz closure, when available, must trivialise this quartic
class together with the compact CY$_3$ Atiyah-class data. On loci with
zero BCOV Euler-characteristic coefficient, such as $K3\times E$ and
$T^6$, the Euler-characteristic contribution vanishes, but this is not a
universal proof that every gauge algebra is unobstructed. On the
quintic, $\chi_{\mathrm{top}}(X_5)=-200$ makes the quartic obstruction
the first compact instance. The Gritsenko $\Delta_5$ sector lives on
the $K3\times E$ formal locus where the Stage-$1$ obstruction classes
used in this correspondence vanish or have zero coefficient.
\end{conjecture}

\begin{remark}[Dimensional origin of the \texorpdfstring{$\En$}{E-n} level]
\label{rem:en-from-dimension}
The $\En$ level is set by the complex dimension of the ambient space, not by the Deligne conjecture on Hochschild cochains. Specifically:
\begin{itemize}
 \item The $\Etwo$ structure on $\C^2$ comes from the $E_2$-operad acting on $\Conf_n(\C^2)$. Topologically, $\Conf_2(\C^2) \simeq \Conf_2(\R^4) \simeq S^3$ has $\pi_1 = 0$, so the topological braiding is \emph{trivial}. The nontrivial braiding of the Yangian $R$-matrix arises not from $\pi_1$ but from the \emph{Omega-background deformation} $(h_1, h_2)$ of the holomorphic factorization algebra: the equivariant parameters break the topological $E_2$-symmetry to a nontrivially braided $E_2$-chiral structure. Without the Omega-background, the braiding is symmetric.
 \item The $E_3$ structure on $\C^3$ comes from the topological $E_3$-operad acting on $\Conf_n(\C^3)$, which has $\pi_1(\Conf_2(\R^6)) \cong \pi_1(S^5) = 0$ (trivial braiding topologically). As in the $\C^2$ case, the nontrivial braiding at the $E_3$ level arises from the Omega-background deformation $(h_1, h_2, h_3)$, not from the topology of the configuration space.
 \item The algebraic $E_3$ from the Deligne conjecture (Remark~\ref{rem:deligne-hochschild}) requires $E_\infty$ input and produces $E_3$ on Hochschild cochains; the topological $E_3$ from $\C^3$ configuration spaces is a priori distinct. Under formality (Kontsevich 1999), the two $E_3$ structures agree rationally.
\end{itemize}
The CY dimension $d$ of the underlying category and the complex dimension $n$ of the ambient holomorphic CS space are related but distinct in general: for $\C^3$, the CY dimension of the quiver category is $d = 3$ and the CS dimension is $n = 3$ (coincidence). For $K3 \times E$, the CY dimension is also $d = 3$ and $\dim_\C(K3 \times E) = 3$, so the coincidence persists. The distinction matters for \emph{non-compact} examples: for the resolved conifold $\cO(-1)^{\oplus 2} \to \bP^1$, the CY dimension is $d = 3$ (it is a non-compact CY threefold) but the holomorphic CS theory is naturally formulated on the total space $\C^4$ with a superpotential constraint, not on a $3$-dimensional ambient space.
\end{remark}

\subsection{The chiral Chevalley--Eilenberg complex}
\label{subsec:chiral-ce}

The bar complex of a chiral algebra is the chiral analogue of the classical Chevalley--Eilenberg complex of a Lie algebra. Making this identification explicit reveals the connection between Costello's holomorphic CS observables and the bar-cobar formalism of Vols~I--II.

\begin{construction}[Chiral CE chains and cochains]
\label{constr:chiral-ce}
Let $A$ be a chiral algebra with augmentation $\varepsilon \colon A \to \Omega_C$ and augmentation ideal $\bar{A} = \ker(\varepsilon)$. The three chiral CE-type complexes:
\begin{enumerate}[label=\textup{(\roman*)}]
 \item The \emph{chiral CE chains} (labeled-ordered) are the ordered bar complex $B^{\mathrm{ord}}(A) = T^c(s^{-1}\bar{A})$, the cofree conilpotent coalgebra on the desuspension of $\bar{A}$, equipped with the deconcatenation coproduct and bar differential built from the chiral product. This is the direct analogue of the Chevalley--Eilenberg chain complex $C_\bullet(\frakg) = \bigwedge^\bullet \frakg$ with its Lie-bracket-induced differential. The labeled-ordered bar retains the full $R$-matrix data.

 \item The \emph{chiral CE chains} (symmetric), under the same locality
 and $R$-twisted descent package, are the symmetric bar complex
 $B^{\Sigma}(A) = \Sym^c(s^{-1}\bar{A})$ with the coshuffle coproduct
 and the symmetrized differential. This is the direct analogue of
 $C_\bullet(\frakg) = \bigwedge^\bullet \frakg$ in its standard
 commutative-coalgebra form. Vol~I's symmetric bar lives here; its
 ordered bar is the labelled complex in \textup{(i)}. Without descent
 this symmetric object is only a graded coinvariant shadow.

 \item The \emph{chiral CE cochains} are the chiral Hochschild cochain
 complex $C^\bullet_{\mathrm{ch}}(A, A)=\RHom(\Omega B(A), A)$, the
 chiral derived center $Z^{\mathrm{der}}_{\mathrm{ch}}(A)$ of Vol~I
 Theorem~H. This is the analogue of
 $C^\bullet(\frakg, \frakg)=\Hom(\bigwedge^\bullet \frakg,\frakg)$,
 the Chevalley--Eilenberg cochains with adjoint coefficients.
\end{enumerate}
The Koszul duality of Vol~I sends the CE chains~(i) to the Koszul
dual $A^!=D_{\Ran}(B(A))$, the Verdier dual of the bar complex. The
Koszul dual $A^!$ is a fourth object, distinct from the three CE
complexes listed above and from the CE cochains~(iii) in particular:
$A^!$ controls the defect, while the CE cochains
$Z^{\mathrm{der}}_{\mathrm{ch}}(A)$ are the algebraic closed-sector
object and become physical bulk observables only after the relevant
open--closed comparison has been supplied
(see Proposition~\ref{prop:three-dualities}). In classical terms,
$A^!$ is the enveloping algebra of the Koszul-dual Lie algebra
$\frakg^\vee$, not the CE cochains
$C^\bullet(\frakg,\frakg)$ with adjoint coefficients.
\end{construction}

\begin{proposition}[Holomorphic CS observables as chiral CE cochains]
\label{prop:hcs-as-ce-cochains}
For $A = V_k(\frakg)$ the Kac--Moody vertex algebra at level $k$ (the
boundary algebra of 3d holomorphic CS), the chiral CE cochains
$C^\bullet_{\mathrm{ch}}(A, A)$ compute the algebraic closed-sector
comparison target for the CS theory; identifying it with physical bulk
observables is the holomorphic-CS open--closed comparison input. At
the critical level $k = -h^\vee$, the zeroth cohomology of the CE
cochains is the Feigin--Frenkel center
$\mathrm{Fun}(\mathrm{Op}_{G^L}(D))$
(Theorem~\ref{thm:feigin-frenkel-center}).
\end{proposition}

\begin{proof}
The identification of the derived center with bulk observables is Vol~I Theorem~H. The Feigin--Frenkel identification is Theorem~\ref{thm:feigin-frenkel-center} (Chapter~\ref{ch:geometric-langlands}).
\end{proof}

\begin{proposition}[Bar complex as CE chains of the Lie conformal algebra]
\label{prop:bar-ce-identification}
Let $L$ be a Lie conformal algebra with $\lambda$-bracket $\{a_\lambda b\}$, and let $A = U^{\mathrm{ch}}(L)$ be its chiral envelope. Then the ordered bar complex of $A$ is the Chevalley--Eilenberg chain complex of $L$:
\[
 B^{\mathrm{ord}}(U^{\mathrm{ch}}(L)) \;\cong\; \mathrm{CE}_\bullet(L).
\]
The graded vector spaces on each side are $T^c(s^{-1}\bar{A}) \cong \bigwedge^\bullet(L)$ with Poincar\'e polynomial $(1+t)^{\dim L}$; the bar differential $d_B$ encodes the $0$\nobreakdash-th product $a_{(0)}b$ (the Lie bracket) as the CE differential:
\[
 d_{\mathrm{CE}}(g_{i_1} \wedge \cdots \wedge g_{i_k}) = \sum_{a < b} (-1)^{a+b} [g_{i_a}, g_{i_b}] \wedge g_{i_1} \wedge \cdots \wedge \widehat{g_{i_a}} \wedge \cdots \wedge \widehat{g_{i_b}} \wedge \cdots \wedge g_{i_k}.
\]
Dually, the chiral CE cochains $\mathrm{CE}^\bullet(L, L) = \Hom(\bigwedge^\bullet L, L)$ are the Hochschild cochains $\HH^\bullet(A, A)$, i.e.\ the derived center $Z^{\mathrm{der}}_{\mathrm{ch}}(A)$. For abelian $L$ (e.g.\ the Heisenberg current algebra), $d_{\mathrm{CE}} = 0$ and $\HH^k(A, A) = \binom{\dim L}{k}$.
\end{proposition}

\begin{proof}
The identification follows from the Poincar\'e--Birkhoff--Witt
filtration on $U^{\mathrm{ch}}(L)$ (Kac 1998, \S\,2.6;
Loday--Vallette 2012, Theorem~10.1.6).  For an ordinary Lie algebra,
the reduced bar resolution of $U(\frak g)$ is compared with the
Chevalley--Eilenberg resolution
$U(\frak g)\otimes\Lambda^\bullet\frak g$.  Applying
$\operatorname{Hom}_{U(\frak g)^e}(-,U(\frak g))$ gives the standard
identity
\[
 HH^\bullet(U(\frak g),U(\frak g))
 \cong H^\bullet(\frak g,U(\frak g)^{\mathrm{ad}}).
\]
The chiral PBW comparison replaces $\frak g$ by the Lie conformal
algebra~$L$ and uses the factorization envelope
$U^{\mathrm{ch}}(L)$.  Its degree-two bar differential is the OPE
residue $a_{(0)}b$, and the degree-three identity is the Jacobi
relation.  A family datum $H_H(U^{\mathrm{ch}}(L);S_L)$ then computes
the support of the chiral derived centre.  A comparison with
$\mathrm{CE}^\bullet(L,U^{\mathrm{ch}}(L)^{\mathrm{ad}})$ is a
separate PBW quasi-isomorphism.

Three elementary models fix the normalisation:
\begin{itemize}
 \item Heisenberg ($L = \langle J \rangle$, abelian): $\mathrm{CE}_\bullet = \bigwedge^\bullet(\langle J \rangle)$, $d_{\mathrm{CE}} = 0$, Poincar\'e~$= (1+t)$, $\HH^\bullet = \{1, 1\}$.
 \item Kac--Moody $\widehat{\fsl}_2$ ($L = \langle e, f, h \rangle$): $\mathrm{CE}_\bullet = \bigwedge^\bullet(\langle e, f, h \rangle)$, $d_{\mathrm{CE}}^2 = 0$ by Jacobi, Poincar\'e~$= (1+t)^3$.
 \item Yangian $Y(\widehat{\fgl}_1)$ (truncated, $3$ generators): strict Lie (class~L), Poincar\'e~$= (1+t)^3$, genus-$3$ bar dimension~$= 2^9 = 512$.
\end{itemize}
\end{proof}

\begin{remark}[\texorpdfstring{$L_\infty$}{L-inf} deformation and shadow tower]
\label{rem:linfinity-shadow-identification}
For a vertex algebra $A$ whose $A_\infty$ structure has higher operations $m_k \neq 0$ ($k \geq 3$), the bar complex $B(A)$ is the CE complex of an $L_\infty$ algebra $L_A$. The $L_\infty$ structure maps are:
\begin{align*}
 l_2 &= \text{Lie bracket (from } m_2\text{, the OPE residue)}, \\
 l_3 &= \text{Jacobiator (from } m_3\text{, the associator)}, \\
 l_k &= \text{higher Jacobiator (from } m_k\text{)}.
\end{align*}
The CE differential of the $L_\infty$ algebra $L_A$ is $d_{\mathrm{CE}}^{L_\infty} = \sum_{k \geq 2} d^{(k)}$, with $d^{(k)}$ encoding $l_k$. The shadow tower of Vol~I is precisely this $L_\infty$ CE tower:
\[
 S_r = l_r \text{ contribution to the CE differential, i.e.\ } \delta^{(r)} \text{ of } \textup{rem:shadow-ainfty-coproduct-vol3}.
\]
For class~G (e.g.\ Heisenberg): $l_k = 0$ for all $k \geq 2$ (abelian), tower terminates at $S_2 = \kappa_{\mathrm{ch}}$. For class~L (e.g.\ Yangian): $l_k = 0$ for $k \geq 3$ (strict Lie), and the CE complex is the standard $\mathrm{CE}_\bullet(L)$ with Poincar\'e $(1+t)^n$; at genus~$3$ this reproduces the $(1+t)^{3g}$ formula of. For class~M (e.g.\ Virasoro): $l_k \neq 0$ for all $k$, and the shadow tower is infinite; computational verification at $c = 1$:
\[
 l_3(T, T, T) = -2, \quad l_4(T, T, T, T) = \tfrac{40}{27}, \quad S_2 = \kappa_{\mathrm{ch}} = \tfrac{1}{2}, \quad S_3 = 2, \quad S_4 = \tfrac{10}{27}.
\]
These values match the $A_\infty$ bar complex and the shadow tower
recursion.
\end{remark}

\subsection{The universal defect and holomorphic Wilson lines}
\label{subsec:universal-defect}

In the holomorphic CS framework, a \emph{defect} is a codimension-$2$ locus along which the gauge field acquires a prescribed singularity. The simplest defect is the holomorphic Wilson line: a line operator valued in a representation $V$ of $\frakg$.

\begin{construction}[Universal defect from Koszul duality]
\label{constr:universal-defect}
Let $A$ be the boundary chiral algebra of holomorphic CS with gauge algebra $\frakg$ on $M$. The \emph{universal defect algebra} is the Koszul dual $A^! = D_{\Ran}(B(A))$. In the 3d case ($M = \Sigma \times \R$), $A^!$ is the chiral CE/bar-cochain object whose reflected current presentation has Feigin--Frenkel level $k' = -k - 2h^\vee$. The bulk-boundary coupling is the canonical map
\[
 \mu_{\mathrm{bb}} \colon Z^{\mathrm{der}}_{\mathrm{ch}}(A) \otimes A^! \;\longrightarrow\; A^!
\]
making $A^!$ a module over the bulk algebra $Z^{\mathrm{der}}_{\mathrm{ch}}(A)$. The holomorphic Wilson line in representation $V$ corresponds to a specific $A^!$-module: the image of $V$ under the equivalence $\Rep^{\Eone}(A^!) \simeq \Rep^{\Eone}(A)^{\mathrm{rev}}$ (Vol~I, Theorem~A, Verdier leg: $D_{\Ran}$ exchanges $A$ and $A^!$ module categories with reversed monoidal structure).

In the 5d case ($M = \C^2 \times \R$), the universal defect algebra is the Koszul dual of the affine Yangian. For $\frakg = \fgl_1$, the defect algebra $A^!$ encodes the Wilson lines of the 5d theory; the RTT formalism applied to the defect $R$-matrix generates the dual Yangian presentation. (The quantum toroidal algebra arises from the \emph{6d} theory on $\C^3$, not from 5d; see Conjecture~\ref{conj:6d-boundary-toroidal}.) The Drinfeld center of the defect module category recovers the $\Etwo$ braided structure:
\[
 \cZ(\Rep^{\Eone}(A^!)) \;\simeq\; \Rep^{\Etwo}(Z^{\mathrm{der}}_{\mathrm{ch}}(A)).
\]
\end{construction}

\begin{conjecture}[6d universal defect and chiral quantum groups on \texorpdfstring{$\C^3$}{C-3}]
\label{conj:6d-defect-c3}
\ClaimStatusConjectured
For the 6d holomorphic theory on $\C^3$ with gauge algebra $\fgl_1$, the universal defect algebra $A^!_{\C^3}$ (the Koszul dual of the quantum toroidal boundary algebra $U_{q,t}(\widehat{\widehat{\fgl}}_1)$) satisfies:
\begin{enumerate}[label=\textup{(\roman*)}]
 \item $A^!_{\C^3}$ carries $E_3$-chiral factorization structure. The underlying topological $E_3$-operad action on $\Conf_n(\C^3)$ has trivial braiding ($\pi_1(\Conf_2(\R^6)) = 0$); the nontrivial $E_3$-chiral braiding arises from the Omega-background deformation $(h_1, h_2, h_3)$ of the holomorphic factorization structure (Remark~\ref{rem:e3-not-symmetric}).
 \item The $\Eone$-projection of $A^!_{\C^3}$ to a curve $C \subset \C^3$ is the Koszul dual of the quantum toroidal algebra, with inverted parameters $q^{-1}, t^{-1}$ (parameter inversion under Koszul duality; see Table of Chapter~\ref{ch:quantum-groups}, Definition~\ref{def:qgf-four-regimes}).
 \item The $\Etwo$-projection to $\C^2 \subset \C^3$ is the $\Etwo$-chiral algebra whose Drinfeld center $\cZ(\Rep^{\Eone}(A^!_{\C^3}|_C))$ is the braided monoidal category of Fock representations of the quantum toroidal algebra.
 \item The \emph{chiral quantum group} on $\C^3$ is the Hopf-algebra-in-braided-category structure on the representation category $\Rep^{\Etwo}(A^!_{\C^3}|_{\C^2})$, equipped with the $R$-matrix from the $E_3$ braiding.
\end{enumerate}
The conjecture composites three constructions (6d $E_3$ factorization, Koszul duality, Drinfeld center), each of which is independently established at lower dimension; the composite arrow at 6d is conjectural.
\end{conjecture}

\begin{conjecture}[6d holomorphic theory on \texorpdfstring{$K3 \times E$}{K3 x E} and the BKM chiral algebra]
\label{conj:6d-k3xe}
\ClaimStatusConjectured
Let $X = K3 \times E$ with $E$ an elliptic curve. On the non-Lagrangian 6d factorization framework, together with the framed K3-fibre specialisation of Theorem~\ref{thm:cy-to-chiral-d3}, the holomorphic theory on $X$ should produce:
\begin{enumerate}[label=\textup{(\roman*)}]
 \item An $\Eone$-chiral algebra $A_{X}^{(K3,E)} = \Phi_3^{(K3,E)}(D^b(\Coh(X)))$ on $E$ on the framed K3-fibre specialisation of Theorem~\ref{thm:cy-to-chiral-d3}, whose character is conjecturally controlled by the Igusa cusp form $\Phi_{10}^{\mathrm{un}}$. The expected chiral algebra is a deformation of the lattice vertex algebra $V_{\Lambda_{K3}}$ of the K3 lattice, with deformation parameters determined by the complex structure of $K3$.
 \item An $\Etwo$-chiral factorization algebra on $K3$ itself (the ``transverse'' directions), whose Drinfeld center gives the braided monoidal category of BPS modules.
 \item The universal defect algebra $A^!_X$ on $E$, whose representation category contains the holomorphic Wilson lines wrapping $E$. These Wilson lines are indexed by the BPS spectrum of $K3 \times E$ (Mukai vectors in $H^*(K3, \Z)$).
 \item The \(\kappa_\bullet\)-spectrum has four coordinates: the Stage-2 Heisenberg reading has $\kappa_{\mathrm{ch}}^{\mathrm{Heis}}(A_X) = 3$, the Borcherds automorphic weight is $\kappa_{\mathrm{BKM}}(\Delta_5) = 5$, the total-space categorical invariant is $\kappa_{\mathrm{cat}}(X) = 0$, and the K3-fibre lattice invariant is $\kappa_{\mathrm{fiber}}(K3) = 24$, consistent with Chapter~\ref{ch:k3-times-e}.
\end{enumerate}
The framed $\Eone$ object is supplied by Theorem~\ref{thm:cy-to-chiral-d3} after the K3-fibre datum is fixed. The conjectural content is the non-Lagrangian 6d algebraic framework for $K3\times E$ (Remark~\ref{rem:6d-not-lagrangian}) and the identification of the BPS root datum with the Borcherds--Kac--Moody root system of $\mathfrak{g}_{\Delta_5}$ (Chapter~\ref{ch:k3-times-e}).
\end{conjecture}

\begin{remark}[What the 6d theory adds to the \texorpdfstring{$K3 \times E$}{K3 x E} chiral-algebra construction]
\label{rem:6d-adds-to-k3xe}
Chapters~\ref{ch:k3-times-e} and~\ref{chap:toroidal-elliptic} develop the $K3 \times E$ chiral algebra construction: the first via the framed K3-fibre specialisation of the CY-to-chiral correspondence $\Phi_3^{(K3,E)}$ (Theorem~\ref{thm:cy-to-chiral-d3}; the further quantum-group identification CY-C remains conjectural), the second via the toroidal/elliptic quantum group presentation (conditional on Conjecture~\ref{conj:toroidal-e1}). Conjecture~\ref{conj:6d-k3xe} adds a third conjectural construction pathway: the holomorphic CS / factorization homology route. The three pathways should agree if all three constructions are realized:
\begin{center}
\small
\begin{tabular}{lll}
 \toprule
 \textbf{Pathway} & \textbf{Input} & \textbf{Output} \\
 \midrule
 CY-to-chiral $\Phi_3^{(K3,E)}$ & $D^b(\Coh(K3 \times E))$ & $A_{K3 \times E}^{(K3,E)}$ (framed object-level specialisation, Theorem~\ref{thm:cy-to-chiral-d3}) \\
 Toroidal presentation & $U_{q,t}(\widehat{\widehat{\fgl}}_1)$ & $\Eone$-chiral on $E$ \\
 6d holomorphic theory & $K3 \times E$ as ambient & $\Eone$-chiral on $E$ \\
 \bottomrule
\end{tabular}
\end{center}
The agreement of the three is itself conjectural. The 6d pathway provides the geometric origin of the quantum toroidal parameters $(q, t)$: they arise from the K3 complex structure moduli via equivariant localisation on the Omega-background of the normal bundle $N_{C/X}$ (the normal bundle grading does not directly identify with $(q,t)$; the intermediary mechanism is Nekrasov-type equivariant localisation).
\end{remark}

\subsection{Codimension-\texorpdfstring{$2$}{2} defect OPE from the 6d action}
\label{subsec:codim2-defect-ope}

The stage-2 specialisation $\SpCh_{N_{C/Y}, C}(\PhiFA_3(\cC))$ for
$Y = \C^3$ and $C = \C_{z_1}$ should be the defect algebra on $C$
obtained by pushing the $E_3$-holomorphic factorization algebra
$\PhiFA_3(\cC)$ forward along the normal bundle
$N_{C/Y} = \C^2_{z_2, z_3}$, with the Omega-background equivariant
weights $(h_2, h_3)$ of $N_{C/Y}$ supplying the specialisation datum
(Definition~\ref{def:phi-fa-and-sp}, Remark~\ref{rem:multiple-e1-specialisations}).
The computation below is a low-spin OPE witness for that expected
identification: restricting the abelian local hCS model to a tubular
neighbourhood of \(C\) and extracting boundary-mode OPEs recovers the
spin-\(1\) current and the Sugawara spin-\(2\) Virasoro field. This
low-spin witness does not construct the full
renormalised 6d hCS factorisation algebra, all $W_{1+\infty}$ spins, or
the hCS-to-Hall comparison of Chapter~\ref{ch:cy3-chain-level-bridge}.

\begin{proposition}[Low-spin defect OPE witness for 6d hCS on \texorpdfstring{$\C^3$}{C-3}]
\label{prop:codim2-defect-ope}
Let $Y = \C^3$ with holomorphic $3$-form $\Omega = dz_1 \wedge dz_2 \wedge dz_3$, and let $C = \C_{z_1} \hookrightarrow Y$ be the curve $\{z_2 = z_3 = 0\}$. The normal bundle is $N_{C/Y} = \C^2_{z_2, z_3}$. In the abelian local hCS model with gauge algebra $\fgl_1$ and Omega-background parameters $(h_1, h_2, h_3)$ satisfying $h_1 + h_2 + h_3 = 0$, restriction to a tubular neighborhood $U = \C_{z_1} \times D^2_{z_2, z_3}$ of $C$ gives the following spin~$1$ and spin~$2$ OPEs expected of the $W_{1+\infty}$ defect algebra at level
\[
 \Psi \;=\; -\sigma_2 \;=\; -(h_1 h_2 + h_1 h_3 + h_2 h_3)
\]
with central charge $c = 1$ for the $\fgl_1$ Sugawara field. In particular, the spin-$1$ current $J$ and spin-$2$ current $T$ on $C$ satisfy:
\begin{enumerate}[label=\textup{(\alph*)}]
 \item \emph{Spin~$1$}: $J(z) J(w) = \Psi/(z - w)^2$, equivalently $\{J_\lambda J\} = \Psi \cdot \lambda$.
 \item \emph{Spin~$2$}: $T(z) T(w) = \tfrac{1}{2}/(z - w)^4 + 2T(w)/(z - w)^2 + \partial T(w)/(z - w)$, equivalently $\{T_\lambda T\} = \tfrac{1}{12}\lambda^3 + 2T\lambda + \partial T$.
\end{enumerate}
\end{proposition}

\begin{proof}
\emph{Arrow~1: restriction to the tubular neighborhood.}
The tubular neighborhood $U = \C_{z_1} \times D^2$ of $C$ in $Y = \C^3$ inherits the holomorphic $3$-form $\Omega|_U = dz_1 \wedge dz_2 \wedge dz_3$. The 6d holomorphic Chern--Simons action on $Y$ with gauge algebra $\fgl_1$:
\[
 S \;=\; \frac{1}{2}\int_Y \Omega \wedge \bigl(A \wedge \bar{\partial} A + \tfrac{2}{3}A \wedge A \wedge A\bigr)
\]
restricts to a free theory on $U$ because the cubic term vanishes for abelian $\fgl_1$.

\emph{Arrow~2: mode expansion in the normal directions.}
Expand the gauge field in modes of the normal bundle:
\[
 A(z_1, z_2, z_3, \bar{z}) \;=\; \sum_{m, n \geq 0} A^{(m,n)}(z_1, \bar{z}_1)\, z_2^m z_3^n.
\]
Each mode $A^{(m,n)}$ is a $(0,1)$-form on $C$ valued in $\fgl_1$. The Omega-background assigns equivariant weight $m \cdot h_2 + n \cdot h_3$ to $A^{(m,n)}$ under the torus $T^2 = \C^*_{z_2} \times \C^*_{z_3}$.

\emph{Arrow~3: integration over the normal fiber.}
Inserting the mode expansion into $S|_U$ and integrating the holomorphic $2$-form $dz_2 \wedge dz_3$ over the formal bidisk $D^2$ (extracting the $(0,0)$ Fourier coefficient of the product of normal modes), the action becomes:
\[
 S|_U \;=\; \sum_{m,n \geq 0} \int_C dz_1 \wedge A^{(m,n)} \bar{\partial} A^{(m,n)} \;+\; \text{(equivariant-weight terms from } \boldsymbol{h}\text{)}.
\]
The propagator of the free-field system on $C$ is
\[
 \langle A^{(m,n)}(z) \, A^{(p,q)}(w) \rangle
 \;=\; \frac{\delta_{m+p,0}\,\delta_{n+q,0}}{(z - w)^2},
\]
where the double pole is the standard holomorphic propagator on $\C$ and the delta-conditions enforce charge conservation in the normal directions.

\emph{Arrow~4: identification of the $W_{1+\infty}$ currents.}
The $\GL(2)$-invariant combinations of the normal modes produce the $W_{1+\infty}$ generators. At spin~$s$:
\[
 J_s(z_1) \;=\; \sum_{m+n = s-1} c_{m,n}^{(s)} \, A^{(m,n)}(z_1),
\]
where $c_{m,n}^{(s)}$ are the $\GL(2)$-invariant coefficients, namely
the minors of the identity matrix matching the invariant polyvector
fields \(\Omega_s\). In particular:
\begin{itemize}
 \item Spin~$1$: $J(z_1) = A^{(0,0)}(z_1)$, the mode at the origin of the normal fiber.
 \item Spin~$2$: $T(z_1) = \frac{1}{2\Psi}\mathopen{:}J(z_1)^2\mathclose{:}$ via the Sugawara construction at level $\Psi$.
\end{itemize}

\emph{Arrow~5: OPE computation.}
\emph{Spin~$1$}: From the propagator in Arrow~3, the contraction of two spin-$1$ modes gives
\[
 J(z) J(w) \;=\; \langle A^{(0,0)}(z)\, A^{(0,0)}(w) \rangle + \mathopen{:}J(z) J(w)\mathclose{:}
 \;=\; \frac{\Psi}{(z - w)^2} + \mathopen{:}J(z) J(w)\mathclose{:},
\]
where the coefficient $\Psi = -\sigma_2 = -(h_1 h_2 + h_1 h_3 + h_2 h_3)$ arises from the equivariant normalization of the propagator: the Omega-background twists the inner product on $\fgl_1$ by $-\sigma_2$ (which is the Kac--Moody level of the 5d theory on $\C^2 \times \R$, as in Costello~2013). In mode language: $J_{(1)} J = \Psi$, all other modes vanish. The classical $r$-matrix is $r^{\mathrm{Heis}}(z) = \Psi/z$ (C10; the level prefix $\Psi$ is retained to distinguish nonzero-level representations). At $\Psi = 0$: the $r$-matrix vanishes.

\emph{Spin~$2$}: The Sugawara construction $T = \frac{1}{2\Psi}\mathopen{:}J^2\mathclose{:}$ at level $\Psi$ for the abelian algebra $\fgl_1$ (with $h^\vee = 0$) produces a Virasoro with central charge
\[
 c \;=\; \frac{\Psi \cdot \dim(\fgl_1)}{\Psi + h^\vee} \;=\; \frac{\Psi \cdot 1}{\Psi + 0} \;=\; 1,
\]
independent of $\Psi$ (a feature of $\fgl_1$; for non-abelian $\frakg$, $c$ depends on the level). The Virasoro OPE is
\[
 T(z)\,T(w) \;=\; \frac{c/2}{(z - w)^4} + \frac{2T(w)}{(z - w)^2} + \frac{\partial T(w)}{z - w}
 \;=\; \frac{1/2}{(z - w)^4} + \frac{2T(w)}{(z - w)^2} + \frac{\partial T(w)}{z - w}.
\]
In mode language: $T_{(3)} T = c/2 = 1/2$, $T_{(1)} T = 2T$, $T_{(0)} T = \partial T$, $T_{(2)} T = 0$. In lambda-bracket form (with divided powers $\lambda^{(n)} = \lambda^n/n!$): $\{T_\lambda T\} = \tfrac{c}{12}\lambda^3 + 2T\lambda + \partial T = \tfrac{1}{12}\lambda^3 + 2T\lambda + \partial T$. The classical $r$-matrix is $r^{\mathrm{Vir}}(z) = \tfrac{1}{2}/z^3 + 2T/z$ (cubic $+$ simple pole; one less than the quartic OPE pole via $d\log$ absorption).

The $T$--$J$ OPE confirms that $J$ is a primary of conformal weight~$1$: $T(z)\,J(w) = J(w)/(z - w)^2 + \partial J(w)/(z - w)$, i.e.\ $T_{(1)}J = J$, $T_{(0)}J = \partial J$.

These are the spin~$1$ and spin~$2$ OPE relations of the expected
$W_{1+\infty}$ defect algebra at level $\Psi$ and central charge
$c = 1$. The all-spin vertex algebra and the renormalised hCS
factorisation-algebra construction remain additional obligations.
\end{proof}

\begin{lemma}[5d-to-6d-on-codim-\texorpdfstring{$2$}{2} effective level bridge]
\label{lem:5d-to-6d-codim2-bridge}
Let $Y = \C^3$ with Omega-background $(h_1, h_2, h_3)$ subject to the CY constraint $h_1 + h_2 + h_3 = 0$, and let $C = \C_{z_1} \hookrightarrow Y$ be a codim-$2$ holomorphic curve with normal bundle $N_{C/Y} \cong \cO_C \oplus \cO_C$ carrying equivariant weights $(h_2, h_3)$ on the two normal coordinates. The 6d holomorphic Chern--Simons action with gauge algebra $\fgl_1$:
\[
 S_{6d} \;=\; \frac{1}{2}\int_Y \Omega \wedge A \wedge \bar{\partial} A
\]
restricted to the tubular neighborhood $U = C \times D^2 \subset Y$ and integrated over the formal bidisk $D^2 \subset N_{C/Y}$ in the equivariant $\bar{\partial}$-regularisation reduces to a 5d holomorphic Chern--Simons action on $C \times \R$ for the $\widehat{\fgl_1}$ Kac--Moody algebra at effective level
\[
 k_{\mathrm{eff}} \;=\; -\sigma_2(h_1, h_2, h_3) \;=\; -(h_1 h_2 + h_1 h_3 + h_2 h_3) \;=\; \tfrac{1}{2}\sum_{i=1}^{3} h_i^2,
\]
where the second equality uses Newton's identity $p_2 = \sigma_1^2 - 2\sigma_2$ under the CY constraint $\sigma_1 = 0$.
\end{lemma}

\begin{proof}
This is the equivariant 1-loop determinant of the $\bar{\partial}$-operator on $D^2$ under the $T^2 = \C^*_{z_2} \times \C^*_{z_3}$ action with weights $(h_2, h_3)$, corrected by the $h_1$ equivariant weight on $C$. The computation is due to Costello (\emph{Supersymmetric gauge theory and the Yangian}, 2013 preprint), with the Costello--Gwilliam factorisation-algebra framework (\emph{Factorization algebras in quantum field theory}, Vol.~II §5) supplying the regularisation axioms. The role of this lemma in Proposition~\ref{prop:codim2-defect-ope} is to justify Arrow~$3$ of the proof: the propagator $\langle A^{(0,0)}(z) A^{(0,0)}(w) \rangle = \Psi/(z-w)^2$ with $\Psi = -\sigma_2$ is the $\widehat{\fgl_1}$ 5d Kac--Moody propagator at level $k_{\mathrm{eff}}$, and the identification $\Psi = k_{\mathrm{eff}}$ is exactly the content of the bridge. The Newton-identity form $\Psi = \tfrac{1}{2}\sum h_i^2$ manifests $\Psi \geq 0$ for real $h_i$ under the CY constraint.
\end{proof}

\begin{remark}[The role of the normal bundle]
\label{rem:normal-bundle-spectral}
The normal bundle $N_{C/Y} = \C^2_{z_2, z_3}$ provides the two spectral parameters $(z_2, z_3)$ of the $E_3$ theory. Under the projection hierarchy $\C^3 \to \C^2 \to C$:
\begin{itemize}
 \item $E_3$ on $\C^3$: two spectral parameters $(z_2, z_3)$ from the full normal bundle;
 \item $\Etwo$ on $\C^2$: one spectral parameter $z_2$ (the Yangian parameter $u$);
 \item $\Eone$ on $C$: no spectral parameter (the chiral algebra on the defect curve).
\end{itemize}
The level $\Psi = -\sigma_2$ is a scalar extracted from the full Omega-background: it is the residue of the classical $r$-matrix $r(u) = -\sigma_2/u$ of the affine Yangian (holomorphic\_cs\_chiral\_engine.py, line~127). The passage from the two-parameter $R$-matrix $R_{\mathrm{ch}}(u, v)$ of Conjecture~\ref{conj:chiral-qg-c3}(ii) to the single-parameter Yangian $R$-matrix is the projection $\C^3 \to \C^2$; the further projection $\C^2 \to C$ forgets the spectral parameter entirely, leaving only the defect algebra with its level $\Psi$.
\end{remark}



\subsection{Consistency comparison: 3d Chern--Simons vs 6d chiral framework}
\label{subsec:cfg25-adversarial}

If the 6d holomorphic theory is a genuine lift of 3d Chern--Simons, every 3d topological result should have a 6d chiral analogue. We compare this expectation against six classical results and catalogue the gaps.

\begin{conjecture}[Chiral analogues of 3d Chern--Simons invariants]
\label{conj:cfg25-adversarial}
\ClaimStatusConjectured
For the 6d holomorphic theory on $\C^3$ with gauge algebra $\fgl_1$ and Omega-background $(h_1, h_2, h_3)$ satisfying $h_1 + h_2 + h_3 = 0$, the following 6d analogues of classical 3d Chern--Simons results should hold:
\begin{enumerate}[label=\textup{(\roman*)}]
 \item \emph{Chiral Kontsevich integral.}
 The perturbative expansion of the 3d CS path integral (the Kontsevich integral $Z^{\mathrm{pert}}(K)$, CFG25 Theorem~A) should lift to a codimension-$2$ defect invariant on $\C^3$: an integral over the configuration space $\Conf_n(C)$ of the holomorphic curve $C \hookrightarrow \C^3$, with the chiral propagator $r_{\mathrm{CY}}(z) = \Psi/z$ replacing the topological propagator. The 3d propagator is smooth; the 6d propagator has poles at $z = -h_i$ from the structure function $g(u) = \prod(u - h_i)/\prod(u + h_i)$, requiring the Omega-background as equivariant regularization.
 \emph{Obstruction}: the codimension-$2$ defect OPE exists (Proposition~\ref{prop:codim2-defect-ope}), but the full chiral Kontsevich integral is not constructed.

 \item \emph{Chiral Jones polynomial.}
 The colored Jones polynomial $J_N(K; q) = \Tr_{V_N}(\cR)$ from quantum traces over the $U_q(\fsl_2)$ $R$-matrix should lift to a ``chiral Jones polynomial'' $J^{\mathrm{ch}}_N(C; q, t) = \Tr_{\mathrm{Fock}_N}(R_{\mathrm{ch}}(u, v))$ using the two-parameter $R$-matrix of the quantum toroidal algebra. The expected output is a triply-graded invariant related to Khovanov--Rozansky homology.
 \emph{Obstruction}: the $R$-matrix exists and satisfies the YBE (Theorem~\ref{thm:yang-r-matrix-ybe}), but the quantum trace is not extracted. The identification with Khovanov--Rozansky is conjectural.

 \item \emph{Chiral volume conjecture.}
 The volume conjecture $\log|J_N(K; e^{2\pi i/N})| \sim N \cdot \mathrm{Vol}(S^3 \setminus K)$ should lift to an asymptotic statement relating the chiral Jones polynomial to the holomorphic Chern--Simons functional $S_{\mathrm{hCS}}$ evaluated on the flat connection associated to $C$.
 \emph{Obstruction}: no chiral volume conjecture is formulated.

 \item \emph{Chiral Verlinde algebra.}
 Quantization of 3d CS on $\Sigma_g \times \R$ produces the Verlinde algebra $V_k(\frakg)$ with $\dim V_k(\fsl_2) = k + 1$. For the 6d theory on $\Sigma_g \times \C^2 \times \R$, the factorization homology $\int_{\Sigma_g} \cA$ gives the genus-$g$ conformal block space. At generic parameters (classes $\mathbf{L}$ and $\mathbf{C}$), $\dim H^*(B_{E_3}(\cA)) = (1+t)^{3g}$, evaluating to $2^{3g}$ at $t = 1$. The factor $3g$ (vs $2g$ at $\Etwo$ or $g$ at $\Eone$) reflects the three factorization directions in $\C^3$. The 3d Verlinde number $V_k = k + 1$ is \emph{not} recovered at generic parameters; it requires root-of-unity truncation of the quantum toroidal algebra.
 \emph{Obstruction}: generic formula exists; root-of-unity truncation is not constructed.

 \item \emph{6-manifold invariant (WRT analogue).}
 The 3d WRT invariant $Z(M^3)$ arises from a $3$-$2$-$1$ extended TFT. The 6d theory should produce a $6$-$5$-$\cdots$-$1$ extended TFT assigning data to manifolds of all dimensions up to~$6$. The current framework assigns data to circles ($\cZ(\Rep^{\Eone}(A))$), tori ($\HH(\HH(A))$), and surfaces ($\int_{\Sigma_g} \cA$), but \emph{not} to $3$-, $4$-, or $5$-manifolds. The partial assignment to $6$-manifolds via $\mathrm{Sp}_4(\Z)$ modularity (Chapter~\ref{ch:k3-times-e}) covers only CY$_3$ targets.
 \emph{Obstruction}: assignments exist for dimensions $1$--$2$; dimensions $3$--$5$ are missing.

 \item \emph{Reshetikhin--Turaev from the quantum toroidal algebra.}
 The 3d RT functor uses $U_q(\frakg)$ at root of unity with the YBE for consistency. At 6d, the Zamolodchikov tetrahedron equation (ZTE) replaces the YBE as the consistency condition. The Yang $R$-matrix satisfies YBE but \emph{fails} ZTE at order $O(\sigma_3^2)$ where $\sigma_3 = h_1 h_2 h_3$ (Theorem~\ref{thm:zte-failure}). The ZTE failure at $\sigma_3 \neq 0$ proves that the $E_3$ structure is genuinely nontrivial and that pairwise $R$-matrices alone do not produce consistent $6$-manifold invariants. Proposition~\ref{prop:zte-explicit-correction} kills the even charge-$2$ obstruction by a rational ternary correction; Proposition~\ref{prop:zte-o-kappa4} leaves rank-$2$ odd-sector residues.
 \emph{Obstruction}: full charge-preserving all-order ZTE completion remains open.
\end{enumerate}
\end{conjecture}

\begin{remark}[The ZTE obstruction as the 6d fingerprint]
\label{rem:zte-6d-fingerprint}
The Zamolodchikov tetrahedron equation failure (item~(vi)) is the deepest structural gap between 3d and 6d. In 3d, the Yang--Baxter equation (the $2$-simplex consistency) suffices for all knot and $3$-manifold invariants. In 6d, the ZTE (the $3$-simplex consistency) is the analogue, and it \emph{fails} for the Yang $R$-matrix. The failure is a quantum effect: it vanishes at $\sigma_3 = 0$ (Kapranov--Voevodsky triviality at the permutation limit) and scales as $O(\sigma_3^2)$. This scaling proves that the $E_3$ corrections live at the \emph{chain level}. The 6d theory is not a ``dimensional upgrade'' of 3d CS; it is a genuinely different structure requiring new algebraic inputs beyond quantum groups.
\end{remark}


\subsection{Systematic obstruction catalogue: what cannot be lifted from 3d to 6d}
\label{subsec:3d-6d-obstruction-catalogue}

The six consistency vectors above decide whether specific 3d
\emph{invariants} extend to 6d.  A deeper question is whether the
underlying \emph{structural features} of 3d CS can be lifted at all.
There are eight structural obstructions: each is either a fundamental
no-go or an unsolved problem, and each has a definite parameter locus
where a 3d feature can be partially recovered.

\medskip

\begin{center}
\small
\renewcommand{\arraystretch}{1.3}
\begin{tabular}{clccl}
\toprule
\# & \textbf{Obstruction} & \textbf{Classification} & \textbf{Depth} & \textbf{Loophole} \\
\midrule
1 & Finite-dimensionality (Verlinde) & Fundamental & Parameter & Root-of-unity locus (codim-$2$) \\
2 & Knot complement topology & Fundamental & Topological & None \\
3 & Surgery formula (Kirby) & Unsolved & Structural & $E_3$ corrections to ZTE \\
4 & Reshetikhin--Turaev functor & Fundamental & Topological & Non-semisimple RT \\
5 & Categorification direction & Unsolved & Structural & M-theory (7d) \\
6 & Unitarity & Fundamental & Parameter & NS limit or imaginary $h_i$ \\
7 & Volume conjecture & Fundamental & Parameter & Rigid CY$_3$ ($h^{2,1} = 0$) \\
8 & Rationality & Fundamental & Parameter & Gepner points \\
\bottomrule
\end{tabular}
\end{center}

\medskip
\noindent The eight obstructions cluster into three types.

\paragraph{Topological obstructions (2, 4).}
These are fundamental: the 6d theory is \emph{holomorphic}, not topological, and the relevant topologies differ irreconcilably. In 3d, the complement $S^3 \setminus K$ of any knot has boundary $T^2$ (the normal bundle is trivial: rank~$2$, oriented, over $S^1$). In 6d, the complement of a holomorphic curve $C$ of genus~$g$ in a CY$_3$ has boundary determined by the normal bundle $N_{C/X}$ with $\deg N_{C/X} = 2g - 2$ (adjunction and $c_1(X) = 0$). Only genus-$1$ curves (elliptic, $\deg N = 0$) have toroidal boundary $T^4$; genus-$0$ curves have $S^3$-type boundaries, and higher-genus curves have still more complex topology. The Reshetikhin--Turaev functor requires a modular tensor category with finitely many simple objects and a non-degenerate $S$-matrix; $\Rep(U_{q,t}(\widehat{\widehat{\fgl}}_1))$ has infinitely many simples (Fock modules $F_n$ for each $n \in \Z$), is non-semisimple for class~$M$, and lacks a known ribbon structure. The non-semisimple RT framework of Costantino--Geer--Patureau-Mirand (2014) provides a conjectural route via modified quantum traces.

\paragraph{Parameter obstructions (1, 6, 7, 8).}
The 3d theory has a \emph{discrete} coupling ($k \in \Z$) while the 6d theory has \emph{continuous} parameters $(h_1, h_2, h_3)$. At special parameter loci, some 3d features are recovered:
\begin{itemize}
 \item \emph{Root of unity}: $q^N = t^M = 1$ (codimension-$2$ in parameter space) should truncate the representation category to finitely many simples. But the truncated quantum toroidal algebra is \emph{not constructed}.
 \item \emph{Unitarity}: algebraic unitarity $R(z)\,R(-z) = \mathrm{Id}$ holds identically (from $g(-z) = 1/g(z)$, a consequence of the CY condition). Hilbert-space unitarity (positive-definite Shapovalov form) \emph{fails} at generic real $h_i$: the Shapovalov norms $\prod_{k=1}^{n} g(k h_1)$ are indefinite. Unitarity is recovered at the topological locus ($h_i$ purely imaginary, giving $|q| = 1$) and at the Nekrasov--Shatashvili limit ($h_2 = 0$, effective 3d theory).
 \item \emph{Volume conjecture}: the 3d conjecture relates quantum invariants to hyperbolic volume via Mostow rigidity (unique hyperbolic structure). CY$_3$ manifolds have positive-dimensional moduli spaces: no Mostow analogue. For \emph{rigid} CY$_3$ ($h^{2,1} = 0$, e.g., the Schoen threefold), the holomorphic volume is unique up to normalization, and a volume conjecture may be formulable.
 \item \emph{Rationality}: $V_k(\frakg)$ at integer level is a rational VOA ($C_2$-cofinite, finitely many simples, modular characters). CY chiral algebras are generically \emph{irrational}: infinitely many charge sectors, non-$C_2$-cofinite, mock-modular characters (class~$M$). Rationality is recovered only at Gepner points (measure zero in CY moduli), where the theory truncates to a tensor product of $N = 2$ minimal models.
\end{itemize}

No single parameter locus recovers \emph{all} 3d features simultaneously: unitarity requires imaginary $h_i$, rationality requires the Gepner locus, and finite-dimensionality requires a root-of-unity specialization. The 3d theory is not ``embedded'' in the 6d theory at any single point.

\paragraph{Structural obstructions (3, 5).}
These are \emph{unsolved problems}: the obstruction may be resolvable by new mathematics.
\begin{itemize}
 \item \emph{Surgery}: 6d handle decomposition exists (Proposition~\ref{prop:handle-decomposition-k3}), but the ZTE failure (Theorem~\ref{thm:zte-failure}) means handle rearrangements are not automatically consistent. ZTE corrections (the $E_3$ correction term $T$) may restore consistency, but sufficiency is not proved.
 \item \emph{Categorification}: Khovanov homology categorifies the Jones polynomial (3d $\to$ 4d). The 6d $\to$ 7d categorification should arise from M-theory on a CY$_2$ (e.g., K3), producing a topological twist of 11d supergravity. This is physically motivated but mathematically unconstructed. The decategorification is consistent: $\dim_{E_3}(g) = \dim_{E_2}(g) \cdot \dim_{E_1}(g)$ by the K\"unneth formula (Dunn additivity).
\end{itemize}

\begin{remark}[Irreducible 6d content]
\label{rem:irreducible-6d-content}
The obstructions are not \emph{defects} of the 6d theory but \emph{signatures} of its richer structure. Features that exist only at 6d, with no 3d analogue:
\begin{enumerate}[label=\textup{(\roman*)}]
 \item Two-parameter $R$-matrix $R(u, v)$ with Zamolodchikov factorization;
 \item Miki automorphism ($\mathbb{Z}/3$ cyclic on $(q, t, s)$ on the positive half / $\mathrm{CoHA}(\C^3) = Y^+$, with ambient $S_3$ on the Drinfeld double; algebra-specific);
 \item ZTE corrections at $O(\sigma_3^2)$ (genuinely new $E_3$ mathematics);
 \item Shadow tower / $A_\infty$ coproduct (infinite corrections $\delta^{(k)} = S_k$);
 \item Mock modularity for class~$M$ characters;
 \item Non-semisimple Drinfeld center (richer than semisimple MTC of 3d).
\end{enumerate}
The 6d theory produces genuinely new mathematical objects (holomorphic curve invariants, CY partition functions, mock modular forms) that do not reduce to classical knot/manifold invariants. The correct relationship between 3d and 6d is not one of ``lift'' but of two outputs of the same $E_3$ Koszul duality formalism, applied to different inputs: topological $E_3$ (framed little $3$-disks on $\R^3$) for 3d, and holomorphic $E_3$ (Omega-background on $\C^3$) for 6d.
\end{remark}

\subsection{The topological--algebraic dichotomy: what the 6d theory produces}
\label{subsec:topological-algebraic-dichotomy}

The obstruction list separates the 3d Chern--Simons structures that
survive 6d holomorphic CS from those that do not.  The answer is a
clean dichotomy: 3d CS produces \emph{topological} invariants, numbers
depending on a manifold or knot, while 6d hCS produces
\emph{algebraic} invariants, rational functions of spectral parameters
encoding the representation theory of quantum algebras. The surviving
structures are precisely the algebraic ones.

\paragraph{The dichotomy.}
We decompose all $34$ substructures across the eight obstruction categories into three types by the mathematical nature of their output:
\begin{itemize}
 \item \emph{Algebraic}: the output is a rational function of spectral parameters, an operator-valued formal series, or a structural datum of an infinite-dimensional algebra. Examples: $R$-matrix satisfying YBE, coproduct $\Delta_z$, structure function $g(z)$, algebraic unitarity $R(z)\,R(-z) = \mathrm{Id}$, fusion product, $T$-matrix.
 \item \emph{Topological}: the output is a number (or finite collection of numbers) that depends on a manifold type, a knot, or a topological invariant, but not on any continuous parameter. Examples: Verlinde dimensions, knot group $\pi_1(S^3 \setminus K)$, Alexander polynomial, Kirby invariance, $S$-matrix (modular, non-degenerate), Khovanov homology, hyperbolic volume.
 \item \emph{Hybrid}: the output has both algebraic and topological aspects. Examples: handle decomposition (topological existence, algebraic gluing data), modular characters (algebraic functions with topological transformation properties).
\end{itemize}
Of the $34$ substructures, $12$ are algebraic, $19$ are topological, and $3$ are hybrid. Among the algebraic structures, the lift rate to 6d is $100\%$: every algebraic structure of 3d CS lifts to a (richer) algebraic structure in 6d. Among the topological structures, the lift rate is $0\%$: no topological structure lifts. The hybrid structures have mixed behavior. The overall lift rate is $12/34 \approx 35\%$; the algebraic fraction of the total catalogue is $12/34$.

\begin{center}
\small
\renewcommand{\arraystretch}{1.3}
\begin{tabular}{lccc}
 \toprule
 \textbf{Invariant type} & \textbf{Count} & \textbf{Lift to 6d} & \textbf{Lift rate} \\
 \midrule
 Algebraic & $12$ & $12$ & $100\%$ \\
 Topological & $19$ & $0$ & $0\%$ \\
 Hybrid & $3$ & $0$ & $0\%$ \\
 \midrule
 \textbf{Total} & $\mathbf{34}$ & $\mathbf{12}$ & $\mathbf{35\%}$ \\
 \bottomrule
\end{tabular}
\end{center}

\paragraph{What the 6d theory produces.}
The 6d holomorphic CS theory is not a failed attempt at knot/manifold invariants. It produces eight algebraic classes of invariants, none of which are knot or manifold invariants:
\begin{enumerate}[label=\textup{(\roman*)}]
 \item \emph{Structure functions} $g(z) = \prod(z - h_i)/\prod(z + h_i)$: rational functions of a spectral parameter encoding the full representation theory. The 3d analogue is the quantum dimension $\dim_q(V)$ -- a \emph{number}, not a function. The entire $z$-dependence is genuinely 6d.
 \item \emph{Two-parameter $R$-matrices} $R(u,v)$ with Zamolodchikov factorization. The 3d analogue is the one-parameter $R(u)$ of $U_q(\frakg)$. The second parameter arises from the Omega-background (normal bundle grading does not directly give $(q,t)$; the passage is through equivariant localization).
 \item \emph{Spectral coproducts} $\Delta_z \colon A \to A \otimes A[\![z]\!]$, with $z$-degree at spin~$s$ equal to~$s$. The 3d coproduct $\Delta \colon A \to A \otimes A$ has no spectral parameter; the $z$-dependence encodes the ordered factorization ($\Eone$).
 \item \emph{Shadow towers} $\{S_k\}_{k \geq 1}$: the $\Ainf$-correction coefficients $\delta^{(k)}$ of the coproduct, equivalently the Maurer--Cartan expansion of the shadow. Class $\mathbf{G}$ (formal) has $S_k = 0$ for $k \geq 2$; class $\mathbf{M}$ (mock modular) has $S_k \neq 0$ for all~$k$. No 3d analogue exists (the 3d coproduct is exact).
 \item \emph{Mock modular partition functions}: for class~$\mathbf{M}$ algebras, the characters are mock modular forms (not genuine modular), arising from the infinite shadow tower. The 3d characters are genuine modular forms (Kac--Peterson).
 \item \emph{ZTE corrections}: the even charge-$2$ obstruction is killed by a rational ternary correction, while the full charge-preserving tetrahedron solution requires the odd sectors as well. No 3d analogue exists, since the YBE suffices in 3d.
 \item \emph{Miki automorphism}: the $\mathbb{Z}/3$ cyclic permutation of $(q,t,s)$ from the CY$_3$ constraint $\sum\epsilon_i = 0$, acting on the positive half $\mathrm{CoHA}(\C^3) = Y^+$; the ambient $S_3$ lives on the Drinfeld double after a $\mathbb{Z}/2$-involution. Algebra-specific, not operadic.
 \item \emph{CY partition functions} (DT/PT/GW): generating functions of curve-counting invariants in CY threefolds. Entirely 6d: 3d has no curve counting.
\end{enumerate}

\paragraph{The universal specialization principle.}
The relationship between 3d and 6d is that of specialization to universality:
\begin{center}
\renewcommand{\arraystretch}{1.3}
\small
\begin{tabular}{lcll}
 \toprule
 \textbf{Dimension} & \textbf{Parameters} & \textbf{Output type} & \textbf{Algebra} \\
 \midrule
 3d ($\C \times \R$) & $1$ (level $k$, discrete) & topological (numbers) & $\widehat{\frakg}$ (Kac--Moody) \\
 5d ($\C^2 \times \R$) & $1$ ($\hbar$, continuous) & algebraic ($f(u)$) & $Y(\frakg)$ (Yangian) \\
 6d ($\C^3$) & $2$ ($(q,t)$, continuous) & algebraic ($f(u,v)$) & $U_{q,t}(\widehat{\widehat{\fgl}}_1)$ \\
 \bottomrule
\end{tabular}
\end{center}
Each step \emph{up} in dimension universalizes the structure: discrete parameters become continuous, numbers become functions, finite-dimensional algebras become infinite-dimensional. The passage \emph{down} recovers the lower-dimensional theory by specialization:
\begin{equation}\label{eq:universal-specialization-bridge}
 \underbrace{U_{q,t}(\widehat{\widehat{\fgl}}_1)}_{\text{6d, generic $(q,t)$}}
 \;\xrightarrow[\text{codim-$2$}]{\;\;q^N = t^M = 1\;\;}
 \underbrace{U_{q,t}^{\mathrm{trunc}}}_{\text{truncated}}
 \;\xrightarrow[\text{codim-$1$}]{\;\;t \to 1\;\;}
 \underbrace{U_q(\fsl_N)\;\text{at root of unity}}_{\text{3d, level }k = N - h^\vee}
\end{equation}
\emph{Neither step is constructed.} Arrow~1 requires the truncated quantum toroidal algebra at joint root of unity (the Hopf ideal, its compatibility with $\Delta_z$, and the $E_3$ factorization of the quotient are all open). Arrow~2 requires understanding the $t \to 1$ limit of the Miki automorphism. The combined bridge is an open problem.

\begin{conjecture}[Root-of-unity bridge]
\label{conj:root-of-unity-bridge}
\ClaimStatusConjectured
The two-step specialization~\eqref{eq:universal-specialization-bridge} recovers the 3d Chern--Simons quantum group $U_q(\frakg)$ at root of unity from the 6d quantum toroidal algebra $U_{q,t}(\widehat{\widehat{\fgl}}_1)$ at generic parameters. The truncated quotient $U_{q,t}^{\mathrm{trunc}}$ at $q^N = t^M = 1$ should carry a finite set of simple modules whose fusion rules, in the $t \to 1$ limit, reproduce the Verlinde algebra $V_k(\frakg)$.
\end{conjecture}

\paragraph{Why the volume conjecture has \texorpdfstring{$0\%$}{0 percent} lift rate.}
The volume conjecture is the most deeply 3d structure in the catalogue: it lives at the intersection of algebra and topology, relating the colored Jones polynomial $J_N(K; e^{2\pi i/N})$ (an algebraic quantity evaluated at a topological point) to the hyperbolic volume $\mathrm{Vol}(S^3 \setminus K)$ (a topological quantity requiring Mostow rigidity). At the 6d level, \emph{neither} ingredient exists. The colored Jones polynomial requires root-of-unity specialization (Arrow~1 of the bridge, unconstructed). The hyperbolic volume requires Mostow rigidity, which fails for CY$_3$ manifolds with positive-dimensional moduli ($h^{2,1} > 0$). The volume conjecture is the canary in the coal mine: its total failure to lift is the clearest evidence that the 6d theory is algebraic, not topological.

\begin{remark}[The correct relationship]
\label{rem:correct-3d-6d-relationship}
The 6d holomorphic CS theory is the \emph{universal algebraic completion} of the 3d--5d tower. It produces the full representation-theoretic data (structure functions, $R$-matrices, coproducts, shadow towers) from which all lower-dimensional specializations should be extractable (Conjecture~\ref{conj:root-of-unity-bridge}). It does \emph{not} produce knot or manifold invariants directly; those arise only after the two-step bridge (truncation $+$ NS limit) to 3d. The analogy is
\[
 \text{6d} : \text{3d} \;\;::\;\; \Z[x] : \Z/n\Z
 \quad\text{(polynomial ring : quotient ring)},
\]
where the polynomial ring $\Z[x]$ contains all the information about all quotients $\Z[x]/(x^n - 1)$, but it is not itself a finite ring, and the specific properties of each quotient (e.g., the number of units) emerge only upon specialization.
\end{remark}

\noindent The topological--algebraic classification, lift rates,
root-of-unity bridge, 6d output catalogue, dimensional hierarchy, and
\(\kappa_\bullet\)-convention split are the data used in this dichotomy.


\section{The Costello--Paquette defect chiral algebra and the \texorpdfstring{$5d \to 6d$}{5d -> 6d} boost}
\label{sec:costello-paquette-defect}

The Costello--Paquette construction (arXiv:2104.14573, arXiv:2208.00354) identifies the boundary chiral algebra of holomorphic Chern--Simons theory as the endomorphism algebra of a \emph{universal defect}: a canonical line defect $D_{\mathrm{univ}}$ whose module category exhausts all Wilson line modules. The construction is established at $5d$ (Costello 2013; Costello--Paquette 2021), producing the affine Yangian $Y(\widehat{\fgl}_1)$, and admits a conjectural boost to $6d$, where the target is the quantum toroidal algebra $U_{q,t}(\widehat{\widehat{\fgl}}_1)$. The $5d$ universal defect and its $6d$ lift culminate in the form factor map $\mathsf{ff} \colon Z^{\mathrm{der}}_{\mathrm{ch}}(A) \to A^!$ bridging the bulk algebra to the defect algebra.

\subsection{The universal defect at each dimension}
\label{subsec:cp-universal-defect-dims}

\begin{construction}[Costello--Paquette universal defect]
\label{constr:cp-universal-defect}
Let $M$ be one of the three ambient spaces in the holomorphic CS hierarchy (Construction~\ref{constr:hol-cs-hierarchy}). The \emph{universal defect} $D_{\mathrm{univ}}$ is a canonical codimension-$2$ defect supported on a holomorphic curve $C \hookrightarrow M$, characterised by the universality property:
\[
 \End(D_{\mathrm{univ}}) \;=\; A, \qquad \Hom(D_{\mathrm{univ}}, W_V) \;=\; V
\]
where $A$ is the boundary chiral algebra and $W_V$ is the Wilson line in representation~$V$. The defect $D_{\mathrm{univ}}$ is the unit object in the category of line defects: every Wilson line factors as $W_V = D_{\mathrm{univ}} \otimes_A V$. The data at each dimension:
\begin{center}
\small
\renewcommand{\arraystretch}{1.2}
\begin{tabular}{clccll}
\toprule
$\dim_\C(M)$ & $M$ & $\rank N_{C/M}$ & $\#$ spectral & $\End(D_{\mathrm{univ}})$ & Construction locus \\
\midrule
$1$ & $\Sigma \times \R$ & $0$ & $0$ & $V_k(\fgl_1)$ & Costello--Gwilliam boundary theorem \\
$2$ & $\C^2 \times \R$ & $1$ & $1$ ($u$) & $Y(\widehat{\fgl}_1)$ & Costello Yangian theorem \\
$3$ & $\C^3$ & $2$ & $2$ ($u, v$) & $U_{q,t}(\widehat{\widehat{\fgl}}_1)$ & requires 6d hCS completion \\
\bottomrule
\end{tabular}
\end{center}
The number of spectral parameters equals the rank of the normal bundle $N_{C/M}$, which is $\dim_\C(M) - 1$. The $\En$ level of the defect algebra equals $\dim_\C(M)$.
\end{construction}

\begin{proposition}[Structural invariants of the universal defect]
\label{prop:cp-defect-invariants}
For the universal defect $D_{\mathrm{univ}}$ in holomorphic CS on $M = \C^n$ with gauge algebra $\fgl_1$ and Omega-background $(h_1, h_2, h_3)$ satisfying $h_1 + h_2 + h_3 = 0$:
\begin{enumerate}[label=\textup{(\roman*)}]
 \item The effective level of the defect algebra is $\Psi = -\sigma_2 = -(h_1 h_2 + h_1 h_3 + h_2 h_3)$, independent of the dimension~$n$. In particular, $\kappa_{\mathrm{ch}}(\End(D_{\mathrm{univ}})) = \Psi$ at all three dimensions.
 \item The Koszul conductor on the free-field branch is $K = \kappa_{\mathrm{ch}}(A) + \kappa_{\mathrm{ch}}(A^!) = 0$ (the Koszul dual has $\kappa_{\mathrm{ch}}(A^!) = -\Psi$).
 \item At the self-dual point ($\sigma_3 = h_1 h_2 h_3 = 0$): $g(u) = 1$, the $R$-matrix is trivial, and $D_{\mathrm{univ}}$ is the trivial defect.
 \item The equivariant weight of the normal mode $A^{(m,n)}$ at $\dim_\C = 3$ is $m \cdot h_2 + n \cdot h_3$, with $m + n + 1$ giving the conformal weight (spin) of the corresponding $W_{1+\infty}$ current.
\end{enumerate}
\end{proposition}

\begin{proof}
The $\kappa_{\mathrm{ch}}$-preservation across dimensions is structural:
$\Psi=-\sigma_2$ depends only on the $\Omega$-background parameters, not
on the ambient dimension.  The same parameters therefore give the same
classical $r$-matrix $r(u)=-\sigma_2/u$ at every $E_n$ level.  The
self-dual triviality follows from
$g(u)=\prod_i(u-h_i)/\prod_i(u+h_i)=1$ when one $h_i$ is zero and
$h_1+h_2+h_3=0$.  The normal-mode weight in dimension three is the
tautological $\Omega$-weight $m h_2+n h_3$, and the shift by one gives
the conformal weight of the corresponding $W_{1+\infty}$ current.
\end{proof}

\subsection{The form factor map: bulk-to-boundary Koszul projection}
\label{subsec:cp-form-factor}

The form factor map of Costello--Paquette encodes the bulk-to-boundary OPE as a map from the derived center (the bulk algebra) to the Koszul dual (the defect algebra). This is the bridge between holomorphic CS observables and quantum group representations.

\begin{construction}[Form factor map]
\label{constr:form-factor-map}
Let $A$ be the boundary chiral algebra of holomorphic CS on $M$. The \emph{form factor map} is the composition
\[
 \mathsf{ff} \colon Z^{\mathrm{der}}_{\mathrm{ch}}(A) \xrightarrow{\;\text{HH-to-bar}\;} B(A) \xrightarrow{\;D_{\Ran}\;} A^!
\]
where the first arrow is the Hochschild-to-bar comparison map (Vol~I Theorem~H, inverse direction) and the second is the Verdier duality on $\Ran(C)$. At the level of modes, for a bulk operator $O$ of conformal weight $\Delta_O$:
\[
 \mathsf{ff}(O)(|v\rangle) \;=\; \Res_{z=0}\, z^{\Delta_O - 1}\, O(z)\, |v\rangle,
\]
extracting the singular part of the bulk-to-boundary OPE: precisely the data captured by the bar complex $B(A)$.
\end{construction}

\begin{proposition}[Form factor axioms]
\label{prop:form-factor-axioms}
The form factor map $\mathsf{ff} \colon Z^{\mathrm{der}}_{\mathrm{ch}}(A) \to A^!$ satisfies:
\begin{enumerate}[label=\textup{(FF\arabic*)}]
 \item $\mathsf{ff}$ is a map of $A$-bimodules (the bulk-boundary OPE is $A$-equivariant from both left and right).
 \item $\mathsf{ff}(\mathbf{1}_{\mathrm{bulk}}) = \mathbf{1}_{A^!}$ (the identity form factor is the Koszul identity).
 \item For $n \geq 2$, $\mathsf{ff}$ intertwines the $E_n$ braiding on the bulk with the Koszul-reflected braiding on $A^!$.
 \item $\kappa_{\mathrm{ch}}$ is preserved: $\kappa_{\mathrm{ch}}(\mathsf{ff}(O)) = \kappa_{\mathrm{ch}}(O)$ for all $O$, and the Koszul conductor on the free-field branch is $K = 0$.
 \item At the self-dual point ($\sigma_3 = 0$), $\mathsf{ff}$ is the identity (the defect is trivial).
\end{enumerate}
The form factor pole structure at each spin:
\begin{center}
\small
\renewcommand{\arraystretch}{1.15}
\begin{tabular}{ccll}
\toprule
\textbf{Spin} & \textbf{Pole order} & \textbf{Residue} & \textbf{Identification} \\
\midrule
$1$ & $1$ & $\Psi = -\sigma_2$ & Classical $r$-matrix $r^{\mathrm{Heis}}(z) = \Psi/z$ \\
$2$ & $3$ & $c/2 = 1/2$ & Classical $r$-matrix $r^{\mathrm{Vir}}(z) = (c/2)/z^3 + 2T/z$ \\
\bottomrule
\end{tabular}
\end{center}
\end{proposition}

\begin{proof}
Axioms~(FF1)--(FF5) follow from the factorisation $\mathsf{ff} = D_{\Ran} \circ B$, the properties of the bar complex (Vol~I), and the Verdier duality (which preserves bimodule structure and $\kappa_{\mathrm{ch}}$). The pole orders and residues agree with Proposition~\ref{prop:codim2-defect-ope}: the spin-$1$ residue $\Psi$ is the $J_{(1)}J$ OPE coefficient, and the spin-$2$ residue $c/2 = 1/2$ is $T_{(3)}T$. The independence of the spin-$2$ form factor from $\Psi$ follows from $c = 1$ for $\fgl_1$ at any level (Proposition~\ref{prop:codim2-defect-ope}).
\end{proof}

\subsection{The \texorpdfstring{$5d \to 6d$}{5d -> 6d} dimensional boost}
\label{subsec:cp-5d-to-6d-boost}

The Costello--Paquette construction at $5d$ is established: the affine Yangian $Y(\widehat{\fgl}_1)$ is the boundary chiral algebra of $5d$ holomorphic CS on $\C^2 \times \R$ (Theorem~\ref{thm:5d-boundary-yangian}). The \emph{boost} to $6d$ is the conjectural lift of this construction to the $6d$ holomorphic theory on $\C^3$, producing the quantum toroidal algebra $U_{q,t}(\widehat{\widehat{\fgl}}_1)$.

\begin{conjecture}[Dimensional boost: \texorpdfstring{$5d \to 6d$}{5d -> 6d}]
\label{conj:cp-5d-to-6d-boost}
\ClaimStatusConjectured
The Costello--Paquette universal defect construction lifts from $5d$ to $6d$ with the following structural changes:
\begin{center}
\small
\renewcommand{\arraystretch}{1.2}
\begin{tabular}{lll}
\toprule
\textbf{Feature} & \textbf{$5d$ theorem} & \textbf{$6d$ lift} \\
\midrule
Spectral parameters & $1$ ($u$) & $2$ ($u, v$) \\
$E_n$ level & $\Etwo$ & $E_3$ \\
Boundary algebra & $Y(\widehat{\fgl}_1)$ & $U_{q,t}(\widehat{\widehat{\fgl}}_1)$ \\
$R$-matrix & $R(u)$, satisfies YBE & $R(u,v)$, should satisfy parametric YBE \\
Consistency eqn & Yang--Baxter (2-simplex) & Zamolodchikov tetrahedron (3-simplex) \\
$\kappa_{\mathrm{ch}}$ & $-\sigma_2$ & $-\sigma_2$ (preserved) \\
Affinization & single ($\widehat{\fgl}_1$) & double ($\widehat{\widehat{\fgl}}_1$) \\
\bottomrule
\end{tabular}
\end{center}
The boost is not dimensional reduction in reverse: the $6d$ theory has genuinely new structure (the second affinization, the Miki automorphism, and the ZTE obstruction) that cannot be obtained from $5d$ alone.

The two-parameter structure function of the $6d$ defect at charge~$1$ is
\begin{equation}\label{eq:cp-two-param-structure}
 G_{\mathrm{ch}}(u, v) \;=\; g(u) \cdot g(v) \cdot g(u - v),
\end{equation}
where $g(u) = \prod_{i=1}^3(u - h_i)/\prod_{i=1}^3(u + h_i)$ is the Yangian structure function. This is the Zamolodchikov factorisation of the charge-$1$ two-parameter $R$-matrix. At charge $\geq 2$, the factorised ansatz $S_{ijk} = R_{ij}R_{ik}R_{jk}$ fails the Zamolodchikov tetrahedron equation at $O(\sigma_3^2)$ where $\sigma_3 = h_1 h_2 h_3$ (Theorem~\ref{thm:zte-failure}), proving the $E_3$ structure is genuinely nontrivial.
\end{conjecture}

\begin{remark}[What is preserved and what changes]
\label{rem:cp-boost-preservation}
The $\kappa_{\mathrm{ch}}$-preservation under the boost ($\kappa_{\mathrm{ch}} = -\sigma_2$ at both $5d$ and $6d$) is a structural consequence: the effective level $\Psi = -\sigma_2$ of the Heisenberg subalgebra on the defect curve~$C$ depends only on the Omega-background parameters, not on the ambient dimension. What \emph{changes} is the structure of the higher-spin sectors:
\begin{itemize}
 \item At $5d$: the spin-$s$ currents form the $\Etwo$-braided algebra $Y(\widehat{\fgl}_1)$, with one spectral parameter $u$ from the normal direction.
 \item At $6d$: the spin-$s$ currents form the $E_3$-braided algebra $U_{q,t}(\widehat{\widehat{\fgl}}_1)$, with two spectral parameters $(u, v)$ from the rank-$2$ normal bundle $N_{C/\C^3} = \C^2$.
\end{itemize}
The second spectral parameter $v$ is the new ingredient at $6d$; it generates the second affinization (the second hat in $\widehat{\widehat{\fgl}}_1$). The $\Etwo \to E_3$ promotion is the holomorphic-geometric incarnation of the derived center construction (Remark~\ref{rem:deligne-hochschild}): at the level of representation categories, $\cZ(\Rep^{\Etwo}(Y)) \simeq \Rep^{E_3}(Z^{\mathrm{der}}_{\mathrm{ch}}(Y))$, where the right side is conjectural.
\end{remark}

\subsection{The K3 universal defect}
\label{subsec:cp-k3-defect}

For $X = K3 \times E$, the holomorphic CS hierarchy specialises to three boundary algebras, each controlled by the Mukai lattice $\Lambda_{K3}$ of signature~$(4, 20)$.

\begin{conjecture}[K3 universal defect hierarchy]
\label{conj:cp-k3-defect-hierarchy}
\ClaimStatusConjectured
The universal defect $D_{\mathrm{univ}}$ for K3 targets produces the following boundary algebras:
\begin{center}
\small
\renewcommand{\arraystretch}{1.2}
\begin{tabular}{clcll}
\toprule
$\dim_\C(M)$ & Ambient & $\kappa_{\mathrm{ch}}$ & Boundary algebra & Construction locus \\
\midrule
$1$ & K3 & $2$ & $V_k(\mathfrak{g}_{K3})$ & CY-A$_2$ \\
$2$ & $K3 \times \C \times \R$ & $3$ & $Y(\mathfrak{g}_{K3})$ & requires 5d K3 hCS construction \\
$3$ & $K3 \times E \times \C$ & $3$ & $U_{q,t}(\widehat{\mathfrak{g}}_{K3})$ & requires 6d non-Lagrangian hCS construction \\
\bottomrule
\end{tabular}
\end{center}
The K3 structure function is degree-$(24, 24)$:
\[
 g_{K3}(u) \;=\; \prod_{\alpha \in \Phi_{K3}^+} \frac{u - \alpha(h)}{u + \alpha(h)},
\]
where $\Phi_{K3}^+$ is the positive root system of the Mukai lattice, with $|\Phi_{K3}^+| = 24$. The structure function satisfies the reflection identity $g_{K3}(u) \cdot g_{K3}(-u) = 1$.

The Wilson lines wrapping curves in $K3$ are indexed by Mukai vectors $v \in H^*(K3, \Z)$, with charge lattice $\Lambda_{K3}$ of signature~$(4, 20)$, and fusion controlled by the Mukai pairing.

The $\kappa_\bullet$-spectrum, distinguishing total-space from fiber contributions:
\[
 \kappa_{\mathrm{cat}} = 0, \quad \kappa_{\mathrm{ch}}^{\mathrm{Hodge}} = 0, \quad \kappa_{\mathrm{ch}}^{\mathrm{Heis}} = 3, \quad \kappa_{\mathrm{BKM}}(\Delta_5) = 5, \quad \kappa_{\mathrm{fiber}} = 24.
\]
Here $\kappa_{\mathrm{cat}}(K3 \times E) = \chi(\cO_{K3 \times E}) = \chi(\cO_{K3}) \cdot \chi(\cO_E) = 2 \cdot 0 = 0$ by K\"unneth; the compact total-space Hodge supertrace $\kappa_{\mathrm{ch}}^{\mathrm{Hodge}}(A_{K3\times E})$ is also $0$, and the K3-fiber value $\chi(\cO_{K3}) = 2$ is a \emph{fiber} invariant, not the total-space one.
The canonical Stage-2 Heisenberg--Cartan additivity is instead $\kappa_{\mathrm{ch}}^{\mathrm{Heis}}(K3 \times E) = \kappa_{\mathrm{ch}}^{\mathrm{Heis}}(K3) + \kappa_{\mathrm{ch}}^{\mathrm{Heis}}(E) = 2 + 1 = 3$. The two Koszul branches are distinct: free-field/Heisenberg ($\fgl_1$-sector, $K = 0$, $\kappa_{\mathrm{ch}}^{\mathrm{Heis}}(A^!) = -3$) and Mukai-doubling/BKM ($\kappa_{\mathrm{ch}}^{\mathrm{Mukai}}(A) = 4$, $\kappa_{\mathrm{ch}}^{\mathrm{Mukai}}(A^!) = 4$, conductor $K^{\kappa_{\mathrm{ch}}^{\mathrm{Mukai}}} = 8 = 2\,c_+(\mathrm{Mukai}(K3))$; see Theorem~\ref{thm:qca-hbar-minus-eighth-lusztig} and Remark~\ref{rem:qca-three-faces-of-eight}). The Borcherds weight $\kappa_{\mathrm{BKM}}(\Delta_5) = c_{\Delta_5}(0)/2 = 5$ is a level-$4$ automorphic invariant of the input Siegel form, not the Koszul conductor. The numerical equality \(5=3+2\) at \(N=1\) mixes a Borcherds weight, a Heisenberg rank, and a fibre Euler characteristic; it is not an invariant identity. The formula $\kappa_{\mathrm{BKM}}(\Phi_N)=\kappa_{\mathrm{ch}}^{\mathrm{Heis}}+\chi(\cO_{\mathrm{fiber}})$ fails as a statement of invariants for every \(N\in\{1,2,3,4,6\}\).
\end{conjecture}



\section{The quantum chiral Koszul dual}
\label{sec:qca-koszul-dual}

The Koszul dual of a boundary chiral algebra is the defect algebra. The $\En$ structure on the defect algebra is inherited from the ambient holomorphic theory: a deformation of the chiral CE cochains by the theory's quantum parameters. The deficiency: the $E_n$ level of $A^!$ is determined by the ambient geometry, not by $A$ alone, and the 6d theory on $\C^3$ produces an $E_3$-chiral deformation that no lower-dimensional argument can access.

\subsection{Chiral Koszul duality: boundary and defect}
\label{subsec:chiral-koszul-bd}

\begin{definition}[Chiral Koszul dual]
\label{def:chiral-koszul-dual}
For a chiral algebra $A$ on a curve $C$ with bar complex $B(A) = T^c(s^{-1}\bar{A})$ (Definition~\ref{def:ordered-bar}, Chapter~\ref{ch:e1-chiral}), the \emph{chiral Koszul dual} is
\[
 A^! \;:=\; D_{\Ran}(B(A)),
\]
the Verdier dual of the bar complex as a factorization coalgebra on $\Ran(C)$. This is the Volume~I Verdier leg (Theorem~A), specialized to the CY setting. The algebra $A^!$ is the \emph{defect algebra}: it controls the category of line operators that can end on the boundary where $A$ lives.
\end{definition}

\begin{proposition}[Three dualities are distinct]
\label{prop:three-dualities}
The chiral Koszul dual $A^!$ is distinct from:
\begin{enumerate}[label=\textup{(\roman*)}]
 \item \emph{Bar-cobar inversion}: $\Omega(B(A)) \simeq A$ recovers the original algebra (the counit face of Vol~I Theorem~A on the Koszul locus; the strict ML weight-completed coderived strengthening is Vol~I Theorem~B on all four classes G/L/C/M -- see Vol~I Corollary~\cite{VolI}). This is $A$, not $A^!$.
 \item \emph{Derived center}: $Z^{\mathrm{der}}_{\mathrm{ch}}(A)
 = C^\bullet_{\mathrm{ch}}(A, A) = \RHom(\Omega B(A), A)$ is the
 closed-sector algebra at level $3$ (Vol~I Theorem~H; Vol~III
 Theorem~\ref{thm:bulk-three-faces-coincidence}). A physical bulk
 identification requires the relevant open--closed comparison datum.
 This is the ``$E_2$ uplift,'' not $A^!$.
 \item \emph{Drinfeld center}: $\cZ(\Rep^{\Eone}(A))$ is a braided monoidal category (Chapter~\ref{ch:drinfeld-center}). This lives on representation categories, not on algebras.
\end{enumerate}
The relationships: $A^!$ is the algebra controlling the defect
(boundary); $Z^{\mathrm{der}}_{\mathrm{ch}}(A)$ is the algebraic
closed-sector object, which becomes a physical bulk only after the
open--closed comparison is supplied; the Drinfeld center is the
categorical passage from $\Eone$ (boundary) to $\Etwo$ (bulk)
($\cZ(\Rep^{\Eone}(A))$ is a braided monoidal \emph{category};
$Z^{\mathrm{der}}_{\mathrm{ch}}(A)$ is a chiral \emph{algebra}; the
Drinfeld center is the categorification of the derived center under
the BZFN hypotheses). The bulk acts on the defect via
$\mu_{\mathrm{bb}}$ (Construction~\ref{constr:universal-defect}).
Three constructions reach the bulk from distinct starting points: the
Hochschild centre $Z^{\mathrm{der}}_{\mathrm{ch}}(A)$ (this
proposition); the positive-half Hall presentation $Y^{+}(X)$
(Definition~\ref{def:bulk-yplus}); the Drinfeld double
$G(X)=\mathcal D_\hbar(Y^{+}(X))$ (Definition~\ref{def:bulk-gx}).
The three meet at one level-$3$ object on the toric stratum in
ordinary chain ambient
(Theorem~\ref{thm:bulk-three-faces-coincidence}).

The distinctness of $A^!$ from $Z^{\mathrm{der}}_{\mathrm{ch}}(A)$ is a fundamental point: for the Kac--Moody family $A = V_k(\frakg)$, the Koszul dual is the completed Verdier-dual bar object $D_{\Ran}B^{\mathrm{ord}}(V_k(\frakg))$, whose reflected current presentation is $V_{k'}(\frakg)$ at $k' = -k - 2h^\vee$, while $Z^{\mathrm{der}}_{\mathrm{ch}}(A) = C^\bullet_{\mathrm{ch}}(V_k(\frakg), V_k(\frakg))$ is the chiral Hochschild cochain complex. These are manifestly different objects: $A^!$ controls defects through the CE/bar-cochain target; $Z^{\mathrm{der}}_{\mathrm{ch}}(A)$ is a derived endomorphism complex.
\end{proposition}

\subsection{The \texorpdfstring{$\En$}{E-n} level of the Koszul dual}
\label{subsec:en-koszul-level}

The $\En$ level of $A^!$ is determined by the ambient theory, not by $A$ alone.

\begin{conjecture}[\texorpdfstring{$\En$}{E-n}-structure on the chiral Koszul dual from holomorphic CS]
\label{conj:en-koszul-from-hcs}
\ClaimStatusConjectured
Let $A_n$ denote the boundary chiral algebra of holomorphic Chern--Simons theory (or its 6d algebraic surrogate) on $\C^n$ projected to a curve $C$, so that the full theory on $\C^n$ carries $E_n$-factorization structure. Then:
\begin{enumerate}[label=\textup{(\roman*)}]
 \item The chiral Koszul dual $A_n^!$ inherits $E_n$-chiral factorization structure from the ambient $E_n$ theory on $\C^n$.
 \item At $n = 1$ (3d CS): $A_1^! = D_{\Ran}B^{\mathrm{ord}}(V_k(\frakg))$, with reflected current presentation $V_{k'}(\frakg)$ at $k' = -k - 2h^\vee$, is an $\Eone$-chiral algebra. The $\Eone$ structure is the standard one.
 \item At $n = 2$ (5d CS): $A_2^!$ carries $\Etwo$-chiral structure on $\C^2$. On a curve, $A_2^!|_C$ is the Koszul dual of the affine Yangian with inverted parameters.
 \item At $n = 3$ (6d theory): $A_3^!$ carries $E_3$-chiral structure on $\C^3$. This is the \emph{large chiral algebra} of the Costello construction: it contains the quantum toroidal algebra (on $C$), the affine Yangian ($\Etwo$ on $\C^2$), and the $E_3$ master structure on $\C^3$.
\end{enumerate}
By item: (i) is the general principle (conjectural at $n \geq 2$). (ii) is a theorem in the completed ordered-bar ambient: the Feigin--Frenkel reflection gives the current presentation of $D_{\Ran}B^{\mathrm{ord}}(V_k(\frakg))$. (iii) is partially established: the affine Yangian Koszul dual is constructed by Schiffmann--Vasserot; the $\Etwo$-chiral structure on the Koszul dual is conjectural. (iv) is fully conjectural and requires the non-Lagrangian 6d algebraic framework of Remark~\ref{rem:6d-not-lagrangian}.
\end{conjecture}

\begin{remark}[\texorpdfstring{$E_3$}{E-3} chiral is not \texorpdfstring{$E_3$}{E-3} symmetric]
\label{rem:e3-not-symmetric}
The $E_3$ structure in Conjecture~\ref{conj:en-koszul-from-hcs}(iv) is \emph{not} the trivially braided $E_3$ of Chapter~\ref{ch:en-factorization} (where $\pi_1(\Conf_2(\R^6)) = 0$ kills the braiding). The holomorphic refinement and the Omega-background deformation produce a nontrivial $E_3$-chiral structure: the factorization algebra on $\Conf_n(\C^3)$ acquires a nontrivially braided structure after deformation by the equivariant parameters $(h_1, h_2, h_3)$. Note that $\Conf_2(\C^3) \simeq \Conf_2(\R^6) \simeq S^5$ as topological spaces, so the braiding is trivial at the topological level; the nontriviality is a feature of the \emph{deformed holomorphic} factorization structure, not of the underlying topology. Without the Omega-background, the braiding is symmetric (the $E_3$-chiral structure has trivial braiding, not that $E_3$ ``becomes'' $E_\infty$: the operadic level remains $E_3$, but the braiding data is trivial).
\end{remark}

\subsection{Deformation of chiral CE cochains to \texorpdfstring{$E_3$}{E-3}}
\label{subsec:ce-deformation-e3}

The chiral CE cochains of an $E_\infty$-chiral algebra $A$ (such as a BD commutative chiral algebra) carry a natural $E_\infty$-algebra structure from the higher Deligne conjecture (Remark~\ref{rem:deligne-hochschild}); in particular, they carry $E_3$-algebra structure. The \emph{deformed} $E_3$-chiral structure relevant to this chapter arises not from the Deligne conjecture alone, but from a deformation of the CE differential by the Omega-background parameters of the holomorphic CS theory.

\begin{construction}[Quantum deformation of chiral CE cochains]
\label{constr:quantum-ce-deformation}
Let $A = V_k(\frakg)$ be a Kac--Moody vertex algebra (an $E_\infty$-chiral algebra). The chiral CE cochains $C^\bullet_{\mathrm{ch}}(A, A) = Z^{\mathrm{der}}_{\mathrm{ch}}(A)$ carry:
\begin{enumerate}[label=\textup{(\roman*)}]
 \item An $\Etwo$-algebra structure from the Deligne conjecture: the cup product and brace operations realize the Dunn additivity decomposition $\Etwo \simeq \Eone \otimes \Eone$ (operadic tensor product, not the coproduct over $E_0$).
 \item A classical Poisson bracket $\{-, -\}_{\mathrm{cl}}$ of degree $-1$ from the Gerstenhaber structure on $\HH^\bullet(A, A)$.
 \item A \emph{quantum deformation} parametrized by the Omega-background $\boldsymbol{h} = (h_1, h_2, h_3)$ with $h_1 + h_2 + h_3 = 0$ (so only two of $(h_1, h_2, h_3)$ are independent). At $\boldsymbol{h} = 0$, the structure is the undeformed $\Etwo$-structure. At generic $\boldsymbol{h}$, the deformed differential
 \[
 d_{\boldsymbol{h}} = d_{\mathrm{CE}} + h_1 \cdot \partial_1 + h_2 \cdot \partial_2 + h_3 \cdot \partial_3
 \]
 (where $\partial_i$ are the equivariant differential operators on $\C^3$ associated to the three coordinate $\C^*$-actions, and the constraint $h_1 + h_2 + h_3 = 0$ ensures $d_{\boldsymbol{h}}^2 = 0$ via the Calabi--Yau condition $\det = 1$ on the torus action) promotes the structure to $E_3$-chiral. The promotion uses the topological $E_3$-operad action on $\Conf_n(\C^3)$, transferred to the algebraic setting via Kontsevich formality; this is the topological $E_3$ of Remark~\ref{rem:two-e3-ap154}(b), not the algebraic $E_3$ of~(a).
\end{enumerate}
\end{construction}

\begin{remark}[Two \texorpdfstring{$E_3$}{E-3} structures: algebraic vs topological]
\label{rem:two-e3-ap154}
The $E_3$ structure in Construction~\ref{constr:quantum-ce-deformation} arises from two independent sources:
\begin{enumerate}[label=\textup{(\alph*)}]
 \item \emph{Algebraic $E_3$}: the higher Deligne conjecture (now a theorem: Lurie, 2012; Costello, 2007) states that for an $E_n$ algebra $A$, the Hochschild cochains $\HH^\bullet(A, A)$ carry an $E_{n+1}$ structure. For $\Eone$ input: $\Etwo$ on $\HH^\bullet$ (standard Deligne conjecture). For $E_\infty$ input: $E_\infty$ on $\HH^\bullet$ (since $\infty + 1 = \infty$). In particular, for $A = V_k(\frakg)$ ($E_\infty$-chiral), $\HH^\bullet(A,A)$ carries \emph{at least} $E_3$; the full $E_\infty$ structure is available but only the $E_3$ sub-structure is relevant to the 6d holomorphic theory on $\C^3$.
 \item \emph{Topological $E_3$}: the $E_3$-operad acts on $\Conf_n(\C^3)$ via the framed little $3$-disks.
\end{enumerate}
These two $E_3$ structures agree under Kontsevich formality over $\Q$ and differ at the chain level. The holomorphic CS construction uses the topological $E_3$ from (b), deformed by the Omega-background; the algebraic $E_3$ from (a) is the classical limit. Integral agreement remains open.
\end{remark}

\subsection{Chiral quantum groups from defect Wilson lines}
\label{subsec:chiral-qg-wilson}

The chiral quantum group is the Hopf-algebra-like structure on the defect algebra $A^!$ extracted from the holomorphic Wilson lines of the ambient CS theory.

\begin{conjecture}[Chiral quantum group on \texorpdfstring{$\C^3$}{C-3}]
\label{conj:chiral-qg-c3}
\ClaimStatusConjectured
For the 6d holomorphic theory on $\C^3$ with $\fgl_1$ gauge algebra and Omega-background $(h_1, h_2, h_3)$:
\begin{enumerate}[label=\textup{(\roman*)}]
 \item The defect algebra $A^!_{\C^3}$ carries a coproduct
 \[
 \Delta_{\mathrm{ch}} \colon A^!_{\C^3} \;\longrightarrow\; A^!_{\C^3} \hat{\otimes} A^!_{\C^3}
 \]
 induced by the collision of two parallel Wilson lines in $\C^3$. The coproduct is coassociative and compatible with the $E_3$-chiral structure.
 \item The $R$-matrix is the universal braiding of the $E_3$ theory:
 \[
 R_{\mathrm{ch}}(u, v) \in A^!_{\C^3} \hat{\otimes} A^!_{\C^3}((u))((v))
 \]
 with $(u, v)$ the two spectral parameters from the two extra directions of $\C^3$ beyond the Wilson line. This $R$-matrix satisfies the parametric Yang--Baxter equation in \emph{two} spectral variables.
 \item The representation category $\Rep^{\Etwo}(A^!_{\C^3}|_{\C^2})$ is the braided monoidal category of quantum toroidal representations. The Drinfeld center $\cZ(\Rep^{\Eone}(A^!_{\C^3}|_C))$ should recover the $\Etwo$-braided category; at 5d (one dimension lower), the analogous statement is Theorem~\ref{thm:drinfeld-center-coha}, which establishes $\cZ(\Rep^{\Eone}(Y^+(\widehat{\fgl}_1))) \simeq \Rep^{\Etwo}(Y(\widehat{\fgl}_1))$. The 6d lift replaces $Y^+$ with $A^!_{\C^3}|_C$ and the full Yangian with the quantum toroidal algebra.
 \item The coproduct deviates from the primitive (Hopf) coproduct by the Dimofte $K$-matrix (Remark~\ref{rem:dimofte-k-matrix-cy}): $\Delta_{\mathrm{ch}}(x) = x \otimes 1 + K_{\cA}(z) \otimes x + \ldots$.
\end{enumerate}
Each item in this conjecture is verified independently at lower dimension: item~(i) at 5d (Schiffmann--Vasserot coproduct on CoHA), item~(ii) at 5d (Yangian $R$-matrix, Theorem~\ref{thm:yang-r-matrix-ybe}), item~(iii) at 5d (Theorem~\ref{thm:drinfeld-center-coha}). Each individual arrow is established at 5d, but the 6d composite is conjectural.
\end{conjecture}

\begin{conjecture}[Chiral quantum group on \texorpdfstring{$K3 \times E$}{K3 x E}]
\label{conj:chiral-qg-k3xe}
\ClaimStatusConjectured
For the 6d holomorphic theory on $K3 \times E$:
\begin{enumerate}[label=\textup{(\roman*)}]
 \item The defect algebra $A^!_{K3 \times E}$ on $E$ is a deformation of the Koszul dual of the BKM vertex algebra $V_{\mathfrak{g}_{\Delta_5}}$, with deformation controlled by the K3 complex structure moduli.
 \item The Wilson lines wrapping $E$ are indexed by Mukai vectors $v \in H^*(K3, \Z)$; the fusion of Wilson lines is controlled by the Mukai pairing.
 \item The coproduct on $A^!_{K3 \times E}$ is the chiral coproduct of Vol~I, specialized to the $K3 \times E$ root datum.
 \item The $\Etwo$-braided representation category of the chiral quantum group is the Verlinde category of the BKM superalgebra $\mathfrak{g}_{\Delta_5}$. This is consistent with Conjecture~\ref{conj:mtc-k3} (Chapter~\ref{ch:e2-chiral}), where $\Phi_{E_2}$ on the \(K3\times E\) input means the chiral Drinfeld centre \(Z^{\mathrm{ch}}(\Phi_{E_1})\), not a native \(E_2\)-output of the \(d=3\) specialisation.
\end{enumerate}
The assertion requires the 6d algebraic framework for $K3 \times E$ (non-Lagrangian; Remark~\ref{rem:6d-not-lagrangian}), the framed object-level \(\Phi_3\) construction of Theorem~\ref{thm:cy-to-chiral-d3}, and the identification of the BKM root datum with the K3 Mukai lattice. These are independent comparison data. The $\kappa_\bullet$-spectrum consistency check is $\kappa_{\mathrm{ch}}^{\mathrm{Heis}} = 3$, $\kappa_{\mathrm{BKM}}(\Delta_5) = 5$, $\kappa_{\mathrm{cat}} = \chi(\cO_{K3 \times E}) = 0$ (K\"unneth: $2 \cdot 0 = 0$; the K3-fiber value $\chi(\cO_{K3}) = 2$ is a separate invariant), and $\kappa_{\mathrm{fiber}} = 24$, with compact total-space $\kappa_{\mathrm{ch}}^{\mathrm{Hodge}}(A_{K3\times E})=0$.
\end{conjecture}

\begin{remark}[Summary: the holomorphic CS\texorpdfstring{$\to$}{->}chiral quantum group construction]
\label{rem:hcs-construction}
The full construction from holomorphic Chern--Simons to chiral quantum group:
\[
\begin{tikzcd}[column sep=1.8cm]
 \text{hol CS on } \C^n \ar[r, "\text{boundary}"] &
 A_n \text{ on } C \ar[r, "B(-)"] &
 B(A_n) \ar[r, "D_{\Ran}"] &
 A_n^! \ar[r, "\Rep^{E_2}"] &
 \text{chiral QG}
\end{tikzcd}
\]
At $n = 1$ (3d CS), the chain
\[
 V_k(\frakg) \to B(V_k) \to
 D_{\Ran}B^{\mathrm{ord}}(V_k(\frakg))
 \to \text{conformal blocks}
\]
is the Feigin--Frenkel reflected-current line \(k'=-k-2h^\vee\), together with
Vol~I Theorem~A for the Kac--Moody family.

At $n = 2$ (5d CS), the chain
\[
 Y(\widehat{\fgl}_1) \to B(Y) \to Y^!
 \to \Etwo\text{-braided representations}
\]
uses the constructed boundary algebra of Theorem~\ref{thm:5d-boundary-yangian},
the CoHA coproduct and Drinfeld centre of
Theorem~\ref{thm:drinfeld-center-coha}, and the MO/\(R\)-matrix
comparison of Theorem~\ref{thm:mo-chiral-rmatrix}.  The remaining
obligation is to identify \(B(Y)\) as a factorisation coalgebra on
\(\Ran(C)\) with the Schiffmann--Vasserot shuffle structure.

At $n = 3$ (6d theory), the expected chain
\[
 U_{q,t}(\widehat{\widehat{\fgl}}_1) \to B(-)
 \to E_3\text{-Koszul dual} \to \text{chiral quantum group}
\]
requires the non-Lagrangian 6d framework and the finite
DWR/Ran Rees comparison before any centre/double operation.

For \(K3\times E\), the expected BKM-related chain
\[
 \text{BKM-related algebra} \to B(-) \to
 \text{defect algebra} \to \text{BKM quantum group}
\]
uses the framed $\Phi_3$ object, the 6d algebraic framework, and the
centre/double passage separated in Proposition~\ref{prop:qca-finite-rees-layer-separation}.
\end{remark}

\subsection{Derived Maurer--Cartan space of chiral CE cochains}
\label{subsec:derived-mc-chiral-ce}

The chiral CE cochains $C^\bullet_{\mathrm{ch}}(L) = \mathrm{CE}^*(L)$ of a Lie conformal algebra $L$ carry a dg Lie algebra structure (the Gerstenhaber bracket). The finite formulas in this subsection use either a finite Rees/DWR truncation or the exterior CE model on a chosen finite generator space \(L_{\mathrm{gen}}\). The full chiral CE complex, with its \(\partial\)-descendant tower and completions, is a filtered or pro-object and is not measured by the binomial dimensions below without this truncation datum. The Maurer--Cartan equation
\begin{equation}\label{eq:mc-chiral-ce}
 d_{\mathrm{CE}}(\alpha) + \tfrac{1}{2}[\alpha, \alpha] = 0, \qquad \alpha \in \mathrm{CE}^1(L),
\end{equation}
parametrises \emph{flat chiral deformations} of the $\lambda$-bracket. The derived Maurer--Cartan space $\mathrm{MC}^{\mathrm{der}}(\mathrm{CE}^{*}(L))$ is the derived stack whose tangent complex at the trivial solution $\alpha = 0$ is the shifted cochain complex $T_0\mathrm{MC} = \mathrm{CE}^{*}(L)[1]$, with $H^0$ controlling infinitesimal automorphisms, $H^1$ controlling deformations, and $H^2$ controlling obstructions.

\begin{proposition}[MC space for the five families]
\label{prop:mc-five-families}
The derived Maurer--Cartan spaces of the chiral CE cochains for the five standard families are:
\begin{enumerate}[label=\textup{(\roman*)}]
 \item \emph{Heisenberg $H_k$ (class~G):} $\mathrm{MC} = \{\mathrm{pt}\}$ (abelian: $d = 0$ and $[\alpha, \alpha] = 0$). The tangent complex has $H^0 = 0$, $H^1 = 1$ (the level shift $k \mapsto k + \varepsilon$), $H^2 = 0$ (unobstructed). The derived structure is trivial: $\mathrm{MC} = \mathrm{MC}^{\mathrm{der}}$.
 \item \emph{Kac--Moody $V_k(\mathfrak{sl}_2)$ (class~L):} $\mathrm{MC} = \mathrm{Conn}_{\mathrm{flat}}(C, \mathfrak{sl}_2)$, the space of flat $\mathfrak{sl}_2$-connections on the curve~$C$. At the abstract level: $H^0 = H^1 = 0$ (Whitehead's theorem for the semisimple Lie algebra $\mathfrak{sl}_2$), $H^2 = 1$ (the top class $H^3(\mathfrak{sl}_2) = k$). The algebra $\mathfrak{sl}_2$ is \emph{rigid}: no infinitesimal deformations.
 \item \emph{Virasoro $\mathrm{Vir}_c$ (class~M):} $\mathrm{MC} = \mathrm{Proj}(C)$, the space of projective structures on~$C$, with $\dim \mathrm{Proj}(C) = 3g - 3$ at genus~$g$. The $L_\infty$ corrections from the shadow tower ($l_3, l_4, \ldots$) make $\mathrm{MC}^{\mathrm{der}}$ genuinely derived: the higher MC equation $d(\alpha) + \tfrac{1}{2}l_2(\alpha, \alpha) + \tfrac{1}{6}l_3(\alpha, \alpha, \alpha) + \cdots = 0$ has infinitely many terms.
 \item \emph{Mukai Heisenberg $H_{\mathrm{Muk}}$ (class~G, rank~$24$):} $\mathrm{MC} = \{\mathrm{pt}\}$ (abelian). The tangent complex has $H^0 = 24$ (the Mukai lattice automorphisms), $H^1 = \binom{24}{2} = 276$ (antisymmetric bilinear forms on the Mukai lattice), $H^2 = \binom{24}{3} = 2024$ (formal obstructions, all of which vanish by abelianity: $[\alpha,\alpha] = 0$). The derived structure is trivial.
 \item \emph{Yangian $Y(\widehat{\mathfrak{gl}}_1)$ (class~L, truncated to $3$ generators):} The MC equation has solutions parametrised by deformations of the Yangian commutation relations. The tangent complex has Nomizu cohomology $H^0 = 2$, $H^1 = 2$, $H^2 = 1$ (for the $3$-dimensional Heisenberg Lie algebra $\mathfrak{h}_3$).
\end{enumerate}
\end{proposition}

\begin{proof}
Items~(i) and~(iv): for an abelian Lie conformal algebra $L$ (all $0$-th products vanish), $d_{\mathrm{CE}} = 0$ and $[\alpha, \alpha] = 0$ for all $\alpha \in \mathrm{CE}^1(L)$. The MC equation is trivially satisfied at $\alpha = 0$ with no obstructions. The shifted CE cohomology is $H^k(\mathrm{CE}^*(L)[1]) = \binom{n}{k+1}$ where $n = \dim(L)$, since the differential vanishes.

Item~(ii): the Lie algebra $\mathfrak{sl}_2$ is semisimple; Whitehead's theorem gives $H^1(\mathfrak{sl}_2) = H^2(\mathfrak{sl}_2) = 0$ (rigidity). The top cohomology $H^3(\mathfrak{sl}_2) = k$ by Poincar\'e duality for the $3$-dimensional Lie algebra.

Item~(iii): the Virasoro Lie conformal algebra has a single chosen
generator $T$ of spin~$2$. In the finite exterior generator quotient,
$\mathrm{CE}^k_{\mathrm{gen}} = 0$ for $k \geq 2$, so
$H^1(\text{tangent}) = 0$ at that quotient level. The class~M
$L_\infty$ corrections $l_3(T,T,T) = -2c \cdot T$ and
$l_4(T,T,T,T) = (40/27) \cdot T$ \textup{(}at $c = 1$\textup{)}
contribute higher-order terms to the MC equation that live on the
ordered bar complex and in the completed chiral CE object, not in the
finite exterior quotient.

Item~(v): the Lie algebra underlying the truncated Yangian is the $3$-dimensional Heisenberg algebra $\mathfrak{h}_3 = \langle e_0, e_1, e_2 \mid [e_0, e_1] = e_2 \rangle$, $2$-step nilpotent. By Nomizu's theorem (1954), $H^k(\mathfrak{h}_3) = (1, 2, 2, 1)$ for $k = 0, 1, 2, 3$.
\end{proof}

\begin{proposition}[Finite exterior tangent Euler characteristic]
\label{prop:mc-tangent-euler-char}
\ClaimStatusProvedHere
Let \(L_{\mathrm{gen}}\) be a finite generator space of rank
\(n\geq1\), and let
\[
  T^{\mathrm{fin}}_0\mathrm{MC}
  :=\mathrm{Hom}\bigl(\Lambda^\bullet L_{\mathrm{gen}},\Bbbk\bigr)[1]
\]
be the finite exterior model for the tangent complex. Then
\[
 \chi(T^{\mathrm{fin}}_0\mathrm{MC})
 = \sum_{k=-1}^{n-1} (-1)^k \dim T^k
 = \sum_{k=-1}^{n-1} (-1)^k \binom{n}{k+1} = 0.
\]
The statement concerns the finite exterior quotient or a finite
Rees/DWR truncation. It does not compute the Euler characteristic of
the uncompleted full chiral CE complex when \(\partial\)-descendants or
infinite shadow towers are present.
\end{proposition}

\begin{proof}
Substituting $j = k + 1$ gives
\[
  \chi(T^{\mathrm{fin}}_0\mathrm{MC})
  = -\sum_{j=0}^{n} (-1)^j \binom{n}{j}
  = -(1 - 1)^n
  = 0
\]
for $n \geq 1$ by the binomial theorem. The proof uses only
\(\dim \Lambda^j L_{\mathrm{gen}}=\binom n j\); it therefore proves the
finite exterior statement and no stronger completed CE assertion.
\end{proof}

\begin{proposition}[ADE deformations of the Mukai Heisenberg]
\label{prop:ade-deformations-mukai}
The abelian Mukai Heisenberg $H_{\mathrm{Muk}}$ (rank~$24$, Mukai pairing of signature~$(4, 20)$) admits ADE deformations to non-abelian current algebras. For each ADE root system $R$ of rank~$r \leq 20$ embeddable in the negative-definite part of the Mukai lattice:
\begin{enumerate}[label=\textup{(\roman*)}]
 \item The deformed algebra has $\dim(\mathfrak{g}_R)$ non-abelian generators (where $\mathfrak{g}_R$ is the ADE Lie algebra: $\dim(\mathfrak{sl}_{n+1}) = n(n+2)$, $\dim(\mathfrak{so}_{2n}) = n(2n-1)$, $\dim(\mathfrak{e}_6) = 78$, $\dim(\mathfrak{e}_7) = 133$, $\dim(\mathfrak{e}_8) = 248$) and $24 - r$ abelian generators.
 \item The deformation space has dimension~$r$ (one parameter per simple root of~$R$).
 \item At level~$k$, the chiral modular characteristic of the non-abelian part is $\kappa_{\mathrm{ch}}(V_k(\mathfrak{g}_R)) = \dim(\mathfrak{g}_R) \cdot (k + h^\vee) / (2h^\vee)$.
\end{enumerate}
The ADE deformations correspond geometrically to ADE singularities of the K3 surface: a K3 with an $R$-singularity has enhanced gauge symmetry~$\mathfrak{g}_R$ from the collapsed rational curves.
\end{proposition}

\begin{proof}
The Mukai lattice $\Lambda = E_8(-1)^{\oplus 2} \oplus U^{\oplus 4}$ has signature $(4, 20)$. The negative-definite part has rank~$20$, and ADE root lattices of rank~$r \leq 20$ embed by Nikulin's embedding theorem (1979). The Lie algebra dimensions are standard (Humphreys, \emph{Introduction to Lie Algebras}, 1972). The $\kappa_{\mathrm{ch}}$ formula is the Kac--Moody modular characteristic from the landscape census. The geometric interpretation of ADE singularities as collapsed $(-2)$-curves follows from the McKay correspondence (Du Val singularities).
\end{proof}

\begin{remark}[Shadow class determines derived MC structure]
\label{rem:shadow-class-mc}
The shadow class (G/L/C/M) of the chiral algebra $A = U^{\mathrm{ch}}(L)$ determines the analytic type of $\mathrm{MC}^{\mathrm{der}}$:
\begin{itemize}
 \item Class~G (Gaussian): $\mathrm{MC}^{\mathrm{der}}$ is classical (no higher homotopies). Example: Heisenberg, Mukai Heisenberg.
 \item Class~L (rational): $\mathrm{MC}^{\mathrm{der}}$ is a derived scheme with finitely many $L_\infty$ corrections. Example: Kac--Moody.
 \item Class~C (convergent): $\mathrm{MC}^{\mathrm{der}}$ has convergent $L_\infty$ series. Example: bc system.
 \item Class~M (mock modular): $\mathrm{MC}^{\mathrm{der}}$ has Gevrey-$1$ divergent $L_\infty$ corrections, requiring Borel summation. The shadow tower $S_k$ contributes to the $k$-th homotopy of $\mathrm{MC}^{\mathrm{der}}$. Example: Virasoro, $W_{1+\infty}$.
\end{itemize}
This is the derived-deformation-theoretic avatar of the shadow tower: the $L_\infty$ structure maps $l_k$ of the bar complex become the higher MC equation terms, and the shadow invariants $S_k$ (Vol~I) are the coefficients of the $k$-th correction $\delta^{(k)}$ (see Vol~I~\cite{VolI}).
\end{remark}

\section{The Costello 5d construction}
\label{sec:costello-5d-construction}

The $d = 5$ regime of holomorphic Chern--Simons theory
(Theorem~\ref{thm:5d-boundary-yangian}) is the middle step of the
\(3\to5\to6\) hierarchy, and the only one where the full construction is
proved: the boundary chiral algebra is the affine Yangian
\(Y(\widehat{\fgl}_1)\). The five ingredients below are the algebraic
faces of Costello's construction (arXiv:1303.2632).

\subsection{The 5d theory and its boundary}
\label{subsec:5d-theory}

\begin{proposition}[5d holomorphic CS specification]
\label{prop:5d-hcs-specification}
The data of 5d holomorphic Chern--Simons theory on $\C^2_{h_1,h_2} \times \R$ with gauge algebra $\fgl_1$ is:
\begin{enumerate}[label=\textup{(\alph*)}]
 \item \emph{Action}: $S = \frac{1}{2}\int \Omega \wedge A \wedge \bar{\partial}A$, where $\Omega = dz_1 \wedge dz_2$ is the holomorphic $2$-form on $\C^2$. The cubic term vanishes for abelian $\fgl_1$.
 \item \emph{Omega-background}: the $T^2$-action $(z_1, z_2) \mapsto (e^{ih_1 t}z_1, e^{ih_2 t}z_2)$ with equivariant parameters $(h_1, h_2)$.
 \item \emph{CY condition}: the one-loop anomaly cancellation of the theory on $\C^2 \times \R$ requires $h_1 + h_2 + h_3 = 0$, where $h_3 = -(h_1 + h_2)$ is the equivariant weight of the $\R$-direction. This is the Calabi--Yau condition on $\C^3 = \C^2 \times \C_R$.
 \item \emph{Boundary algebra}: the observables on the boundary $\{R = 0\}$, projected to a holomorphic curve $C \subset \C^2$, form the affine Yangian $Y(\widehat{\fgl}_1)$ with structure function $g(u) = \prod_{i=1}^{3}(u - h_i)/(u + h_i)$.
\end{enumerate}
\end{proposition}

\subsection{Tree-level OPE and Yangian relations}
\label{subsec:tree-level-ope}

\begin{proposition}[Tree-level structure function]
\label{prop:tree-level-g}
The structure function $g(u) = \prod_{i=1}^{3}(u-h_i)/(u+h_i)$, arising from tree-level Feynman diagrams of the 5d theory, satisfies:
\begin{enumerate}[label=\textup{(\roman*)}]
 \item \emph{Reflection identity}: $g(u)\,g(-u)=1$. Indeed
 \[
 g(-u)=\prod_{i=1}^{3}\frac{-u-h_i}{-u+h_i}
      =\prod_{i=1}^{3}\frac{u+h_i}{u-h_i}=g(u)^{-1}.
 \]
 This ensures the antisymmetry of the $ee$ relation $g(\Delta)\, e(z)\, e(w) = -g(-\Delta)\, e(w)\, e(z)$.
 \item \emph{Coefficient vanishing}: expand $g(z) = \sum_{j \geq 0} \phi_j z^{-j}$. Then:
 \[
 \phi_0 = 1,\quad \phi_1 = 0,\quad \phi_2 = 0,\quad \phi_3 = -2\sigma_3,\quad \phi_4 = 0.
 \]
 The vanishing $\phi_1 = 0$ is equivalent to $h_1 + h_2 + h_3 = 0$ (the CY condition). The vanishing of all even $\phi_{2k}$ for $2k < 6$ follows from parity: $\log g(z) = \sum_{k \text{ odd}} \alpha_k z^{-k}$, so $\phi_{2k} = 0$ unless $2k$ admits a partition into odd parts $\geq 3$. The first nonzero even coefficient is $\phi_6 = \alpha_3^2/2$, from the partition $6 = 3 + 3$.
 \item \emph{Classical $r$-matrix}: the $R$-matrix on $\mathrm{Fock}_2$ has leading behaviour $R(u) = 1 - \sigma_2/u + O(1/u^2)$ as $u \to \infty$. The classical $r$-matrix is $r(u) = -\sigma_2/u$, with residue $-\sigma_2$ equal to the Kac--Moody level $\Psi = -\sigma_2$.
\end{enumerate}
\end{proposition}

\begin{proof}
Part~(i): the reflection identity $g(u) g(-u) = 1$ is an algebraic consequence of the rational function form. Explicitly, $g(-u) = \prod (-u-h_i)/(-u+h_i) = \prod (u+h_i)/(u-h_i) = 1/g(u)$, with each factor applying the identity $(-u-h)/(h-u) = (u+h)/(u-h)$. This holds for \emph{any} $(h_1,h_2,h_3)$, including the CY locus.

Part~(ii): $\log g(z) = \sum_a [\log(1-h_a/z) - \log(1+h_a/z)] = -2\sum_{k \text{ odd}} p_k/(k z^k)$, where $p_k = h_1^k + h_2^k + h_3^k$. By the CY condition, $p_1 = 0$, giving $\phi_1 = \alpha_1 = 0$. Even $\alpha_k$ vanish identically from the $\log(1-x) - \log(1+x)$ expansion. The leading nontrivial coefficient is $\phi_3 = \alpha_3 = -2p_3/3 = -2\sigma_3$ (using Newton's identity $p_3 = 3\sigma_3$ under the CY condition $\sigma_1 = 0$). The $\phi_j$ for $j \geq 6$ even are generically nonzero, arising from products of odd $\alpha_k$'s.

Part~(iii) is the classical limit computed in
Theorem~\ref{thm:yang-r-matrix-ybe}.  At $(h_1,h_2,h_3)=(1,-2,1)$ it
gives $-\sigma_2=3$; at $(1,-3,2)$ it gives $-\sigma_2=7$.
\end{proof}

\subsection{Tree-level OPE at spins \texorpdfstring{$1$}{1} and \texorpdfstring{$2$}{2}}
\label{subsec:tree-level-spin12}

\begin{proposition}[Spin-\texorpdfstring{$1$}{1} and spin-\texorpdfstring{$2$}{2} OPE from the 5d theory]
\label{prop:5d-spin12-ope}
The boundary algebra of 5d holomorphic CS with $\fgl_1$ gauge algebra has:
\begin{enumerate}[label=\textup{(\roman*)}]
 \item \emph{Spin-$1$ OPE}: $J(z) J(w) = \Psi/(z-w)^2$, where $\Psi = -\sigma_2 = -(h_1 h_2 + h_1 h_3 + h_2 h_3)$ is the Kac--Moody level of the 5d theory. In lambda-bracket form: $\{J_\lambda J\} = \Psi \cdot \lambda$. This matches Proposition~\ref{prop:codim2-defect-ope}(a).
 \item \emph{Spin-$2$ OPE}: the Sugawara construction $T = \mathopen{:}JJ\mathclose{:}/(2\Psi)$ gives a Virasoro with central charge $c = 1$ (for $\fgl_1$, $c = \dim(\fgl_1) \cdot \Psi/(\Psi + h^\vee) = 1$, independent of $\Psi$). The OPE is $T(z)T(w) = \frac{1/2}{(z-w)^4} + \frac{2T(w)}{(z-w)^2} + \frac{\partial T(w)}{z-w}$. In lambda-bracket form: $\{T_\lambda T\} = \frac{1}{12}\lambda^3 + 2T\lambda + \partial T$. This matches Proposition~\ref{prop:codim2-defect-ope}(b).
\end{enumerate}
\end{proposition}

\begin{proof}
This is the Sugawara construction, as in
Proposition~\ref{prop:codim2-defect-ope}.
\end{proof}

\subsection{One-loop correction: the CY constraint}
\label{subsec:one-loop-cy}

\begin{proposition}[One-loop enforcement of the CY condition]
\label{prop:one-loop-cy-constraint}
The one-loop Feynman diagram of 5d holomorphic CS on $\C^2 \times \R$ produces a correction to the boundary OPE that is proportional to $h_1 + h_2 + h_3$. Specifically:
\begin{enumerate}[label=\textup{(\roman*)}]
 \item The structure function coefficient $\phi_1 = -2(h_1 + h_2 + h_3)$. The CY condition $h_1 + h_2 + h_3 = 0$ is equivalent to $\phi_1 = 0$.
 \item When $\phi_1 \neq 0$, the Serre relation of the would-be Yangian acquires a correction of order $\phi_1 \cdot \phi_3$, which breaks the Yangian structure.
 \item The anomaly cancellation mechanism: the total equivariant weight of the holomorphic $3$-form on $\C^3 = \C^2 \times \C_R$ is $\mathrm{wt}(\Omega_3) = h_1 + h_2 + h_3$. Setting this to zero ensures $\Omega_3$ is $T^3$-invariant, which is the CY condition $\det(T^3|_{\Omega^{3,0}}) = 1$.
 \item The reflection identity $g(u) g(-u) = 1$ holds for ALL $(h_1,h_2,h_3)$ (it is an algebraic property of the rational function, not a consequence of the CY condition). The CY condition enters through $\phi_1$, which controls the Serre relations.
\end{enumerate}
\end{proposition}

\begin{proof}
Part~(i): \(\phi_1=\alpha_1=-2p_1=-2(h_1+h_2+h_3)\) from the
log-series expansion. Part~(ii): the Serre correction
\(\delta_{\mathrm{Serre}}=\phi_1\phi_3/3\) vanishes at CY
(\(\phi_1=0\)) and is nonzero away from it. Parts~(iii)--(iv) are
Costello 2013, Proposition~3.2, together with the displayed algebraic
consequences.
\end{proof}

\subsection{The Omega-background deformation}
\label{subsec:omega-deformation}

\begin{proposition}[Deformation hierarchy from the Omega-background]
\label{prop:omega-deformation}
The Omega-background $(h_1, h_2)$ on $\C^2$ (with $h_3 = -(h_1+h_2)$) deforms the boundary algebra through a single effective parameter $\sigma_3 = h_1 h_2 h_3$:
\begin{enumerate}[label=\textup{(\roman*)}]
 \item At $\sigma_3 = 0$ (the self-dual point, where one $h_i = 0$): $g(u) = 1$ identically, the $R$-matrix is trivial, and $Y(\widehat{\fgl}_1)$ degenerates to the Heisenberg algebra $H_1$. The central charge remains $c = 1$.
 \item At generic $\sigma_3 \neq 0$: the structure function $g(u)$ is a nontrivial rational function with poles at $u = -h_i$ and zeros at $u = h_i$. The $R$-matrix on $\mathrm{Fock}_1$ is $R^{(1)}(u) = g(u)$ (the scalar structure function). On $\mathrm{Fock}_2 = \mathrm{span}\{|\mathrm{row}\rangle, |\mathrm{col}\rangle\}$, the diagonal entries of the $R$-matrix are:
 \[
 R_{(\mathrm{row},\mathrm{row})}(u) = g(u) \cdot g(u + h_1), \qquad
 R_{(\mathrm{col},\mathrm{col})}(u) = g(u) \cdot g(u + h_2),
 \]
 as products over box contents $c(s)$ of the partitions $(2)$ and $(1,1)$ respectively.
 \item The charge-$1$ unitarity $g(u) g(-u) = 1$ holds at all parameter values. The charge-$2$ unitarity $R_{12}(u) R_{21}(-u) = \mathrm{id}$ is a matrix identity involving off-diagonal entries (Theorem~\ref{thm:yang-r-matrix-ybe}).
 \item The invariant $\sigma_2 = h_1 h_2 + h_1 h_3 + h_2 h_3$ generates a rescaling, not a true deformation: the Yangian at parameters $(h_1,h_2,h_3)$ is isomorphic to the one at $(\lambda h_1, \lambda h_2, \lambda h_3)$ for $\lambda \neq 0$.
\end{enumerate}
\end{proposition}

\begin{proof}
Part~(i): at \(\sigma_3=0\), say \(h_2=0\), then \(h_3=-h_1\) and
\[
g(u)=\frac{(u-h_1)u(u+h_1)}{(u+h_1)u(u-h_1)}=1 .
\]
Parts~(ii)--(iv) are the charge-\(1\) and charge-\(2\) product
identities, with the charge-\(2\) case matching the
Maulik--Okounkov diagonal eigenvalue formula
(Theorem~\ref{thm:mo-chiral-rmatrix}).
\end{proof}

\subsection{The Koszul dual: boundary-to-bulk}
\label{subsec:5d-koszul-dual}

\begin{proposition}[Koszul dual of the affine Yangian from 5d CS]
\label{prop:5d-koszul-dual}
The Koszul dual $A^! = D_{\Ran}(B(Y(\widehat{\fgl}_1)))$ of the 5d boundary algebra satisfies:
\begin{enumerate}[label=\textup{(\roman*)}]
 \item \emph{Parameter inversion}: $A^!$ lives at the inverted Omega-background $(-h_1, -h_2, -h_3)$. Its structure function is $g^!(u) = g(-u) = 1/g(u)$.
 \item \emph{$\kappa_{\mathrm{ch}}$ complementarity}: $\kappa_{\mathrm{ch}}(A) + \kappa_{\mathrm{ch}}(A^!) = \rho_K$, where $\kappa_{\mathrm{ch}}(A) = -\sigma_2$, $\kappa_{\mathrm{ch}}(A^!) = \sigma_2$, and $\rho_K = 0$ for $\fgl_1$ (class $\mathbf{L}$, $h^\vee = 0$). The Koszul conductor $\rho_K$ is family-dependent: $\rho_K = 0$ for free-field/$\fgl_1$ and class $\mathbf{L}$; $\rho_K = 13$ for the Virasoro family (Vol~I).
 \item \emph{$\sigma$-invariant behaviour}: under parameter inversion, $\sigma_2$ is preserved ($\sigma_2(-h) = \sigma_2(h)$) and $\sigma_3$ is negated ($\sigma_3(-h) = -\sigma_3(h)$).
 \item \emph{Shadow class preservation}: $A^!$ is also class $\mathbf{L}$ (shadow depth $3$). The pole structure of $g^!(u) = 1/g(u)$ has the same simple poles (at $u = h_i$ instead of $u = -h_i$), preserving the shadow depth.
 \item \emph{Reflected level}: the current presentation of $A^!$ has Kac--Moody level $k' = -k - 2h^\vee = \sigma_2$ (for $\fgl_1$, $h^\vee = 0$, so $k' = -k$). At $n = 1$ (3d CS), this specializes to the Feigin--Frenkel reflected-current surface of the CE/bar-cochain Koszul target.
\end{enumerate}
\end{proposition}

\begin{proof}
The parameter inversion $g^!(u)=1/g(u)$ follows from the reflection
identity: $g^!(u)=g(-u)$, since evaluating $g$ at the inverted
parameters is equivalent to evaluating the original $g$ at $-u$, and
$g(-u)=1/g(u)$.  The $\kappa_{\mathrm{ch}}$ complementarity is the
identity
$\kappa_{\mathrm{ch}}(A)+\kappa_{\mathrm{ch}}(A^!)=0$ on the free-field
branch.
\end{proof}

\begin{remark}[Costello 5d structural checks]
\label{rem:costello-5d-summary}
The construction is controlled by the CY condition, the invariants
$\sigma_2,\sigma_3$, the reflection identity $g(u)g(-u)=1$, the
coefficient pattern $\phi_3=-2\sigma_3$, the spin-$1$ level
$\Psi=-\sigma_2$, the spin-$2$ Virasoro coefficient
$T_{(3)}T=1/2$, the one-loop Serre correction, the
$\Omega$-background deformation, the charge-$2$ diagonal $R$-matrix,
and the Koszul parameter inversion.
\end{remark}


%% The K3 holomorphic CS hierarchy: 3d -> 5d -> 6d

\section{The holomorphic CS hierarchy for K3}
\label{sec:hcs-k3-hierarchy}

The flat-space hierarchy of Section~\ref{sec:qca-from-qg} (Construction~\ref{constr:hol-cs-hierarchy}) has a K3 specialisation in which the ambient spaces $\C^n$ are replaced by K3-fibered geometries. The three levels of the hierarchy, their boundary algebras, and the dimensional reductions connecting them are the subject of this section.

\begin{conjecture}[K3 holomorphic CS hierarchy]
\label{conj:hcs-k3-hierarchy}
\ClaimStatusConjectured
The holomorphic Chern--Simons hierarchy specialised to K3 consists of three levels, with boundary algebras and dimensional reductions as follows:
\begin{center}
\small
\renewcommand{\arraystretch}{1.3}
\begin{tabular}{lllll}
\toprule
\textbf{Level} & \textbf{Geometry} & \textbf{Boundary algebra} & $E_n$ & \textbf{Construction locus} \\
\midrule
3d & $K3$ sigma model on $C$ & $H_{\mathrm{Muk}}$ (Mukai Heisenberg, rank $24$) & $\Eone$ & CY-A$_2$ and Mukai lattice \\
5d & $K3 \times C \times \R$ & $Y(\fg_{K3})$ (K3 Yangian) & $\Etwo$ & requires 5d K3 hCS construction \\
6d & $K3 \times E \times C$ & $U_{q,t}(\fg_{K3}^{\mathrm{tor}})$ (K3 quantum toroidal) & $E_3$ & requires 6d K3$\times E$ hCS construction \\
\bottomrule
\end{tabular}
\end{center}
The dimensional reductions are:
\begin{enumerate}[label=\textup{(\roman*)}]
 \item \emph{6d $\to$ 5d: shrink $E$}. As $\tau \to i\infty$ (equivalently $q = e^{2\pi i \tau} \to 0$), the elliptic curve $E$ degenerates and the quantum toroidal algebra reduces to the Yangian:
 \[
 U_{q,t}(\fg_{K3}^{\mathrm{tor}}) \;\xrightarrow{q \to 1}\; Y(\fg_{K3}).
 \]
 The trigonometric structure function $G_{K3}(x; \{q_a\})$ reduces to the rational structure function $g_{K3}(u; \{h_a\})$ via $x = e^u$, $q_a = e^{h_a}$.

 \item \emph{5d $\to$ 3d: shrink $C$}. As all Omega-background parameters $h_a \to 0$ (equivalently $\hbar \to 0$), the Yangian deformation disappears:
 \[
 Y(\fg_{K3}) \;\xrightarrow{h_a \to 0}\; U(\widehat{\fg}_{K3}) \;\simeq\; H_{\mathrm{Muk}}.
 \]
 The rational structure function $g_{K3}(u) = \prod_{a=1}^{24} (u - h_a)/(u + h_a) \to 1$ (trivial $R$-matrix), and the Yangian degenerates to the current algebra of the Mukai Heisenberg.
\end{enumerate}
The 3d level is established by CY-A$_2$ (Theorem~CY-A, $d = 2$). CY-A$_3$ supplies the $d = 3$ framed object-level assignment on its H1--H4 locus (Theorem~\ref{thm:derived-framing-m3b2} and Theorem~\ref{thm:cy-to-chiral-d3}); it does not produce a global $\Phi_3$ functor on all smooth proper CY$_3$ categories. The 5d and 6d levels remain conjectural because they depend on the constructions of the K3 Yangian (Conjecture~\ref{conj:k3-yangian}, Section~\ref{sec:rtt-ope-dictionary}) and K3 quantum toroidal algebra (Conjecture~\ref{conj:6d-k3xe}) respectively, which require the 5d/6d holomorphic CS frameworks (not merely CY-A$_3$).
\end{conjecture}

\subsection{The Omega-background at each level}
\label{subsec:k3-omega-background}

The Omega-background parameters control the deformation from classical to quantum at each level.

\begin{construction}[K3 Omega-background hierarchy]
\label{constr:k3-omega-hierarchy}
The Omega-background data at each level of the K3 hierarchy:
\begin{enumerate}[label=\textup{(\roman*)}]
 \item \emph{3d ($K3$ sigma model)}: no Omega-background. The boundary algebra $H_{\mathrm{Muk}}$ is classical (free field, $\hbar = 0$). The $R$-matrix is trivial and $\kappa_{\mathrm{ch}}(H_{\mathrm{Muk}}) = 2$ (from the CY-A$_2$ construction; $\kappa_{\mathrm{ch}} = \chi^{\mathrm{CY}}_2(K3) = 2$).

 \item \emph{5d ($K3 \times C \times \R$)}: Omega-background given by $24$ parameters $(h_1, \ldots, h_{24})$ in the complexified Mukai lattice, subject to the CY$_2$ constraint
 \[
 \sum_{a=1}^{24} h_a = 0.
 \]
 This leaves $23$ free parameters (the Mukai lattice $\widetilde{\Lambda}_{K3}$ has rank $24$, minus one constraint). The effective deformation parameter is the Newton power sum $p_3 = \sum h_a^3$; when $p_3 = 0$, the structure function degenerates to $g_{K3}(u) = 1$.

 For the flat-space $\C^3$, the $24$ parameters specialise to $(h_1, h_2, h_3, 0, \ldots, 0)$ with $h_1 + h_2 + h_3 = 0$, recovering the standard Omega-background of Section~\ref{sec:costello-5d-construction}.

 \item \emph{6d ($K3 \times E \times C$)}: the $24$ Mukai parameters $(h_1, \ldots, h_{24})$ are supplemented by the modular parameter $\tau$ of the elliptic curve $E$, giving $24$ free parameters in total. The nome $q = e^{2\pi i \tau}$ satisfies $|q| < 1$ since $\Imag(\tau) > 0$, ensuring convergence of all $q$-expansions. The multiplicative parameters are $q_a = e^{h_a}$, and the CY$_2$ constraint becomes
 \[
 \prod_{a=1}^{24} q_a = 1.
 \]
\end{enumerate}
\end{construction}

\subsection{Boundary conditions and the bulk-defect structure}
\label{subsec:k3-boundary-defect}

At each level, the theory on a manifold of the form $M \times \R$ or $M \times S^1$ has two types of boundary conditions: Dirichlet (fixing boundary values) and Neumann (fixing normal derivatives). These produce different algebras.

\begin{construction}[K3 boundary and defect algebras]
\label{constr:k3-boundary-defect}
At each level of the K3 hierarchy:
\begin{center}
\small
\renewcommand{\arraystretch}{1.3}
\begin{tabular}{llll}
\toprule
\textbf{Level} & \textbf{Boundary (Dirichlet)} & \textbf{Bulk (derived center)} & \textbf{Defect (Koszul dual)} \\
\midrule
3d & $H_{\mathrm{Muk}}$ & $Z^{\mathrm{der}}_{\mathrm{ch}}(H_{\mathrm{Muk}})$ & $H_{\mathrm{Muk}}^!$ \\
5d & $Y(\fg_{K3})$ & $Z^{\mathrm{der}}_{\mathrm{ch}}(Y(\fg_{K3}))$ & $Y(\fg_{K3})^!$ \\
6d & $U_{q,t}(\fg_{K3}^{\mathrm{tor}})$ & $Z^{\mathrm{der}}_{\mathrm{ch}}(U_{q,t})$ & $U_{q,t}^!$ \\
\bottomrule
\end{tabular}
\end{center}
The bulk-boundary coupling at each level is the canonical map $\mu_{\mathrm{bb}} \colon Z^{\mathrm{der}}_{\mathrm{ch}}(A) \otimes A^! \to A^!$ of Construction~\ref{constr:universal-defect}, making the defect algebra $A^!$ a module over the bulk. The defect algebra at the 5d level has inverted parameters: $Y(\fg_{K3})^!$ lives at $(-h_1, \ldots, -h_{24})$ with structure function $g_{K3}^!(u) = g_{K3}(-u) = 1/g_{K3}(u)$ (Proposition~\ref{prop:5d-koszul-dual}, generalised from $\C^3$ to K3). The $\kappa_{\mathrm{ch}}$-complementarity gives $\kappa_{\mathrm{ch}}(A) + \kappa_{\mathrm{ch}}(A^!) = 0$ for the free-field/$\fgl_1$ branch ($\rho_K = 0$; Koszul conductor vanishes for class $\mathbf{L}$).
\end{construction}

\subsection{Wilson lines and coproducts}
\label{subsec:k3-wilson-lines}

The Wilson lines (defects) at each level produce modules over the boundary algebra. Their fusion is controlled by the coproduct.

\begin{construction}[K3 Wilson line hierarchy]
\label{constr:k3-wilson-hierarchy}
The Wilson lines at each level of the K3 hierarchy:
\begin{enumerate}[label=\textup{(\roman*)}]
 \item \emph{3d: point operators on $C$}. Indexed by Mukai vectors $v \in H^*(K3, \Z)$, these produce Fock modules $\cF_v$ of $H_{\mathrm{Muk}}$. The coproduct is \emph{primitive}:
 \[
 \Delta_0(J_{a,n}) = J_{a,n}^L + J_{a,n}^R,
 \]
 as established in Proposition~\ref{prop:heisenberg-primitive}.

 \item \emph{5d: line operators on $C$}. The line operators wrapping $C$ in $K3 \times C \times \R$ are indexed by Fock modules $\cF_\lambda$ of $Y(\fg_{K3})$. The coproduct is the \emph{Miura factorisation}:
 \[
 \Delta_z(T(u)) = T^L(u) \cdot T^R(u - z),
 \]
 where $z$ is the spectral parameter and $T(u) = 1 + \sum_{s \geq 1} \psi_s u^{-s}$ is the transfer matrix. Fusion of Wilson lines is controlled by the $R$-matrix $R_{K3}(u)$ (Conjecture~\ref{conj:rtt-ope-dictionary}).

 \item \emph{6d: surface operators wrapping $E$}. The Wilson surfaces wrapping $E$ in $K3 \times E \times C$ are indexed by BPS states labelled by Mukai vectors together with $E$-winding numbers. The coproduct is a $q$-deformed Miura factorisation:
 \[
 \Delta_x^{(q)}(\Psi(w)) = \Psi^L(w) \cdot \Psi^R(w/x),
 \]
 where $x = e^z$ is the multiplicative spectral parameter. This is conjectural.
\end{enumerate}
\end{construction}

\subsection{The structure function at each level}
\label{subsec:k3-structure-fn}

The structure function encodes the $R$-matrix data and determines the Yangian/quantum group relations. Its type changes across the hierarchy.

\begin{proposition}[K3 structure function hierarchy]
\label{prop:k3-structure-fn-hierarchy}
The structure function at each level of the K3 hierarchy:
\begin{enumerate}[label=\textup{(\roman*)}]
 \item \emph{3d (classical)}: $g_{K3}^{(3d)}(u) = 1$ identically (trivial $R$-matrix).

 \item \emph{5d (rational)}: $g_{K3}^{(5d)}(u) = \prod_{a=1}^{24} \frac{u - h_a}{u + h_a}$, a degree-$(24,24)$ rational function satisfying:
 \begin{itemize}
 \item Reflection: $g_{K3}(u) \cdot g_{K3}(-u) = 1$.
 \item CY$_2$ constraint: $\phi_1 = -2\sum h_a = 0$.
 \item Classical limit: $g_{K3}(u) \to 1$ as $u \to \infty$.
 \end{itemize}

 \item \emph{6d (trigonometric)}: $G_{K3}^{(6d)}(x) = \prod_{a=1}^{24} \frac{1 - q_a x}{1 - q_a^{-1} x}$, where $q_a = e^{h_a}$, satisfying:
 \begin{itemize}
 \item Inversion: $G_{K3}(x) \cdot G_{K3}(1/x) = 1$ (requires CY$_2$: $\prod q_a = 1$).
 \item Reduction: $G_{K3}(e^{\varepsilon u}; \{e^{\varepsilon h_a}\}) \to g_{K3}(u; \{h_a\})$ as $\varepsilon \to 0$.
 \item Classical limit: $G_{K3}(x) \to 1$ as all $q_a \to 1$.
 \end{itemize}
\end{enumerate}
Parts~(i) and~(ii) are direct substitutions. Part~(iii) is an
algebraic identity for the product formula; the inversion identity
requires the CY$_2$ constraint \(\prod q_a=1\). The \(6d\to5d\)
reduction is the \(\varepsilon\to0\) expansion of the same product.
\end{proposition}

\begin{proof}
Part~(i): at the 3d level, all Omega-background parameters are zero, so $g = \prod_{a=1}^{24} u/u = 1$.

Part~(ii): the reflection identity $g(u) g(-u) = 1$ follows from factorwise cancellation: each factor $(u - h_a)/(u + h_a)$ paired with $(-u - h_a)/(-u + h_a) = (u + h_a)/(u - h_a)$ gives product~$1$. The CY$_2$ vanishing of $\phi_1$ follows from the log-expansion $\log g(u) = -2 \sum_{k \text{ odd}} p_k / (k u^k)$, so $\phi_1 = -2 p_1 = -2 \sum h_a = 0$.

Part~(iii): the inversion identity $G(x) G(1/x) = \prod q_a^2 = (\prod q_a)^2 = 1$ (where the last equality uses $\prod q_a = 1$). The reduction is the standard passage from trigonometric to rational: expand $e^{\varepsilon h_a} = 1 + \varepsilon h_a + O(\varepsilon^2)$ and $e^{\varepsilon u} = 1 + \varepsilon u + O(\varepsilon^2)$; the leading terms in the product $\prod (1 - e^{\varepsilon h_a} e^{\varepsilon u})/(1 - e^{-\varepsilon h_a} e^{\varepsilon u})$ match $\prod (u - h_a)/(u + h_a)$ at $O(1)$ with \(O(\varepsilon^2)\) corrections. At fifteen generic points the relative error is bounded by \(6 \times 10^{-3}\) for \(\varepsilon = 0.01\).
\end{proof}

\subsection{The \texorpdfstring{$\kappa_\bullet$}{kappa-.}-spectrum across the hierarchy}
\label{subsec:k3-kappa-hierarchy}

\begin{remark}[\texorpdfstring{$\kappa_\bullet$}{kappa-.}-spectrum of the K3 hierarchy]
\label{rem:kappa-k3-hierarchy}
Each level of the K3 hierarchy contributes to the $\kappa_\bullet$-spectrum of $K3 \times E$:
\begin{center}
\small
\renewcommand{\arraystretch}{1.3}
\begin{tabular}{llrl}
\toprule
\textbf{Subscript} & \textbf{Origin} & \textbf{Value} & \textbf{Level} \\
\midrule
$\kappa_{\mathrm{ch}}^{\mathrm{Hodge}}(K3)$ & CY-A$_2$ correspondence $\Phi_2$ & $2$ & 3d (proved) \\
$\kappa_{\mathrm{ch}}^{\mathrm{Heis}}(K3 \times E)$ & Beauville-additive Heisenberg--Cartan: $2 + 1$ & $3$ & 5d/6d \\
$\kappa_{\mathrm{BKM}}(\Delta_5)$ & Igusa cusp form $\Delta_5$ weight & $5$ & 6d \\
$\kappa_{\mathrm{cat}}(K3)$ & $\chi(\cO_{K3}) = 2$ (fibre) & $2$ & 3d \\
$\kappa_{\mathrm{cat}}(K3 \times E)$ & K\"unneth: $\chi(\cO_{K3}) \cdot \chi(\cO_E) = 0$ & $0$ & 5d/6d \\
$\kappa_{\mathrm{fiber}}(K3)$ & Mukai lattice rank & $24$ & All levels \\
\bottomrule
\end{tabular}
\end{center}
The five distinct values $\{\kappa_{\mathrm{cat}}, \kappa_{\mathrm{ch}}^{\mathrm{Hodge}}, \kappa_{\mathrm{ch}}^{\mathrm{Heis}}, \kappa_{\mathrm{BKM}}(\Delta_5), \kappa_{\mathrm{fiber}}\} = \{0, 0, 3, 5, 24\}$ on $K3 \times E$ are the outputs of five distinct constructions: the Tier-(i) CY-datum intrinsic $\kappa_{\mathrm{cat}}(K3 \times E) = \chi(\cO_{K3 \times E}) = \chi(\cO_{K3}) \cdot \chi(\cO_E) = 2 \cdot 0 = 0$ (K\"unneth-multiplicative); the compact total-space Hodge supertrace $\kappa_{\mathrm{ch}}^{\mathrm{Hodge}}(K3 \times E) = 0$ by Serre at odd $d = 3$; the Stage-2 Heisenberg rank $\kappa_{\mathrm{ch}}^{\mathrm{Heis}} = 3$; the Borcherds weight $\kappa_{\mathrm{BKM}}(\Delta_5) = 5$; and the Mukai lattice rank $\kappa_{\mathrm{fiber}} = 24$.
\end{remark}

\subsection{The defect algebra DAG: Koszul duality construction verification}
\label{subsec:defect-koszul-dag}

The defect algebra (chiral quantum group) arises from the five-stage construction of Remark~\ref{rem:hcs-construction}, instantiated for each CY target as a directed acyclic graph (DAG) of computations. The five stages are:
\begin{center}
\small
\renewcommand{\arraystretch}{1.3}
\begin{tabular}{clll}
\toprule
\textbf{Stage} & \textbf{Input} & \textbf{Output} & \textbf{Constraint} \\
\midrule
1 & CY$_d$ category & Bulk algebra $A$ with $\kappa_{\mathrm{ch}}(A)$ & Shadow class G/L/C/M \\
2 & Bulk $A$ & Bar complex $B(A)$, differentials & $d_{\mathrm{bar}} = 0$ for class G \\
3 & Bar $B(A)$ & Defect $A^! = D_{\Ran}(B(A))$, $\kappa_{\mathrm{ch}}(A^!)$ & Structure function inversion \\
4 & Bulk + Defect & Conductor $K = \kappa_{\mathrm{ch}} + \kappa_{\mathrm{ch}}^!$ & $K = 0$ (free-field) or $K > 0$ (class M) \\
5 & Bar + Defect & Morita $\Omega(B(A)) \simeq A$ & Koszul locus check \\
\bottomrule
\end{tabular}
\end{center}

\begin{conjecture}[Defect algebra of \texorpdfstring{$K3 \times E$}{K3 x E} via Koszul duality]
\label{conj:defect-koszul-k3xe}
\ClaimStatusConjectured
Let $A_{K3 \times E}^{(K3,E)}$ denote the bulk chiral algebra on $E$ from the framed K3-fibre specialisation $\Phi_3^{(K3,E)}(D^b(\Coh(K3 \times E)))$ (Theorem~CY-A$_3$, Theorem~\ref{thm:cy-to-chiral-d3}; the 6d framework of Conjecture~\ref{conj:6d-k3xe} provides an independent route). Then the five-stage construction produces two branches:
\begin{enumerate}[label=\textup{(\roman*)}]
 \item \emph{Free-field/Heisenberg branch} (class~G, $\fgl_1$):
 $\kappa_{\mathrm{ch}}^{\mathrm{Heis}}(A) = 3$,
 $\kappa_{\mathrm{ch}}^{\mathrm{Heis}}(A^!) = -3$,
 conductor $K = 0$. The compact total-space algebra
 $A_{K3\times E}$ has $\kappa_{\mathrm{ch}}=0$; the value $3$
 is the Heisenberg/free-field specialization. The defect algebra
 has inverted structure function $g^!_{K3}(u) = 1/g_{K3}(u)$
 and reversed braiding
 $\Rep^{\Etwo}(A^!) \simeq \Rep^{\Etwo}(A)^{\mathrm{rev}}$.
 The bar-cobar adjunction recovers $A$ on the Koszul locus:
 $\Omega(B(A)) \simeq A$.
 \item \emph{Mukai-doubling branch} (class~M, Theorem-C row~$\mathsf{B}$):
 $\kappa_{\mathrm{ch}}^{\mathrm{Mukai}}(A) = 4$,
 $\kappa_{\mathrm{ch}}^{\mathrm{Mukai}}(A^!) = 4$,
 conductor $K^{\kappa_{\mathrm{ch}}^{\mathrm{Mukai}}} = 8 = 2\,c_+(\mathrm{Mukai}(K3))$
 (Theorem~\ref{thm:qca-hbar-minus-eighth-lusztig};
 Remark~\ref{rem:qca-three-faces-of-eight}). The defect algebra is
 a deformation of the chiral Koszul dual of the BKM vertex algebra
 $V_{\mathfrak{g}_{\Delta_5}}$, with infinite shadow depth. The
 Mukai integer \(8\) places this branch outside the free-field Koszul
 class and gives the order-eight comparison datum at the Lusztig
 specialisation \(\ell=8\).
\end{enumerate}
The two branches coexist: the Heisenberg branch captures the abelian/$\fgl_1$ free-field sector, while the Mukai-doubling branch captures the non-abelian/interacting BKM sector controlled by $V_{\mathfrak{g}_{\Delta_5}}$. The Borcherds weight $\kappa_{\mathrm{BKM}}(\Delta_5) = c_{\Delta_5}(0)/2 = 5$ is the universal level-$4$ automorphic invariant of the input Siegel form; it is structurally distinct from the Mukai integer \(8\). The two numbers \(5\) and \(8\) are separate diagnostics.

Conjecture~\ref{conj:defect-koszul-k3xe} lies over the framed $\Phi_3$
object and over the 6d construction of Conjecture~\ref{conj:6d-k3xe};
the defect branch inherits those hypotheses.
\end{conjecture}

\begin{proposition}[K3 defect algebra: the \texorpdfstring{$d = 2$}{d = 2} verified case]
\label{prop:defect-koszul-k3}
For $K3$ at $d = 2$ (CY-A$_2$ proved), the five-stage construction is fully verified:
\begin{enumerate}[label=\textup{(\roman*)}]
 \item Bulk: $A_{K3} = \mathcal H_{\mathrm{Muk}}$ on the free Mukai--Heisenberg branch; $\kappa_{\mathrm{ch}}^{\mathrm{Hodge}}(K3)=2$, $\kappa_{\mathrm{ch}}^{\mathrm{Heis}}(\mathcal H_{\mathrm{Muk}})=24$, class~G.
 \item Bar: $B(A_{K3})$ has $24$ generators (Mukai rank), $d_{\mathrm{bar}} = 0$ (free-field), no $A_\infty$ corrections.
 \item Defect: $A_{K3}^! = \mathcal H_{\mathrm{Muk}}^!$ with $\kappa_{\mathrm{ch}}^{\mathrm{Hodge}}(A^!) = -2$, $g^!(u) = 1/g_{K3}(u)$.
 \item Conductor: $K_{\mathrm{Hodge}} = 2 + (-2) = 0$ (free-field branch).
 \item Morita: $\Omega(B(A_{K3})) \simeq A_{K3}$ at $\Etwo$-level on the Koszul locus.
\end{enumerate}
The structure function inversion $g_{K3}^!(u) = 1/g_{K3}(u)$ follows from the algebraic unitarity identity $g_{K3}(u) \cdot g_{K3}(-u) = 1$ (Proposition~\ref{prop:k3-structure-fn-hierarchy}(ii)) together with the Verdier parameter inversion $h_a \mapsto -h_a$ under Koszul duality (Construction~\ref{constr:k3-boundary-defect}). Braiding reversal: $\Rep^{\Etwo}(A_{K3}^!) \simeq \Rep^{\Etwo}(A_{K3})^{\mathrm{rev}}$ (the sign flip $\kappa_{\mathrm{ch}}^{\mathrm{Hodge}} \to -\kappa_{\mathrm{ch}}^{\mathrm{Hodge}}$ inverts the leading $R$-matrix coefficient on this scalar line).
\end{proposition}

\begin{proof}
Parts~(i)--(iv): the Mukai--Heisenberg branch $\mathcal H_{\mathrm{Muk}}$ is the free-field output of $\Phi$ at $d = 2$ (Section~\ref{sec:cy-d2-functor}); its source Hodge scalar is $\kappa_{\mathrm{ch}}^{\mathrm{Hodge}}(K3) = \chi^{\mathrm{CY}}_2(K3) = 2$ (Proposition~\ref{prop:cy-kappa-d2}), and its Heisenberg scalar is the Mukai rank $24$. The branch is class~G because the Heisenberg bracket $J^{(a)}(z) J^{(b)}(w) \sim \omega^{ab}/(z-w)^2$ has no first-order pole, so $\mu = 0$ and $d_{\mathrm{bar}} = 0$. Koszul duality for free-field algebras gives $\kappa_{\mathrm{ch}}^{\mathrm{Hodge},!} = -\kappa_{\mathrm{ch}}^{\mathrm{Hodge}}$ with conductor $K_{\mathrm{Hodge}} = 0$ (matching $e1\_koszul\_three\_families.py$ for the Heisenberg family at rank~$24$).

Part~(v): for $A_{K3} = \mathcal H_{\mathrm{Muk}}$ on the chirally Koszul locus, the raw-direct-sum scope applies because $A_{K3}$ is class~G. The bar-cobar adjunction $\Omega \dashv B$ (cobar left, bar right) is a Quillen equivalence (Vol~I, Theorem~B, raw-direct-sum scope on classes G and L; see Corollary~\cite{VolI}), so $\Omega(B(A_{K3})) \simeq A_{K3}$ as $\Etwo$-algebras. The $\Etwo$-level follows from CY dimension $d = 2$ (Dunn additivity: $\Eone \otimes \Eone \simeq \Etwo$).

\end{proof}

\begin{remark}[Three dualities on K3 are distinct]
\label{rem:three-dualities-k3}
The K3 defect algebra $A_{K3}^!$ sits alongside three other K3-attached objects, each with its own role:
\begin{enumerate}[label=\textup{(\alph*)}]
 \item the \emph{mirror} K3 algebra $A_{K3^v}$ has $\kappa_{\mathrm{ch}}^{\mathrm{Hodge}} = 2$ (HMS preserves the Hodge scalar, so the mirror reads the same chiral characteristic);
 \item the \emph{derived center} $Z^{\mathrm{der}}_{\mathrm{ch}}(A_{K3})$ is the bulk algebra controlling the $49$-dimensional Drinfeld double;
 \item the \emph{Drinfeld center} $\cZ(\Rep^{\Eone}(A_{K3}))$ is the braided monoidal \emph{category} carrying the $\Etwo$ promotion.
\end{enumerate}
The $\kappa_{\mathrm{ch}}^{\mathrm{Hodge}}$-sign diagnostic separates Koszul duality from HMS: $\kappa_{\mathrm{ch}}^{\mathrm{Hodge}}(A_{K3}^!) = -2$ (Koszul dual flips the sign), while $\kappa_{\mathrm{ch}}^{\mathrm{Hodge}}(A_{K3^v}) = 2$ (mirror preserves the sign).
\end{remark}

\subsection{The master diagram}
\label{subsec:k3-master-diagram}

\begin{conjecture}[Master diagram of the K3 holomorphic CS hierarchy]
\label{conj:k3-master-diagram}
\ClaimStatusConjectured
The complete K3 holomorphic CS hierarchy connects three levels via dimensional reduction, with the structure function transitioning from trigonometric (6d) to rational (5d) to classical (3d):
\[
\begin{tikzcd}[column sep=1.2cm, row sep=1.8cm]
 K3 \times E \times C \ar[d, "\text{6d hol theory}"] \ar[r, "\text{shrink } E"]
 & K3 \times C \times \R \ar[d, "\text{5d hol CS}"] \ar[r, "\text{shrink } C"]
 & K3 \text{ on } C \ar[d, "\text{3d } \sigma\text{-model}"] \\
 U_{q,t}(\fg_{K3}^{\mathrm{tor}}) \ar[r, "q \to 1"]
 \ar[d, "B(-)"]
 & Y(\fg_{K3}) \ar[r, "h_a \to 0"]
 \ar[d, "B(-)"]
 & H_{\mathrm{Muk}} \ar[d, "B(-)"] \\
 B(U_{q,t}) \ar[d, "D_{\Ran}"]
 & B(Y(\fg_{K3})) \ar[d, "D_{\Ran}"]
 & B(H_{\mathrm{Muk}}) \ar[d, "D_{\Ran}"] \\
 U_{q,t}^! \ar[d, "\Rep^{E_3}"]
 & Y(\fg_{K3})^! \ar[d, "\Rep^{\Etwo}"]
 & H_{\mathrm{Muk}}^! \ar[d, "\Rep^{\Eone}"] \\
 \text{chiral QG}_{K3}^{(6d)}
 & \text{chiral QG}_{K3}^{(5d)}
 & \text{conformal blocks}_{K3}
\end{tikzcd}
\]
Each column is the holomorphic CS $\to$ chiral quantum group construction of Remark~\ref{rem:hcs-construction}, specialised to the K3 geometry at the given dimension. Each horizontal arrow is a dimensional reduction. The vertical arrows are:
\begin{itemize}
 \item $B(-)$: the ordered bar complex (chiral CE chains; Construction~\ref{constr:chiral-ce});
 \item $D_{\Ran}$: the Verdier dual on $\Ran(C)$ (Koszul duality; Construction~\ref{constr:universal-defect});
 \item $\Rep^{E_k}$: the braided representation category at $E_k$-level $k = n$ (the complex dimension).
\end{itemize}
The 3d column is established (CY-A$_2$ proved). The 5d column is partially established: the flat-space boundary algebra is proved (Theorem~\ref{thm:5d-boundary-yangian}), and the K3 specialisation is conjectural. The 6d column is conditional on the non-Lagrangian 6d algebraic framework (Remark~\ref{rem:6d-not-lagrangian}), the compact \(K3\times E\) hCS--Hall comparison, and the K3 quantum-toroidal construction. Theorem~\ref{thm:cy-to-chiral-d3} supplies the framed object-level \(\Phi_3\) input; the remaining condition is that dimensional reduction commute with bar, Koszul duality, and representation categories in this compact 6d geometry.
\end{conjecture}



%% Comparison: hCS route vs sigma model route

\subsection{Comparison: holomorphic CS vs.\ sigma model}
\label{subsec:hcs-vs-sigma-adversarial}
\index{holomorphic CS!vs sigma model}
\index{sigma model!vs holomorphic CS}

The K3 chiral algebra arises from three constructions on the same K3 geometry: the holomorphic CS/Costello--Li boundary algebra $A_E$ (Route~A), the CY-to-chiral correspondence output $H_{\mathrm{Muk}} = \Phi_2(D^b(\Coh(K3)))$ (Route~B), and the K3 sigma model VOA $V_{K3}$ (Route~C). Each is a distinct algebraisation of the same input, with its own central charge, shadow class, and operadic level; the four-vector comparison below identifies which invariant each construction reads.

\begin{proposition}[Route comparison across four vectors]
\label{prop:hcs-vs-sigma}
\index{modular characteristic!comparison across routes}
The following apparent contradictions between the holomorphic CS and sigma model approaches are resolved:
\begin{enumerate}[label=\textup{(\roman*)}]
 \item \emph{Central charge.} The Costello--Li KS boundary algebra $A_E$ (Construction~\textup{\ref{constr:costello-li-k3xe}}) has $c = 24$ ($24$~free bosons from $H^*(K3)$). The $\cN = 4$ sigma model VOA $V_{K3}$ has $c = 6$ (Definition~\textup{\ref{def:n4-sca-c6}}). These are \emph{different algebras}: $A_E$ is the B-model reduction (Hochschild data), $V_{K3}$ is the full SCFT. The CY-A$_2$ functor output $\Phi(D^b(\Coh(K3))) = H_{\mathrm{Muk}}$ coincides with $A_E$ as an algebra ($c = 24$), not with $V_{K3}$ ($c = 6$). The common number $2$ is the K3 Hodge scalar $\kappa_{\mathrm{ch}}^{\mathrm{Hodge}}(K3)=\chi(\cO_{K3})$, not the central charge of either algebra.

 \item \emph{Shadow class.} $H_{\mathrm{Muk}}$ is class~$\mathbf{G}$ (Gaussian, shadow depth $r = 2$): it is a free-field VOA with $m_k = 0$ for $k \geq 3$ (K3 is formal by DGMS). The $\cN = 4$ SCA $V_{K3}$ is class~$\mathbf{M}$ (infinite tower, $r_{\max} = \infty$): the stress tensor introduces non-trivial $A_\infty$ corrections with critical discriminant $\Delta = 8\kappa_{\mathrm{ch}}^{\mathrm{SCFT}} S_4 = 20/39 \neq 0$ (Computation~\textup{\ref{comp:shadow-tower-k3}}). Resolution: class~$\mathbf{G}$ is the shadow class of the $\Phi$-output $H_{\mathrm{Muk}}$; class~$\mathbf{M}$ is the shadow class of the physical SCFT $V_{K3}$. Same geometry, different algebraizations.

 \item \emph{$E_n$~structure.} K3 alone ($d = 2$) has native $\Etwo$ from the $\bS^2$-framing (Arrow~3 of the cyclic-to-chiral passage). $K3 \times E$ ($d = 3$) has native $\Eone$ from ordered factorization. Resolution: these apply to \emph{different CY dimensions}. The $\Etwo$ on $H_{\mathrm{Muk}}$ is automatic (free fields are $E_\infty$). The $\Eone$ on $Y(\fg_{K3})$ is the native level at $d = 3$; $\Etwo$ is derived via the Drinfeld center (Conjecture~\textup{\ref{conj:k3-master-diagram}}). The classical limit $h_a \to 0$ consistently sends $Y(\fg_{K3})$ ($\Eone$) to $H_{\mathrm{Muk}}$ ($\Etwo \hookrightarrow \Eone$, forgetful).

 \item \emph{Yangian.} The 5d holomorphic CS with Omega-background $(h_1, \ldots, h_{24})$ produces $Y(\fg_{K3})$; the sigma model is the classical $h_a \to 0$ limit, with image the abelian Heisenberg $H_{\mathrm{Muk}}$ on the Mukai lattice. The Yangian deformation is the Omega-background lift: $Y(\fg_{K3}) \xrightarrow{h_a \to 0} H_{\mathrm{Muk}}$ sends $g_{K3}(u) \to 1$ (trivial $R$-matrix) and $\Delta_z \to \Delta_0$ (primitive coproduct). The leading non-trivial term in $\log g_{K3}(u)$ is $O(h^3)$, controlled by $p_3 = \sum h_a^3$; the CY$_2$ constraint $\sum h_a = 0$ kills the $O(h)$ contribution.
\end{enumerate}
All four comparison vectors index distinct algebraisations of the same K3 geometry, with the central charge, Heisenberg/sigma-model split, $E_n$-level shift, and Omega-background deformation reading the K3 data through different construction layers. The comparison table:
\begin{center}
\small
\renewcommand{\arraystretch}{1.3}
\begin{tabular}{lccccl}
\toprule
\textbf{Route} & $c$ & \textbf{distinguished scalar} & \textbf{Shadow class} & $E_n$ & \textbf{Construction locus} \\
\midrule
$A_E$ (KS boundary) & $24$ & $\rho_{\mathrm{KS}}=24$ & G & $E_\infty$ & requires KS-boundary construction \\
$H_{\mathrm{Muk}} = \Phi(D^b(\Coh(K3)))$ & $24$ & $\kappa_{\mathrm{ch}}^{\mathrm{Hodge}}=2$; $\rho_{\mathrm{Muk}}=24$ & G & $\Etwo$ & CY-A$_2$ \\
$V_{K3}$ ($\cN = 4$ SCA) & $6$ & $\kappa_{\mathrm{ch}}^{\mathrm{SCFT}}=2$ & M & $\Etwo$ & $\cN=4$ SCA construction \\
$Y(\fg_{K3})$ (Yangian) & -- & $\kappa_{\mathrm{ch}}^{\mathrm{Heis}}=3$ & G & $\Eone$ & requires K3 Yangian construction \\
\bottomrule
\end{tabular}
\end{center}
The scalar entries $24$ and $2$ do not name one invariant. Route~A uses the OPE-level rank trace $\rho_{\mathrm{KS}}=\tr(\delta^{ab}) = 24$, while Routes~B and~C use respectively the K3 Hodge scalar and the normalised small-$\cN=4$ scalar, both equal to $\chi(\cO_{K3}) = 2$ on this surface. Route~D gives $\kappa_{\mathrm{ch}}^{\mathrm{Heis}} = 3$ by additivity: $\kappa_{\mathrm{ch}}^{\mathrm{Heis}}(K3 \times E) = 2 + 1$. The boundary-to-sigma ratio is $24/2 = 12$ (Proposition~\textup{\ref{prop:boundary-sigma-ratio}}).

\end{proposition}

\begin{proof}
Item~(i): $c(A_E) = 24$ by generator count (Construction~\ref{constr:costello-li-k3xe}); $c(V_{K3}) = 6$ by definition (Definition~\ref{def:n4-sca-c6}); $\kappa_{\mathrm{ch}}^{\mathrm{SCFT}}(V_{K3}) = 2$ by Proposition~\ref{prop:kappa-k3}; $\kappa_{\mathrm{ch}}^{\mathrm{Hodge}}(H_{\mathrm{Muk}}) = 2$ by Proposition~\ref{prop:cy-kappa-d2} (the Serre duality $\mathbb{S}_{\cC} = [2]$ kills the one-loop correction). Item~(ii): class~G for $H_{\mathrm{Muk}}$ by formality (DGMS) and shadow depth computation (all $S_r = 0$ for $r \geq 3$); class~M for $V_{K3}$ by Computation~\ref{comp:shadow-tower-k3} (the computed support packet is mixed and has scalar witness $\Delta = 20/39 \neq 0$). Item~(iii): $\Etwo$ at $d = 2$ from the $\bS^2$-framing (Theorem~\ref{thm:cy-to-chiral}, Arrow~3); $\Eone$ at $d = 3$ from ordered factorization; the forgetful functor $\Etwo \hookrightarrow \Eone$ gives compatibility. Item~(iv): the Yangian classical limit $g_{K3}(u; \{h_a\}) = \prod(u - h_a)/(u + h_a) \to 1$ as $h_a \to 0$ is verified directly; the CY$_2$ constraint $\sum h_a = 0$ kills the $O(h)$ term in $\log g_{K3}$; the leading deformation is $O(h^3)$ with coefficient $(-2/3) p_3 / u^3$ where $p_3 = \sum h_a^3$.
\end{proof}

\begin{remark}[Hodge decomposition and the generator discrepancy]
\label{rem:hcs-sigma-deepened}
The $24$~vs.~$8$~generator discrepancy between $H_{\mathrm{Muk}}$ and $V_{K3}$ is resolved by the HKR decomposition. The $24$~currents of $H_{\mathrm{Muk}}$ decompose as $h^{0,0} + h^{2,0} + h^{0,2} + h^{1,1} + h^{2,2} = 1 + 1 + 1 + 20 + 1 = 24$ (the full Hodge diamond of K3). The $8$~generators of $V_{K3}$ are $(T, G^\pm, \widetilde{G}^\pm, J^{++}, J^{--}, J^3)$ from the $\cN = 4$ superconformal structure. These count generators of \emph{different algebras}: the ratio $24/8 = 3 = \dim_\C(K3) + 1$.

The shadow tower of $V_{K3}$ exhibits sign alternation from depth~$4$ onward ($S_4 = 5/156 > 0$, $S_5 = -1/26 < 0$, $S_6 > 0$, \ldots) with critical discriminant $\Delta = 20/39$, confirming class~$\mathbf{M}$. All $S_r$ in the computed range $2 \leq r \leq 12$ are nonzero. In contrast, $H_{\mathrm{Muk}}$ has $S_r = 0$ for all $r \geq 3$ (class~$\mathbf{G}$, formal by DGMS).

Two further distinctions are used: (v)~The CoHA of a CY$_3$ category is an \emph{associative algebra} (Schiffmann--Vasserot--Yang--Zhao multiplication), \emph{not} a chiral algebra. The connection CoHA $\to$ $\Eone$-chiral algebra is via the CY-to-chiral correspondence $\{\Phi_d\}$ (CY-A), not by identification. (vi)~The $\Omega$-background $(h_1, \ldots, h_{24})$ on $\C$ in $K3 \times \C \times \R$ is the intermediary mechanism producing $Y(\fg_{K3})$ from 5d holomorphic CS; the normal-bundle grading $N_{C/Y}$ does \emph{not} directly identify with the $(q,t)$~parameters. The labelled scalars are $\kappa_{\mathrm{ch}}^{\mathrm{Hodge}}(K3)=2$ and $\kappa_{\mathrm{ch}}^{\mathrm{Heis}}(K3\times E)=3$ (Remark~\ref{rem:beauville-kappa-formula-subscript-split}).
\end{remark}

\begin{remark}[Four pairs of distinct objects in the hCS vs.\ sigma model comparison]
\label{rem:hcs-sigma-conflations}
The hCS / sigma-model comparison involves four pairs of distinct objects connected by named arrows; the arrow direction and intermediary in each case fixes which side of the pair is in play.
\begin{enumerate}[label=(\alph*)]
 \item \emph{CoHA versus $\Eone$-chiral algebra.} The CoHA carries an associative (Hall) product; the $\Eone$-chiral algebra is a factorization algebra on $\C \times \R$. The connecting arrow is the CY-to-chiral correspondence $\{\Phi_d\}$.
 \item \emph{$N_{C/Y}$ grading versus $(q,t)$ parameters.} The intermediary mechanism (equivariant localisation, Nekrasov partition function, refinement) supplies the dictionary between the two.
 \item \emph{Spectral $z$ versus worldsheet $z$.} The spectral $z$ enters the shift in $\Delta_z$; the worldsheet $z$ enters the OPE pole structure. Setting spectral $z = 0$ removes the shift; sending worldsheet $z \to w$ produces the OPE poles.
 \item \emph{OPE-level rank trace $24$ versus CY $\kappa_{\mathrm{ch}}^{\mathrm{Hodge}} = 2$.} The two are different invariants of the same K3: OPE trace $\rho_{\mathrm{OPE}}=\tr(\delta^{ab}) = 24$ counts free-field modes, CY supertrace $\chi(\cO_{K3}) = 2$ reads the Hodge data. The subscript names which.
\end{enumerate}
\end{remark}

%% RTT-OPE dictionary for the K3 Yangian

\section{The RTT--OPE dictionary for the K3 Yangian}
\label{sec:rtt-ope-dictionary}

The (conjectural) K3 Yangian $Y(\fg_{K3})$ for $\fg = \gl_1$ admits two
compatible presentation surfaces at the Heisenberg mode-algebra shadow,
not a proved equivalence of full chiral-Yangian definitions. The RTT
presentation below is weaker than the full chiral quadruple with modular
MC tower; the dictionary records the data visible at this abelian mode
level:
\begin{enumerate}[label=(\Alph*)]
 \item \textbf{RTT presentation}: generators are modes of the transfer matrix $T(u) = 1 + \sum_{s \geq 1} \psi_s u^{-s}$, with the defining relation
 \[
 R_{12}(u - v)\, T_1(u)\, T_2(v) \;=\; T_2(v)\, T_1(u)\, R_{12}(u - v).
 \]
 \item \textbf{OPE presentation}: generators are $24$ Heisenberg currents $J_i(z) = \sum_n J_{i,n}\, z^{-n-1}$, $i = 0, \ldots, 23$, with the OPE
 \[
 J_i(z)\, J_j(w) \;\sim\; \frac{\omega^{ij}}{(z - w)^2},
 \]
 where $\omega^{ij}$ is the Mukai pairing of signature~$(4, 20)$.
\end{enumerate}

\begin{remark}[Spectral parameter vs.\ worldsheet coordinate]
\label{rem:spectral-vs-worldsheet}
The variable $u$ in the RTT presentation is a \emph{Yangian spectral parameter}, living in the complexified weight space. The variable $z$ in the OPE presentation is a \emph{worldsheet coordinate} on the chiral Riemann surface. These are different mathematical objects:
\begin{itemize}
 \item Setting $u = 0$ in $\Delta_u$ yields the cocommutative coproduct $\Delta_0(T(v)) = T^L(v) \cdot T^R(v)$: no singularity.
 \item Setting $z \to w$ in $J_i(z)\,J_j(w)$ produces the OPE poles $\sim (z-w)^{-2}$, a physical collision singularity.
\end{itemize}
Both encode the \emph{same mode algebra} $[J_{i,m},\, J_{j,n}] = \omega^{ij}\, m\, \delta_{m+n,0}$; the bridge is through the modes.
\end{remark}

\subsection{The 24 Heisenberg currents and the Sugawara construction}
\label{subsec:k3-heisenberg-ope}

\begin{proposition}[Mukai Heisenberg OPE]
\label{prop:mukai-heisenberg-ope}
The $24$ currents $J_i(z)$ associated to the Mukai lattice $\widetilde{\Lambda}_{K3} = U^4 \oplus E_8(-1)^2$ satisfy the Heisenberg OPE
\begin{equation}\label{eq:mukai-heisenberg-ope}
 J_i(z)\, J_j(w) \;\sim\; \frac{\omega^{ij}}{(z - w)^2},
 \qquad i, j = 0, \ldots, 23,
\end{equation}
where $\omega^{ij} = \mathrm{diag}(\underbrace{+1, \ldots, +1}_{4},\, \underbrace{-1, \ldots, -1}_{20})$ in the diagonal basis. The corresponding $\lambda$-bracket is
\begin{equation}\label{eq:mukai-lambda-bracket}
 \{J_{i\,\lambda}\, J_j\} \;=\; \omega^{ij}\, \lambda.
\end{equation}
\end{proposition}

\begin{proof}
This is the standard Heisenberg OPE for $24$ free bosons with the Mukai pairing, as constructed in Section~\ref{subsec:k3-fact-tree} of Chapter~\ref{ch:k3-times-e}. The $\lambda$-bracket is the Fourier transform of the OPE: $\{a_\lambda\, b\} = \mathrm{Res}_{z=0}\, e^{\lambda z}\, a(z)\, b(0)$, giving $\mathrm{Res}_{z=0}\, e^{\lambda z}\, \omega^{ij}/z^2 = \omega^{ij}\, \lambda$.
\end{proof}

\begin{proposition}[Sugawara stress tensor at \texorpdfstring{$c = 24$}{c = 24}]
\label{prop:sugawara-c24}
The Sugawara construction
\begin{equation}\label{eq:sugawara-mukai}
 T(z) \;=\; \frac{1}{2}\sum_{i,j=0}^{23} \omega_{ij}\, {:}J_i(z)\, J_j(z){:}
\end{equation}
produces a Virasoro field with central charge $c = 24 = \kappa_{\mathrm{fiber}}$. Each free boson contributes $c_i = 1$ independent of the sign of $\omega^{ii}$, giving $c = \sum_{i=0}^{23} 1 = 24$. The self-OPE is
\begin{equation}\label{eq:tt-ope-c24}
 T(z)\, T(w) \;\sim\; \frac{12}{(z - w)^4} \;+\; \frac{2\, T(w)}{(z - w)^2} \;+\; \frac{\partial T(w)}{z - w}.
\end{equation}
\end{proposition}

\begin{proof}
The central charge of a rank-$1$ Heisenberg algebra at level $k$ ($[J_m, J_n] = k\, m\, \delta_{m+n,0}$) with Sugawara $T_n = (1/(2k)) \sum_a {:}J_a\, J_{n-a}{:}$ is $c = 1$ for any nonzero $k$ (the factor $k$ in the commutator cancels with the $1/k$ normalisation). For the rank-$24$ Mukai Heisenberg with $k_i = \omega^{ii} = \pm 1$, the total is $c = 24$. The fourth-order pole coefficient $c/2 = 12$ in~\eqref{eq:tt-ope-c24} confirms this.
\end{proof}

\begin{remark}[\texorpdfstring{$\kappa_\bullet$}{kappa-.}-spectrum]
The central charge $c = 24$ is $\kappa_{\mathrm{fiber}}$ (the Mukai lattice rank on the K3 fibre), distinct from $\kappa_{\mathrm{ch}}^{\mathrm{Hodge}}(K3) = 2$ (the K3 chiral algebraisation), $\kappa_{\mathrm{ch}}^{\mathrm{Heis}}(K3 \times E) = 3$ (Beauville-additive Heisenberg--Cartan rank: $2 + 1$), $\kappa_{\mathrm{BKM}}(\Delta_5) = 5$ (the Igusa weight), and $\kappa_{\mathrm{cat}}(K3) = \chi(\cO_{K3}) = 2$ (K3-fibre per-fibre value; $\kappa_{\mathrm{cat}}(K3 \times E) = 0$ by K\"unneth on the total space). The five $K3 \times E$ tower values are $\{\kappa_{\mathrm{cat}}, \kappa_{\mathrm{ch}}^{\mathrm{Hodge}}, \kappa_{\mathrm{ch}}^{\mathrm{Heis}}, \kappa_{\mathrm{BKM}}(\Delta_5), \kappa_{\mathrm{fiber}}\} = \{0, 0, 3, 5, 24\}$, with the first two reading the total-space Hodge layer.
\end{remark}

\subsection{The RTT--OPE bridge: classical \texorpdfstring{$r$}{r}-matrix = OPE coefficient}
\label{subsec:rtt-ope-bridge}

\begin{conjecture}[RTT--OPE dictionary for \texorpdfstring{$Y(\fg_{K3})$}{Y(g-K3)} at \texorpdfstring{$\gl_1$}{gl-1}]
\label{conj:rtt-ope-dictionary}
\ClaimStatusConjectured
The K3 Yangian $Y(\fg_{K3})$ for $\fg = \gl_1$ (Conjecture~\ref{conj:k3-yangian}) admits an RTT presentation with $R$-matrix
\begin{equation}\label{eq:rk3-diagonal}
 R_{K3}(u) \;=\; \mathrm{diag}\!\left(\frac{u - h_1}{u + h_1},\; \ldots,\; \frac{u - h_{24}}{u + h_{24}}\right),
\end{equation}
a degree-$(24,24)$ rational function with $24$ parameters $h_i$ satisfying $\sum h_i = 0$ (CY$_2$ constraint). The classical limit as $u \to \infty$ gives:
\begin{equation}\label{eq:classical-r-mukai}
 R_{K3}(u) \;=\; I_{24} \;-\; \frac{2}{u}\, \mathrm{diag}(h_1, \ldots, h_{24}) \;+\; O(u^{-2}).
\end{equation}
At unit deformation $h_i = \omega^{ii}$, the classical $r$-matrix is $r = -2\, \omega^{-1}$ (the inverse Mukai pairing, up to normalisation). This $r$-matrix is dual to the OPE second-order pole: the coefficient $\omega^{ij}$ in~\eqref{eq:mukai-heisenberg-ope} equals the leading $1/u$ coefficient of $R_{K3}(u)$.
\end{conjecture}

The RTT--OPE dictionary is summarised in Table~\ref{tab:rtt-ope}.

\begin{table}[ht]
\centering
\small
\renewcommand{\arraystretch}{1.3}
\begin{tabular}{lll}
\toprule
\textbf{Object} & \textbf{RTT side} & \textbf{OPE side} \\
\midrule
Generators & $T(u) = 1 + \sum \psi_s\, u^{-s}$ & $J_i(z) = \sum J_{i,n}\, z^{-n-1}$, $i=0,\ldots,23$ \\
Relations & $R_{12}(u{-}v)\, T_1 T_2 = T_2 T_1\, R_{12}$ & $J_i(z)\, J_j(w) \sim \omega^{ij}/(z{-}w)^2$ \\
$R$-matrix & $R(u) = \mathrm{diag}(\frac{u-h_i}{u+h_i})$ & OPE coefficient $\omega^{ij}$ \\
Stress tensor & Transfer matrix $T(u)$ & $T(z) = \tfrac{1}{2}\sum \omega_{ij}\,{:}J_i J_j{:}$, $c = 24$ \\
Coproduct & $\Delta_u(T(v)) = T^L(v) \cdot T^R(v{-}u)$ & $E_1$-factorisation on $\C \times \R$ \\
Variable & $u$ = spectral (weight space) & $z$ = worldsheet (Riemann surface) \\
Lie conformal & $Y(\fg_{K3})$ via RTT & $U^{\mathrm{ch}}(L_{\mathrm{Muk}}) = V_{H_{\mathrm{Muk}}}$ \\
Unitarity & $R(u)\, R(-u) = I$ & $J_i$ self-conjugate \\
Central charge & Encoded in quantum determinant & $c = 24$ from $(z{-}w)^{-4}$ pole \\
\bottomrule
\end{tabular}
\caption{RTT--OPE dictionary for the K3 Yangian at $\gl_1$. The bridge between the two presentations passes through the common mode algebra $[J_{i,m},\, J_{j,n}] = \omega^{ij}\, m\, \delta_{m+n,0}$.}
\label{tab:rtt-ope}
\end{table}

\subsection{The Mukai Lie conformal algebra and chiral envelope}
\label{subsec:mukai-lie-conformal}

\begin{proposition}[Chiral envelope of \texorpdfstring{$L_{\mathrm{Muk}}$}{L-Muk}]
\label{prop:chiral-envelope-mukai}
Let $L_{\mathrm{Muk}} = \bC[\partial] \otimes \bC^{24}$ be the rank-$24$ Lie conformal algebra with $\lambda$-bracket~\eqref{eq:mukai-lambda-bracket}. Then:
\begin{enumerate}[label=\textup{(\roman*)}]
 \item $L_{\mathrm{Muk}}$ is abelian (no structure constants beyond the pairing). Shadow class~$\mathbf{G}$ (free field, depth~$2$).
 \item The chiral (factorisation) envelope equals the rank-$24$ Heisenberg VOA:
 \begin{equation}\label{eq:chiral-envelope-mukai}
 U^{\mathrm{ch}}(L_{\mathrm{Muk}}) \;\simeq\; V_{H_{\mathrm{Muk}}}.
 \end{equation}
 \item Generator count: $24$ strong generators at conformal weight~$1$.
 \item Character:
 $\mathrm{ch}(V_{H_{\mathrm{Muk}}}) = q^{-1} \prod_{n \geq 1} 1/(1 - q^n)^{24} = q^{-1}/\eta(q)^{24}$, with $q^1$ coefficient $24$ and $q^2$ coefficient $324$.
 \item $V_{H_{\mathrm{Muk}}}$ is the \emph{classical limit} ($\sigma_3 = 0$) of $Y(\fg_{K3})$: the undeformed Heisenberg sector before the Yangian deformation is turned on.
\end{enumerate}
\end{proposition}

\begin{proof}
(i)~The $\lambda$-bracket $\{J_{i\,\lambda}\, J_j\} = \omega^{ij}\, \lambda$ is linear in $\lambda$ (no $\lambda^0$ or higher terms), so $L_{\mathrm{Muk}}$ is the Lie conformal algebra of a $2$-step nilpotent Lie algebra (class~$\mathbf{G}$). (ii)~For an abelian Lie conformal algebra, the factorisation envelope adds no new generators or relations beyond the PBW basis; the result is the free-field (Heisenberg) vertex algebra. (iii)~Generator count: $24 = \mathrm{rank}(\widetilde{\Lambda}_{K3})$. (iv)~The PBW basis $J_{i_1,-n_1} \cdots J_{i_k,-n_k} |0\rangle$ ($n_j > 0$, $24$ colours) gives the Fock space character $\prod_{n \geq 1} 1/(1-q^n)^{24}$. At $q^1$: $24$ oscillators. At $q^2$: $\binom{25}{2} + 24 = 300 + 24 = 324$. (v)~At $\sigma_3 = 0$ the structure function $g_{K3}(z) = 1$ identically (all factors cancel), the $R$-matrix is trivial, and $Y(\fg_{K3})$ degenerates to $H_{\mathrm{Muk}}$.
\end{proof}

\subsection{The coproduct in OPE language}
\label{subsec:coproduct-ope-language}

\begin{proposition}[Primitivity of the Heisenberg coproduct]
\label{prop:heisenberg-primitive}
The Heisenberg generators $J_{i,n}$ are \emph{primitive} in the $E_1$-chiral bialgebra:
\begin{equation}\label{eq:heisenberg-primitive}
 \Delta_0(J_{i,n}) \;=\; J_{i,n}^L \;+\; J_{i,n}^R.
\end{equation}
The coproduct preserves the Heisenberg commutator on the tensor product:
\[
 [\Delta_0(J_{i,m}),\, \Delta_0(J_{j,n})]
 \;=\; 2\, \omega^{ij}\, m\, \delta_{m+n,0}
\]
(the factor~$2$ arises from the two independent copies $H_L$ and $H_R$). The spectral shift $u$ appears only at the \emph{composite} level: the Sugawara coproduct $\Delta_u(T(v)) = T^L(v) \cdot T^R(v - u)$ acquires $u$-dependent cross-terms via the Miura factorisation (Section~\ref{sec:e1-chiral-bialgebras}).
\end{proposition}

\begin{proof}
Direct Fock space computation on the tensor product $\cF_L \otimes \cF_R$. The cross-commutator $[J_{i,m}^L, J_{j,n}^R] = 0$ (operators on different tensor factors), so $[\Delta_0(J_{i,m}), \Delta_0(J_{j,n})] = [J_{i,m}^L, J_{j,n}^L] + [J_{i,m}^R, J_{j,n}^R] = 2\, \omega^{ij}\, m\, \delta_{m+n,0}$.
\end{proof}

\begin{remark}[RTT--OPE dictionary: verification summary]
\label{rem:rtt-ope-summary}
The RTT--OPE dictionary is governed by the following nine identities:
\begin{center}
\small
\renewcommand{\arraystretch}{1.15}
\begin{tabular}{llr}
\toprule
\textbf{Block} & \textbf{Content} & \textbf{Identities} \\
\midrule
Heisenberg OPE & $\omega^{ij}$ coefficients, rank, signature, pole order & 7 \\
Sugawara & $c = 24$, conformal weights, $\kappa_\bullet$-spectrum & 8 \\
Fock space Virasoro & $c = r$ at ranks $1$--$4$, positive/negative pairings & 6 \\
RTT--OPE bridge & Classical $r$-matrix, $R(u) = I + \omega^{-1}/u$, proportionality & 7 \\
Spectral vs worldsheet & distinction, $\Delta_0$ no singularity & 4 \\
Lie conformal envelope & $L_{\mathrm{Muk}}$ abelian, $U^{\mathrm{ch}} = V_{H_{\mathrm{Muk}}}$, character & 8 \\
Coproduct in OPE & Primitivity $\Delta_0(J) = J^L + J^R$, commutator & 6 \\
Dictionary consistency & $9$ entries, cross-references & 4 \\
RTT--OPE domain & Conjectural bridge domain & 2 \\
\bottomrule
\end{tabular}
\end{center}
\end{remark}


\section{The CFG25 \texorpdfstring{$E_1$}{E-1}-chiral lift: from 3d CS to 6d holomorphic theory}
\label{sec:cfg25-e1-chiral-lift}

Costello--Francis--Gwilliam (2025) establish five structural results for 3d Chern--Simons theory on $\R^3$ using factorization homology. Each result has a precise $E_1$-chiral analogue in the 6d holomorphic theory on $\C^3$ of the Costello construction. This section makes each analogue explicit and identifies what works, what breaks, and what requires genuinely new ideas.

\begin{center}
\renewcommand{\arraystretch}{1.3}
\small
\begin{tabular}{lp{4.5cm}p{5.5cm}l}
 \toprule
 \textbf{CFG25 result (3d on $\R^3$)} & \textbf{$E_1$-chiral lift (6d on $\C^3$)} & \textbf{Obstruction / new idea} & \textbf{Construction locus} \\
 \midrule
 1. $\mathrm{Obs}^q(\R^3)$ is $E_3$-fact. & $E_3$-chiral fact.\ on $\C^3$ (holomorphic refinement $E_6 \to E_3$) & Omega-background provides nontrivial braiding & works \\
 2. Boundary $= V_k(\frakg)$ & Boundary $= Y(\widehat{\fgl}_1)$ & $E_\infty \to E_1$: Deligne level drops $E_3 \to E_2$ & works \\
 3. Link invariants & Holomorphic submanifold invariants & $\pi_1(\text{complement}) \to$ vanishing cycles; normal bundle $N_{\Sigma/\C^3}$ & partial \\
 4. $\int_{M \setminus L} = $ QG invariant & $\int_{\C^3 \setminus \Sigma} = $ quantum toroidal invariant & Holomorphic dependence; Ran space required for compact CY$_3$ & requires compact Ran construction \\
 5. $\mathrm{Obs}^q(\R^3)^! = U_q(\frakg)\text{-mod}$ & $\mathrm{Obs}^q(\C^3)^! = U_{q,t}(\widehat{\widehat{\fgl}}_1)\text{-mod}$ & ZTE failure (Theorem~\ref{thm:zte-failure}); even-sector ternary correction & requires full charge-preserving $E_3$ complex \\
 \bottomrule
\end{tabular}
\end{center}

\subsection{Result 1: \texorpdfstring{$E_n$}{E-n}-factorization structure}
\label{subsec:cfg25-result1}

\begin{proposition}[Holomorphic refinement of the \texorpdfstring{$E_n$}{E-n} level]
\label{prop:holomorphic-en-refinement}
The observables of holomorphic Chern--Simons theory on an $n$-dimensional complex manifold $M$ carry $E_{2n}$-factorization structure topologically (from $\dim_\R M = 2n$) and $E_n$-chiral factorization structure holomorphically (each complex direction contributes one chiral level). The relationship:
\begin{center}
\renewcommand{\arraystretch}{1.2}
\small
\begin{tabular}{llll}
 \toprule
 \textbf{Theory} & \textbf{Manifold} & \textbf{Topological $\En$} & \textbf{Holomorphic $\En$} \\
 \midrule
 CFG25 (3d CS) & $\R^3$ & $E_3$ & (not holomorphic) \\
 5d CS & $\C^2 \times \R$ & $E_4$ & $\Etwo$-chiral on $\C^2$ \\
 6d theory & $\C^3$ & $E_6$ & $E_3$-chiral on $\C^3$ \\
 \bottomrule
\end{tabular}
\end{center}
\noindent The $E_3$-chiral structure on $\C^3$ is the direct lift of the $E_3$ structure on $\R^3$. The holomorphic refinement is not a weakening: it is a \emph{different} $E_3$ structure, one that depends on the Omega-background parameters $(h_1, h_2, h_3)$ with $h_1 + h_2 + h_3 = 0$ for the nontrivial braiding (the topological braiding on $\Conf_2(\R^6)$ is trivial since $\pi_1(\Conf_2(\R^6)) \cong \pi_1(S^5) = 0$).
\end{proposition}

\begin{proof}
Costello--Gwilliam (2021), Costello--Francis--Gwilliam (2025). The holomorphic refinement is Proposition~\ref{prop:holomorphic-e1}; the $E_3$ on $\C^3$ is Conjecture~\ref{conj:topological-e3-comparison}(a).
\end{proof}

\subsection{Result 2: the boundary algebra upgrade}
\label{subsec:cfg25-result2}

The CFG25 boundary algebra is the Kac--Moody VOA $V_k(\frakg)$, which is $E_\infty$-chiral (a commutative vertex algebra). In the 6d lift, the boundary algebra becomes the affine Yangian $Y(\widehat{\fgl}_1)$, which is $E_1$-chiral (non-commutative, ordered).

\begin{proposition}[Deligne conjecture level drop in the CFG25 lift]
\label{prop:deligne-level-drop-cfg25}
The higher Deligne conjecture (Lurie, \emph{Higher Algebra}, Theorem~5.3.1.30) assigns to the Hochschild cochains $\HH^\bullet(A, A)$ of an $E_n$-algebra $A$ an $E_{n+1}$-algebra structure. In the CFG25 lift:
\begin{enumerate}[label=\textup{(\roman*)}]
 \item \emph{3d (CFG25)}: The boundary $V_k(\frakg)$ is $E_\infty$-chiral. Its Hochschild cochains carry $E_\infty$ structure (and in particular $E_3$, generated by cup product, Gerstenhaber bracket, and BV operator). The BV operator requires $E_\infty$ input.
 \item \emph{6d (lift)}: The boundary $Y(\widehat{\fgl}_1)$ is $E_1$-chiral. By Dunn additivity ($E_1 \otimes_{E_0} E_1 = E_2$), $\HH^\bullet(Y, Y)$ carries only $E_2$-algebra structure. No BV operator is available.
\end{enumerate}
The Deligne conjecture level drops from $E_3$ to $E_2$ in the lift. This is a genuine structural consequence: the commutativity of $V_k(\frakg)$ is \emph{lost} in passing to the Yangian.
\end{proposition}

\begin{proof}
Part~(i) is the higher Deligne conjecture applied to $E_\infty$ input (Lurie, \emph{loc.\ cit.}). Part~(ii) follows from Dunn additivity: $\HH^\bullet$ of an $E_n$-algebra is $E_{n+1}$, so $\HH^\bullet$ of $E_1$ is $E_2$. The BV operator arises from the $S^1$-action on the cyclic bar complex, which requires the commutativity of $A$.
\end{proof}

The passage from $E_1$ to $E_2$ is achieved not by the Deligne conjecture but by the \emph{Drinfeld center} construction:
\begin{equation}\label{eq:drinfeld-e1-to-e2-cfg25}
 \cZ(\Rep^{E_1}(Y^+(\widehat{\fgl}_1))) \;\simeq\; \Rep^{E_2}(Y(\widehat{\fgl}_1)),
\end{equation}
as proved in Theorem~\ref{thm:drinfeld-center-coha}. The full Yangian $Y(\widehat{\fgl}_1) = Y^+ \bowtie (Y^+)^{*\mathrm{cop}}$ is the Drinfeld double of the positive half $Y^+$ (the CoHA; note the CoHA itself is associative, not chiral).

\subsection{Result 3: submanifold invariants replace link invariants}
\label{subsec:cfg25-result3}

\begin{conjecture}[Holomorphic submanifold invariants from 6d factorization homology]
\label{conj:cfg25-submanifold-invariants}
\ClaimStatusConjectured
Let $\Sigma \subset \C^3$ be a smooth holomorphic submanifold of complex dimension $k$ ($k = 1$ for curves, $k = 2$ for surfaces). The factorization homology
\[
 \int_{\C^3 \setminus \Sigma} \cF
\]
of the $E_3$-chiral factorization algebra $\cF$ on the complement $\C^3 \setminus \Sigma$ produces an invariant of the pair $(\C^3, \Sigma)$ that depends on the holomorphic data of $\Sigma$, including its complex structure and topology. For codimension-$2$ curves ($k = 1$), the defect algebra on $\Sigma$ is $W_{1+\infty}$ at level $\Psi = -\sigma_2$ (Proposition~\ref{prop:codim2-defect-ope}).

The lift from 3d to 6d changes the invariant theory in three ways:
\begin{enumerate}[label=\textup{(\alph*)}]
 \item \emph{Classification}: links in $\R^3$ are classified by isotopy (finite combinatorial data); holomorphic submanifolds in $\C^3$ are classified by algebraic geometry (moduli spaces of potentially infinite dimension).
 \item \emph{Fundamental group}: $\pi_1(\R^3 \setminus L)$ is the knot group (nontrivial for every nontrivial link); $\pi_1(\C^3 \setminus \Sigma)$ may be trivial for smooth holomorphic curves. The invariant data migrates from $\pi_1$ (knot group representations) to the \emph{normal bundle} $N_{\Sigma / \C^3}$.
 \item \emph{Normal bundle}: for a curve $C \subset \C^3$ of genus $g$, the CY adjunction formula (trivial canonical bundle of $\C^3$) gives $\det(N_{C/\C^3}) = K_C^{-1} = \cO(2-2g)$. The normal bundle replaces the framing data of the link.
\end{enumerate}
\end{conjecture}

\begin{remark}[Why \texorpdfstring{$\pi_1$}{pi-1} migrates to the normal bundle]
\label{rem:pi1-to-normal-bundle}
In 3d CS, the Wilson line observable $W_\rho(K) = \mathrm{tr}_\rho(\mathrm{Hol}_K(A))$ is the trace in a representation $\rho$ of the holonomy around the knot $K$. The holonomy is a $\pi_1$-representation. In the 6d lift, the holonomy around a codimension-$2$ submanifold $\Sigma$ is trivial when $\pi_1(\C^3 \setminus \Sigma) = 0$, but the normal bundle $N_{\Sigma/\C^3}$ provides a \emph{replacement}: the equivariant parameters $(h_1, h_2)$ of the normal $\C^2$-fiber give the spectral parameters of the Yangian/toroidal $R$-matrix. The connection to quantum group parameters is through the Omega-background equivariant localisation on the normal bundle.
\end{remark}

\subsection{Result 4: factorization homology and quantum toroidal invariants}
\label{subsec:cfg25-result4}

\begin{conjecture}[Factorization homology on \texorpdfstring{$\C^3$}{C-3} complements]
\label{conj:cfg25-fact-hom-complement}
\ClaimStatusConjectured
For a holomorphic submanifold $\Sigma \subset \C^3$ in the toric CY$_3$ setting, the factorization homology $\int_{\C^3 \setminus \Sigma} \cF$ computes quantum toroidal $U_{q,t}(\widehat{\widehat{\fgl}}_1)$-module data. The identification with the Nekrasov partition function
\begin{equation}\label{eq:nekrasov-fact-hom-cfg25}
 Z_{\mathrm{Nek}}(\varepsilon_1, \varepsilon_2; q) \;=\; \int_{\C^2} \cF\big|_{\C^2}
\end{equation}
provides the computational bridge: $Z_{\mathrm{Nek}}^{U(1)} = \prod_{n \geq 1} 1/(1-q^n)$ with plethystic logarithm $q/(1-q)$ (one bosonic BPS mode per instanton number). This is the $\Etwo$-projection of the $E_3$ integral over $\C^3$.

For compact CY$_3$ (quintic, $K3 \times E$): the holomorphic factorization homology integral requires the Ran-space formulation of chiral homology, and the available CY-A$_3$ input is the framed object-level theorem under H1--H4 and fixed admissible specialisation \textup{(}Theorem~\ref{thm:cy-to-chiral-d3}\textup{)}. The 6d factorization homology route does \emph{not} bypass that apparatus: each subproblem (local $E_3$ algebra for compact targets, handle decomposition, VOA identification of output) requires the same chain-level data that CY-A$_3$ assumes or verifies on toric/formal loci. Explicit chain-level convergence and framing for non-formal compact CY$_3$ remain residual analytic content.
\end{conjecture}

\subsection{Result 5: Koszul duality and the ZTE obstruction}
\label{subsec:cfg25-result5}

The CFG25 Koszul duality $\mathrm{Obs}^q(\R^3)^! \simeq U_q(\frakg)\text{-mod}$ is \emph{proved}. The 6d lift to $\mathrm{Obs}^q(\C^3)^! \simeq U_{q,t}(\widehat{\widehat{\fgl}}_1)\text{-mod}$ is \emph{conjectural} (Conjecture~\ref{conj:e3-koszul-duality}).

\begin{proposition}[Parameter inversion under \texorpdfstring{$E_3$}{E-3}-chiral Koszul duality]
\label{prop:parameter-inversion-cfg25}
Under the Verdier duality on $\C^3$ factorization coalgebras, the Omega-background parameters invert: $(h_1, h_2, h_3) \mapsto (-h_1, -h_2, -h_3)$. On the elementary symmetric polynomials:
\begin{enumerate}[label=\textup{(\roman*)}]
 \item $\sigma_2(-\boldsymbol{h}) = \sigma_2(\boldsymbol{h})$: the level $\Psi = -\sigma_2$ is \emph{preserved} by Koszul duality (even function).
 \item $\sigma_3(-\boldsymbol{h}) = -\sigma_3(\boldsymbol{h})$: the deformation parameter $\sigma_3 = h_1 h_2 h_3$ is \emph{negated} (odd function).
 \item In the $(q, t)$ language: $q = e^{h_1} \mapsto e^{-h_1} = q^{-1}$, $t = e^{-h_2} \mapsto e^{h_2} = t^{-1}$.
\end{enumerate}
The Koszul conductor for the free-field (KM/$\fgl_1$) case is $\rho_K = 0$: $\kappa_{\mathrm{ch}}(A) + \kappa_{\mathrm{ch}}(A^!) = (-\sigma_2) + \sigma_2 = 0 = \rho_K$.
\end{proposition}

\begin{proof}
Direct computation. Let $f(\boldsymbol{h}) = e_k(h_1, h_2, h_3)$ be the $k$-th elementary symmetric polynomial. Under $h_i \mapsto -h_i$: $e_k(-\boldsymbol{h}) = (-1)^k e_k(\boldsymbol{h})$. So $e_2(-\boldsymbol{h}) = e_2(\boldsymbol{h})$ (even) and $e_3(-\boldsymbol{h}) = -e_3(\boldsymbol{h})$ (odd). For the modular characteristic, write $\kappa_{\mathrm{ch}}^{\mathrm{KM}} = -\sigma_2$; then $\kappa_{\mathrm{ch}}^{\mathrm{KM}}(-\boldsymbol{h}) = -\sigma_2(-\boldsymbol{h}) = -\sigma_2(\boldsymbol{h}) = \kappa_{\mathrm{ch}}^{\mathrm{KM}}(\boldsymbol{h})$. The Koszul dual has $\kappa_{\mathrm{ch}}^{!,\mathrm{KM}} = \sigma_2$, so $\kappa_{\mathrm{ch}}^{\mathrm{KM}} + \kappa_{\mathrm{ch}}^{!,\mathrm{KM}} = -\sigma_2 + \sigma_2 = 0 = \rho_K$ for the KM/$\fgl_1$ family (class~$\mathbf{G}$, $\rho_K = 0$). For other families the Koszul conductor differs: $\rho_K = 13$ for Virasoro, $\rho_K = 24$ for lattice rank (Vol~I Theorem~C).
\end{proof}

The key obstruction to proving the 6d Koszul duality is the \emph{Zamolodchikov tetrahedron equation} (ZTE). In 3d, the quantum Yang--Baxter equation (YBE) governs the $E_2$ coherence of the $R$-matrix, and the CFG25 result uses the YBE solution to build the Koszul dual. In 6d, the $E_3$ coherence requires the ZTE (the $3$-simplex analogue of the YBE $2$-simplex relation). Theorem~\ref{thm:zte-failure} proves that the na\"ive factorisation $S_{ijk} = R_{ij} R_{ik} R_{jk}$ from pairwise YBE solutions does \emph{not} satisfy the ZTE: the obstruction is $O(\sigma_3^2)$ and antisymmetric. Proposition~\ref{prop:zte-explicit-correction} kills the charge-$2$ obstruction by a genuinely ternary correction. Proposition~\ref{prop:zte-o-kappa4} shows that the odd charge sectors still carry rank-$2$ residues. Thus the local even-sector obstruction is solved, but the full charge-preserving all-order ZTE is not yet constructed.

\subsection{The obstruction hierarchy}
\label{subsec:cfg25-obstruction-hierarchy}

The three obstructions to the CFG25 lift form a partial order by severity:

\begin{center}
\renewcommand{\arraystretch}{1.3}
\small
\begin{tabular}{cllp{6cm}}
 \toprule
 \textbf{Level} & \textbf{Obstruction} & \textbf{Type} & \textbf{Required datum} \\
 \midrule
 1 & ZTE failure & Local algebraic & charge-$2$ obstruction killed by a ternary correction; full charge-preserving completion open \\
 2 & CY-A$_3$ ($\bS^3$-framing) & Global geometric & requires the corrected TCFT comparison or the complete $\HH^{-2}_{E_1}$ filtration theorem (Theorem~\ref{thm:derived-framing-m3b2}); chain-level explicit framing open for non-formal \\
 3 & 6d non-Lagrangian & Foundational & construct route~(c) of Remark~\ref{rem:6d-not-lagrangian} \\
 \bottomrule
\end{tabular}
\end{center}

\noindent Each later obstruction depends on resolving the earlier ones. The charge-$2$ ZTE obstruction is algebraic and has a ternary resolution; the odd-sector and all-order ZTE completion remain part of the local problem. CY-A$_3$ is controlled only after the corrected TCFT carrier or the complete $\HH^{-2}_{E_1}$ filtration comparison is supplied, and chain-level explicit framing data for non-formal algebras remains open; the non-Lagrangian nature of the 6d theory is the foundational question of whether the algebraic formulation of Costello--Francis--Gwilliam extends to the 6d setting.


% ============================================================
% Perturbative chiral invariants from 6d hCS Feynman diagrams
% ============================================================

\section{Perturbative chiral invariants from Feynman diagrams}
\label{sec:perturbative-chiral-feynman}

The perturbative expansion of the local hCS model is expected to see the
shadow tower of $\cW_{1+\infty}$ through the identification of loop order
with shadow depth. The proved content available here is finite:
the Virasoro channel has been checked through $L\leq 3$ against the
shadow recursion and the $\Ainf$ operations. The all-loop
Feynman/shadow dictionary remains a conjectural extension until the
regularised hCS graph integrals and their normalisation are constructed
uniformly.

\subsection{The perturbative expansion of 6d hCS}
\label{subsec:perturbative-expansion-6d}

\begin{proposition}[Perturbative chiral expansion]
\label{prop:perturbative-chiral-expansion}
For the local abelian hCS model on $\C^3_{h_1,h_2,h_3}$ restricted to
the codimension-$2$ defect $C \subset \C^3$, after a choice of
regularisation and effective Omega-background vertices, the perturbative
expansion has the formal graph shape
\begin{equation}\label{eq:perturbative-expansion}
 Z_{\mathrm{chiral}}(A_C) \;=\; \sum_{L \geq 0} \sigma_3^L \; Z^{(L)},
 \qquad
 Z^{(L)} \;=\; \sum_{\Gamma:\, b_1(\Gamma) = L}
 \frac{1}{|\!\Aut(\Gamma)|}\, w(\Gamma)\, I(\Gamma),
\end{equation}
where the sum ranges over connected Feynman graphs~$\Gamma$ at loop order~$L = b_1(\Gamma)$ (first Betti number), $w(\Gamma)$ is the Lie algebra weight from the Yangian structure constants, and $I(\Gamma)$ is the regularised configuration-space integral over $\Conf_{|V(\Gamma)|}(C)$.

For $\fgl_1$ (abelian gauge group), the cubic vertex of the holomorphic CS action vanishes, and the effective vertices arise from the Omega-background mode couplings at each spin level.
\end{proposition}

\begin{proof}
The displayed formula is the standard graph expansion one obtains after
choosing a propagator, regularisation, and effective vertices. For
$\fgl_1$, the cubic term in the bare hCS action vanishes identically, so
the nontrivial vertices in this witness come from the Omega-background
deformation rather than from a nonabelian cubic bracket. The finite graph
normalisations below are not an all-loop construction of renormalised
hCS observables.
\end{proof}

\subsection{Standard graphs and their contributions}
\label{subsec:standard-graphs}

The standard Feynman graphs at each loop order, for the spin-$s$ channel of the $\cW_{1+\infty}$ algebra at $c = 1$, are:

\begin{enumerate}[label=\textup{(\roman*)}]
 \item \emph{Tree level ($L = 0$): propagator exchange.} Two external vertices connected by one edge. The graph contributes $Z^{(0)}_s = \kappa_{\mathrm{ch},s} = c/s$, the modular characteristic of the spin-$s$ channel.

 \item \emph{One loop ($L = 1$): bubble diagram.} Two vertices connected by two parallel edges. Symmetry factor $1/|\!\Aut| = 1/2$. The regularised integral gives $Z^{(1)}_s = S_2^{(s)} = \kappa_{\mathrm{ch},s}$ (the shadow depth-$2$ invariant equals $\kappa_{\mathrm{ch}}$; this is the content of Theorem~A at genus~$1$).

 \item \emph{Two loops ($L = 2$): theta graph.} Three vertices forming a triangle ($3$ edges, $b_1 = 1$). This graph carries the cubic shadow $S_3 = \alpha$. For spin~$2$ (Virasoro at $c = 1$): $\alpha = 2$, arising from the $\Ainf$-operation $m_3(T, T, T) = -2T$. For spin~$1$ (Heisenberg, class~$\mathbf{G}$): $\alpha = 0$.

 \item \emph{Three loops ($L = 3$).} Multiple graph topologies contribute to $S_4$ (the quartic contact term). For spin~$2$ at $c = 1$: $S_4 = 10/27$, consistent with $m_4(T, T, T, T) = \tfrac{40}{27}\, T$.
\end{enumerate}

\subsection{The shadow--Feynman dictionary}
\label{subsec:shadow-feynman-dictionary}

\begin{proposition}[Finite loop/shadow witness]
\label{prop:loop-shadow-identification}
For the Virasoro spin-$2$ channel with the graph normalisation of
Proposition~\ref{prop:perturbative-chiral-expansion}, the finite
checked dictionary is
\begin{equation}\label{eq:shadow-feynman-dictionary}
 Z^{(L)} \;=\; S_{L+1}, \qquad L = 0, 1, 2, 3.
\end{equation}
Explicitly, for the Virasoro channel at $c = 1$:
\begin{center}
\small
\renewcommand{\arraystretch}{1.15}
\begin{tabular}{ccccl}
\toprule
\textbf{Loop $L$} & \textbf{Shadow $S_{L+1}$} & \textbf{Value} & $\boldsymbol{\Ainf}$ \textbf{operation} & \textbf{Graph} \\
\midrule
$0$ & $S_1 \equiv \kappa_{\mathrm{ch}}$ & $1/2$ & $m_2$ (OPE bracket) & propagator \\
$1$ & $S_2$ & $1/2$ & $m_2$ (self-energy) & bubble \\
$2$ & $S_3 = \alpha$ & $2$ & $m_3(T,T,T) = -2T$ & theta \\
$3$ & $S_4$ & $10/27$ & $m_4(T,T,T,T) = \tfrac{40}{27}T$ & sunset + others \\
\bottomrule
\end{tabular}
\end{center}
The all-loop identity $Z^{(L)} = S_{L+1}$ for every $L$ is the natural
shadow--$\Ainf$--coproduct conjecture suggested by the Vol~I
correspondence, not a consequence of the finite computation above.
\end{proposition}

\begin{proof}
For $L\leq3$, the Feynman contribution $Z^{(L)}$ is computed from the
chosen regularised configuration-space integrals on
$\overline{\Conf}_n(C)$ using holomorphic residue calculus. The shadow
invariant $S_{L+1}$ is computed from the Vol~I convolution recursion
(the quadratic $Q_L(t) = 4\kappa_{\mathrm{ch}}^2 + 12\kappa_{\mathrm{ch}}\,\alpha\, t + (9\alpha^2 + 16\kappa_{\mathrm{ch}}\, S_4)\, t^2$).

The identification proceeds by matching both sides at each order:
\begin{itemize}
 \item $L = 1$: $Z^{(1)} = \kappa_{\mathrm{ch}} = S_2$. The one-loop bubble computes the self-energy correction, which equals $\kappa_{\mathrm{ch}}$ by Theorem~A.
 \item $L = 2$: $Z^{(2)} = \alpha = S_3 = 2$. The two-loop theta graph computes the cubic shadow, matching $|m_3(T,T,T)| = 2$.
 \item $L = 3$: $Z^{(3)} = S_4 = 10/27$. The three-loop graphs give the quartic contact, matching $m_4(T,T,T,T) = \tfrac{40}{27}\, T$ via the normalisation $S_4 = m_4\text{-coefficient} / 4$.
\end{itemize}
The finite dictionary matches the shadow tower recursion and the
$\Ainf$ operations at these four loop orders.
\end{proof}

\subsection{Shadow class from the loop expansion}
\label{subsec:shadow-class-loops}

The shadow class (Definition~\ref{def:shadow-class}, Vol~I) is determined by the loop expansion structure:
\begin{itemize}
 \item Class~$\mathbf{G}$: all loops $L \geq 1$ vanish (tree level only). The Heisenberg algebra (spin~$1$) is class~$\mathbf{G}$: the perturbative expansion terminates at $L = 0$.
 \item Class~$\mathbf{L}$: finitely many nonzero loop orders.
 \item Class~$\mathbf{C}$: convergent loop expansion (hence Borel summable).
 \item Class~$\mathbf{M}$: divergent loop expansion (Gevrey-$1$); Borel resummation is available only in the positivity/discriminant regime $\kappa_{\mathrm{ch}}^3S_4>0$ of Proposition~\ref{prop:class-m-borel-summability}. The Virasoro channel (spin~$2$) at $c = 1$ lies in this class~$\mathbf{M}$ regime: the shadow tower is infinite.
\end{itemize}

The spin~\(1\) Heisenberg channel is class~\(\mathbf{G}\). The
Virasoro subalgebra at spin~\(2\) already has an infinite shadow tower,
so a shadow theory retaining that sector is not the rank-one Heisenberg
scalar shadow.

\subsection{The MacMahon connection}
\label{subsec:macmahon-connection-feynman}

The MacMahon function records primitive BPS multiplicities. It is not
obtained by summing spin-channel modular characteristics.

\begin{proposition}[MacMahon primitive multiplicities]
\label{prop:constant-effective-kappa}
\ClaimStatusProvedHere
Let \(V_{\mathrm{BPS}}=\oplus_{n\geq 1}V_n\) with \(\dim V_n=n\).
Then
\[
 M(q)=\prod_{n\geq 1}(1-q^n)^{-n}
 =\mathrm{ch}\,\Sym(V_{\mathrm{BPS}})
 =\operatorname{PE}\!\left(\sum_{n\geq 1}nq^n\right).
\]
Thus the exponent \(n\) in the \(n\)-th factor of \(M(q)\) is the
primitive DT multiplicity \(\Omega(n)\), not a scalar
\(\kappa_{\mathrm{ch}}\)-channel.
\end{proposition}

\begin{proof}
The plethystic exponential identity gives the displayed product. The
plethystic logarithm is
\[
 \operatorname{PLog}M(q)=\sum_{n\geq 1}nq^n,
\]
which is precisely the primitive BPS spectrum of the \(\C^3\) Hall
algebra.
\end{proof}

\begin{remark}[Perturbative chiral invariants]
\label{rem:perturbative-feynman-summary}
The perturbative expansion is governed by graph combinatorics, Lie
algebra weights, holomorphic configuration integrals, the tree-level
identity $\kappa_{\mathrm{ch},s}=c/s$, the one-loop equality
$S_2=\kappa_{\mathrm{ch}}$, the two-loop value $S_3=\alpha=2$, the
three-loop value $S_4=10/27$, and the dictionary
$L\leftrightarrow S_{L+1}$ between loop order and shadow depth.
\end{remark}

% ============================================================
% E_3 Hochschild deformation theory: the 6d analogue of CFG25
% ============================================================

\section{\texorpdfstring{$E_3$}{E-3} Hochschild cohomology as chiral deformation theory}
\label{sec:e3-hochschild-deformation-theory}

In CFG25, the deformation quantisation of 3d Chern--Simons theory uses the $E_3$ bracket on $\HH^\bullet(\mathrm{Obs}_{\mathrm{cl}})$: the Maurer--Cartan element $\alpha \in \HH^2_{E_3}$ is the quantum correction, and the deformed observables $\mathrm{Obs}^q = \mathrm{CE}_*({\frak g})[[\hbar]]$ are controlled by this MC element. The 6d analogue is the $E_3$ Hochschild cohomology $\HH^\bullet_{E_3}(A, A)$ of a chiral algebra~$A$ as the deformation theory controlling $A \to A[[\varepsilon_1, \varepsilon_2, \varepsilon_3]]$. We open with the chain-level foundation: the brace-operad proof of Deligne's conjecture at $n = 2$ on Hochschild cochains, and its extension to $n = 3$ under the CY cyclic hypothesis via the Menichi BV operator.

\subsection{Deligne at \texorpdfstring{$n = 2$}{n=2}: the brace operad on Hochschild cochains}
\label{subsec:deligne-n2-brace}

The $E_2$-structure on $\CC^\bullet(A,A)$ is the Deligne conjecture at $n = 2$: McClure--Smith, Tamarkin, and Voronov. The chain-level witness is the brace operad.

\begin{definition}[Brace operations on Hochschild cochains;
\ClaimStatusDefinitional]
\label{def:brace-operations}
Let $A$ be an $A_\infty$-algebra over a field of characteristic zero. Write $\CC^n(A,A) = \Hom_k(A^{\otimes n}, A)$ for the Hochschild cochains (with a degree shift by $n$; our degree conventions follow Gerstenhaber). For $x \in \CC^m, y_1 \in \CC^{k_1}, \ldots, y_n \in \CC^{k_n}$, the \emph{$n$-brace} of $x$ by $y_1, \ldots, y_n$ is
\begin{equation}\label{eq:hoch-brace-explicit}
 \{x; y_1, \ldots, y_n\}(a_1, \ldots, a_N)
 \;=\;
 \sum_{0 \leq i_1 \leq i_2 \leq \cdots \leq i_n} (-1)^\eta\,
 x\bigl(a_1, \ldots, a_{i_1}, y_1(a_{i_1+1}, \ldots, a_{i_1+k_1}),
 a_{i_1+k_1+1}, \ldots, y_n(\ldots), \ldots, a_N\bigr),
\end{equation}
with $N = m + \sum k_j - n$, summation over non-overlapping insertions of $y_j$ starting at position $i_j$ with $i_j + k_j \leq i_{j+1}$, and the Koszul sign
\[
 \eta \;=\; \sum_{j=1}^n |y_j|\!\left(\sum_{\ell = 1}^{i_j} |a_\ell|\right)
 \;+\; \sum_{j < j'} |y_j| |y_{j'}|.
\]
The arity-$0$ brace is $\{x; \,\} = x$; the arity-$2$ brace with no insertions (the \emph{cup product}) is $(f \smile g)(a_1, \ldots, a_{m+n}) = (-1)^{|g| \sum_{\ell \leq m} |a_\ell|} f(a_1, \ldots, a_m) \cdot g(a_{m+1}, \ldots, a_{m+n})$, where the product uses $\mu_2$ of $A$.
\end{definition}

\begin{theorem}[Deligne at $n = 2$: brace model]
\label{thm:deligne-n2-brace-model}
Let $\Br$ denote the non-symmetric dg operad generated by the brace operations of Definition~\ref{def:brace-operations} with the brace-associativity relations
\[
 \bigl\{\{x; y_1, \ldots, y_m\};\, z_1, \ldots, z_n\bigr\}
 \;=\;
 \sum (-1)^{\eta'}\, \bigl\{x;\, \ldots,\, \{y_j; z_?, \ldots, z_?\},\, \ldots\bigr\},
\]
where the right-hand sum ranges over insertion patterns that assign each $z_k$ either to a free slot of $x$ or to one of the $y_j$. Then:
\begin{enumerate}[label=\textup{(\roman*)}]
\item The operations~\eqref{eq:hoch-brace-explicit} equip $\CC^\bullet(A,A)$ with a $\Br$-algebra structure for any $A_\infty$-algebra $A$.
\item There is a quasi-isomorphism of dg operads over $\Q$
\[
 \Br \;\xrightarrow{\;\sim\;}\; C_\bullet(E_2;\Q),
\]
canonical up to the $\mathrm{GRT}_1(\Q)$-torsor of Drinfeld associators (Tamarkin 1998; the choice of associator pins the rational homotopy type uniformly).
\item The Hochschild differential $b = \{\mu_A;-\} - \{-;\mu_A\}$ is a $\Br$-derivation, so $\CC^\bullet(A,A)$ is an $E_2$-algebra in dg vector spaces.
\end{enumerate}
Passing to cohomology, $(\HH^\bullet(A,A), \smile, [-,-]_\Ger)$ is a Gerstenhaber algebra with the Gerstenhaber bracket $[f,g]_\Ger = \{f;g\} - (-1)^{(|f|-1)(|g|-1)} \{g;f\}$.
\end{theorem}

\begin{proof}
(i) The brace identities are verified by direct chain-level computation (Voronov \cite{VoronovHomotopyGer} Theorem~3.3): each brace identity expands to a sum over shuffles of insertion positions with Koszul signs that cancel in pairs. (ii) McClure--Smith \cite{McClureSmithDeligne} Theorem~A constructs an explicit zigzag of dg operads through the Smith filtration of the sequence operad; Tamarkin \cite{TamarkinAnotherProof} \S3 gives the alternative construction via a Drinfeld associator, with the choice space a $\mathrm{GRT}_1(\Q)$-torsor (compare Remark~\ref{rem:deligne-kt-grt1-e2}). (iii) The identity $[b, \{-;-, \ldots, -\}] = 0$ reduces, after expansion, to the $A_\infty$-relation at $n = 2$: $b^2 = 0$ is the $A_\infty$-relation for $\mu_A$.
\end{proof}

\begin{remark}[Gerstenhaber pre-Lie and cohomological Poisson]
\label{rem:gerstenhaber-pre-lie}
The Gerstenhaber pre-Lie product $f \circ g = \{f;g\}$ is not associative; its failure of associativity is an explicit brace:
\[
 (f \circ g) \circ h \;-\; f \circ (g \circ h)
 \;=\; \{f; g, h\} \;+\; (-1)^{(|g|-1)(|h|-1)} \{f; h, g\}.
\]
The Gerstenhaber bracket, being the antisymmetrisation of $\circ$, is Lie on cohomology. The Poisson identity $[f, g \smile h]_\Ger = [f,g]_\Ger \smile h + (-1)^{(|f|-1)|g|} g \smile [f,h]_\Ger$ holds strictly on $\HH^\bullet(A,A)$ and up to explicit brace homotopies on $\CC^\bullet(A,A)$.
\end{remark}

\subsection{The Menichi BV operator on cyclic Hochschild cochains}
\label{subsec:menichi-bv-operator}

The cyclic pairing of a CY $A_\infty$-algebra upgrades the $E_2$-structure of Theorem~\ref{thm:deligne-n2-brace-model} to a chain-level framed $E_{d+1}$ structure. The upgrade operator is the Menichi BV operator, dual to the Connes cyclic $B$-operator on chains.

\begin{proposition}[Chain-level BV operator on cyclic Hochschild cochains]
\label{prop:menichi-delta-explicit}
Let $(A, \{\mu_n\}, \langle-,-\rangle)$ be a unital cyclic $A_\infty$-algebra of CY dimension $d$, with pairing $\langle-,-\rangle \colon A \otimes A \to k[-d]$. Fix a pair of dual bases $\{e_\alpha\}, \{e^\alpha\}$ satisfying $\langle e_\alpha, e^\beta\rangle = \delta_\alpha^\beta$. For $f \in \CC^n(A,A)$, define
\begin{equation}\label{eq:menichi-delta-explicit}
 \Delta(f)(a_1, \ldots, a_{n-1})
 \;=\;
 \sum_{i = 1}^{n}\, (-1)^{\sigma_i}\,
 \sum_\alpha \langle f(a_i, \ldots, a_{n-1}, e_\alpha, a_1, \ldots, a_{i-1}), e^\alpha\rangle^{\mathrm{PD}},
\end{equation}
where $\langle-,-\rangle^{\mathrm{PD}}$ denotes the Poincar\'e-dual pairing that uses $\langle-,-\rangle$ to convert a scalar back into an element of $A$ via the isomorphism $A \cong A^\vee[-d]$, and the sign $\sigma_i$ is the Koszul sign from cyclically permuting $n - i$ arguments of total degree past $i$ arguments with degrees $|a_j|$, plus the $d$-shift from the pairing. Then:
\begin{enumerate}[label=\textup{(\roman*)}]
\item $\Delta$ has degree $1 - d$: $\Delta \colon \CC^n(A,A) \to \CC^{n-1}(A,A)[1-d]$.
\item $\Delta^2 = 0$ and $[b, \Delta] = 0$.
\item \textup{(BV identity.)} For $f \in \CC^m, g \in \CC^n$,
\begin{equation}\label{eq:bv-identity}
 [f, g]_\Ger
 \;=\; (-1)^{|f|}\bigl(\Delta(f \smile g) - \Delta(f)\smile g - (-1)^{|f|} f \smile \Delta(g)\bigr).
\end{equation}
\end{enumerate}
Equivalently, $(\CC^\bullet(A,A), b, \smile, \Delta)$ is a BV algebra up to homotopy, and strictly a BV algebra on $\HH^\bullet(A,A)$, with BV-degree shift $d - 1$.
\end{proposition}

\begin{proof}
Menichi \cite{MenichiBV} Proposition~2.1 establishes $\Delta^2 = 0$ by direct chain-level expansion: the double cyclic rotation decomposes into pairs $(i, n-i)$ that cancel by cyclic symmetry of $\langle-,-\rangle$. The BV identity~\eqref{eq:bv-identity} is Menichi \cite{MenichiBV} Theorem~A (via $\infty$-inner-product transport of Tradler--Zeinalian \cite{TradlerZeinalian} Theorem~3.4): expanding $\Delta(f \smile g)$ in the cup and brace shows that the $\mu_2$-insertions match the Gerstenhaber pre-Lie up to the stated symmetriser, and Koszul signs match. The chain-homotopy $[b, \Delta] = 0$ follows because $b$ and $\Delta$ both factor through $\mu_A$-contraction and cyclic rotation, which commute as endomorphisms of $A^{\otimes \bullet}$.
\end{proof}

\begin{remark}[Degree shift comes from the cyclic pairing]
\label{rem:delta-degree-shift}
The degree $1 - d$ of $\Delta$ is the pairing-degree shift: the dual $e^\alpha$ lives in $A[-d]$ because $\langle-,-\rangle$ has degree $-d$, so reassembling $\CC^{n-1}$ from a scalar costs a $-d$ shift, compensated by the $+1$ from removing an argument slot. At $d = 0$ (ordinary BV): $\Delta$ has degree $+1$, the classical Connes $B$-dual. At $d = 1$ (topological string surfaces): $\Delta$ has degree $0$, the Kontsevich cyclic operator. At $d = 2$ (K3 / CY$_2$): $\Delta$ has degree $-1$, the classical BV generator. At $d = 3$ (CY$_3$): $\Delta$ has degree $-2$, the CY$_3$-twisted BV generator used below.
\end{remark}

\subsection{Deligne at \texorpdfstring{$n = 3$}{n=3}: three compatible \texorpdfstring{$E_3$}{E-3}-structures}
\label{subsec:deligne-n3-three-structures}

The naive statement ``Deligne at $n = 3$ for a CY$_3$ category'' conflates three distinct $E_3$-structures on Hochschild-type complexes. We name them, prove their chartwise compatibility, and separate the finite Rees gluing theorem from the remaining compact completion and realisation obligations.

\begin{theorem}[Three compatible \texorpdfstring{$E_3$}{E-3}-structures at CY\texorpdfstring{$_3$}{-3}]
\label{thm:deligne-n3-three-e3}
Let $X$ be a smooth Calabi--Yau threefold with cyclic $A_\infty$-enhancement $A_X$ of $D^b(\Coh(X))$. Three $E_3$-algebra structures arise on Hochschild-type complexes:
\begin{description}
\item[\textup{(A) Topological (CFG 2026).}] $\Obs^{\mathrm{cl}}_{\mathrm{CS}}(\R^3)$ is $E_3$ from the little $3$-disks operad on $\R^3$; the structure is topological, with Deligne at $n = 3$ supplying an $E_4$-algebra on its Hochschild cochains by Lurie \cite{LurieHA} Theorem~5.3.1.30.

\item[\textup{(B) Holomorphic (hCS on $\C^3$).}] $\Obs_{\hCS}(\C^3)$ is holomorphic-$E_3$ (Gwilliam--Williams \cite{GwilliamWilliams}) from the three complex directions of $\C^3$; the structure is real $E_3$ after Kontsevich--Tamarkin formality of $\cD_3$ (Theorem~\ref{thm:plat-hCS-quantum}).

\item[\textup{(C) Algebraic (CY$_3$-cyclic on Hochschild).}] $\CC^\bullet(A_X, A_X)$ is $E_3$ via the combined $(\Br, \Delta)$-action: the brace operad of Theorem~\ref{thm:deligne-n2-brace-model} supplies the $E_2$-structure; the Menichi BV operator of Proposition~\ref{prop:menichi-delta-explicit} supplies the $\Delta$ of degree $-2$ (since $d = 3$); their combined operadic level is $fE_4$ on cohomology, and the restricted $E_3$-sub-action (forgetting one $SO(2)$-plane) is canonical up to the contractible $S^2 = SO(3)/SO(2)$ of forgetful maps.
\end{description}
On each $\C^3$-chart $U \subseteq X$, structures \textup{(B)} and \textup{(C)} agree via Kontsevich--Tamarkin formality: the HKR isomorphism $\CC^\bullet(\cO_U, \cO_U) \simeq \mathrm{PV}^\bullet(U)$ intertwines the two $E_3$-actions up to a $\mathrm{GRT}_1(\Q)$-torsor.  On a finite DWR/Ran cover with bounds \((N,r,L,m)\), the hCS--Hall comparison glues at the Rees ordered Hall layer under the hypotheses of Proposition~\ref{prop:qca-finite-rees-layer-separation}.  The ordinary compact critical CoHA comparison requires \(R_{\mathrm{vc}}\) and the vanishing of the pro-descent \(\varprojlim^1\)-obstruction.
\end{theorem}

\begin{proof}
Structure \textup{(A)} is direct: $\R^3$ has a little $3$-disks operad action on observables, and Lurie \cite{LurieHA} Theorem~5.3.1.30 gives $E_{n+1}$ on Hochschild of $E_n$. Structure \textup{(B)} is Theorem~\ref{thm:plat-hCS-quantum}: shuffle integrals against the Bochner--Martinelli kernel realise the Fulton--MacPherson little $3$-disks action on $\Obs_{\hCS}(\C^3)$.

Structure \textup{(C)}: Theorem~\ref{thm:deligne-n2-brace-model} gives the $E_2$-action on $\CC^\bullet(A_X, A_X)$ via brace; Proposition~\ref{prop:menichi-delta-explicit} at $d = 3$ gives $\Delta$ of degree $-2$; Getzler's theorem identifies the homology of the combined operad with $\Ger \otimes \Lambda[\Delta]$ in the $d$-shifted grading, which is the homology of $fE_{d+1} = fE_4$. Forgetting the $SO(3) \supset SO(2)$-rotation gives an $E_3$-sub-action canonical up to $SO(3)/SO(2) = S^2$. The CY$_3$ volume form $\Omega_X = dz_1 \wedge dz_2 \wedge dz_3$ pins the forgetful map by distinguishing the third complex direction, making the $E_3$-action canonical.

Compatibility of \textup{(B)} and \textup{(C)} on a $\C^3$-chart: the HKR quasi-isomorphism $\CC^\bullet(\cO_{\C^3}, \cO_{\C^3}) \simeq \mathrm{PV}^\bullet(\C^3)$ (Hochschild--Kostant--Rosenberg, Tamarkin \cite{TamarkinAnotherProof}; see Remark~\ref{rem:three-atiyah-cocycles-qca} below for the Kapranov $\ell_\infty$-structure) is an $E_2$-quasi-isomorphism. The Kontsevich formality theorem for the $n$-disks operad (Kontsevich \cite{KontsevichOperads}; Tamarkin) lifts this to an $E_3$-quasi-isomorphism after choice of an associator. The formality bridge is canonical up to $\mathrm{GRT}_1(\Q)$.

Globally on compact $X$: the chartwise $E_3$-isomorphisms glue to the finite Rees Hall layer precisely under the finite DWR/Ran hypotheses named in Proposition~\ref{prop:qca-finite-rees-layer-separation}.  This does not yet identify the completed compact hCS observables with the ordinary critical CoHA: that stronger statement needs the monoidal vanishing-cycle realisation and pro-compatible transition data in the truncation variables.

Compatibility of \textup{(A)} with \textup{(B)}--\textup{(C)}: the CFG topological $E_3$ lives on $\R^3$; the hCS holomorphic $E_3$ lives on $\C^3 = \R^6$. The comparison is through the Wick-rotation / holomorphic-topological transition, a further layer that the CY$_3$-cyclic $E_3$ of \textup{(C)} does not need to traverse. Remark~\ref{rem:cy3-no-cfg-shortcut} of the CY$_3$ chain-level bridge chapter records this separation: CFG 2026's $E_3$ is not by itself a proof that the hCS $E_3$ on $\C^3$ agrees with the CY$_3$-cyclic $E_3$ on Hochschild; the two are related through the finite Rees comparison only after the DWR/Ran cyclic-atlas data have been fixed.
\end{proof}

\begin{remark}[Operadic separations]
\label{rem:deligne-n3-ap-cleared}
Three confusions are frequent on this axis and are cleared by Theorem~\ref{thm:deligne-n3-three-e3}. First, ``Deligne at $n = 3$ applied to a CY$_3$ category gives $E_3$ on $\CC^\bullet$'' is ambiguous: it may mean Lurie's iterated $\HH^\bullet$ of an $E_2$-input (giving $E_3$ without CY) or it may mean the BV-lift of the $E_2$-brace action under a CY$_3$ cyclic pairing. The two produce $E_3$-structures on different complexes; they coincide only after the Kontsevich--Tamarkin formality chartwise identification above. Second, ``BV is $E_3$'' is imprecise: BV is $fE_2$ on cohomology (Getzler); the $E_3$-promotion requires the CY dimension as an operadic shift, so ``$E_3$'' without naming the shift is type-ambiguous. Third, ``$R$-matrix equals the $E_3$-bracket on $\CC^\bullet$'' is false: the $R$-matrix at $d = 3$ arises from $E_2$-braiding on $\Rep(D(Y^+))$ after the Drinfeld-double step; the $E_3$-bracket on $\CC^\bullet$ controls the Maurer--Cartan deformation theory upstream of the double, not the braiding downstream.
\end{remark}

\begin{remark}[Three Atiyah cocycles and the CY$_3$-cyclic \texorpdfstring{$E_3$}{E-3}]
\label{rem:three-atiyah-cocycles-qca}
On a smooth Calabi--Yau threefold $X$ the Atiyah class $\At(T_X) \in \HH^1(X, T_X[-1])$ sources cohomologically distinct structures (Kapranov \cite{Kapranov1999}; Costello--Li \cite{CostelloLi2016}): the binary bracket $\ell_2 = \At(T_X)$; the tree-level Yukawa $Y_3 \in H^{0,3}(X)$, identified with Kapranov's minimal $\ell_3$; and the one-loop BCOV determinant-line or gravitational datum.  The latter is not a class \((\chi_{\mathrm{top}}(X)/24)[\Omega_X]^{0,1}\), since a holomorphic CY$_3$ volume form has no \((0,1)\)-component.  These structures live in different complexes, so none is the algebraic $E_3$-bracket of Theorem~\ref{thm:deligne-n3-three-e3}(C); the Menichi $\Delta$ of Proposition~\ref{prop:menichi-delta-explicit} sits in a different Hochschild tridegree and is pinned by the cyclic pairing, not by the Atiyah class. At $K3 \times E$ the Euler-characteristic coefficient of the BCOV one-loop datum vanishes because $\chi_{\mathrm{top}}(K3 \times E) = 24 \cdot 0 = 0$, while the product Yukawa and the algebraic $E_3$-bracket remain separate inputs.
\end{remark}

\begin{theorem}[Restricted \texorpdfstring{$E_3$}{E-3} feeds \texorpdfstring{$\Phi^{FA}_3$}{Phi-FA-3}]
\label{thm:restricted-e3-feeds-phi-fa-3}
Let $\cC$ be a smooth proper cyclic $A_\infty$-category of CY dimension $3$ on an admissible Stage-$1$ chart. The restricted $E_3$-sub-action on $\CC^\bullet(\cC,\cC)$ of Theorem~\ref{thm:deligne-n3-three-e3}(C) is the algebraic operadic input for the chartwise Stage-$1$ output $\PhiFA_3(\cC)$ of the two-stage factorisation $\Phi_3 = \SpCh_{\Sigma_2, C} \circ \PhiFA_3$. Specifically:
\begin{enumerate}[label=\textup{(\roman*)}]
\item The $E_3$-Maurer--Cartan element $\alpha \in \HH^2_{E_3}(A,A)$ of Subsection~\ref{subsec:mc-hh-e3} is the algebraic avatar of the hCS Maurer--Cartan solution, read through the $\Br \otimes \Lambda[\Delta]$ chain-level model.
\item The $E_4$-bracket $[-,-]_{E_3}$ of equation~\eqref{eq:mc-e3} is the higher Deligne-iteration bracket: Lurie \cite{LurieHA} Theorem~5.3.1.30 applied to the $E_3$-algebra $\CC^\bullet(\cC,\cC)$ gives an $E_4$-structure on $\HH^\bullet_{E_3}(\CC^\bullet(\cC,\cC))$, whose degree-$(-1)$ Lie bracket is $[-,-]_{E_3}$.
\item The classical limit $\alpha_1$ of the Maurer--Cartan element is the Gerstenhaber bracket of the tree-level Yukawa $Y_3$ of Remark~\ref{rem:three-atiyah-cocycles-qca} with itself; the one-loop coefficient is controlled by the BCOV determinant-line datum; the full tower $\{\alpha_k\}$ is the shadow tower of Conjecture~\ref{conj:mc-shadow-tower}.
\end{enumerate}
\end{theorem}

\begin{proof}
(i) The $E_3$-structure on $\CC^\bullet(\cC,\cC)$ of Theorem~\ref{thm:deligne-n3-three-e3}(C) produces the $E_3$-Hochschild cohomology $\HH^\bullet_{E_3}(A,A) = \mathrm{R}\Hom_{\Mod_A^{E_3}}(A,A)$ used in Subsection~\ref{subsec:e3-tricomplex}. Its Maurer--Cartan locus $\mathrm{MC}(\HH^2_{E_3})$ parametrises $E_3$-deformations of the input $A_\infty$-structure; by the chartwise formality Theorem~\ref{thm:deligne-n3-three-e3}, these match hCS deformations on $\C^3$-charts.

(ii) Direct: Lurie \cite{LurieHA} Theorem~5.3.1.30 gives $E_{n+1}$ on Hochschild of $E_n$; at $n = 3$, $E_4$; the $E_4$-structure contains a degree-$(-1)$ Lie bracket by Dunn--Lurie additivity $E_4 \simeq E_1 \otimes E_3$.

(iii) Classical-limit identification: Witten's $\partial\bar\partial$-action for hCS on $X$ has Maurer--Cartan solution at leading order the Yukawa coupling $Y_3$ (cubic tensor in the gauge-algebra directions), matching the brace $\{Y_3; Y_3\}$ in the Gerstenhaber pre-Lie grading. At one loop, Costello--Li \cite{CostelloLi2016} inserts the BCOV determinant-line contribution; this is the source of the $\alpha_2$ coefficient. The shadow tower recursion of Conjecture~\ref{conj:mc-shadow-tower} extends the pattern through finite-depth chain models.
\end{proof}

\begin{problem}[Global compact-CY$_3$ compatibility of the restricted $E_3$ input]
\label{op:restricted-e3-feeds-phi-fa-3-global}
Upgrade Theorem~\ref{thm:restricted-e3-feeds-phi-fa-3} from the
admissible Stage-$1$ chartwise locus to a global compact CY$_3$ statement by
constructing compatible finite DWR/Ran Rees comparisons at all
\((N,r,L,m)\), killing the resulting \(\varprojlim^1\)-obstruction, and
supplying the monoidal vanishing-cycle realisation
\(R_{\mathrm{vc}}\) from the Rees critical Hall layer to the ordinary
compact critical CoHA.
\end{problem}

\begin{remark}[CFG 2026 and the chain-level bridge]
\label{rem:cfg2026-chain-level-bridge}
CFG 2026's Deligne-at-$n=3$ derivation uses the little $3$-disks operad on the topological space $\R^3$ directly: classical observables of Chern--Simons with spectral parameter carry an $E_3$-algebra structure by factorisation, and Lurie's theorem produces the $E_4$-upgrade on their Hochschild cochains. Our derivation of Theorem~\ref{thm:deligne-n3-three-e3}(C) instead uses the brace operad on Hochschild cochains of a cyclic $A_\infty$-category, with the Menichi BV operator supplying the extra layer from the CY cyclic pairing. The two derivations produce compatible $E_3$-algebras chartwise on $\C^3$-charts, via the Kontsevich--Tamarkin formality zigzag $\Br \otimes \Lambda[\Delta] \xrightarrow{\sim} C_\bullet(fE_4;\Q) \to C_\bullet(E_3;\Q)$.  On compact CY$_3$ the finite DWR/Ran Rees comparison gives the current Hall-valued descent layer; completion and ordinary critical CoHA realisation are the additional obligations named in Problem~\ref{op:restricted-e3-feeds-phi-fa-3-global}.
\end{remark}

\subsection{The \texorpdfstring{$E_3$}{E-3} Hochschild tricomplex}
\label{subsec:e3-tricomplex}
\label{subsec:e3-bar-class-m}

For an $E_3$-algebra~$A$ on $\C^3$, the $E_3$ Hochschild cohomology $\HH^\bullet_{E_3}(A, A) = R\Hom_{\mathrm{Mod}_A^{E_3}}(A, A)$ carries an $E_4$-algebra structure by the higher Deligne conjecture (Lurie, \emph{Higher Algebra}, Theorem~5.3.1.30). The underlying cochain complex decomposes as a \emph{tricomplex} $C^{p,q,r}(A)$ with three commuting differentials $d_1, d_2, d_3$ (one per complex direction on $\C^3$):
\begin{equation}\label{eq:e3-tricomplex}
 d_i \colon C^{p,q,r} \to C^{p + \delta_{i1},\, q + \delta_{i2},\, r + \delta_{i3}},
 \qquad d_i^2 = 0,\quad [d_i, d_j] = 0.
\end{equation}

At the chain level, for a chiral algebra with $n$ generators:
\begin{equation}\label{eq:e3-tricomplex-chain-dim}
 \dim C^{p,q,r} = \binom{n}{p}\binom{n}{q}\binom{n}{r},
 \qquad
 \dim_{\mathrm{total}} = 8^n,
 \qquad
 P(s,t,u) = ((1+s)(1+t)(1+u))^n.
\end{equation}

The equivariant weights of the three differentials are the Omega-background parameters:
\[
 \mathrm{wt}(d_1) = h_1, \quad \mathrm{wt}(d_2) = h_2, \quad \mathrm{wt}(d_3) = h_3,
 \qquad h_1 + h_2 + h_3 = 0 \;\text{(CY)}.
\]
The CY condition forces $\mathrm{wt}(d_{\mathrm{total}}) = h_1 + h_2 + h_3 = 0$, ensuring $d_{\mathrm{total}}^2 = 0$.

\begin{proposition}[Finite-mode tricomplex pages by shadow class]
\label{prop:tricomplex-dimensions}
Let \(A\) be a chiral algebra with \(n\) finite-mode generators and shadow class \(\mathbf X\in\{G,L,C,M\}\). The finite-mode \(E_3\) tricomplex has the following terminal page when no later shadow differential is present:
\begin{center}
\small
\renewcommand{\arraystretch}{1.15}
\begin{tabular}{ccccc}
\toprule
\textbf{Class} & \textbf{Chain dim} & \textbf{Page dim} & \textbf{Poincar\'e} & \textbf{Deficit} \\
\midrule
$\mathbf{G}$ & $8^n$ & $8^n$ & $((1+s)(1+t)(1+u))^n$ & $0$ \\
$\mathbf{L}$ & $8^n$ & $8^n$ & $(1+t)^{3n}$ & $0$ \\
$\mathbf{C}$ & $8^n$ & $8^n$ & $(1+t)^{3n}$ & $0$ \\
$\mathbf{M}$ & $8^n$ & $\dim E_4=6^n$ & $(3t(1+t))^n$ & $8^n - 6^n$ at \(E_4\) \\
\bottomrule
\end{tabular}
\end{center}
The deficit for class~$\mathbf{M}$ at \(E_4\) is the image-kernel pair of the quartic shadow differential (see Proposition~\ref{prop:e3-spectral-degeneration} below). For \(n\leq3\) no later differential can act, hence this page is \(E_\infty\). For \(n\geq4\), \(d_5,d_6,\ldots,d_{n+1}\) are additional finite-mode operators. For class~$\mathbf{G}$, the cohomology Poincar\'e polynomial retains the three-variable form because all differentials vanish; for classes~$\mathbf{L}$ and~$\mathbf{C}$, the three variables collapse to the total degree~$t = stu$ after the differentials act.
\end{proposition}

\begin{proof}
The chain-level dimension \(8^n=(2^n)^3\) follows from the product decomposition \(C^{p,q,r}=\Lambda^p(V)\otimes\Lambda^q(V)\otimes\Lambda^r(V)\) with \(\dim V=n\). The cohomology for class~\(\mathbf G\) equals the chain level because all \(d_i\) vanish. For class~\(\mathbf M\), \(d_4\) contracts the top exterior class in each finite-mode factor: per factor, the \(d_4\)-homology of \(\Lambda^*(k^3)\) is \([0,3,3,0]\) with Poincar\'e \(3t+3t^2\) (Proposition~\ref{prop:virasoro-d4} in~\S\ref{subsec:e3-bar-class-m}); by K\"unneth, the \(n\)-fold tensor gives \(E_4\)-page \((3t(1+t))^n\) with total dimension \(6^n\). Classes~\(\mathbf L\) and~\(\mathbf C\) have \(d_4=0\) because the shadow tower terminates before depth~\(4\), respectively because charge conservation kills the quartic operator, so their cohomology equals the \(E_3\)-page dimension \(2^{3n}=8^n\).

This is the exterior-algebra count in the \(E_3\) tricomplex.
\end{proof}

\subsection{The spectral sequence \texorpdfstring{$E_3 \Rightarrow E_2 \Rightarrow E_1$}{E-3 => E-2 => E-1}}
\label{subsec:e3-spectral-sequence}

The tricomplex admits a spectral sequence by filtering on the third direction:
\begin{align*}
 E_1^{p,q,r} &= H^r(C^{p,q,\bullet}, d_3) & &\text{($d_3$-cohomology)} \\
 E_2^{p,q,r} &= H^q(E_1^{p,\bullet,r}, d_2) & &\text{($d_2$-cohomology of $d_3$-cohomology)} \\
 E_\infty^{p,q,r} &= H^p(E_2^{\bullet,q,r}, d_1) & &\text{($= \HH^\bullet_{E_3}(A,A)$).}
\end{align*}

The pages record the successive Hochschild reductions: $E_1 \approx E_2$-Hochschild (one direction computed), $E_2 \approx E_1$-Hochschild (two directions), $E_\infty = E_3$-Hochschild (all three).

\begin{proposition}[Degeneration page by shadow class]
\label{prop:e3-spectral-degeneration}
The spectral sequence of the $E_3$ tricomplex degenerates as follows:
\begin{center}
\small
\renewcommand{\arraystretch}{1.15}
\begin{tabular}{clcl}
\toprule
\textbf{Class} & \textbf{Degeneration} & \textbf{$d_4$ coefficient $\Delta$} & \textbf{Reason} \\
\midrule
$\mathbf{G}$ & $E^{\mathrm{ss}}_1 = E^{\mathrm{ss}}_\infty$ & $0$ & All $d_i = 0$ (abelian) \\
$\mathbf{L}$ & $E^{\mathrm{ss}}_3 = E^{\mathrm{ss}}_\infty$ & $0$ & Shadow terminates at depth~$2$ \\
$\mathbf{C}$ & $E^{\mathrm{ss}}_3 = E^{\mathrm{ss}}_\infty$ & $0$ & Charge conservation kills $d_4$ \\
$\mathbf{M}$, $n \leq 3$ & $E^{\mathrm{ss}}_4 = E^{\mathrm{ss}}_\infty$ & $8\kappa_{\mathrm{ch}} S_4$ & Degree reasons (support too narrow for $d_5$) \\
$\mathbf{M}$, $n \geq 4$ & $E_{n+2}$ at latest & $8\kappa_{\mathrm{ch}} S_4$ & $d_r = 0$ for $r \geq n+2$ by degree \\
\bottomrule
\end{tabular}
\end{center}
For the Virasoro at $c = 1$: $\Delta = 8 \cdot \tfrac{1}{2} \cdot \tfrac{10}{27} = \tfrac{40}{27}$, agreeing with the direct computation in Proposition~\ref{prop:virasoro-d4}.
\end{proposition}

\begin{proof}
For class~$\mathbf{G}$: all differentials $d_i$ vanish on an abelian algebra, so $E^{\mathrm{ss}}_1 = E^{\mathrm{ss}}_\infty$ immediately. For class~$\mathbf{L}$: the shadow tower has $S_k = 0$ for $k \geq 3$, so the higher differential $d_4$ (which is controlled by $S_4$) vanishes; $E^{\mathrm{ss}}_3 = E^{\mathrm{ss}}_\infty$. For class~$\mathbf{C}$: the $U(1)$-charge grading forces $d_4 = 0$ by the same mechanism as the $E_3$ bar for the bc system.

For class~$\mathbf{M}$: the $E_4$ page is supported in tridegrees $(p,q,r)$ with each entry in $\{0, 1\}$ (after the $d_1, d_2, d_3$ contractions), giving a band of total degree $[n, 2n]$ (width $n+1$). The differential $d_r$ has shift $r-1$. For $d_5$ to act nontrivially, both source degree $k$ and target degree $k - 4$ must lie in $[n, 2n]$, requiring $4 \leq n$. So $d_5 = 0$ for $n \leq 3$, giving $E^{\mathrm{ss}}_4 = E^{\mathrm{ss}}_\infty$. For $n \geq 4$: the spectral sequence terminates at $E_{n+2}$ at the latest (since $d_r$ with $r - 1 > n$ has shift exceeding the band width).

\end{proof}

\begin{proposition}[Class $\mathbf{M}$ bar growth past \texorpdfstring{$n = 3$}{n = 3}]
\label{prop:e3-bar-class-m-asymptotic}
Let $A$ be a class-$\mathbf{M}$ chiral algebra with $n$ finite-mode generators and shadow invariants \(S_r\). On the finite one-primitive quotient, the \(d_4\)-homology has Poincar\'e polynomial
\[
 P(E_4;t)=(3t+3t^2)^n,
 \qquad
 \dim E_4=6^n .
\]
The later differentials are precisely
\[
 d_r:E_r^k\longrightarrow E_r^{k-r+1},
 \qquad 5\leq r\leq n+1,\qquad n+r-1\leq k\leq 2n.
\]
Thus \(E_\infty=E_4\) for \(n\leq3\). For \(n=4\),
\[
 \dim E_\infty=1296-2\,\rho_5,
 \qquad
 \rho_5=\operatorname{rank}\!\bigl(d_5:E_5^8\to E_5^4\bigr),
 \qquad
 0\leq\rho_5\leq81 .
\]
In general
\[
 \dim E_\infty
 =
 \dim H^\bullet\!\left(E_4,\ d_5+d_6+\cdots+d_{n+1}\right),
\]
where \(d_r\) is the finite-mode operator induced by the \(r\)-th shadow \(S_r\).
The ordinary PBW bar complex is infinite-dimensional before this finite-mode completion is imposed.
\end{proposition}

\begin{proof}
After the three Hochschild contractions, the \(d_4\)-homology of one finite-mode generator is \([0,3,3,0]\), with support in degrees \(1\) and \(2\). K\"unneth gives \(E_4=[0,3,3,0]^{\otimes n}\), hence \(P(E_4;t)=(3t+3t^2)^n\) and \(\dim E_4=6^n\). The support of \(E_4\) is the degree band \([n,2n]\). A differential \(d_r\) lowers total degree by \(r-1\); it can act only from degrees \(k\) satisfying \(k\leq2n\) and \(k-r+1\geq n\), equivalently \(n+r-1\leq k\leq2n\). This is empty for every \(r\geq5\) when \(n\leq3\). For \(n=4\) the only non-empty case is \(r=5,k=8\); the source and target are both \(3^4=81\)-dimensional, giving the displayed rank formula. The last possible differential has \(r=n+1\). The spectral sequence identity gives the final formula.
\end{proof}

\begin{remark}[Admissible shadow instances]
\label{rem:e3-bar-class-m-admissible-instances}
The number of possible higher-shadow positions is
\[
 N(n) \;:=\; \sum_{r = 5}^{n+1} \bigl(n - r + 2\bigr) \;=\; \frac{(n-3)(n-2)}{2} \quad (n \geq 3).
\]
For \(n\leq3\) this number is \(0\). For \(n=4\) it is \(1\), the single map \(d_5:E_5^8\to E_5^4\). For larger \(n\) the ranks of these maps are additional shadow data, not consequences of the \(d_4\) calculation alone.
\[
\begin{array}{c|c|c}
n & \dim E_4 & N(n)\\
\hline
3 & 216 & 0\\
4 & 1296 & 1\\
5 & 7776 & 3\\
6 & 46656 & 6\\
10 & 60466176 & 28
\end{array}
\]
\end{remark}

\subsection{Maurer--Cartan deformation on \texorpdfstring{$\HH^2_{E_3}$}{HH-2-E-3}}
\label{subsec:mc-hh-e3}

The Maurer--Cartan element governing the chiral deformation quantisation lives in $\HH^2_{E_3}(A, A)$. The MC equation is:
\begin{equation}\label{eq:mc-e3}
 d(\alpha) + \tfrac{1}{2}[\alpha, \alpha]_{E_3} = 0,
 \qquad \alpha \in \HH^2_{E_3}(A, A),
\end{equation}
where $[-, -]_{E_3}$ is the degree-$(-1)$ Lie bracket from the $E_4$-structure on $\HH^\bullet_{E_3}$ (the higher Deligne conjecture provides $E_{n+1}$ structure on $\HH^\bullet_{E_n}$; here $n = 3$, giving $E_4$, and the $E_4$ structure includes a Lie bracket of degree $-1$).

The MC element expands in powers of $\sigma_3 = h_1 h_2 h_3$:
\begin{equation}\label{eq:mc-expansion}
 \alpha = \sum_{k \geq 1} \sigma_3^{k-1}\, \alpha_k,
 \qquad
 \alpha_k \in \HH^2_{E_3}(A, A).
\end{equation}

\begin{conjecture}[MC element = shadow tower]
\label{conj:mc-shadow-tower}
\ClaimStatusConjectured
The $k$-th MC coefficient $\alpha_k$ is identified with the shadow invariant $S_k$ (Vol~I, Definition~\ref{def:shadow-invariants}):
\begin{equation}\label{eq:mc-shadow-identification}
 \alpha_k \;\longleftrightarrow\; S_k \;\longleftrightarrow\; Z^{(k-1)} \;\longleftrightarrow\; m_k,
\end{equation}
where $Z^{(L)}$ is the regularised $L$-loop Feynman contribution when
the graph-integral package is defined (finite evidence is recorded in
Proposition~\ref{prop:loop-shadow-identification}), and $m_k$ is the
$\Ainf$~operation of the bar complex (Vol~I~\cite{VolI}). The conjectural content is that the MC equation~\eqref{eq:mc-e3}, order by order in $\sigma_3$, reproduces the shadow tower recursion.
\end{conjecture}

\begin{remark}[Chain-level tricomplex obligation]
\label{rem:mc-shadow-why-conjecture}
The identification~\eqref{eq:mc-shadow-identification} is conjectural because it requires the \emph{explicit} $E_3$-chiral tricomplex structure on $A$ for CY$_3$ targets. The \((\infty,1)\)-categorical CY-A$_3$ input of Theorem~\ref{thm:derived-framing-m3b2} gives the framed existence statement; the explicit chain-level tricomplex data for non-formal algebras, needed to compute the MC coefficients $\alpha_k$, remains open. The correspondence $S_k \leftrightarrow Z^{(k-1)}$ is currently a finite Virasoro-channel witness through $k\leq4$ (Proposition~\ref{prop:loop-shadow-identification}), while the $S_k \leftrightarrow m_k$ comparison belongs to the Vol~I bar-complex lane. The composite arrow $\alpha_k \leftrightarrow S_k$ is the new content: it identifies the MC coefficients on $\HH^2_{E_3}$ with the shadow tower. The conjecture asserts that the deformation-theoretic packaging (the MC equation on $\HH^2_{E_3}$) and the combinatorial packaging (the shadow tower recursion) describe the same underlying structure.
\end{remark}

The MC series has convergence controlled by the shadow class:
\begin{center}
\small
\renewcommand{\arraystretch}{1.15}
\begin{tabular}{clll}
\toprule
\textbf{Class} & \textbf{MC series type} & \textbf{Convergence} & \textbf{Example} \\
\midrule
$\mathbf{G}$ & Polynomial (one term) & Radius $= \infty$ & Heisenberg: $\alpha = \alpha_1$ \\
$\mathbf{L}$ & Polynomial (finite) & Radius $= \infty$ & Yangian: $\alpha = \alpha_1 + \sigma_3 \alpha_2$ \\
$\mathbf{C}$ & Convergent & Radius $> 0$ & bc system \\
$\mathbf{M}$ & Gevrey-$1$ divergent & Radius $= 0$ & Virasoro: Borel on $\kappa_{\mathrm{ch}}^3S_4>0$ \\
\bottomrule
\end{tabular}
\end{center}

\subsection{The CFG25 dictionary: 3d vs 6d deformation quantisation}
\label{subsec:cfg25-dictionary}

The relationship between CFG25 (3d CS on $\R^3$) and 6d hCS on $\C^3$ is captured by a dictionary that translates each element of the 3d deformation quantisation to its 6d chiral analogue:

\begin{center}
\small
\renewcommand{\arraystretch}{1.3}
\begin{tabular}{p{4.5cm}p{6cm}}
\toprule
\textbf{CFG25 (3d CS on $\R^3$)} & \textbf{6d hCS on $\C^3$} \\
\midrule
Single differential $d$ on $\mathrm{CE}_*(\frak{g})$ & Three commuting differentials $d_1, d_2, d_3$ on $\HH^\bullet_{E_3}(A, A)$ \\
One parameter $\hbar$ & $(h_1, h_2, h_3)$ with $h_1 + h_2 + h_3 = 0$; essential parameter $\sigma_3$ \\
MC in $\HH^2(\mathrm{Obs}_{\mathrm{cl}})$: $E_3$ bracket from $\R^3$ & MC in $\HH^2_{E_3}(A, A)$: $E_3$ bracket from $\C^3$ tricomplex \\
$\alpha = \hbar\, r + \hbar^2 \cdot \ldots$ (classical $r$-matrix $+$ corrections) & $\alpha = S_1 + \sigma_3 S_2 + \sigma_3^2 S_3 + \cdots$ (shadow tower) \\
$d(\alpha) + \tfrac{1}{2}[\alpha,\alpha] = 0$ (CYBE $+$ quantum) & $d(\alpha) + \tfrac{1}{2}[\alpha,\alpha]_{E_3} = 0$ (shadow recursion) \\
Koszul dual $\mathrm{CE}_*(\frak{g})^! = U_q(\frak{g})$ & Koszul dual $\HH^\bullet_{E_3}(A)^! =$ chiral quantum group [\textsc{conj}] \\
Boundary $V_k(\frak{g})$ ($E_\infty$-chiral) & Boundary $Y(\widehat{\fgl}_1)$ ($E_1$-chiral) \\
Deligne: $E_3$ on $\HH^\bullet(V_k)$ (from $E_\infty$ input) & Deligne: $E_2$ on $\HH^\bullet(Y)$ (from $E_1$ input) \\
\bottomrule
\end{tabular}
\end{center}

The Deligne conjecture level drop (the last row) is the structural
reason that the 6d theory requires the \emph{external} \(E_3\) structure
from holomorphic factorisation on \(\C^3\), rather than the
\emph{internal} \(E_3\) from the higher Deligne conjecture: the boundary
algebra \(Y(\widehat{\fgl}_1)\) is only \(E_1\)-chiral, so the Deligne
conjecture produces only \(E_2\) on its Hochschild cochains
(Proposition~\ref{prop:deligne-level-drop-cfg25}).  The
Omega-background supplies the equivariant weights of this factorisation
structure; it is not an additional operadic direction.

\subsection{The tangent-obstruction sequence}
\label{subsec:e3-tangent-obstruction}

The \(E_3\) deformation complex of a chiral algebra~\(A\) has tangent
and obstruction spaces
\begin{equation}\label{eq:e3-tangent-obstruction}
 \HH^0_{E_3}(A,A) \to \HH^1_{E_3}(A,A) \to \HH^2_{E_3}(A,A) \to \cdots
\end{equation}
depending on the differential, relations, coefficients, and
completion.  For the associated-graded free tricomplex with \(n\)
generators and zero internal differential one obtains the model
dimensions
\begin{align}
 \dim \HH^0_{E_3} &= n & &\text{(infinitesimal $E_3$-automorphisms)}, \label{eq:hh0-dim} \\
 \dim \HH^1_{E_3} &= 3n & &\text{(first-order $E_3$ deformations: $n$ per direction)}, \label{eq:hh1-dim} \\
 \dim \HH^2_{E_3} &= \tfrac{3n(3n-1)}{2} & &\text{(obstructions: tridegree-$2$ contributions)}, \label{eq:hh2-dim}
\end{align}
giving the free-model virtual dimension
\(\mathrm{vdim}_{\mathrm{gr}} = 3n - \tfrac{3n(3n-1)}{2}\), which is
\(0\) for \(n = 1\) and negative for \(n \geq 2\).  No assertion is made
that these are the Hochschild dimensions of an arbitrary chiral algebra.

\begin{proposition}[\texorpdfstring{$\HH^2_{E_3}$}{HH-2-E-3} in the free associated-graded tricomplex]
\label{prop:hh2-tridegree}
In the free associated-graded tricomplex with \(n\) generators and zero
internal differential, the tridegree \(p + q + r = 2\) part receives
contributions:
\begin{enumerate}[label=\textup{(\roman*)}]
 \item Tridegrees $(2,0,0)$, $(0,2,0)$, $(0,0,2)$: each $\binom{n}{2}$, total $3 \cdot \tfrac{n(n-1)}{2}$.
 \item Tridegrees $(1,1,0)$, $(1,0,1)$, $(0,1,1)$: each $n^2$, total $3n^2$.
\end{enumerate}
Sum: $\tfrac{3n(n-1)}{2} + 3n^2 = \tfrac{3n(3n-1)}{2}$. Verified at $n = 1$ (dim $= 3$), $n = 3$ (dim $= 36$), $n = 24$ (dim $= 2556$).
\end{proposition}

\begin{proof}
Direct count of the tridegree-$(p,q,r)$ contributions to the exterior-algebra tricomplex. For item~(i): $\binom{n}{2} \cdot \binom{n}{0} \cdot \binom{n}{0} = \binom{n}{2}$ for each of $3$ orderings. For item~(ii): $\binom{n}{1}^2 \cdot \binom{n}{0} = n^2$ for each of $3$ orderings. The Mukai value $n = 24$ gives \(3\cdot276+3\cdot576=2556\).
\end{proof}



%% Four-manifold invariants from 6d holomorphic CS

\subsection{Four-manifold invariants from 6d holomorphic theory}
\label{subsec:four-manifold-6d}

The dimensional hierarchy of holomorphic CS produces invariants of
manifolds in defect dimension $d_{\mathrm{defect}} = \dim(\text{theory}) - 2$:
$3d$ CS gives WRT invariants of $3$-manifolds; $5d$ holomorphic CS gives
conjectural invariants via affine Yangians; $6d$ $(2,0)$ theory gives
$4$-manifold invariants via quantum toroidal algebras. Compactification
of the $6d$ $(2,0)$ theory on $X^4 \times \Sigma_g$ produces an effective
theory on $X^4$; for $\Sigma_g = T^2$, this is $4d$ $\mathcal{N} = 4$
SYM under the Vafa--Witten twist.

\begin{conjecture}[Four-manifold invariants from the \texorpdfstring{$6d$}{6d} chiral construction]
\label{conj:four-manifold-6d-invariants}
\ClaimStatusConjectured
\index{four-manifold invariants!6d holomorphic theory}%
For a compact oriented $4$-manifold $X$ with $b_1(X) = 0$, the $6d$
invariant $Z_{6d}(X \times T^2)$ equals the Vafa--Witten partition
function $Z_{\mathrm{VW}}(X; \tau)$, a modular form
\textup{(}or mock modular form\textup{)} of weight $w = -\chi(X)/2$
under $\mathrm{SL}_2(\Z)$. The modularity character depends on the
intersection form $Q_X$:
\begin{center}
\small
\renewcommand{\arraystretch}{1.15}
\begin{tabular}{lcccl}
\toprule
$X$ & $\chi(X)$ & $b_2^+(X)$ & $w$ & $Z_{\mathrm{VW}}$ \\
\midrule
$K3$ & $24$ & $3$ & $-12$ & $1/\eta^{24}$ \textup{(}genuine modular\textup{)} \\
$T^4$ & $0$ & $3$ & $0$ & $1$ \textup{(}trivial\textup{)} \\
$\C\bP^2$ & $3$ & $1$ & $-3/2$ & mock modular \\
$S^2 \times S^2$ & $4$ & $1$ & $-2$ & mock modular \\
$E(n)$ & $12n$ & $2n{-}1$ & $-6n$ & $1/\eta^{12n}$ \textup{(}$b_2^+ > 1$ for $n \geq 2$\textup{)} \\
\bottomrule
\end{tabular}
\end{center}

The dichotomy is: $b_2^+(X) > 1 \Rightarrow$ genuine modular form;
$b_2^+(X) = 1 \Rightarrow$ mock modular \textup{(}requiring
non-holomorphic completion\textup{)}.

The scalar entering the rank-one Vafa--Witten Hilbert-scheme product is
the surface Euler characteristic $\chi_{\mathrm{top}}(X)$; it is not, by
itself, the compact chiral invariant of a CY$_3$ object. A comparison
with a chiral characteristic of $A_{\mathrm{Tot}(K_X)}$ requires the
framed $\Phi_3^{(\Sigma_2,C)}$ specialisation and the local
non-compact boundary condition on $\mathrm{Tot}(K_X)$,
$\kappa_{\mathrm{cat}} = \chi(\mathcal{O}_X)$ \textup{(}Noether:
$(\chi + \sigma)/4$\textup{)},
$\kappa_{\mathrm{fiber}} = b_2(X)$ \textup{(}lattice rank\textup{)}.

For $K3$: $\kappa_{\mathrm{ch}}^{\mathrm{OPE}} = 24$ (the free-field
generator count at the OPE-level on $\mathrm{Tot}(K_{K3})$;
distinct from $\kappa_{\mathrm{ch}}^{\mathrm{Hodge}}(D^b\mathrm{Coh}(K3)) = 2$),
$\kappa_{\mathrm{cat}} = 2$,
$\kappa_{\mathrm{fiber}} = 22$, and $Z_{\mathrm{VW}} = 1/\eta^{24}$
is unconditional \textup{(}CY-A$_2$\textup{)}. For $\C\bP^2$:
$\kappa_{\mathrm{cat}} = 1$, and the mock modularity of $Z_{\mathrm{VW}}$
is established by Manschot--Thomas.

The construction via $\mathrm{Tot}(K_X)$ requires the framed CY$_3$
comparison data for the non-compact total space. The factorization homology route
$\int_{X^4} A_{E_2}$ is conjectural.
\end{conjecture}

\begin{remark}[Noether formula and Hilbert scheme consistency]
\label{rem:four-manifold-noether-hilb}
For complex surfaces, the holomorphic Euler characteristic
$\chi(\mathcal{O}_X) = (\chi(X) + \sigma(X))/4$ (Noether's formula)
provides a cross-check: $\chi(\mathcal{O}_{K3}) = (24 - 16)/4 = 2$,
$\chi(\mathcal{O}_{\C\bP^2}) = (3 + 1)/4 = 1$. The rank-$1$ VW
generating function $\prod_{k \geq 1} (1 - q^k)^{-\chi(X)}$ reproduces
Hilbert scheme Euler characteristics $\chi(\mathrm{Hilb}^n(X))$
(Gottsche 1990): for $K3$, the sequence $1, 24, 324, 3200, 25650,
176256$ matches $\prod(1 - q^k)^{-24}$ exactly.

\noindent The relevant four-manifold invariants
evaluates the surface catalogue, Noether formula, Hilbert scheme Euler
characteristics through $n = 10$, VW modular weights, mock/genuine
modularity dichotomy at $b_2^+ = 1$, $\kappa_\bullet$-spectrum, and
$S$-duality cross-checks.
\end{remark}

%% CHIRAL WRT INVARIANTS FROM THE K3 QUANTUM TOROIDAL ALGEBRA

\subsection{Chiral Witten--Reshetikhin--Turaev invariants}
\label{subsec:chiral-wrt}

The $3d$ WRT invariant $Z_{\mathrm{WRT}}(M^3;\, \frakg,\, k)$ of a closed $3$-manifold $M$ is constructed from three ingredients:
\begin{enumerate}[label=\textup{(\roman*)}]
 \item a quantum group $U_q(\frakg)$ at a root of unity $q = e^{2\pi i/(k + h^\vee)}$, whose representation category is a modular tensor category;
 \item a surgery presentation $M = S^3(L)$ for a framed link $L$ with linking matrix $L_{ij}$;
 \item a state sum over colourings of the link components.
\end{enumerate}
The $6d$ chiral WRT invariant $Z^{\mathrm{ch}}(K3;\, N)$ replaces each ingredient as follows.

\begin{center}
\renewcommand{\arraystretch}{1.3}
\small
\begin{tabular}{llp{5.5cm}}
 \toprule
 \textbf{$3d$ CS ingredient} & \textbf{$6d$ chiral analogue} & \textbf{Construction locus} \\
 \midrule
 $U_q(\frakg)$ at root of unity & $U_{q,t}(\fg_{K3}^{\mathrm{tor}})$ at $\omega = e^{2\pi i/N}$ & requires root-of-unity bridge \\
 Surgery $M = S^3(L)$ & Handle decomposition: $K3 = D^4 \cup 22\,(D^2 \times D^2) \cup D^4$ & Freedman--Quinn handle theorem \\
 Linking matrix $L_{ij}$ & Intersection form $Q_{K3} = 3U \oplus 2E_8(-1)$, sig.\,$(3, 19)$ & lattice computation \\
 State sum & Factorization homology on handles & requires compact handle factorisation \\
 $\sigma(L) = \mathrm{sig}(L_{ij})$ & $\sigma(K3) = 3 - 19 = -16$ & signature computation \\
 \bottomrule
\end{tabular}
\end{center}

\begin{conjecture}[Abelian chiral WRT invariant of K3]
\label{conj:chiral-wrt-k3}
\ClaimStatusConjectured
For the abelian truncation of the K3 quantum toroidal algebra at level $N \geq 2$, the chiral WRT invariant is
\begin{equation}\label{eq:chiral-wrt-k3}
 Z^{\mathrm{ch,ab}}(K3;\, N)
 \;=\; N^{\chi(K3)/2} \cdot \varphi(N)
 \;=\; N^{12} \cdot \varphi(N),
\end{equation}
where $\varphi(N)$ is the Gauss phase product over the Mukai lattice $\widetilde{\Lambda}_{K3}$ of signature~$(4, 20)$:
\[
 \varphi(N)
 \;=\; \prod_{a=1}^{24}
 \frac{g(\sigma_a, N)}{|g(\sigma_a, N)|},
 \qquad
 g(\sigma, N)
 \;=\; \sum_{n=0}^{N-1} \exp\!\bigl(\pi i \sigma n^2 / N\bigr).
\]
Here $\sigma_a = +1$ for $a = 1, \ldots, 4$ and $\sigma_a = -1$ for $a = 5, \ldots, 24$.

The phase $\varphi(N) = 1$ for all $N$, because $\sigma(K3) = 4 - 20 = -16 \equiv 0 \pmod{8}$ and an even unimodular lattice with signature divisible by~$8$ has trivial Gauss phase. Therefore $Z^{\mathrm{ch,ab}}(K3;\, N) = N^{12}$, an integer.
\end{conjecture}

The exponent $12 = \chi(K3)/2$ arises from the handle decomposition: each Betti direction contributes $\sqrt{N}$ to the quantum dimension, and the total is $N^{(b_0 + b_2 + b_4)/2} = N^{(1 + 22 + 1)/2} = N^{12}$. At $N = 2$:
\[
 Z^{\mathrm{ch,ab}}(K3;\, 2) \;=\; 2^{12} \;=\; 4096.
\]
The same integer is obtained from the Gauss-sum product formula, the
handle-by-handle decomposition, and the Betti-number product.

\begin{remark}[Connection to G\"ottsche and Vafa--Witten]
\label{rem:chiral-wrt-gottsche-vw}
The chiral WRT invariant $Z^{\mathrm{ch}}(K3;\, N) = N^{12}$ is \emph{not} the same as:
\begin{enumerate}[label=\textup{(\alph*)}]
 \item The G\"ottsche Euler characteristics $\chi(\mathrm{Hilb}^n(K3)) = p_{24}(n)$ with generating function $\prod_{k \geq 1}(1 - q^k)^{-24}$: these are instanton moduli space invariants, while $Z^{\mathrm{ch}}$ is a topological state count.
 \item The Vafa--Witten partition function $Z_{\mathrm{VW}}(K3, \mathrm{SU}(N);\, \tau)$, which is a modular form of weight $-\chi(K3)/2 = -12$.
\end{enumerate}
The relationship: $Z_{\mathrm{VW}}(K3, \mathrm{SU}(N);\, \tau) = Z^{\mathrm{ch}}(K3;\, N) \cdot F(\tau;\, N)$ where $F$ is the one-loop fluctuation determinant around the truncated vacuum. The chiral WRT extracts the zero-mode (topological) component.
\end{remark}

\begin{conjecture}[Chiral WRT of \texorpdfstring{$K3 \times E$}{K3 x E} and \texorpdfstring{$\Delta_5$}{Delta-5}]
\label{conj:chiral-wrt-k3xe}
\ClaimStatusConjectured
For $K3 \times E$, the product formula gives
\begin{equation}\label{eq:chiral-wrt-k3xe}
 Z^{\mathrm{ch}}(K3 \times E;\, N) \;=\; Z^{\mathrm{ch}}(K3;\, N) \cdot N \;=\; N^{13},
\end{equation}
where the factor $N$ is the chiral Verlinde dimension $V^{\mathrm{ch}}_N(E)$ on the elliptic curve~$E$. The exponent $13 = \chi(K3)/2 + 1$ is related to the weight $\kappa_{\mathrm{BKM}}(\Delta_5) = 5$ of the Igusa cusp form $\Delta_5$: in the unnormalised Igusa-square convention, the DMVV formula
\[
 \frac{1}{\Phi_{10}^{\mathrm{un}}}=\Delta_5^{-2}
 = \sum_N p^N \chi(S^N K3;\, q,\, y)
\]
has $N$-th coefficients whose polynomial prefactor matches $N^{13}\).
The identification of \(Z^{\mathrm{ch}}(K3 \times E;\, N)\) with the
asymptotics of \(\Delta_5\) Fourier coefficients becomes a theorem only
after CY-A and the Vol~I Borcherds-lift comparison are supplied.
\end{conjecture}

%% THE CHIRAL VOLUME CONJECTURE

\subsection{The chiral volume conjecture}
\label{subsec:chiral-volume-conjecture}

The $3d$ volume conjecture (Kashaev 1997; Murakami--Murakami 2001) asserts that
for a hyperbolic knot $K$ in $S^3$,
\begin{equation}\label{eq:3d-volume-conjecture}
\lim_{N \to \infty} \frac{2\pi}{N}\,
 \log\bigl| J_N\!\bigl(K;\, e^{2\pi i / N}\bigr)\bigr|
 \;=\; \operatorname{Vol}(S^3 \setminus K),
\end{equation}
where $J_N(K; q)$ is the $N$-th coloured Jones polynomial and
$\operatorname{Vol}(S^3 \setminus K)$ is the hyperbolic volume of the
knot complement. The $6d$ holomorphic theory does not have knots
(holomorphic curves are not knots: Obstruction~$2$ in
Section~\ref{subsec:cfg25-obstruction-hierarchy}), and the lift-rate
assessment of Section~\ref{subsec:cfg25-result3} classified the volume
conjecture as a \emph{complete gap} with $0\%$ lift rate. This
assessment is correct for a literal lift. What follows is an
\emph{analogue}, not a lift: it occupies the same structural
position (relating quantum invariants to classical geometry) but
uses different mathematics on both sides.

\medskip\noindent\textbf{The $3d$-to-$6d$ dictionary.}
\begin{center}
\small
\renewcommand{\arraystretch}{1.2}
\begin{tabular}{lll}
\toprule
& \textbf{3d Chern--Simons} & \textbf{6d chiral construction} \\
\midrule
Ambient space & $S^3$ & CY$_3$ $X$ \\
Submanifold & knot $K$ (real $1$-manifold) & holomorphic curve $C$ (complex $1$-manifold) \\
Quantum invariant & $J_N(K;\, e^{2\pi i/N})$ & $Z_N^{\mathrm{ch}}(C, X)$ \\
Classical geometry & $\operatorname{Vol}(S^3 \setminus K)$ & $\lvert\mathrm{AJ}(C)\rvert$ \\
Rigidity & Mostow (unique metric) & attractor mechanism (unique at attractor) \\
\bottomrule
\end{tabular}
\end{center}

\begin{conjecture}[Chiral volume conjecture]
\label{conj:chiral-volume-conjecture}
\ClaimStatusConjectured
Let $X$ be a Calabi--Yau threefold with Hodge-normalised holomorphic
$3$-form $\Omega$ ($\int_X \Omega \wedge \bar\Omega = 1$). Let
$C \subset X$ be a holomorphic curve representing the class
$\beta \in H_2(X, \Z)$. Define the \emph{charge-$N$ chiral partition
function}
\[
 Z_N^{\mathrm{ch}}(C, X)
 \;:=\;
 \sum_{\substack{d \,\geq\, 1 \\ d\,[\hspace{0.2pt}C\hspace{0.2pt}]\,=\, N\beta}}
 \Omega_{\mathrm{DT}}(d \cdot \beta,\, X)\, q^{\hspace{0.3pt}d}
\]
where $\Omega_{\mathrm{DT}}(\gamma, X)$ is the Donaldson--Thomas
invariant in charge class $\gamma$. Define the \emph{Abel--Jacobi
period}
\[
 \mathrm{AJ}(C)
 \;=\; \int_\Gamma \Omega
 \;\in\; J^2(X) \;=\; H^{2,1}(X)^* \big/ H_3(X, \Z),
\]
where $\Gamma$ is a $3$-chain with $\partial\Gamma = C - C_0$ for a
reference cycle $C_0$ \textup{(}the Griffiths normal function\textup{)}.
Then:
\begin{equation}\label{eq:chiral-volume-conjecture}
 \lim_{N \to \infty} \frac{1}{N}\,
 \log\bigl\lvert Z_N^{\mathrm{ch}}(C, X) \bigr\rvert
 \;=\;
 \frac{1}{2\pi}\,\bigl\lvert \mathrm{AJ}(C) \bigr\rvert,
\end{equation}
where $\lvert\mathrm{AJ}(C)\rvert$ is the norm in the Hodge metric on
$J^2(X)$.
\end{conjecture}

\begin{remark}[Three regimes]
\label{rem:chiral-volume-three-regimes}
The conjecture specialises to three regimes determined by the Hodge
number $h^{2,1}(X)$.

\emph{Toric regime} ($X$ toric, e.g.\ resolved conifold). The
Abel--Jacobi period reduces to the complexified K\"ahler parameter
$t_C = \int_C \omega$, and the conjecture becomes
$\lim (1/N)\,\log\lvert Z_N^{\mathrm{DT}}\rvert = t_C$. This is
consistent with the Okounkov--Reshetikhin--Vafa asymptotics: for the
conifold, $Z_N^{\mathrm{DT}}(N\beta) \sim p(N)\cdot Q^N$ where
$Q = e^{-t}$ and $p(N)$ is the partition function, giving growth rate
$t + O(1/\!\sqrt{N}\,)$.

\emph{Compact regime} ($h^{2,1} > 0$, e.g.\ quintic). The period is a
genuine Griffiths normal function depending on the complex structure
$\tau \in \mathcal{M}_{cs}(X)$. The conjecture predicts a specific
function of $\tau$ (the Abel--Jacobi image) from the
asymptotics of DT invariants. Under mirror symmetry $X \leftrightarrow
X^\vee$, this reduces to the statement that the mirror map
$t(z) = w_1(z)/w_0(z)$ controls DT growth on the mirror side.

\emph{Rigid regime} ($h^{2,1} = 0$, e.g.\ Schoen manifold). The
intermediate Jacobian $J^2(X)$ is trivial, so $\mathrm{AJ}(C) = 0$ for
all curves $C$. The conjecture predicts \emph{subexponential} growth:
$(1/N)\,\log\lvert Z_N^{\mathrm{ch}}\rvert \to 0$. Proxy check: for the
conifold at $Q = 1$, $Z_N = p(N)$ and
$(1/N)\,\log p(N) \sim \pi/\!\sqrt{N} \to 0$ (Hardy--Ramanujan).
\end{remark}

\begin{remark}[Relation to the attractor mechanism]
\label{rem:chiral-volume-attractor}
For BPS black holes in Type~IIB on $X$, the attractor mechanism
(Ferrara--Kallosh--Strominger 1995) fixes the complex structure at a
point $\tau_*(\gamma)$ determined by the charge $\gamma$. At the
attractor point, the BPS entropy is
$S_{\mathrm{BH}}(\gamma) = \pi\,\lvert Z(\gamma, \tau_*)\rvert^2$
where $Z(\gamma, \tau) = \int_\gamma \Omega(\tau)$ is the central
charge (Ooguri--Strominger--Vafa 2004). This is the period integral in
the charge class $\gamma$. The attractor mechanism provides the
analogue of Mostow rigidity: at the attractor point, the complex
structure is \emph{unique}, and the chiral volume conjecture reduces to
a single number. The conjecture is strongest at attractors.
\end{remark}

\begin{remark}[Period growth and obstructions]
\label{rem:chiral-volume-evidence}
\emph{Period growth}: (i) the toric regime is verified to leading order by
the Okounkov--Reshetikhin--Vafa asymptotics; (ii) mirror symmetry
relates DT invariants of $X$ in class $\beta$ to periods of $X^\vee$,
so the mirror map \emph{is} the Abel--Jacobi map on the mirror side;
(iii) the attractor mechanism gives
$\log\Omega_{\mathrm{BPS}}(\gamma) \sim \pi\lvert Z(\gamma)\rvert^2$,
supporting the period--growth link.

The chiral-algebraic formulation, the trace over Fock modules of $A_C$,
lives on the framed object-level locus of
Theorem~\textup{\ref{thm:cy-to-chiral-d3}}. Chain-level convergence of
the trace for non-formal compact CY$_3$ is analytic input; the
normalisation of $\Omega$ is fixed by the Hodge metric; the absence of a
CY analogue of Mostow rigidity makes the statement a family over moduli;
and the large-$N$ limit of DT invariants in a fixed curve class requires
compact-CY$_3$ asymptotics.

\noindent The curve geometry, CY$_3$ data, DT asymptotics,
Abel--Jacobi periods, three-regime analysis, \(3d\) comparison,
obstruction classes, and conjectural hypotheses are the data entering
the large-\(N\) claim.
\end{remark}

\begin{proposition}[Toric small-$N$ convergence for the resolved conifold]
\label{prop:chiral-volume-small-N-conifold}
For the resolved conifold $X = \cO(-1) \oplus \cO(-1) \to \P^1$ with K\"ahler parameter $t = 2\pi s$, $s > 0$, and $Q = e^{-t}$, the charge-$N$ chiral partition function in the unique compact curve class is $Z_N(X, Q) = p(N) \cdot Q^N$ where $p(N)$ is the integer partition function. Then
\[
 \frac{1}{N}\,\log\bigl\lvert Z_N(X, Q) \bigr\rvert \;=\; -t + \frac{\log p(N)}{N},
\]
and $(1/N) \log p(N) \sim \pi \sqrt{2/(3N)} \to 0$ by Hardy--Ramanujan, so
\[
 \lim_{N \to \infty} \frac{1}{N}\,\log\bigl\lvert Z_N(X, Q) \bigr\rvert \;=\; -t \;=\; -\bigl\lvert \mathrm{AJ}(C) \bigr\rvert.
\]
At finite $N$ the convergence rate is $|{-}t - N^{-1}\log p(N)|$:
\begin{center}
\small\renewcommand{\arraystretch}{1.1}
\begin{tabular}{rrrr}
\toprule
$N$ & $N^{-1} \log p(N)$ & Hardy--Ramanujan $\pi\sqrt{2/(3N)}$ & ratio \\
\midrule
$10$ & $0.3738$ & $0.8112$ & $0.46$ \\
$100$ & $0.1907$ & $0.2565$ & $0.74$ \\
$1000$ & $0.0723$ & $0.0811$ & $0.89$ \\
\bottomrule
\end{tabular}
\end{center}
The rate matches the conjectured $|\mathrm{AJ}(C)|/(2\pi) \cdot 2\pi = t$ (absolute value) to leading order at every $N$, with sub-leading $O(N^{-1/2})$ Hardy--Ramanujan correction.
\end{proposition}

\begin{proof}
The $\mathrm{DT}$ partition function of the resolved conifold in curve class $N[\P^1]$ is $\mathrm{DT}_N(Q) = p(N) Q^N$ (Okounkov--Reshetikhin--Vafa, \emph{Quantum Calabi--Yau and Classical Crystals}, Theorem~2). Taking logarithms gives the stated formula. The Hardy--Ramanujan expansion $\log p(N) = \pi \sqrt{2N/3} + O(\log N)$ yields $(1/N) \log p(N) \sim \pi \sqrt{2/(3N)}$. The complexified K\"ahler parameter $t = \int_{[\P^1]} \omega$ coincides with $|\mathrm{AJ}(C)|$ for the unique toric curve class since $h^{2,1}(X) = 0$.
\end{proof}

\begin{proposition}[Rigid-vertical small-$N$: subexponential rate for K3 classes]
\label{prop:chiral-volume-rigid-vertical}
For $X = \mathrm{K3}$ with curve class $\beta \in H_2(\mathrm{K3}, \Z)$ of self-intersection $\beta^2 = 2h - 2$, the vertical $\mathrm{PT}$ partition function is $\mathrm{PT}(X, N\beta; q) = q^{hN}/\prod_{k \geq 1} (1 - q^k)^{24}$ (Kawai--Yoshioka / Maulik--Pandharipande--Thomas). Its Fourier coefficients $c_N$ are the coefficients of $\prod(1-q^k)^{-24}$; Cardy asymptotics give $\log c_N \sim 4\pi \sqrt{N}$, so
\[
 \frac{1}{N} \log c_N \;\to\; 0 \quad (N \to \infty),
\]
matching the conjecture's \emph{rigid regime} prediction $|\mathrm{AJ}(C)| = 0$ on vertical classes. Small-$N$ values: $c_1 = 24$, $c_2 = 324$, $c_3 = 3200$, $c_5 = 176256$, with $(1/N)\log c_N$ decreasing monotonically from $3.18$ at $N = 1$ to $1.21$ at $N = 50$, and Cardy ratio converging to $1$ as $N^{1/2} \to \infty$.
\end{proposition}

\begin{proof}
The generating-function identity is the Maulik--Pandharipande--Thomas theorem (2010, \emph{Topology}); the Cardy asymptotic $\log c_N \sim 4\pi\sqrt{N}$ follows from the circle-method expansion of $\eta^{-24}$ (Rademacher formula, applied to weight $-12$ eta-quotient). Vertical $\mathrm{AJ}(\beta) \in J^2(\mathrm{K3} \times E)$ lies in the $E$-factor and integrates to zero on $\beta$ (by the K\"unneth decomposition $\beta \in H_2(\mathrm{K3}) \otimes 1$), giving $|\mathrm{AJ}(\beta)| = 0$. Consistency with the conjecture is the subexponential rate.
\end{proof}

\section{The K3 chiral quantum-group \texorpdfstring{$R$}{R}-matrix, associator, and classification}
\label{sec:qca-k3-chiral-quantum-group}

The K3 chiral Hall--Drinfeld double
$\mathcal{D}_\hbar(\mathcal{Y}^{\mathrm{Hall}}_\hbar(\mathrm{CoHA}_{K3\times E}))$
of Chapter~\ref{ch:k3-chiral-bialgebra-platonic} occupies a fourth quasi-Hopf
taxonomic slot. Five structural data pin its quantum-group content:
the Siegel-elliptic dynamical $R$-matrix,
the twisted Siegel--Borcherds associator,
a super-quasi-Hopf compatibility,
an $A_\infty$-quasi-Hopf upgrade at the Humbert walls,
and the classification of its gauge-equivalence class by a $1$-dimensional
Lie-bialgebra cohomology group.

The ambient-qualifier discipline is twofold. At chain level
(explicit complexes, named differentials, Chevalley--Eilenberg cochains,
pentagon 3-cocycle witness $d_{\mathrm{CE}}\phi^{(3)} = -\frac{1}{2}[\phi^{(2)},\phi^{(2)}]_{\mathrm{Ger}}$),
the structure computes; at $(\infty,1)$-categorical level (derived
categories of $D$-modules on $\overline{\mathcal{A}_2}$, Lurie $\mathrm{HA}$
formalism of $E_2$-algebras in a presentable $\infty$-category), the structure
universalises. Both lanes appear simultaneously: the explicit Pasol--Zagier
Kronecker--Eisenstein series on $\mathbb{H}_2$ (chain-level), the
Etingof--Kazhdan super-category extension (categorical), the Bruinier 2002
Chern-class reciprocity (cohomological), the Felder dynamical YBE on the
paramodular quotient (geometric), and the Markl--Shnider--Stasheff
$A_\infty$-pentagon at Humbert walls (operadic). Each appearance is one face
of the same five-datum object.

The five theorems and five bridging propositions that follow make each datum
precise, with arguments citing primary literature.

\begin{theorem}[Siegel-elliptic dynamical \texorpdfstring{$R$}{R}-matrix: a fourth class in the Drinfeld--Belavin taxonomy]
\label{thm:qca-R-Sieg-dynamical}
Let $A_\tau$ be the principally polarised abelian surface with period matrix
$\bigl(\begin{smallmatrix}\tau&z\\z&\rho\end{smallmatrix}\bigr)\in\mathbb{H}_2$.
The Siegel-elliptic dynamical $R$-matrix
$R_{\mathrm{Sieg,dyn}}(u-v;\tau,\rho,z)\in\mathrm{End}(V_{24}\otimes V_{24})$
of Chapter~\ref{ch:k3-chiral-bialgebra-platonic} Definition~\ref{def:R-Sieg-dynamical},
constructed from the genus-$2$ Riemann theta function
$\theta[{\alpha\atop\beta}](w;\tau,z,\rho)$ via the Pasol--Zagier 2013 extension of
Kronecker--Eisenstein to $\mathbb{H}_2$, occupies a fourth structural slot beyond
Drinfeld 1985's rational/trigonometric/elliptic classification: it is neither
rational (no $1/(u-v)$-pole structure), nor trigonometric (no $\mathbb{C}^*$-loop),
nor purely elliptic (no single elliptic parameter $\tau$). The fourth class is
Siegel-elliptic-dynamical: dynamical parameter in $\mathbb{H}_2$, spectral parameter
separate, with the three-tier degeneration
$R_{\mathrm{Sieg,dyn}}\xrightarrow{\rho\to i\infty} R^{\mathrm{ell,dyn}}_{\mathrm{Felder}}
\xrightarrow{\tau\to i\infty} R^{\mathrm{rat}}(u)$
recovering the Felder 1994 elliptic dynamical and Yang rational $R$-matrices as
Humbert-cusp and double-Humbert-cusp degenerations.

On the paramodular locus $K(1)\subset\mathbb{H}_2$,
$R_{\mathrm{Sieg,dyn}}$ satisfies the dynamical Yang--Baxter equation
\[
 R_{12}(u_1-u_2)\,R_{13}(u_1-u_3)\,R_{23}(u_2-u_3)
 \;=\;
 R_{23}(u_2-u_3)\,R_{13}(u_1-u_3)\,R_{12}(u_1-u_2)
\]
modulo shifts of dynamical parameter $(\tau,\rho,z)\mapsto(\tau,\rho,z)+h_i$
at Cartan-generator insertions (the Felder dynamical shift, extended to
$\mathrm{Sp}_4$-compatible period matrix). Off $K(1)$, the YBE acquires a
half-integer Jacobi-index anomaly proportional to the Maass multiplier
$v_{\Delta_5}$ of $\Delta_5$ as a half-integral-weight paramodular form.
\end{theorem}

\begin{proof}
The fourth-class statement is Chapter~\ref{ch:k3-chiral-bialgebra-platonic},
Theorem~\ref{thm:R-sieg-fourth-class}; the dynamical YBE on $K(1)$ is
Theorem~\ref{thm:R-Sieg-YBE} of the same chapter. The three-tier degeneration is
Remark~\ref{rem:R-Sieg-felder-genus1} (genus-$1$ Felder limit at $\rho\to i\infty$,
rational limit at $\tau\to i\infty$). The classical $\hbar^1$-YBE follows from
Pasol--Zagier 2013 Theorem~3.2 cyclic Siegel-Kronecker identity; the quantum
$\hbar^2$-YBE follows from Etingof--Kazhdan quantisation of the Manin pair, as
in Chapter~\ref{ch:k3-chiral-bialgebra-platonic}
Proposition~\ref{prop:pasol-zagier-to-dynamical-YBE}. The Felder dynamical
mechanism--dynamical parameter shift by simple root $h_i$ at each Cartan
insertion--is the genus-$2$ lift of Felder 1994 \S 4, with the Cartan eigenvalue
$(\alpha,\alpha)$ supplying the shift magnitude through the cross-bracket
$[t_{12},e_\alpha]=(\alpha,\alpha)\,e_\alpha$.
\end{proof}

\begin{theorem}[Pentagon and hexagon for the Siegel--Borcherds associator]
\label{thm:qca-pentagon-hexagon-sieg-bor}
The twisted associator
$\widetilde{\Phi}^{\mathrm{Sieg\text{-}Bor}}_{\mathrm{Sp}_4}[\Phi_{10}^{\mathrm{un}}/\eta^{24}]$
lives in the pro-nilpotent pro-$\hbar$-filtered completion
\[
 \widetilde{\Phi}^{\mathrm{Sieg\text{-}Bor}}_{\mathrm{Sp}_4}[\Phi_{10}^{\mathrm{un}}/\eta^{24}]
 \;\in\;
 \exp\bigl(\widehat{\mathfrak{t}^{\mathrm{Sieg}}_{2,[2]}\oplus\mathfrak{n}_+^{\mathrm{imag}}}\bigr)^{\mathrm{grouplike}},
\]
with $\mathfrak{t}^{\mathrm{Sieg}}_{2,[2]}$ the genus-$2$ infinitesimal pure-braid
Lie algebra (Enriquez--Gomez-Gonzalez--Maassarani 2022) and
$\mathfrak{n}_+^{\mathrm{imag}}$ the BKM imaginary-root nilpotent.

The pentagon equation holds at orders $\hbar^{\leq 3}$. The order-$\hbar^2$ part
reduces to the closedness condition $d_{\mathrm{CE}}\phi^{(2)} = 0$, which holds
by the Kohno--Drinfeld four-term relation on the symmetric part and the BKM
Jacobi identity on the imaginary part. At order $\hbar^3$, the obstruction
correction
\[
 \phi^{(3)}
 \;=\;
 \zeta(3)\,c_{\mathrm{symm}}
 \;+\;
 \tfrac{26}{3}\,c_{\mathrm{timelike}}
 \;+\;
 \bigl(\Phi_{10}^{\mathrm{un}}/\eta^{24}\bigr)\,c_{\Phi_{10}^{\mathrm{un}}},
\]
with $c_{\mathrm{symm}},c_{\mathrm{timelike}},c_{\Phi_{10}^{\mathrm{un}}}$ three
$3$-cochains, linearly independent in the coboundary space
$B^3 \subset C^3(\mathfrak{t}^{\mathrm{Sieg}}_{2,[2]}\oplus\mathfrak{n}_+^{\mathrm{imag}};\,\mathcal{O}(\mathbb{D}))$
with function coefficients over the period domain $\mathbb{D}$ (the third
coefficient is the modular function $\Phi_{10}^{\mathrm{un}}/\eta^{24}$;
restricted to a fixed chamber it is read as its scalar value),
balances the Gerstenhaber self-bracket $-\frac{1}{2}[\phi^{(2)},\phi^{(2)}]_{\mathrm{Ger}}$
term by term in the Borcherds product expansion of $\Phi_{10}^{\mathrm{un}}$ at Fourier level
$c(N,\ell)$ (Borcherds 1998, Theorem~13.3; Gritsenko--Nikulin 1997, Theorem~1.2).
The constant $26/3$ decomposes as $\mathrm{rk}(\mathrm{II}_{25,1})/3 = 26/3$ (Fake Monster
Cartan rank, $\mathrm{rk}(\mathrm{II}_{25,1}) = 26$, divided by cyclic triple-leg
averaging), not a Virasoro central charge.

Both hexagons at $\hbar^{\leq 2}$ hold for the pair
$(\widetilde{\Phi}^{\mathrm{Sieg\text{-}Bor}},R_{\mathrm{Sieg,dyn}})$ on
$K(1)\subset\overline{\mathcal{A}_2}$; restriction to $K(1)$ is forced by the
half-integer Jacobi-index anomaly of $\Delta_5$ on $\mathrm{Sp}_4(\mathbb{Z})$.
\end{theorem}

\begin{proof}
The pentagon is Chapter~\ref{ch:k3-chiral-bialgebra-platonic},
Theorem~\ref{thm:pentagon-sieg-bor}, together with the KZ coboundary lift and
Borcherds identity propositions (the order-$\hbar^{\leq 2}$ Knizhnik--Zamolodchikov pentagon proposition
and the order-$\hbar^3$ Borcherds-pentagon proposition). The hexagon is
Theorem~\ref{thm:hexagon-R-Sieg}, decomposing at order
$\hbar^2$ into: (H-I) Pasol--Zagier 2013 cyclic Siegel-Kronecker compatibility;
(H-II) Maass-form compatibility on $K(1)$ via Lorgat 2020 Proposition~4.1; and
(H-III) the Casimir--imaginary cross-term grouplike-exponential compatibility.
The $26/3$ arises in the order-$\hbar^3$ Borcherds-pentagon proposition via Nikulin
1979 Lorentzianisation $\Gamma^{4,20}\hookrightarrow\mathrm{II}_{25,1}$
(of rank $26$) together with cyclic triple-leg $|\mathbb{Z}_3|^{-1}$
averaging.
\end{proof}

\begin{theorem}[Super-quasi-Hopf structure and \texorpdfstring{$A_\infty$}{A-inf} upgrade at Humbert walls]
\label{thm:qca-super-A-infty-upgrade}
The $\mathbb{Z}/2$-grading on the K3 chiral Hall--Drinfeld double is genuine
super-grading (Koszul sign rule on the pentagon, hexagon, and antipode axioms),
sourced from the K3 sigma-model's Ramond--Ramond fermion number and lifted
through the $\widetilde{M}_{24}$ cover to BKM imaginary-root parity: odd imaginary
simple roots $\alpha$ with $(\alpha,\alpha)<0$ are fermionic generators, bosonic
when $(\alpha,\alpha)=0$. The Igusa form $\Delta_5$ lives in odd weight $5$;
$\Phi_{10}^{\mathrm{un}}=\Delta_5^2$ in even weight $10$.

Along the Humbert divisors $H_1\cup H_4\subset\overline{\mathcal{A}_2}$, the
Siegel--Borcherds associator degenerates; the strict quasi-Hopf pentagon at
$\hbar^{\leq 3}$ upgrades to an $A_\infty$-quasi-Hopf pentagon with higher-coherence
data $\phi^{(n)},\,n\geq 4$, satisfying the Markl--Shnider--Stasheff $A_\infty$
pentagon equation
\[
 d_{\mathrm{CE}}\,\phi^{(n)}
 \;=\;
 -\,\tfrac{1}{2}\sum_{\substack{k+l=n+1\\ k,l\geq 2}}
 \bigl[\phi^{(k)},\phi^{(l)}\bigr]_{\mathrm{Ger}},
\]
with the $n$-th higher-coherence datum carrying the $n$-th Fourier-Jacobi
coefficient of the expansion
$\Phi_{10}^{\mathrm{un}}(\rho,\tau,z) = \sum_{m\geq 1}\phi_{10,m}(\tau,z)e^{2\pi im\rho}$
(Eichler--Zagier Fourier-Jacobi theory).
\end{theorem}

\begin{proof}
The genuineness of the super-grading is Chapter~\ref{ch:k3-chiral-bialgebra-platonic}
Remark~\ref{rem:super-genuine-fermion}; the lift through the K3 sigma model and
$\widetilde{M}_{24}$ cover uses Gaberdiel--Hohenegger--Volpato 2012 for the
$M_{24}$-equivariance of the K3 elliptic genus, and Gritsenko--Nikulin 1995
\S 2 for the BKM imaginary-root fermionic/bosonic classification by
$(\alpha,\alpha)$-sign. The $A_\infty$ upgrade at Humbert walls is
Theorem~\ref{thm:A-infty-upgrade} together with
the $A_\infty$ higher-coherence Humbert proposition. The recursion at depth
$n$ is Markl--Shnider--Stasheff 2002 Chapter~8, with the Maurer--Cartan equation
in the Chevalley--Eilenberg DG Lie algebra of
$(\mathfrak{t}^{\mathrm{Sieg}}_{2,[2]}\oplus\mathfrak{n}_+^{\mathrm{imag}})[[\hbar]]$
supplying the higher-pentagon constraint.
\end{proof}

\begin{theorem}[Pentagon coboundary decomposition at order \texorpdfstring{$\hbar^{\leq 3}$}{hbar-<= 3}]
\label{thm:qca-pentagon-hbar3-explicit}
Let $\mathfrak{L} := \mathfrak{t}^{\mathrm{Sieg}}_{2,[2]} \oplus \mathfrak{n}_+^{\mathrm{imag}}$ denote the home Lie algebra of the Siegel--Borcherds associator. The three independent $3$-coboundaries on $\mathfrak{L}$ are
\[
 c_{\mathrm{symm}}
 \;:=\;
 [t_{12},t_{23}] \otimes [t_{23}, t_{13}],
 \qquad
 c_{\mathrm{timelike}}
 \;:=\;
 \sum_{\alpha\in\mathrm{II}_{25,1}^{\mathrm{timelike}}} \alpha\otimes\alpha\otimes\alpha,
 \qquad
 c_{\Phi_{10}^{\mathrm{un}}}
 \;:=\;
 \sum_{\substack{(N,\ell,M)\\ 4NM-\ell^2=-1}} c_{\phi_{0,1}}(N,\ell)\, e_\alpha\otimes e_\beta\otimes e_\gamma,
\]
with $t_{ij}$ the Drinfeld--Kohno infinitesimal braiding, $\mathrm{II}_{25,1}^{\mathrm{timelike}}$ the timelike component of the Fake Monster Cartan, and $c_{\phi_{0,1}}$ the $(N,\ell)$-th Fourier coefficient of the K3 weak Jacobi form. The pentagon equation at order $\hbar^{\leq 3}$ reads
\[
 \phi^{(3)}
 \;=\;
 \zeta(3)\,c_{\mathrm{symm}}
 \;+\;
 \tfrac{26}{3}\,c_{\mathrm{timelike}}
 \;+\;
 (\Phi_{10}^{\mathrm{un}}/\eta^{24})\,c_{\Phi_{10}^{\mathrm{un}}},
\]
where the constant $26/3 = \mathrm{rk}(\mathrm{II}_{25,1})/3$ is the Fake Monster Cartan rank $26$ divided by cyclic triple-leg averaging, and the third coefficient is the modular function $\Phi_{10}^{\mathrm{un}}/\eta^{24}$ on the period domain: the decomposition is read in $C^3(\mathfrak{L};\,\mathcal{O}(\mathbb{D}))$ with function coefficients, its scalar form being the evaluation in a fixed chamber. The three cochains are linearly independent in the coboundary space $B^3(\mathfrak{L};\,\mathcal{O}(\mathbb{D})) \subset C^3$ (as coboundaries they vanish in cohomology; independence holds at the cochain level), so the decomposition is unique.
\end{theorem}

\begin{theorem}[Conditional Etingof--Kazhdan characteristic comparison]
\label{thm:qca-etingof-kazhdan-super-classification}
Assume two compact Hall--Drinfeld source packages over $K3\times E$ have
the same finite Hall/CoHA source window, positive and negative halves,
Cartan part, continuous Hopf pairing, source-to-target comparison maps,
radical descent, no-extra-relations kernel equality, PBW comparison, and
exact inverse-limit passage. Assume moreover that their primitive maps
land in the Manin-pair target
$(\mathfrak{g}_{\Delta_5},\,\mathfrak{n}_+^{\mathrm{imag}}\oplus
\mathfrak{h}^{\mathrm{imag,rk\,23}})$ with the super structure and the
paramodular $K(1)$-restriction of Theorem~\ref{thm:hexagon-R-Sieg}.
Then the comparison data define a characteristic map
\[
 \chi^{\mathrm{cmp}}_{\Delta_5}\colon
 \Def_{\mathrm{src}\to\Delta_5}(\mathbf H)
 \longrightarrow
 H^2_{\mathrm{Lie}}\bigl(\mathfrak q_{K(1)};\mathrm{grt}_1\bigr)
 \;\cong\;
 \mathbb{C}\cdot\Delta_5.
\]
Here $\mathfrak q_{K(1)}$ is the $K(1)$-restricted target deformation
Lie algebra of the recognised Manin-pair package. The two quantisations
are compared by their finite Hall--Borcherds matrices and by their image
under $\chi^{\mathrm{cmp}}_{\Delta_5}$; the scalar Igusa line is the
automorphic target of the comparison, not a direct computation of
the intrinsic second Lie-bialgebra cohomology of
$\mathfrak{g}_{\Delta_5}$ and not a construction of $\mathbf H_{\Delta_5}$
from automorphy alone.
\end{theorem}

\begin{proof}
Etingof--Kazhdan Parts~IV--V classify quasi-Hopf quantisations once the
completed super Manin-pair input and associator class have been fixed.
They do not construct the compact Hall source. The finite Hall--Borcherds
recognition gates supply that source: comparison matrices, radical
descent, no-extra-relations kernel equality, PBW compatibility, and strict
completion transitions. Bruinier 2002 Proposition~5.1 supplies the
Heegner-divisor Chern-class reciprocity which places the automorphic
class of $\Delta_5$ in the transgressed paramodular obstruction line.
Gritsenko--Nikulin and Borcherds supply the denominator arithmetic and
weight. These automorphic facts are target conditions for the recognised
source package, not a Lie-bialgebra cohomology computation by
automorphy.
\end{proof}

\begin{remark}[Three faces of the number \texorpdfstring{$8$}{8}]
\label{rem:qca-three-faces-of-eight}
The number $8$ appears in three structurally distinct roles: $(\mathrm{a})$ categorical Mukai doubling $2\,c_+(\mathrm{Mukai}(K3))=2\cdot 4 = 8$; $(\mathrm{b})$ the branch-normalised local-exponent denominator of the automorphic line of $\Delta_5$ around the Humbert divisor $H_1\subset\overline{\mathcal{A}_2}$; $(\mathrm{c})$ the Lusztig specialisation order $\ell=8$ of the Hall--Drinfeld double, when that specialisation is constructed. Bruinier~$2002$ Proposition~$5.1$ controls the automorphic Chern class. It does not identify a Chow class with a small-quantum-group cohomology class, and it does not imply a canonical equation for $\hbar^2$.
\end{remark}

\begin{theorem}[Root-of-unity comparison at \texorpdfstring{$\ell=8$}{l=8}]
\label{thm:qca-hbar-minus-eighth-lusztig}
The Lusztig specialisation of the Hall--Drinfeld double at $\ell=8$ produces a
distinguished root-of-unity object at $\zeta = e^{2\pi i/8}$. Bruinier 2002, Proposition~5.1
(Chern-class formula for Heegner divisors) computes the automorphic Chern class
\[
 c_1\bigl(\mathcal{L}^{\Delta_5}\bigr)\big|_{H_1}
 \;=\;
 [\xi]
 \;\in\;
 \mathrm{CH}^1(H_1)\otimes\mathbb{Q}.
\]
After a primitive branch is chosen, its local exponent may have denominator $8$. This produces a comparison with the Lusztig order $\ell=8$ and with the Mukai integer $2c_+=8$. The formal parameter itself is not fixed by this comparison: if $q=\exp(\hbar)$, then $\hbar$ is a logarithm of $\zeta$ and is defined only modulo $2\pi i\mathbb{Z}$.
\end{theorem}

\begin{proof}
The Bruinier calculation is a calculation on the automorphic parameter space. The Lusztig construction is a root-of-unity specialisation of the exponentiated quantum parameter. A comparison between them must be supplied by the Hall--Borcherds recognition data. Without this comparison map, the two order-eight records are independent records with the same integer.
\end{proof}

\begin{theorem}[Recognised characteristic invariant: the Igusa cusp form]
\label{thm:qca-classification-invariant}
On the finite Hall--Borcherds recognition locus of
Theorem~\ref{thm:qca-etingof-kazhdan-super-classification}, the
completed source deformation complex has a Bruinier--Heegner
characteristic line
\[
 \chi^{\mathrm{cmp}}_{\Delta_5}\bigl(\Def_{\mathrm{src}\to\Delta_5}(\mathbf H)\bigr)
 \subset
 H^2_{\mathrm{Lie}}\bigl(\mathfrak q_{K(1)};\mathrm{grt}_1\bigr)
 \;\cong\;
 M^{!,\mathrm{odd}}_{5}(K(1))\,/\,(\text{periods})
 \;\cong\;
 \mathbb{C}\cdot\Delta_5.
\]
Two recognised Hall--Drinfeld packages with the same finite comparison
matrices and the same image in this line give gauge-equivalent
Etingof--Kazhdan quantisations. The Igusa cusp form is the scalar
automorphic characteristic of the recognised source; it is not, by
itself, an intrinsic classification of all quasi-Hopf structures on
$\mathfrak{g}_{\Delta_5}$.
\end{theorem}

\begin{proof}
Apply the conditional comparison theorem of
Chapter~\ref{ch:k3-chiral-bialgebra-platonic} to a completed source
package which has already passed the finite Hall--Borcherds gates. The
Etingof--Kazhdan super-category extension (Etingof 2002, \S 6.5;
Gurevich--Majid 1994 for super-Hopf axioms; Etingof--Kazhdan 1996--2000
Parts~I--V of the quantisation theorem) transports the fixed
super-Manin input to a quasi-Hopf output. Bruinier--Heegner
Chern-class reciprocity and the Gritsenko--Nikulin denominator identify
the scalar automorphic target with the $\Delta_5$ line. The finite
source matrices and PBW/no-extra-relations checks are the source
recognition; the automorphic line is only the characteristic target.
\end{proof}

\begin{remark}[Seven lenses on one object]
\label{rem:qca-seven-lenses}
The K3 chiral Hall--Drinfeld double is not seven separate facts but one
fact seen through seven lenses; they are enumerated here and displayed
simultaneously.
\begin{enumerate}[label=\textup{(L\arabic*)}]
\item \textbf{Gritsenko additive $\leftrightarrow$ Weyl--Kac--Borcherds denominator.}
 $\mathrm{Grit}(\eta^9\vartheta_1) = \Delta_5$ is the Fourier--Jacobi face of the
 Weyl--Kac character formula for $\mathfrak{g}_{\Delta_5}$
 (Chapter~\ref{ch:k3e-bkm}, Remark~\ref{rem:k3e-gritsenko-seeds-WKB}).
\item \textbf{Maass $\mathbb{Z}/2$-spin cover $\leftrightarrow$ Jacobi half-integral-weight paramodular.}
 $\widehat{\mathrm{Sp}_4(\mathbb{Z})}^{v_{\Delta_5}}$ is the paramodular metaplectic
 cover where half-integral Jacobi-index theory lives
 (Theorem~\ref{thm:Maass-spin-cover}).
\item \textbf{Arthur packet $\psi_{\Delta_{10}}$ $\leftrightarrow$ $R_{\mathrm{Sieg,dyn}}$ via Langlands dual group.}
 The Arthur packet's Langlands dual group at the spin cover is $\mathrm{Sp}_4$ itself,
 and the local $p$-adic Hecke category is the Kazhdan--Lusztig category of the
 $p$-adic loop group, with $R_{\mathrm{Sieg,dyn}}$ its KL $R$-matrix
 (Chapter~\ref{ch:k3-chiral-bialgebra-platonic},
 Remark~\ref{rem:arthur-packet-langlands-dual-R-matrix}).
\item \textbf{Pasol--Zagier $H_2$ Kronecker--Eisenstein $\leftrightarrow$ elliptic dynamical YBE on $K(1)$.}
 The Siegel Kronecker--Eisenstein cyclic identity is simultaneously the classical $r$-matrix
 Yang--Baxter equation and the pentagon Maurer--Cartan defect
 (Chapter~\ref{ch:k3-chiral-bialgebra-platonic},
 Remarks~\ref{rem:pasol-zagier-H2-associator}
 and~\ref{rem:elliptic-dynamical-YBE-K1}).
\item \textbf{Shimura--Waldspurger $\leftrightarrow$ Ikeda-vs-Gritsenko distinction.}
 The Gritsenko and Borcherds lifts produce the same $\Delta_5$; the Shimura--Waldspurger
 correspondence is the automorphic witness. Ikeda produces $\Delta_{10}$ only; the
 Gritsenko half-integer-index lift has no Ikeda analogue
 (Chapter~\ref{ch:k3-chiral-bialgebra-platonic},
 Remark~\ref{rem:ikeda-gritsenko-disambiguation}).
\item \textbf{Humbert monodromy $\leftrightarrow$ Lusztig root-of-unity specialisation.}
 The order-eight Humbert datum, the Mukai integer \(2c_+(\mathrm{Mukai}(K3))=8\),
 and the Lusztig order \(\ell=8\) are compared on the fixed K3 branch.
 Bruinier's Heegner Chern-class formula supplies the automorphic side only.
\item \textbf{Conditional characteristic $\leftrightarrow$ 1-loop-anomaly output.}
 After finite Hall--Borcherds source recognition, the comparison class maps to
 the one-dimensional Bruinier--Heegner characteristic line
 $\mathbb{C}\cdot\Delta_5$. The same scalar $\Delta_5$ is the
 1-loop anomaly-cancellation output of twisted 11D supergravity on
 $K3\times T^2$ under the protected comparison hypothesis
 (Chapter~\ref{ch:k3-chiral-bialgebra-platonic},
 Theorems~\ref{thm:k3-classification-etingof}
 and~\ref{thm:k3-1loop-output}).
\end{enumerate}
Each lens refracts the same underlying object; the interlocks between them
are mathematical. Lens~(L1) seeds imaginary-root multiplicity
(Humbert monodromy, lens~L6); lens~(L2) seeds the paramodular
$K(1)$-restriction that lens~(L4) requires for the hexagon to close;
lens~(L3) seeds the Langlands-reciprocity characteristic comparison
that lens~(L7) completes after source recognition. No single lens is
sufficient; all seven are necessary. The
unity of the object is the joint coherence of the seven lenses.
\end{remark}

\section{Source-recognised \texorpdfstring{$\mathbf H_{\Delta_5}$}{H-Delta-5} as quantum chiral algebra}
\label{sec:qca-hdelta5-quantum-chiral-algebra}

\begin{remark}[Source-recognised \texorpdfstring{$\mathbf H_{\Delta_5}$}{H-Delta-5} as the fourth taxonomic slot]
\label{rem:qca-hdelta5-fourth-slot}
The quantum chiral algebra construction surveyed above (Drinfeld--Jimbo,
Drinfeld Yangian, RTT in the chiral setting) reaches its K3 master
form only on the source-recognised locus where the finite
Hall--Borcherds gates construct the non-abelian chiral bialgebra
$\mathbf H_{\Delta_5}$. As a quantum chiral algebra on that locus,
$\mathbf H_{\Delta_5}$ has:
\begin{itemize}
\item \emph{Coalgebra side:} bar complex $B_{\mathrm{ch}}(\mathbf H_{\Delta_5})$
 in the Hall subcategory of $\mathrm{CoHA}_{K3\times E}$;
\item \emph{Algebra side:} cobar $\Omega_X(B_{\mathrm{ch}}(\mathbf H_{\Delta_5}))\simeq\mathbf H_{\Delta_5}$
 on the chosen order-eight root-of-unity branch;
\item \emph{$R$-matrix side:} Siegel-elliptic dynamical $R_{\mathrm{Sieg,dyn}}(z_1,z_2;\tau,w)$
 satisfying the dynamical Yang-Baxter equation on $\bar{\mathcal A}_2$;
\item \emph{Associator side:} twisted Siegel-Borcherds associator
 $\widetilde\Phi^{\mathrm{Sieg\text{-}Bor}}_{\mathrm{Sp}_4}[\Phi_{10}^{\mathrm{un}}/\eta^{24}]$
 satisfying the pentagon at order $\hbar^{\le 3}$ with coboundary
 $\phi^{(3)} = \zeta(3) c_{\mathrm{symm}} + \tfrac{25}{3} c_{\mathrm{timelike}} + (\Phi_{10}^{\mathrm{un}}/\eta^{24}) c_{\Phi_{10}^{\mathrm{un}}}$.
\end{itemize}
These are the four target-facing data of the recognised package for
$\mathbf H_{\Delta_5}$. The order-eight comparison fixes a branch after
the source package has been supplied; it does not replace the Hall
source construction. Cross-ref
Chapter~\ref{ch:k3-chiral-bialgebra-platonic}.
\end{remark}

\begin{remark}[\texorpdfstring{$\Psi$}{Psi} as a functor of operads: operadic compatibility]
\label{rem:qca-psi-functor-of-operads}
\index{$\Psi$!functor of operads}
\index{operadic compatibility!Psi}
The universal-trace functor $\Psi$ of
Chapter~\ref{ch:phi-universal-trace-platonic},
Theorem~\ref{thm:phiutp-psi-functor-lax-symmon}, extends the recognised
source-package rule $(L, \phi_L) \mapsto \mathbf{H}_{\phi_L}$ to a lax
symmetric-monoidal functor of symmetric monoidal $(\infty,1)$-categories
\[
 \Psi\colon (\mathrm{CY}^{\mathrm{Siegel\text{-}aut}}_2, \oplus)
 \longrightarrow
 (\mathrm{QHopf}^{\mathrm{BKM}}, \otimes),
\]
with Jacobi-form K\"unneth $\phi_1 \boxtimes \phi_2$ on the domain, and
Drinfeld--Reshetikhin--Turaev tensor product $R_1 \otimes R_2$,
$\Phi_1 \otimes \Phi_2$ on the codomain. Applied at the K3 input,
$\Psi(\Lambda^{3,2}, \phi_{K3}) = \mathbf{H}_{\Delta_5}$ only after
the finite Hall--Borcherds source gates and Bruinier--Heegner
characteristic comparison have been supplied; this is the quantum
chiral algebra of
Remark~\ref{rem:qca-hdelta5-fourth-slot}. The operadic compatibility of
$\Psi$ implies that $\mathbf{H}_{\Delta_5}$ is not an isolated construction
but the recognised value at one object of a functor respecting composition:
the four data (coalgebra bar, algebra cobar, $R$-matrix, associator) listed
in Remark~\ref{rem:qca-hdelta5-fourth-slot} are functorial in $L$ after
source recognition,
with direct-sum $L_1 \oplus L_2$ mapping to tensor product
$\mathbf{H}_{\phi_1} \otimes \mathbf{H}_{\phi_2}$ of the four-fold data
up to natural isomorphism. The rank-$4$ BKM instance
(Proposition~\ref{prop:phiutp-rank-four}) identifies this at
$\mathrm{II}_{1,1} \oplus \mathrm{II}_{1,1} = \mathrm{II}_{2,2}$ giving
$\mathbf{H}_{\mathrm{Monster}} \otimes \mathbf{H}_{\mathrm{Monster}}$.
Primary-lit: Borcherds 1998 \S14 (K\"unneth for singular theta
correspondence on lattice direct sums), Dunn 1988 Theorem~3.4
(additivity $E_1 \otimes_{\mathrm{BV}} E_1 \simeq E_2$),
Reshetikhin--Turaev 1991 Proposition~1.5, Lurie \emph{Higher Algebra}
5.3.2.1.
\end{remark}

\begin{remark}[\texorpdfstring{$E_2 \to E_1$}{E-2 -> E-1} Dunn stabilisation of \texorpdfstring{$\Psi$}{Psi} on the \texorpdfstring{$\mathbf H_{\Delta_5}$}{H-Delta-5} branch]
\label{rem:qca-psi-dunn-stabilisation-hd5}
\index{Dunn stabilisation!Psi E_2 to E_1}
\index{Psi!Dunn stabilisation}
On the source-recognised quantum-chiral-algebra branch that hosts
$\mathbf H_{\Delta_5}$, the operadic compatibility of $\Psi$ takes the
specific form of an $E_2 \to E_1$ Dunn stabilisation
(Theorem~\ref{thm:phiutp-e2-to-e1-dunn}):
the domain carries an $E_2$-structure via Fubini on the two commuting
pairings $(\oplus, \boxtimes)$, while the codomain
$(\mathrm{QHopf}^{\mathrm{BKM}}, \otimes)$ carries only an
$E_1$-structure because Eckmann--Hilton would force any second
compatible pairing on $\mathrm{QHopf}$ to collapse each object to a
commutative algebra, contradicting non-abelian BKM. The forgetful
$E_2 \to E_1$ collapse is \emph{not} information loss: the coalgebra
structure on $\mathbf H_{\Delta_5}$ (equivalently, the Jacobi-index
coupling on $\phi_{K3}$) is absorbed \emph{internal} to each BKM output
object, rather than surviving as a second external pairing on
$\mathrm{QHopf}^{\mathrm{BKM}}$. The operadic compatibility therefore
translates: external $E_2$ on the input (lattices + Jacobi forms) becomes
internal coalgebra + external $E_1$ on the output (quantum chiral
algebras, tensor product). This is the precise operadic content of the
Koszul reflection acting through $\Psi$: the reflection halves the
operadic degree as it pushes from the CY-Siegel side to the BKM-chiral
side, consistent with the $[d]$-shift on $\Phi$ of
Theorem~\ref{thm:phi-output-shift-k3-2}.
\end{remark}

\begin{remark}[Etingof-Kazhdan Part V after source recognition]
\label{rem:qca-2}
The Etingof--Kazhdan 2007 Part~V super-category extension of the
quantisation theorem (for Lie bialgebras over a Grothendieck--Witt
ring) applies at the K3 base point only after the compact source package
has supplied its finite Hall--Borcherds recognition data. It then gives
a characteristic comparison
\[
H^2_{\mathrm{Lie}}\bigl(\mathfrak q_{K(1)};\mathrm{grt}_1\bigr)
\;\cong\; \mathbb C\cdot \Delta_5,
\]
where the right side is the scalar automorphic target attached to
$\Delta_5$, the primitive square root of the unnormalised Igusa form
$\Phi_{10}^{\mathrm{un}}$. This statement does not compute
the intrinsic second Lie-bialgebra cohomology of
$\mathfrak g_{\Delta_5}$ and does not make $\mathbf H_{\Delta_5}$
unique from automorphy alone. The uniqueness
assertion is conditional: same source window, same comparison matrices,
same radical/PBW/no-extra-relations gates, and same
Bruinier--Heegner characteristic image. Primary-lit:
Etingof-Kazhdan 2007 Part~V, Morava 1985, Hopkins-Kuhn-Ravenel 2000.
\end{remark}


%% 6D holomorphic Chern--Simons: classical BV datum, quantum observables,
%% anomaly, field-strength renormalisation, minimal model, dualizability, and the
%% all-orders non-abelian Yangian compound.

\section{6D holomorphic Chern--Simons: BV datum, quantum observables, dualizability}
\label{sec:hcs-bv-quantum-dualizability}

The holomorphic Chern--Simons theory on a Calabi--Yau $d$-fold $X$ carries a
classical BV datum whose field-pairing shift is $\mathrm{shift}=d-4$,
a quantum BV observable algebra whose $E_3$-structure on $\C^3$ is
witnessed by an explicit Bochner--Martinelli propagator, a corrected
one-loop anomaly split in which $C_2(\fg)$ controls field-strength
renormalisation while the dimension-$3$ holomorphic gauge anomaly is
quartic, and a dualizability profile governed by $E_3$-Koszul
self-duality. These structures give the all-orders identification
$\partial\hCS_5(\fg) \simeq Y_{\epsilon_1,\epsilon_2,\epsilon_3}(\widehat\fg)$
for simply-laced $\fg$, together with the corner identification at finite
$n$ and the disambiguation of the three chiral-Yangian variants. The
seven results below are stated in the lane in which their proof
works; $(\infty,1)$-categorical inputs and chain-level computations are
both load-bearing.

\subsection{Classical BV datum on a Calabi--Yau \texorpdfstring{$d$}{d}-fold}
\label{subsec:hcs-classical-bv}

\begin{theorem}[Classical BV structure of \texorpdfstring{$\hCS$}{hCS} on a CY\texorpdfstring{$_d$}{-d}]
\label{thm:plat-hCS-classical}
Let $X$ be a smooth Calabi--Yau $d$-fold with holomorphic volume form
$\Omega_X \in H^{d,0}(X)$, and let $\fg$ be a finite-dimensional Lie
algebra with nondegenerate invariant pairing
$\langle-,-\rangle_\fg$. The classical BV datum of holomorphic
Chern--Simons theory on $X$ with gauge Lie algebra $\fg$ consists of:
\begin{enumerate}[label=\textup{(\alph*)}]
\item the super-field
\[
 \cA \;\in\; \Omega^{0,\bullet}(X,\fg)[1],
\]
whose $\Omega^{0,p}$ component has BV-degree $1-p$. In complex
dimension $3$ this is
$c + A_{0,1} + A^*_{0,2} + c^*_{0,3}$ with degrees
$+1,0,-1,-2$;
\item the classical action
\[
 S_{\mathrm{cl}}(\cA)
 \;=\; \int_X \Omega_X \wedge
 \Big\langle \cA,\; \bar\partial \cA + \tfrac{1}{3}[\cA,\cA]\Big\rangle_\fg;
\]
\item a $(d{-}4)$-shifted symplectic structure on the space of BV
fields, induced from Serre duality pairing
$\Omega^{0,p} \otimes \Omega^{0, d-p} \to \Omega^{0,d} \to \C$.
\end{enumerate}
The shift law
\[
 \mathrm{shift}(d) \;=\; d - 4,
 \qquad
 E_n(d) \;=\; E_{d},
\]
organises the operadic level of the BV observable algebra by dimension:
\[
 (d,\,\mathrm{shift},\,E_n)
 \;\in\;
 \bigl\{(2,-2,E_2),\ (3,-1,E_3),\ (4,0,E_4),\ (5,+1,E_5)\bigr\}.
\]
This is the BV-field pairing shift, not the PTVV shift on the moduli
stack of objects. Thus at $d=5$ the BV field space has a
$+1$-shifted symplectic pairing, while the PTVV form on
$\cM(D^b\Coh(X))$ has shift $2-d=-3$.
\end{theorem}

\begin{proof}
The field complex $\Omega^{0,\bullet}(X,\fg)[1]$ and the cubic
functional are Costello's holomorphic Chern--Simons BV datum; the
Serre pairing
$\Omega^{0,p}\otimes\Omega^{0,d-p}\to\Omega^{0,d}\to\C$ has total
field degree $d-4$ after the shift by $[1]$ on each input. Local
observables of a holomorphic theory on a complex $d$-fold form an
$E_d$-algebra by the Costello--Gwilliam factorisation formalism, and
the $d=3$ entry is the usual $(-1)$-shifted BV symplectic structure
used in the quantisation of hCS on $\C^3$.
\end{proof}

\subsection{Quantum observables on \texorpdfstring{$\C^3$}{C-3}: \texorpdfstring{$E_3$}{E-3}-structure and propagator}
\label{subsec:hcs-quantum-observables}

\begin{theorem}[Quantum observables of \texorpdfstring{$\hCS$}{hCS} on \texorpdfstring{$\C^3$}{C-3} and the \texorpdfstring{$E_3$}{E-3}-algebra structure]
\label{thm:plat-hCS-quantum}
Let $Y = \C^3$ with holomorphic volume form $\Omega = dz_1 \wedge dz_2 \wedge dz_3$ and let $\fg$ be semisimple.
Costello renormalisation (arXiv:$1303.2632$; Costello--Gwilliam~$2017$ Vol~II~\S$5$)
produces a factorisation algebra of quantum observables
\[
 \Obs_{\hCS}(\C^3)
 \;=\; \bigl(\Sym(\cE^\vee[1])[[\hbar]],\; Q + \hbar\Delta\bigr)
\]
with propagator
\[
 P_{\mathrm{BM}}(z,w)
 \;=\; \frac{2}{(2\pi i)^3}
 \sum_{k=1}^{3} (-1)^{k-1}\, \overline{(z_k - w_k)}\; \|z - w\|^{-6}\;
 \widehat{d\bar{z}_k} \wedge dw_1 \wedge dw_2 \wedge dw_3,
\]
the Bochner--Martinelli kernel on $\C^3 \setminus \Delta$, and OPE
$\cA(z)\cA(w) \sim \hbar\, P_{\mathrm{BM}}(z,w) \cdot \mathbf{1}$
modulo $Q$-exact terms. The observables carry a genuine
$E_3$-algebra structure in the Dolbeault chain complex:
\[
 \Obs_{\hCS}(\C^3)
 \;\simeq\;
 \CE^\bullet_{\bar\partial,\chir}(\cE_{\hCS},\,\cO_{\C^3}),
\]
with $E_3$-composition realised by a sum-over-shuffles prescription with
Koszul signs, associativity enforced by \v{C}ech--Dolbeault Mayer--Vietoris
on $\overline{\mathrm{Conf}}_n(\C^3)$, and commutativity by
$\pi_1(\mathrm{Conf}_2(\C^3)) = \pi_1(S^5) = 0$.
\end{theorem}

\begin{proof}
\emph{Propagator and OPE.}
The kernel $P_{\mathrm{BM}}$ is the unique
$\partial$-harmonic representative of the diagonal class in
$H^{d-1, d-1}(\C^d \setminus \Delta)$ at $d = 3$ (Range~$1986$ \S$6.1$);
direct computation gives
$\bar\partial P_{\mathrm{BM}} = \delta_\Delta$. The Costello--Li
$\bar\partial$-regularisation (arXiv:$1601.04040$ \S$4$) realises
$P_{\mathrm{BM}}$ as the heat-kernel limit
$\lim_{L \to \infty} \int_0^L K_t\, dt$ of the Dolbeault heat kernel on $\C^3$,
which in turn supplies the OPE through the counterterm construction of
Costello--Gwilliam Vol~II Thm.~$9.3.1$.

\emph{$E_3$-structure.}
Shuffles on ordered insertions of points in $\mathrm{Conf}_n(\C^3)$ with
Koszul signs produce a chain-level action of the Fulton--MacPherson little
$3$-disks operad on $\Obs_{\hCS}(\C^3)$, by Kontsevich--Tamarkin formality
of $\cD_3$ (Tamarkin~$2003$; Kontsevich~$1999$ Thm.~$1$) composed with the
Gwilliam--Williams~$2021$ comparison $E_d^{\mathrm{hol}} \simeq E_d$
(arXiv:$2009.05037$ Thm.~$2.5.5$).
Mayer--Vietoris on the strata of the Fulton--MacPherson compactification
$\overline{\mathrm{Conf}}_n(\C^3)$ supplies the associator cocycle; the
resulting pairing
$\mu_n \colon \Obs^{\otimes n} \to \Obs$ is strictly associative on
cohomology and $A_\infty$-associative on chains with higher
$\mu_{n \geq 3}$ computed by shuffle-type integrals against the
Bochner--Martinelli kernel.

\emph{Commutativity.}
The link of the diagonal in $\mathrm{Conf}_2(\C^3)$ is $S^5$ (removing
a point from $\C^3 \setminus \{0\}$ and deformation retracting);
$\pi_1(S^5) = 0$, so the $E_3$-operad trivialises braid monodromy at the
topological level. The nontriviality of the chain-level $E_3$-structure
is supplied by the holomorphic deformation: the $\bar\partial$-grading
breaks the symmetric action of $S_n$ on $\mu_n$ to the $E_3$-action of the
framed little $3$-disks.

\emph{CE-chiral identification.}
Costello--Gwilliam Vol~II Thm.~$6.9.0.1$ realises
$\Obs_{\hCS}(\C^3)$ as the holomorphic Chevalley--Eilenberg complex of the
$L_\infty$-algebra $\cE_{\hCS} = \Omega^{0,\bullet}(\C^3, \fg)[1]$ valued
in the commutative ring $\cO_{\C^3}$; the isomorphism is the chain-level
content of the BV quantisation.
\end{proof}

\subsection{Anomaly factorisation}
\label{subsec:hcs-anomaly}

\begin{theorem}[Corrected anomaly split for \texorpdfstring{$\hCS$}{hCS} on CY\texorpdfstring{$_3$}{-3}]
\label{thm:plat-anomaly}
Let $X$ be a compact Calabi--Yau $3$-fold with holomorphic volume form
$\Omega_X$ and topological Euler characteristic $\chi_{\mathrm{top}}(X)$,
and let $\fg$ be a semisimple Lie algebra. In complex dimension $d$, the
local holomorphic gauge anomaly is controlled by an invariant polynomial
of degree $d+1$ in the gauge curvature. Thus for $d=3$ the cohomological
anomaly is quartic, schematically
\[
 \kanom(X,\fg)
 \;=\;
 \hbar \int_X \Omega_X \wedge
 I_4(\cA,F_\cA,F_\cA,F_\cA),
\]
with $I_4 \in \mathrm{Sym}^4(\fg^\vee)^{\fg}$ in the chosen
representation. The quadratic-Casimir coefficient $C_2(\fg)$ belongs to
field-strength renormalisation and to BV-trivial counterterms; it is not
the anomaly class.
Consequences:
\begin{enumerate}[label=\textup{(\roman*)}]
\item Vanishing of the cubic tensor $d^{abc}$ for
$\mathfrak{su}(2)$, $\mathfrak{so}(N)$, $E_6,E_7,E_8,F_4,G_2$ is not an
anomaly-cancellation theorem for $6$D $\hCS$; the quartic invariant
$I_4$ must be evaluated.
\item On flat $\C^3$ the local anomaly is absent in the translation
invariant perturbative model. On $K3\times E$ the
Euler-characteristic/BCOV coefficient vanishes because
$\chi_{\mathrm{top}}(K3\times E)=24\cdot0=0$.
\item The quintic, with $\chi_{\mathrm{top}}(X_5)=-200$, remains the
first compact instance for the quartic anomaly. The cubic symmetric
tensor $d^{abc}$ does not determine this coefficient.
\end{enumerate}
\end{theorem}

\begin{proof}
Costello's heat-kernel regularisation supplies the perturbative
renormalisation framework. The four-dimensional gauge-anomaly pattern
(triangle diagram, cubic Casimir) does not transfer to holomorphic
dimension~$3$: the local Chern--Weil degree of the hCS anomaly is
$d+1=4$. A full compact CY$_3$ quantisation theorem with the quartic
coefficient fixed is not asserted here; the statement records the
correct obstruction degree and separates it from the $C_2(\fg)$
field-strength counterterm.
\end{proof}

\subsection{The bar complex as BV-BRST complex: identification and anomaly universality}
\label{subsec:phys-bv-bar-identification}

%% Absorption from physics_bv_brst_cy.tex:
%% The identifications in this subsection--QME = MC for Theta_A,
%% bar differential = Q_BRST, curvature = BRST anomaly--are not
%% present elsewhere in this chapter. They are recorded here as
%% structural facts, with the physics note as primary source and
%% primary-literature attribution.

The holomorphic Chern--Simons BV datum of Theorem~\ref{thm:plat-hCS-classical} carries an open-string interpretation: for a brane $\mathcal{F} \in \mathrm{Ob}(\mathcal{C})$ in the CY category $\mathcal{C}$, the open string field is
\[
 \alpha \;\in\; \mathfrak{F} \;:=\; \mathrm{RHom}_\mathcal{C}(\mathcal{F}, \mathcal{F})[1],
\]
and the BV action is the CY-trace functional
\begin{equation}
\label{eq:phys-bv-bv-action}
 S(\alpha) \;=\; \sum_{n \geq 2} \frac{1}{n}\, \mathrm{Tr}\bigl(\mu_n(\alpha, \ldots, \alpha)\bigr),
\end{equation}
where $\mu_n$ are the $A_\infty$-operations and $\mathrm{Tr} \colon A \to k[-d]$ is the CY trace. The CY pairing $\langle -,- \rangle \colon A \otimes A \to k[-d]$ supplies the BV symplectic form.

\begin{proposition}[Classical master equation = \texorpdfstring{$A_\infty$}{A-inf} relations]
\label{prop:phys-bv-classical-me}
\ClaimStatusProvedHere
The BV action \eqref{eq:phys-bv-bv-action} satisfies the classical master equation
\[
 \{S, S\} = 0
\]
with respect to the BV antibracket $\{-,-\}$ defined by inverting the CY pairing. This is equivalent to the $A_\infty$-relations for $(\mathrm{End}_\mathcal{C}(\mathcal{F}), \{\mu_n\})$.
\end{proposition}

\begin{proof}
The classical master equation $\{S, S\} = 0$ expands, using the BV antibracket from the CY pairing, as
\[
 \sum_{p+q = n+1} \mathrm{Tr}\bigl(\mu_p(\alpha^{\otimes (p-1)}, \mu_q(\alpha^{\otimes q}))\bigr) = 0
 \qquad \text{for all } n \geq 1.
\]
This is the $A_\infty$-relation $\sum_{p+q=n+1} \mu_p \circ (\mathrm{id}^{\otimes i} \otimes \mu_q \otimes \mathrm{id}^{\otimes (p-1-i)}) = 0$ contracted against $\alpha^{\otimes n}$ and traced. Non-degeneracy of the CY pairing shows the traced and untraced $A_\infty$-relations are equivalent.
\end{proof}

\begin{theorem}[Quantum master equation = Maurer--Cartan equation for \texorpdfstring{$\Theta_A$}{Theta-A}]
\label{thm:phys-bv-qme-mc}
Let $A = A_\mathcal{C} = \Phi(\mathcal{C})$ be the quantum chiral algebra associated to the CY category $\mathcal{C}$, and let $\Theta_A$ be the universal Maurer--Cartan element of $A$ (Vol~I). The quantum master equation
\[
 \hbar\,\Delta S + \tfrac{1}{2}\{S, S\} = 0
\]
for the BV action $S$ of the open string field theory on any brane $\mathcal{F} \in \mathcal{C}$ is equivalent to the Maurer--Cartan equation
\[
 d\Theta_A + \tfrac{1}{2}[\Theta_A, \Theta_A] \;=\; \mathcal{O}_{\mathrm{full}}(A)
\]
for $\Theta_A$ in the chiral deformation complex of $A$. Its scalar
Hodge projection is $\operatorname{pr}_{\mathrm{sc}}\mathcal{O}_{\mathrm{full},g}(A)
= \kappa_{\mathrm{ch}}(A)\lambda_g$: unconditionally at genus~$1$, and at
all genera on the uniform-weight scalar lane. The BV-to-chiral dictionary is:
\begin{center}
\renewcommand{\arraystretch}{1.3}
\begin{tabular}{lll}
 \textbf{BV-BRST} & \textbf{Chiral bar complex} & \textbf{Identification} \\
 \hline
 Field $\alpha \in \mathfrak{F}$ & Bar element $\alpha \in B(A)$ & Same underlying vector space \\
 BV action $S(\alpha)$ & MC element $\Theta_A$ & $S = \mathrm{Tr}(\Theta_A)$ \\
 CY pairing $\langle -,- \rangle$ & Factorisation coproduct & $\omega_{\mathrm{BV}} = \Delta_B^*\omega_{\mathrm{CY}}$ \\
 $\{S,S\} = 0$ (CME) & $[\Theta_A, \Theta_A] = 0$ at genus $0$ & Classical $A_\infty$ relations \\
 $\hbar\,\Delta S$ & $d\Theta_A$ at genus $\geq 1$ & One-loop anomaly \\
 QME obstruction & $\mathcal{O}_{\mathrm{full},g}(A)$, scalar projection $\kappa_{\mathrm{ch}}(A)\cdot\omega_g$ & Genus-$g$ anomaly
\end{tabular}
\end{center}
Expanding in powers of $\hbar$ (genus):
\begin{itemize}
 \item \emph{Genus $0$}: $\{S,S\} = 0$, the $A_\infty$-relations (Proposition~\ref{prop:phys-bv-classical-me}). On the chiral side, $[\Theta_A, \Theta_A]|_{g=0} = 0$: the bar complex is a differential coalgebra at tree level.
 \item \emph{Genus $1$}: the one-loop anomaly $\Delta S_2 + \{S_3, S_2\} + \cdots$ equals $\kappa_{\mathrm{ch}}(A)\cdot\lambda_1$, where $\lambda_1 = c_1(\mathbb{E})$ is the Hodge class on $\overline{\mathcal{M}}_{1,1}$.
 \item \emph{Genus $g$}: the full genus-$g$ obstruction is
 $\mathcal{O}_{\mathrm{full},g}(A)$. Its scalar component is
 $\mathrm{obs}^{\mathrm{sc}}_g(A)=\kappa_{\mathrm{ch}}(A)\cdot\lambda_g$
 on the uniform-weight scalar lane (Vol~I; cf.~equation~\eqref{eq:phys-anom-obs-g}
 below), while multi-weight families at $g\geq2$ carry cross-channel
 summands.
\end{itemize}
\end{theorem}

\begin{proof}
The identification $S = \mathrm{Tr}(\Theta_A)$ and $\omega_{\mathrm{BV}} = \Delta_B^*\omega_{\mathrm{CY}}$ are term-by-term: the degree-$n$ term of $\Theta_A$ is $\frac{1}{n}\mu_n \in A^{\otimes n}$, and $\mathrm{Tr}$ on $A^{\otimes n}$ applied to $\alpha^{\otimes n}$ gives $\frac{1}{n}\mathrm{Tr}(\mu_n(\alpha,\ldots,\alpha)) = S_n$, the $n$-th term of the BV action. The BV Laplacian $\Delta = \sum_{i,j}\omega^{ij}\partial^2/\partial\phi^i\partial\phi^j$ acting on $S$ contracts a pair of fields via $\omega_{\mathrm{BV}} = \Delta_B^*\omega_{\mathrm{CY}}$; on the chiral side this is $d\Theta_A|_{g=1}$, the one-loop piece of the deformation differential. The genus-$g$ anomaly coefficient $\kappa_{\mathrm{ch}}(A)$ is the scalar projection of $\Theta_A$ onto the Hodge class $\lambda_g$ (Vol~I, Theorem~D), not the full multi-weight component. The equivalence of the full QME with the MC equation follows by induction on genus: at each genus $g$, the BV equation at order $\hbar^{2g-2}$ is precisely the $g$-th component of the full MC equation; applying the scalar projection gives the displayed Theorem~D component, and the remaining non-scalar part is the cross-channel correction.
\end{proof}

\begin{proposition}[Bar differential = BRST operator]
\label{prop:phys-bv-bar-brst}
Under the identification $\mathfrak{F} = A[1]$ and $S = \mathrm{Tr}(\Theta_A)$, the BRST operator $Q_{\mathrm{BRST}} = \{S, -\}$ acting on $\mathfrak{F}$ coincides with the bar differential $d_B$ on $B(A)$. Explicitly, for $\alpha_1 \otimes \cdots \otimes \alpha_n \in A[1]^{\otimes n}$:
\[
 d_B(\alpha_1 \otimes \cdots \otimes \alpha_n)
 = \sum_{i < j} \pm\, \alpha_1 \otimes \cdots \otimes \mu_{j-i+1}(\alpha_i, \ldots, \alpha_j) \otimes \cdots \otimes \alpha_n.
\]
\end{proposition}

\begin{proof}
On the BV side, $Q_{\mathrm{BRST}} = \{S,-\}$ acts on $\mathfrak{F}$ via the antibracket; expanding $\{S_n, -\}$ for each $n$ using the BV antibracket from the CY pairing produces a sum over insertions of $\mu_n$ at consecutive positions, which is precisely the bar differential $d_B$. The Koszul sign rule matches the graded-commutative conventions of the bar complex.
\end{proof}

\begin{theorem}[Curvature = BRST anomaly; Adler--Bardeen universality]
\label{thm:phys-bv-curvature-anomaly}
At genus $g \geq 1$, the bar complex $B(A)$ acquires curvature from the Hodge bundle over $\overline{\mathcal{M}}_{g,n}$. Its scalar Hodge projection is:
\begin{equation}
\label{eq:phys-bv-curvature}
 \operatorname{pr}_{\mathrm{sc}}(d_B^2) \;=\; \kappa_{\mathrm{ch}}(A) \cdot \omega_g,
\end{equation}
where $\omega_g \in H^2(\overline{\mathcal{M}}_{g,1})$ is the Hodge class. The three consequences:
\begin{enumerate}[label=\textup{(\roman*)}]
 \item At genus $g$, the scalar failure of nilpotency
 $\operatorname{pr}_{\mathrm{sc}}(Q_{\mathrm{BRST}}^2)
 = \kappa_{\mathrm{ch}}(A)\cdot\omega_g$ is the scalar BRST anomaly at
 $g$ loops.
 \item The scalar anomaly is \emph{universal}: it depends only on
 $\kappa_{\mathrm{ch}}(A)$ and the Hodge class, not on the particular
 brane $\mathcal{F}$. At $g\geq2$, the full multi-weight obstruction can
 also contain cross-channel summands.
 \item Scalar anomaly cancellation $\kappa_{\mathrm{ch}}(A) = 0$ restores
 $\operatorname{pr}_{\mathrm{sc}}(d_B^2)=0$ at all genera. Full
 nilpotency additionally requires vanishing of the cross-channel
 summand, for example on the uniform-weight lane.
\end{enumerate}
\end{theorem}

\begin{proof}
The genus-$0$ nilpotency $d_B^2 = 0$ is the content of the $A_\infty$-relations (Proposition~\ref{prop:phys-bv-classical-me}). At genus $g \geq 1$, the bar complex is a curved factorisation coalgebra: the scalar curvature is computed by the trace $\mathrm{Tr}(Q_{\mathrm{BRST}}^2)$ on the endomorphism algebra, which by the CY trace axiom picks up $\kappa_{\mathrm{ch}}(A)$; the Hodge class $\omega_g$ is the cohomological class of the Hodge bundle insertion. The brane-independence of $\kappa_{\mathrm{ch}}$ follows from its definition as a scalar projection of $\Theta_A$, which is defined on $A$ (not on any individual brane module): the trace $\mathrm{Tr}$ on $\mathrm{End}(\mathcal{F})$ is fixed by the CY structure, and $\kappa_{\mathrm{ch}}(A)$ is its genus-$1$ shadow. Universality is Theorem~D of Vol~I on the scalar lane; the full $g\geq2$ multi-weight curvature also records cross-channel terms.
\end{proof}

\begin{proposition}[\texorpdfstring{$E_2$}{E-2} enhancement of the bar coalgebra: native at \texorpdfstring{$d=2$}{d=2}, centre-side at \texorpdfstring{$d\geq3$}{d>=3}]
\label{prop:phys-bv-e2-bar}
For a CY$_d$ category $\mathcal{C}$ with $d \geq 2$, the bar construction has the following \(E_2\) scope:
\begin{enumerate}[label=\textup{(\roman*)}]
 \item The $E_1$ direction is the BRST direction: the bar differential $d_B=Q_{\mathrm{BRST}}$ and the deconcatenation coproduct $\Delta$.
 \item At \(d=2\), \(A_\mathcal{C}\) is natively \(E_2\) (Definition~\ref{def:e2-chiral-algebra}), so \(B(A_\mathcal{C})\) is the factorisation-dual \(E_2\) coalgebra.
 \item At \(d\geq3\), \(A_\mathcal{C}\) is natively \(E_1\). The CY \(R\)-matrix \(R_\mathcal{C}(z)=\mathrm{Res}^{\mathrm{coll}}_{0,2}(\Theta_{A_\mathcal{C}})\) equips the Drinfeld-centre side \(\mathcal{Z}(\mathrm{Rep}^{E_1}(A_\mathcal{C}))\), and its bar object, with the braided \(E_2\) structure. It does not make \(B(A_\mathcal{C})\) an intrinsic \(E_2\) coalgebra.
\end{enumerate}
The open-closed datum lives on the \emph{pair}
$(C^\bullet_{\mathrm{ch}}(A,A), A)$: the chiral Hochschild cochain
complex $C^\bullet_{\mathrm{ch}}(A,A)$ is the algebraic closed-sector
object, $A$ is the open/boundary sector, and a physical bulk reading
requires the open--closed comparison coupling
$\mathrm{HH}^\bullet(\mathcal{C})$ to $A$ by deformations of the
$A_\infty$-structure.
\end{proposition}

\begin{proof}
For $d = 2$: the native $E_2$ structure is Theorem~CY-A$_2$
(Chapter~\ref{ch:cy-to-chiral}). The $E_2$ bar coalgebra is the
factorisation dual of the $E_2$-chiral algebra $A_\mathcal{C}$. For
$d \geq 3$: the CY $R$-matrix $R_\mathcal{C}(z)$ is extracted from
$\Theta_{A_\mathcal{C}}$ at degree~$2$
(Construction~\ref{constr:cy-r-matrix}), and its YBE property
(Theorem~\ref{thm:yang-r-matrix-ybe}) gives the half-braiding on the
Drinfeld centre \(\mathcal{Z}(\mathrm{Rep}^{E_1}(A_\mathcal{C}))\).
The centre-side bar object is therefore \(E_2\); the boundary bar
\(B(A_\mathcal{C})\) remains \(E_1\). The open-closed pair structure
is the Swiss-cheese datum $(C^\bullet_{\mathrm{ch}}(A,A), A)$ from
Vol~II; comparison with BV open-closed string field theory requires
the Kontsevich--Soibelman open-closed map.
\end{proof}

\begin{conjecture}[BV-BRST complex of the topological B-model and the bar complex of \texorpdfstring{$G(X)$}{G(X)}]
\label{conj:phys-bv-cy3-bar}
\ClaimStatusConjectured
For a CY$_3$ manifold $X$, assume the compact Hall/centre package defining the quantum vertex chiral group \(G(X)\) has been constructed. The full BV-BRST complex of the topological B-model is then expected to compare with the bar complex \(B(G(X))\). Specifically:
\begin{enumerate}[label=\textup{(\roman*)}]
 \item The bar differential $d_B$ on $B(G(X))$ is the BRST operator $Q_{\mathrm{BRST}}$ of the B-model.
 \item The genus-$0$ cohomology $H^\bullet(B(G(X)), d_B)|_{g=0}$ computes the classical B-model open string states modulo gauge equivalence.
 \item At genus $g \geq 1$, the scalar Hodge projection
 $\operatorname{pr}_{\mathrm{sc}}(d_B^2)
 = \kappa_{\mathrm{ch}}(G(X))\cdot\omega_g$ is the scalar $g$-loop
 BRST anomaly on the chiral lane; at $g\geq2$ the full multi-weight
 obstruction can retain cross-channel summands. The Borcherds scalar is a distinct datum: for $K3 \times E$, $\chi_{\mathrm{top}} = 0$ but $\kappa_{\mathrm{BKM}}(\Delta_5) = 5$, so the BKM weight is not computed by $\chi_{\mathrm{top}}(X)/24$ (Remark~\ref{rem:phys-anom-chi-vs-weight} below).
 \item The bar-cobar adjunction $\Omega(B(G(X))) \simeq G(X)$ (on the Koszul locus) is the statement that the open string field theory determines, and is determined by, the quantum vertex chiral group.
 \item The BKM denominator identity of $\mathfrak{g}_X$ is the BV partition function:
 \[
 Z_{\mathrm{BRST}} \;=\;
 \frac{\det Q_{\mathrm{BRST}}|_{\mathrm{ghost}}}{\det Q_{\mathrm{BRST}}|_{\mathrm{field}}}
 \;=\; \prod_{\alpha \in \Delta_+} (1 - q^\alpha)^{\mathrm{mult}(\alpha)}.
 \]
 For $K3 \times E$, this reproduces the Borcherds product expansion of the Igusa cusp form $\Delta_5$.
\end{enumerate}
\end{conjecture}

\subsection{Modular characteristics: values, anomaly ratios, and the three-candidate problem}
\label{subsec:phys-anom-kappa-table}

%% Absorption from physics_anomaly_cancellation.tex:
%% The master table of kappa_ch values with rho-ratios, the three-candidate
%% problem (chi/24, chi/12, automorphic weight), the anomaly non-cancellation
%% theorem, and the toric genus-1 matching proposition are new to this chapter.

The modular characteristic $\kappa_{\mathrm{ch}}(\mathcal{A})$ of a chiral algebra $\mathcal{A}$ controls the scalar Hodge projection of the genus tower:
\begin{equation}
\label{eq:phys-anom-obs-g}
 \mathrm{obs}^{\mathrm{sc}}_g(\mathcal{A}) \;=\; \kappa_{\mathrm{ch}}(\mathcal{A}) \cdot \lambda_g
 \;\in\; H^{2g}(\overline{\mathcal{M}}_g, \mathbb{Q}),
\end{equation}
where $\lambda_g = c_g(\mathbb{E})$ is the top Chern class of the Hodge bundle. At genus~$1$ this scalar identity is unconditional; at all genera it is the uniform-weight scalar lane of Vol~I, Theorem~D. For multi-weight algebras and $g\geq2$, the full obstruction is $\mathrm{obs}^{\mathrm{full}}_g(\mathcal{A})=\mathrm{obs}^{\mathrm{sc}}_g(\mathcal{A})+\mathrm{obs}^{\mathrm{cross}}_g(\mathcal{A})$. The genus-$1$ free energy is $F_1(\mathcal{A}) = \kappa_{\mathrm{ch}}(\mathcal{A})/24$ (Faber--Pandharipande evaluation $\lambda_1^{\mathrm{FP}} = 1/24$).

The \emph{anomaly ratio} $\varrho(\mathcal{A}) := \kappa_{\mathrm{ch}}(\mathcal{A})/c(\mathcal{A})$ encodes how much of the conformal anomaly survives in the genus tower. The master table of standard chiral algebras:
\begin{center}
\renewcommand{\arraystretch}{1.3}
\begin{tabular}{lccc}
\textbf{Algebra} & $c$ & $\kappa_{\mathrm{ch}}$ & $\varrho$ \\ \hline
$\mathcal{H}^{\oplus d}$ (Heisenberg, rank $d$, level $1$) & $d$ & $d$ & $1$ \\
$\widehat{\mathfrak{g}}_k$ (Kac--Moody) &
 $k\dim\mathfrak{g}/(k+h^\vee)$ &
 $(k+h^\vee)\dim\mathfrak{g}/(2h^\vee)$ &
 $(k+h^\vee)^2/(2kh^\vee)$ \\
$\mathrm{Vir}_c$ (Virasoro) & $c$ & $c/2$ & $1/2$ \\
$\mathcal{W}_N^k$ ($\mathfrak{sl}_N$-type) & $c$ & $c(H_N - 1)$ & $H_N - 1$ \\
$bc$ (ghosts, weights $(2,-1)$) & $-26$ & $-13$ & $1/2$ \\
$\beta\gamma$ (symplectic bosons) & $c$ & $c/2$ & $1/2$ \\
$V_\Lambda$ (lattice, rank $d$) & $d$ & $d$ & $1$
\end{tabular}
\end{center}
where $H_N = \sum_{j=1}^N 1/j$ is the $N$-th harmonic number. The Heisenberg algebra is the unique standard family with $\varrho = 1$: $\kappa_{\mathrm{ch}} = k = c$ at rank $d$, level $1$. For the bosonic string: $\kappa_{\mathrm{eff}} = \kappa_{\mathrm{ch}}(\mathrm{Vir}_c) + \kappa_{\mathrm{ch}}(\mathrm{ghost}) = c/2 - 13 = 0$ at $c = 26$.

\begin{remark}[Complementarity: Koszul pairs]
\label{rem:phys-anom-complementarity}
For a Koszul pair $(\mathcal{A}, \mathcal{A}^!)$, the complementarity conductor $\kappa_{\mathrm{ch}}(\mathcal{A}) + \kappa_{\mathrm{ch}}(\mathcal{A}^!)$ depends on the family: it vanishes for Heisenberg, Kac--Moody, and free-field algebras (antisymmetric under Feigin--Frenkel $k \mapsto -k-2h^\vee$); it equals $13$ for Virasoro (Vol~I, Theorem~C); $250/3$ for $\mathcal{W}_3$; $(H_N-1)K_N/2$ for $\mathcal{W}_N^k$. For the $\fgl_1$/$\C^3$ Yangian branch: $\kappa_{\mathrm{ch}} + \kappa_{\mathrm{ch}}^! = -\sigma_2 + \sigma_2 = 0$ (Proposition~\ref{prop:5d-koszul-dual}).
\end{remark}

\begin{proposition}[Toric CY\texorpdfstring{$_3$}{-3} scalar and refined characters]
\label{prop:phys-anom-toric-genus1}
For the toric CY$_3$ \(X=\mathbb C^3\):
\begin{enumerate}[label=\textup{(\roman*)}]
 \item The evaluated rank-one Heisenberg sector has
 \(\kappa_{\mathrm{ch}}^{\mathrm{Heis}}(\mathbb C^3)=1\) and vacuum
 character \(P(q)=\prod_{n\geq 1}(1-q^n)^{-1}\).
 \item The Virasoro subalgebra at \(c=1\) has scalar
 \(\kappa_{\mathrm{ch}}^{\mathrm{Vir}}=1/2\). This is a subalgebra
 invariant, not the scalar CY trace of the \(\C^3\) Hall output.
 \item The DT partition function is the refined character
 \[
  Z^{\mathrm{DT}}(\mathbb C^3)=M(q)
  =\operatorname{PE}\!\left(\sum_{n\geq 1}nq^n\right),
 \]
 so \(\Omega(n)=n\).
 \item The finite \(W_N\) conductor \(c(H_N-1)\) has no limit which
 defines a scalar \(\kappa_{\mathrm{ch}}\) for the \(\C^3\) Hall output.
\end{enumerate}
\end{proposition}

\begin{proof}
The first two assertions are the standard OPE normalisations:
\[
 J(z)J(w)\sim (z-w)^{-2},
 \qquad
 T(z)T(w)\sim \frac{1/2}{(z-w)^4}+\cdots .
\]
The third assertion is the product formula for the Hilbert scheme of
points on \(\C^3\), equivalently the plethystic exponential of the
primitive BPS spectrum. The last assertion follows because
\(H_N-1\to\infty\), while the scalar Heisenberg sector remains rank one.
\end{proof}

\begin{theorem}[Anomaly non-cancellation: independent scalar lanes]
\label{thm:phys-anom-non-cancellation}
For a CY$_3$ $X$, the chiral modular characteristic of a constructed quantum vertex group, the BCOV genus-one scalar, and the Borcherds automorphic weight are independent pieces of data until a comparison theorem identifies them:
\begin{enumerate}[label=\textup{(\roman*)}]
 \item $\kappa_{\mathrm{ch}}(G(X))$ is defined only where the Hall positive half, Drinfeld double, and centre comparison have been constructed.
 \item The BCOV genus-one scalar is controlled by the topological and holomorphic anomaly data of $X$.
 \item The Borcherds weight $w(\Phi_X)$ is an invariant of a chosen automorphic denominator.
 \item The anomaly cancellation condition $\kappa_{\mathrm{eff}} = \kappa_{\mathrm{ch}}(G(X)) + \kappa_{\mathrm{ghost}} = 0$ determines a \emph{critical weight} $w_{\mathrm{crit}} = -\kappa_{\mathrm{ghost}}$, not a critical central charge.
\end{enumerate}
There are three scalar lanes which may agree only on a constructed comparison locus:
\begin{center}
\begin{tabular}{lcl}
\textbf{Lane} & \textbf{Scalar} & \textbf{Agreement condition} \\ \hline
Heisenberg analogy & $\chi_{\mathrm{top}}(X)/24$ & framed free-field comparison \\
BCOV genus-$1$ & BCOV anomaly scalar & genus-one comparison map \\
Automorphic denominator & $w(\Phi_X)$ & BKM recognition theorem
\end{tabular}
\end{center}
No equality among these three numbers is a formal consequence of the symbol \(G(X)\).
\end{theorem}

\begin{proof}[Argument]
The three scalars live on different levels of the construction. The Euler and BCOV scalars are target-space anomaly data. The Borcherds weight is a scalar shadow of an automorphic denominator. The chiral modular characteristic is a property of a constructed vertex object. A theorem may identify two of them only after specifying the comparison arrow and its hypotheses.
\end{proof}

\begin{remark}[Why \texorpdfstring{$\chi(K3 \times E) = 0$}{chi(K3 x E) = 0} but \texorpdfstring{$\kappa_{\mathrm{BKM}}(\Delta_5) = 5 \neq 0$}{kappa-BKM(Delta-5) = 5 != 0}]
\label{rem:phys-anom-chi-vs-weight}
The Euler characteristic $\chi_{\mathrm{top}}(K3 \times E) = \chi_{\mathrm{top}}(K3)\chi_{\mathrm{top}}(E) = 24 \cdot 0 = 0$ belongs to the target-space scalar lane. The Borcherds invariant $\kappa_{\mathrm{BKM}}(\Delta_5) = 5$ is the weight of the denominator $\Delta_5$ in the Gritsenko normalisation. These are different invariants. The sentence \(\chi_{\mathrm{top}}=0\) therefore gives no contradiction and no derivation of \(\kappa_{\mathrm{BKM}}(\Delta_5)=5\).
\end{remark}

\subsection{Non-abelian field-strength renormalisation}
\label{subsec:hcs-field-strength-renormalisation}

\begin{theorem}[One-loop non-abelian field-strength renormalisation of \texorpdfstring{$\hCS$}{hCS} on \texorpdfstring{$\C^3$}{C-3}]
\label{thm:plat-Z-counterterm}
The one-loop bubble of $\hCS$ on $\C^3$ with semisimple gauge algebra
$\fg$ requires the counterterm
\[
 S^{(1)}_{\mathrm{c.t.}}(\cA)
 \;=\;
 -\hbar\, C_2(\fg)\, (4\pi)^{-3}\, \log(L/\varepsilon)
 \int_{\C^3} \Omega \wedge \Tr(\cA\, \bar\partial \cA),
\]
where $C_2(\fg)$ is the quadratic Casimir, $L$ is the IR scale, and
$\varepsilon$ is the UV cutoff. The associated field-strength
renormalisation factor is
\[
 Z^{(1)}_\cA
 \;=\;
 1 - \hbar\, C_2(\fg)\, (4\pi)^{-3}\, \log(L/\varepsilon).
\]
For $\fg = \mathfrak{su}(N)$: $C_2 = N$, and the coefficient of
$\log(L/\varepsilon)$ equals $-\hbar N/(32\pi^3)$.
\end{theorem}

\begin{remark}[Field-strength renormalisation vs anomaly]
\label{rem:plat-Z-vs-anomaly}
The coefficient $-\hbar C_2(\fg)(4\pi)^{-3}\log(L/\varepsilon)$ of the
kinetic counterterm is a \emph{field-strength} renormalisation of $\cA$,
not an anomaly: it rescales the kinetic term but preserves the BV master
equation $\{S, S\} = 0$ once the counterterm is included. The anomaly
$\kanom$ of Theorem~\ref{thm:plat-anomaly} is the obstruction to
solving the quantum master equation at one loop. In complex dimension
$3$ that cohomological anomaly is controlled by the quartic invariant
$I_4\in\mathrm{Sym}^4(\fg^\vee)^\fg$ in the chosen representation, not
by $C_2(\fg)$ and not by the cubic tensor $d^{abc}$. Thus $C_2$ sets
field renormalisation; the quartic Chern--Weil slot sets the local hCS
anomaly. Conflating these invariants misattributes the one-loop
dressing.
\end{remark}

\begin{proof}
The counterterm is extracted by Costello's BV-cohomology regularisation
(Costello~$2011$ \S$5.4$; Costello--Gwilliam Vol~II Thm.~$9.5.0.6$) on
the Bochner--Martinelli propagator of
Theorem~\ref{thm:plat-hCS-quantum}:
the coefficient of the kinetic term receives a
$\log(L/\varepsilon)$ correction proportional to the
quadratic-Casimir trace
$C_2(\fg)\cdot \Tr(\cA \bar\partial \cA)$, with the $(4\pi)^{-3}$
normalisation fixed by the three-dimensional Dolbeault heat-kernel
measure. For $\mathrm{SU}(N)$: $C_2 = N$ in the normalisation
$\Tr(T^a T^b) = \tfrac{1}{2}\delta^{ab}$ (Peskin--Schroeder Appendix~A,
eq.~$(A.94)$).
\end{proof}

\subsection{Minimal \texorpdfstring{$L_\infty$}{L-inf}-model on flat and compact CY\texorpdfstring{$_3$}{-3}}
\label{subsec:hcs-minimal-model}

\begin{theorem}[Minimal \texorpdfstring{$L_\infty$}{L-inf}-model of \texorpdfstring{$\hCS$}{hCS} on \texorpdfstring{$\C^3$}{C-3}]
\label{thm:plat-Linf-minimal}
On flat $\C^3$ with harmonic subspace $\mathcal{H} = \C[z_1, z_2, z_3]
\otimes \fg$, Kontsevich--Soibelman homotopy transfer along the
Bochner--Martinelli propagator of
Theorem~\ref{thm:plat-hCS-quantum} gives a minimal $L_\infty$-model
$(\mathcal{H}, \{\ell_n^{\min}\}_{n \geq 1})$ for
$\cE_{\hCS}(\C^3)$ with
\[
 \ell_1^{\min} \;=\; 0,
 \qquad
 \ell_2^{\min}([f, g]) \;=\; [f, g]_\fg,
 \qquad
 \ell_n^{\min} \;=\; 0 \quad \text{for all } n \geq 3.
\]
On a compact Calabi--Yau $3$-fold $X$, the Atiyah class
\[
 \mathrm{At}(TX) \;\in\;
 H^1\!\bigl(X,\,\Omega^1_X \otimes \mathrm{End}(TX)\bigr)
\]
is the sole formality obstruction: $\mathrm{At}(TX) = 0$ implies
$\ell_n^{\min} = 0$ for $n \geq 3$. In particular:
\begin{enumerate}[label=\textup{(\roman*)}]
\item On $K3 \times E$: $\mathrm{At}(TE) = 0$ (the elliptic curve is a
complex Lie group, so $TE$ is trivial), but the total-space cubic
receptacle does not vanish. Indeed
$H^3(K3\times E,\Omega^3_{K3\times E})$ contains the K\"unneth summand
$H^2(K3,\Omega^2_{K3})\otimes H^1(E,\Omega^1_E)\cong\C$.
Therefore strict Lie formality on $K3\times E$ is not a consequence of
this Hodge group vanishing.
\item On the quintic: $\mathrm{At}(TX) \neq 0$; the Kuranishi cubic
receptacle $H^3(X, \Omega^3_X) = \C$ is nonzero, so the minimal model
carries a nontrivial $\ell_3^{\min}$ representing the CHSW obstruction.
\end{enumerate}
\end{theorem}

\begin{proof}
\emph{Flat case.}
Every rooted binary tree with an internal edge contributes to
$\ell_n^{\min}$ via the composition
$p \circ \mu \circ (h \otimes h) \circ \mu \circ \cdots$
where $h$ is the chain homotopy inverting $\bar\partial$ via
$P_{\mathrm{BM}}$. On the polynomial harmonic subspace $\mathcal{H}$,
$h$ acts as zero (by the Hodge decomposition
$\Omega^{0,\bullet}(\C^3) = \mathcal{H} \oplus \mathrm{im}(\bar\partial)
\oplus \mathrm{im}(\bar\partial^*)$
restricted to polynomial forms: the propagator kills harmonic inputs).
Every tree with an internal edge vanishes; only the binary root tree
$\ell_2$ survives, reproducing the Lie bracket of $\fg$.

\emph{Compact case.}
The Atiyah class $\mathrm{At}(TX)$ is the class of the extension
$0 \to \Omega^1_X \otimes TX \to J^1(TX) \to TX \to 0$; by
Kapranov~$1999$ Thm.~$2.8.1$ (\emph{Compositio}~$115$), the
$L_\infty$-structure of $\Omega^{0,\bullet}(X, TX)$ at the chain level
has cubic operation $\ell_3 = \mathrm{At}(TX) \cdot$ (anti-symmetrised
triple bracket). On $K3 \times E$: the K\"unneth decomposition
$T(K3 \times E) = \pi_1^* TK3 \oplus \pi_2^* TE$ and
$\mathrm{At}(TE) = 0$ (Biswas--Nag~$2002$ for the elliptic-curve case)
reduces the K3-side Atiyah class to a class in
$H^1(K3, \Omega^1 \otimes \mathrm{End})$, and the cubic receptacle
 $H^3(K3 \times E, \Omega^3)$ is nonzero by K\"unneth:
 $H^3(K3 \times E, \Omega^3) = \bigoplus_{p+q=3}
 H^p(K3, \Omega^p_{K3}) \otimes H^q(E, \Omega^q_E)$, and the summand
 $H^2(K3,\Omega^2_{K3})\otimes H^1(E,\Omega^1_E)$ is $\C$.
 Consequently the strict-Lie conclusion for $K3\times E$ requires an
 additional vanishing of the transferred operation, not merely Hodge
 degree counting.
On the quintic: $H^3(X, \Omega^3_X) = \C$ by Serre duality to
$H^0(X, \cO_X) = \C$, and the Atiyah class is nonzero by
Griffiths--Harris~$1978$ Ch.~$5$.
\end{proof}

\subsection{Dualizability, \texorpdfstring{$3$}{3}-duality, and \texorpdfstring{$E_3$}{E-3}-Koszul self-duality}
\label{subsec:hcs-dualizability}

\begin{theorem}[Dualizability of \texorpdfstring{$\Obs_{\hCS}$}{Obs-hCS} and \texorpdfstring{$3$}{3}-duality]
\label{thm:plat-dualizability}
Let $Y = \C^3$ with the quantum observable $E_3$-algebra
$\Obs_{\hCS}(\C^3)$ of Theorem~\ref{thm:plat-hCS-quantum}, and let
$\partial\overline{\mathrm{Conf}}_2(\C^3) \simeq S^5$ be the link of
the fat diagonal in the Fulton--MacPherson compactification. Then:
\begin{enumerate}[label=\textup{(\roman*)}]
\item The $E_3$-trace defined by integration over
$S^5 \subset \partial\overline{\mathrm{Conf}}_2(\C^3)$ is
nondegenerate on $\Obs_{\hCS}(\C^3)$.
\item The abelian sector (gauge algebra $\fg = \fgl_1$) is
$3$-dualisable in the $(\infty,3)$-category of $E_3$-algebras.
\item On flat $\C^3$ with non-abelian semisimple $\fg$, the
$3$-dualizability is obstructed by an infinite-dimensional
Hochschild cohomology: by Gwilliam--Williams~$2021$
(arXiv:$2009.05037$) Prop.~$5.3.2$,
\[
 \HH^0_{E_3}(\Obs_{\hCS}(\C^3),\, \Obs_{\hCS}(\C^3))
 \;=\; \C[[\tau_1, \tau_2, \tau_3]],
\]
which has infinite cohomological dimension, violating the dualizability
finiteness axiom.
\item On a compact Calabi--Yau $3$-fold $X$, the $3$-dualizability is
restored: the Hochschild cohomology becomes finite-dimensional
under the compactness-plus-properness of $\Coh(X)$.
\item The $E_3$-operad is Koszul self-dual up to the shift:
\[
 \cD_3^! \;\simeq\; \mathrm{Lie}[2],
\]
by Fresse~$2017$ \emph{Homotopy of Operads \& Grothendieck--Teichm\"uller Groups}
Vol.~I Thm.~$14.1.A$, recovering the Calaque--Pantev--To\"en--Vaqui\'e
Koszul duality on the Poisson side. The strict
Koszul duality (Gwilliam--Williams~$2021$ \S$5.3$) and the
homotopy Koszul duality (Francis--Gaitsgory~$2012$) agree through
Fresse Thm.~$12.3.A$ and the Positselski coderived/contraderived
transfer (Positselski~$2011$~\S$6$).
\end{enumerate}
\end{theorem}

\begin{proof}
Item~(i) is the nondegeneracy of the $E_3$-trace on the Fulton--MacPherson
boundary, Costello--Gwilliam Vol~II~\S$6$. Item~(ii) is the abelian
specialisation where
$\Obs_{\hCS}(\C^3)|_{\fgl_1} \simeq \Sym(\Omega^{0,\bullet}(\C^3)[1])[[\hbar]]$
is a free $E_3$-algebra on a dualizable input; its
$3$-dualizability is Lurie~\emph{HA}~$5.3.2$ via the cobordism hypothesis
applied to the free $E_3$-algebra.
Item~(iii) is Gwilliam--Williams~$2021$ Prop.~$5.3.2$: the
$E_3$-Hochschild cohomology of the non-abelian free
holomorphic factorization algebra on $\C^3$ is
$\C[[\tau_1,\tau_2,\tau_3]]$, infinite-dimensional, ruling out
$3$-dualizability on open domains.
Item~(iv) is the proper-Calabi--Yau restoration: on compact $X$,
Hochschild cohomology agrees with derived global sections of an
internal-$\RHom$ sheaf on $X \times X$, finite-dimensional by proper
coherent duality.
Item~(v) is Fresse~$2017$ Vol.~I Thm.~$14.1.A$. The strict-vs-homotopy
compatibility is Fresse Thm.~$12.3.A$ in the Koszul-dual direction
together with Positselski's coderived/contraderived equivalence.
\end{proof}

\begin{remark}[Obstruction locus is \texorpdfstring{$\HH^0$}{HH-0}, not \texorpdfstring{$\HH^1$}{HH-1}]
\label{rem:plat-dualizability-HH0}
The $3$-dualizability obstruction in
Theorem~\ref{thm:plat-dualizability}(iii) lives in $\HH^0$, not $\HH^1$:
the infinite-dimensionality of $\HH^0_{E_3}(\Obs_{\hCS}(\C^3),
\Obs_{\hCS}(\C^3)) = \C[[\tau_1, \tau_2, \tau_3]]$ is the receptacle
of the formal deformations
$(\tau_1, \tau_2, \tau_3)$ of the Omega-background;
it does not vanish on flat $\C^3$ because the three equivariant
parameters $h_1, h_2, h_3$ provide three independent formal directions
in the unconstrained deformation space. The Calabi--Yau condition
$\sum h_i = 0$ cuts out a codimension-$1$ slice
$\HH^0_{\mathrm{CY}} = \C[[\tau_1, \tau_2, \tau_3]]/(\tau_1 + \tau_2 + \tau_3)$,
still infinite-dimensional. Compactness of $X$ collapses the formal
power series to a finite-dimensional quotient; properness is what
supplies the dualizability.
\end{remark}

\subsection{All-orders non-abelian compound: the simply-laced Yangian VOA}
\label{subsec:hcs-nonab-compound}

\begin{theorem}[All-orders non-abelian \texorpdfstring{$5$}{5}D hCS boundary for simply-laced \texorpdfstring{$\fg$}{g}]
\label{thm:plat-nonab-compound}
For every simply-laced semisimple Lie algebra \(\fg\) in the
Costello--Gaiotto--Yagi \(5\)d holomorphic Chern--Simons setting, the
boundary vertex algebra on
\(\C^2_{\epsilon_1,\epsilon_2}\times\R\), with Calabi--Yau slice
\(\epsilon_1+\epsilon_2+\epsilon_3=0\), is the Yangian VOA attached to
\(\fg\):
\[
 \partial \hCS_5(\fg)
 \;\simeq\;
 Y^{\mathrm{VOA}}_{\epsilon_1,\epsilon_2,\epsilon_3}(\fg),
\]
as a vertex algebra to all orders in the perturbative parameter.  An
affine or toroidal algebra \(Y_{\epsilon_1,\epsilon_2,\epsilon_3}(\widehat\fg)\)
requires additional affine/toroidal Hall or stable-envelope input and is
not a consequence of the \(5\)d theorem alone.
\end{theorem}

\begin{proof}
The proof concatenates twelve ingredients, each a theorem in the
primary literature:
\begin{enumerate}[label=\textup{(P\arabic*)}]
\item Costello~$2013$ (arXiv:$1303.2632$) sets up $5$D hCS on
$\C^2 \times \R$ with $\Omega$-background and identifies tree-level
boundary observables with the Yangian generators.
\item Costello--Gwilliam~$2017$ Vol~II supplies the factorisation-algebra
formalism that realises boundary observables as a chiral algebra.
\item Costello--Yagi~$2018$ (arXiv:$1810.01970$) proves the rank-$1$
case at all orders via $4$D Chern--Simons $\to$ classical Yangian
reduction.
\item Costello--Gaiotto~$2018$ (arXiv:$1810.10024$) establishes the
universal category-theoretic framework (vertex algebra at a corner).
\item Costello--Dimofte--Paquette~$2020$ (arXiv:$2005.00083$) proves the
rank-$1$ boundary chiral algebra all-orders theorem for $\fgl_1$ (the
$A_0$ case of the present statement).
\item Chan--Paton brane fibre identification: for ADE-types $A_n$,
$D_n$, $E_6$, $E_7$, the minuscule representation
$V_{\mathrm{min}} = \C^{N_\fg}$ realises the universal
Wilson-line module (Fulton--Harris~$1991$ \S$15.1$, Bourbaki
\emph{Groupes et alg\`ebres de Lie}~Ch.~$8$~\S$7.3$).
For $E_8$ (no minuscule representation), the Kostant--Slodowy
transverse slice to the subregular nilpotent orbit supplies the
Wilson-line fibre (Slodowy~$1980$ \S$8.3$; Kostant--Wallach~$2006$).
\item Francis~$2013$ (arXiv:$1107.0087$) Thm.~$2.29$ reduces the
chiral Chevalley--Eilenberg complex of an $E_1$-algebra to the
ordinary CE complex of its underlying Lie algebra, supplying the
passage from chain-level observables to Lie-algebra cohomology.
\item Whitehead's second lemma: $H^2_{\mathrm{Lie}}(\fg, \fg) = 0$ for
every semisimple $\fg$ (Chevalley--Eilenberg~$1948$ \S$6$;
Humphreys~$1972$ Thm.~$7.8.2$), uniformly killing quadratic
obstructions to deformation.
\item Frenkel--Ben-Zvi~$2004$ \emph{Vertex Algebras and Algebraic
Curves} Thm.~$3.4.3$ establishes simplicity of the vacuum module at
generic level.
\item Frenkel--Ben-Zvi~$2004$ Thm.~$18.4.2$ establishes the smoothness
of the space of $\fg^L$-opers at the critical level, supplying the
moduli of flat $\fg^L$-connections.
\item On Nakajima quiver-variety charts, Maulik--Okounkov~$2012$
(arXiv:$1211.1287$) supplies stable-envelope \(R\)-matrices matching
the Yangian \(R\)-matrix in the symplectic-resolution sector.  This is
the surface/quiver chart of the comparison, not a construction of a
generic affine/toroidal CY$_3$ algebra.
\item Proch\'azka--Rap\v{c}\'ak~$2018$ (arXiv:$1711.09993$) level-$n$
truncation supplies the matching of corner structure at finite $n$
(used in Corollary~\ref{cor:plat-corner}).
\end{enumerate}

\emph{Synthesis.}
P$1$--P$5$ establish the statement for $\fgl_1$ at all orders.
P$7$--P$9$ lift the chain-level identification to the vertex-algebra
level by reducing the chiral CE complex to ordinary Lie cohomology and
killing the quadratic obstruction uniformly in $\fg$.
P$6$ supplies the Chan--Paton brane fibre that threads simply-laced
ADE data through the universal construction of P$3$: for $A, D, E_6,
E_7$ via minuscule, for $E_8$ via Kostant--Slodowy.
P$10$ supplies the moduli of boundary conditions needed for the
critical-level compatibility (extending the generic-level vacuum of
P$9$).
P$11$ supplies the explicit \(R\)-matrix on the symplectic-resolution
charts on which the stable-envelope construction is defined.
The resulting composite is a morphism of vertex algebras that is an
isomorphism at tree level (by P$1$) and remains an isomorphism at all
orders in \(\hbar\) under the simply-laced hCS hypotheses of
Costello--Gaiotto--Yagi.  Whitehead vanishing and the vacuum simplicity
control the named deformation problem; they do not by themselves
construct an affine or toroidal ADE Hall algebra.
\end{proof}

\begin{corollary}[Corner identification at finite \texorpdfstring{$n$}{n}]
\label{cor:plat-corner}
At finite rank $n$, the boundary corner of
Theorem~\ref{thm:plat-nonab-compound} is
$Y_{0, 0, n}\!\bigl(\widehat{\mathfrak{sl}}_n\bigr)$, not the
Calabi--Yau-symmetric $Y_{n, n, n}$. The symmetric corner emerges only
in the $n \to \infty$ limit. Agreement between the two Yangian
presentations is the Proch\'azka--Rap\v{c}\'ak level-$n$ embedding
\[
 Y_{0,0,n}\!\bigl(\widehat{\mathfrak{sl}}_n\bigr)
 \;\hookrightarrow\;
 Y_{0,0,1}\!\bigl(\widehat{\mathfrak{gl}}_1\bigr)^{\otimes n}
 /\;(\text{level-}n\text{ identification})
\]
from Proch\'azka--Rap\v{c}\'ak~$2018$ (arXiv:$1711.09993$) \S$5$.
\end{corollary}

\begin{proof}
The asymmetric corner $Y_{0, 0, n}$ arises from the boundary
condition on $\partial \C^2_{\epsilon_1, \epsilon_2} \times \R$: the
$\R$-direction is not deformed by the Omega-background, so two of the
three Omega parameters are zero at finite $n$. The symmetric corner
$Y_{n,n,n}$ corresponds to the triple-$\Omega$-background of the $6$D
lift; it is reached only in the $n \to \infty$ limit where the
normal bundle of the boundary curve acquires the full three-parameter
structure. The level-$n$ embedding is Proch\'azka--Rap\v{c}\'ak~$2018$
\S$5$.
\end{proof}

\begin{remark}[Three variants of the chiral Yangian]
\label{rmk:plat-chiral-yangian}
Three distinct algebras travel under the name ``Yangian'' in the
holomorphic Chern--Simons hierarchy:
\begin{enumerate}[label=\textup{(\alph*)}]
\item \emph{$4$D CS boundary}: the classical Yangian $Y(\fg)$ of
Drinfeld~$1985$; it arises at the boundary of $4$D Chern--Simons on
$\Sigma \times \C$ (Costello--Yagi~$2018$) as a quasi-triangular Hopf
algebra with one spectral parameter, no conformal weight.
\item \emph{$5$D hCS boundary}: the affine Yangian
$Y_{\epsilon_1,\epsilon_2,\epsilon_3}(\widehat\fg)$ with conformal
weight (Theorem~\ref{thm:plat-nonab-compound}); it is a vertex
algebra, not merely a Hopf algebra, carrying two spectral parameters
and the chiral field content of the $5$D boundary theory.
\item \emph{Chiral Yangian $\mathcal{Y}^{\chir}(\fg)$}: the zero-mode
sector of (b) for $\Omega$-background partition functions, identified
with the locally finite-dimensional sub-vertex-algebra of
$Y_{\epsilon_1,\epsilon_2,\epsilon_3}(\widehat\fg)$ spanned by
$\hbar^0$ modes. This is the object that couples to Nekrasov
partition functions on $\mathbb{C}^2 \times \mathbb{C}^*$ and computes
instanton counting weights.
\end{enumerate}
The three algebras are mathematically distinct: (a) has no conformal
structure, (b) has both the $R$-matrix and the conformal weight,
and (c) has the conformal weight but only the zero-mode Hopf
structure. Identifications across (a)--(c) arise by specialisation:
(b) $\to$ (a) by forgetting conformal weight (the $\epsilon_3 \to 0$
limit of the third Omega parameter); (b) $\to$ (c) by projection to
the zero-mode sector. The all-orders compound of
Theorem~\ref{thm:plat-nonab-compound} is the statement at variant
(b).
\end{remark}
